% !TeX program = pdflatex
% Samodzielny skrypt: Teoria kategorii — podstawy i przyklady.
\documentclass[11pt,a4paper]{article}
\usepackage[T1]{fontenc}
\usepackage[utf8]{inputenc}
\usepackage[polish]{babel}
\usepackage{lmodern}
\usepackage{amsmath,amssymb,amsthm,mathtools,amscd}
\usepackage{mathrsfs}
\usepackage{microtype}
\usepackage{xcolor}
\usepackage{geometry}
\usepackage{enumitem}
\usepackage[unicode]{hyperref}
\hypersetup{colorlinks=true,linkcolor=blue!45!black,urlcolor=blue!45!black}
\geometry{margin=2.5cm}
\setlength{\parindent}{0pt}
\setlength{\parskip}{4pt}
\setlength{\emergencystretch}{2em}
\setlist[enumerate]{topsep=4pt,itemsep=2pt,leftmargin=*}
\setlist[itemize]{topsep=4pt,itemsep=2pt,leftmargin=*}
\numberwithin{equation}{section}
\theoremstyle{plain}
\newtheorem{twierdzenie}{Twierdzenie}[section]
\newtheorem{stwierdzenie}[twierdzenie]{Stwierdzenie}
\newtheorem{lemat}[twierdzenie]{Lemat}
\newtheorem{wniosek}[twierdzenie]{Wniosek}
\theoremstyle{definition}
\newtheorem{definicja}[twierdzenie]{Definicja}
\newtheorem{przyklad}[twierdzenie]{Przykład}
\theoremstyle{remark}
\newtheorem{uwaga}[twierdzenie]{Uwaga}
\newcommand{\cat}[1]{\mathsf{#1}}
\newcommand{\Set}{\cat{Set}}
\newcommand{\Grp}{\cat{Grp}}
\newcommand{\Mon}{\cat{Mon}}
\newcommand{\Ring}{\cat{Ring}}
\newcommand{\Top}{\cat{Top}}
\newcommand{\Vect}{\cat{Vect}}
\newcommand{\Ab}{\cat{Ab}}
\newcommand{\Pos}{\cat{Pos}}
\newcommand{\Rel}{\cat{Rel}}
\newcommand{\id}{\mathrm{id}}
\newcommand{\Hom}{\operatorname{Hom}}
\newcommand{\Ob}{\operatorname{Ob}}
\newcommand{\Nat}{\operatorname{Nat}}
\newcommand{\op}{\mathrm{op}}
\newcommand{\im}{\operatorname{im}}
\newcommand{\Ker}{\operatorname{ker}}
\newcommand{\colim}{\operatorname{colim}}
\DeclareMathOperator{\End}{End}
\DeclareMathOperator{\Aut}{Aut}
\DeclareMathOperator{\ev}{ev}
% Odsyłacze Atlasu: w samodzielnym PDF-ie bez zewnętrznego indeksu
% zachowujemy czytelny opis odsyłacza. Identyfikator pozostaje w źródle rozdziału.
\providecommand{\atlasref}[2]{#2}
\title{\textbf{Teoria kategorii od podstaw}\\\large Definicje, intuicje, przykłady i najważniejsze konstrukcje}
\author{}
\date{}

% BEGIN GLOBAL CODE LISTINGS (generated from latex/pl/code-listings.tex)
\usepackage{xcolor}
\usepackage{listings}
\definecolor{codebg}{RGB}{248,249,251}
\definecolor{codeborder}{RGB}{180,191,203}
\definecolor{codecomment}{RGB}{61,113,79}
\definecolor{codekeyword}{RGB}{41,78,132}
\lstdefinelanguage{Julia}{
  morekeywords={abstract,baremodule,begin,break,catch,const,continue,do,else,elseif,
    end,export,false,finally,for,function,global,if,import,in,isa,let,local,macro,
    module,mutable,primitive,quote,return,struct,true,try,type,using,where,while,
    missing,nothing,println,print},
  sensitive=true,morecomment=[l]{\#},morecomment=[n]{\#=}{=\#},
  morestring=[b]",
}
\lstdefinestyle{skrypt}{
  backgroundcolor=\color{codebg},frame=single,rulecolor=\color{codeborder},
  framesep=5pt,frameround=tttt,framerule=.5pt,
  basicstyle=\ttfamily\footnotesize,keywordstyle=\color{codekeyword}\bfseries,
  commentstyle=\color{codecomment}\itshape,stringstyle=\color{purple!65!black},
  numbers=left,numberstyle=\tiny\color{gray!75!black},numbersep=8pt,
  stepnumber=1,breaklines=true,breakatwhitespace=false,
  columns=fullflexible,keepspaces=true,showstringspaces=false,tabsize=4,
  linewidth=.90\linewidth,captionpos=b,aboveskip=11pt,belowskip=12pt,
  xleftmargin=12pt,xrightmargin=2pt,
}
\lstdefinestyle{pythonpde}{style=skrypt,language=Python}
\lstdefinestyle{python}{style=skrypt,language=Python}
\lstdefinestyle{julia}{style=skrypt,language=Julia}
\lstset{style=skrypt}
\renewcommand{\lstlistingname}{Kod}
\renewcommand{\lstlistlistingname}{Spis kodów}
% END GLOBAL CODE LISTINGS
\begin{document}
\maketitle
\begin{abstract}
Skrypt wprowadza język teorii kategorii, zakładając znajomość podstaw algebry liniowej, teorii mnogości i elementarnej algebry. Punktem wyjścia są konkretne przykłady: zbiory i funkcje, grupy i homomorfizmy, przestrzenie wektorowe i odwzorowania liniowe oraz zbiory uporządkowane. Następnie omawiamy funktory, transformacje naturalne, konstrukcje uniwersalne, granice i kogranice, lemat Yonedy, sprzężenia oraz podstawowe idee monad. Przy każdym nowym pojęciu wyjaśniamy zarówno jego formalną treść, jak i to, jakie pytanie matematyczne rozwiązuje. Ćwiczenia przeznaczone do samodzielnego rozwiązania znajdują się w osobnym zbiorze zadań; przykłady z objaśnieniami pozostają w skrypcie.
\end{abstract}
\tableofcontents
\newpage

\section{Po co teoria kategorii?}\label{sec:motywacja}

W wielu działach matematyki pojawia się podobna sytuacja. Mamy pewnego rodzaju obiekty i odwzorowania, które zachowują strukturę tych obiektów. Dla zbiorów są to funkcje, dla grup homomorfizmy, dla przestrzeni topologicznych odwzorowania ciągłe, a dla przestrzeni wektorowych odwzorowania liniowe. Odwzorowania można składać, a każdy obiekt ma odwzorowanie identycznościowe.

Teoria kategorii bada własności, które można sformułować wyłącznie za pomocą obiektów, odwzorowań i ich składania. Nie usiłuje zastąpić algebry lub topologii. Pozwala natomiast rozpoznawać te same konstrukcje w różnych działach matematyki i wyjaśnia, dlaczego są one do siebie podobne.

\begin{przyklad}[Iloczyn kartezjański i iloczyn grup]
Dla zbiorów $A,B$ definiujemy $A\times B$ i rzuty $p_A:A\times B\to A$, $p_B:A\times B\to B$. Jeśli mamy funkcje $f:X\to A$ i $g:X\to B$, istnieje dokładnie jedna funkcja $h:X\to A\times B$ spełniająca $p_A\circ h=f$ oraz $p_B\circ h=g$, mianowicie $h(x)=(f(x),g(x))$. Ta sama własność charakteryzuje iloczyn grup z rzutami będącymi homomorfizmami. Rzeczywiście, dla homomorfizmów $f:X\to A$ i $g:X\to B$ oraz $x,y\in X$ mamy
\[
h(xy)=(f(xy),g(xy))=(f(x)f(y),g(x)g(y))=h(x)h(y).
\]
Zatem $h$ jest homomorfizmem. Jeżeli inny homomorfizm $u:X\to A\times B$ ma te same rzuty, to dla każdego $x\in X$ obie współrzędne $u(x)$ są równe odpowiednim współrzędnym $h(x)$, więc $u=h$. Własność tę nazwiemy \emph{uniwersalną}. Jej sformułowanie, w przeciwieństwie do dowodu istnienia, nie odwołuje się do elementów iloczynu. Rozwiniemy ją w rozdziale~\ref{sec:produkty}.
\end{przyklad}

\begin{przyklad}[Ta sama własność w zbiorze uporządkowanym]\label{cat:ex:motywacja-porzadek}
Rozważmy zbiór częściowo uporządkowany $(P,\leq)$. Między $a$ i $b$ dopuszczamy dokładnie jeden morfizm $a\to b$, jeśli $a\leq b$, a w przeciwnym razie żadnego. Zwrotność zapewnia identyczności, przechodniość pozwala składać morfizmy, a łączność wynika z tego, że między ustalonymi obiektami istnieje co najwyżej jeden morfizm.

Odpowiednikiem iloczynu $a$ i $b$ jest ich kres dolny $m$, o ile istnieje. Istnienie rzutów $m\to a$ i $m\to b$ oznacza $m\leq a$ i $m\leq b$. Jeżeli $x\leq a$ oraz $x\leq b$, własność uniwersalna wymaga dokładnie jednego morfizmu $x\to m$, czyli $x\leq m$. Są to dokładnie warunki definiujące największe ograniczenie dolne. W $\mathbb R$ z naturalnym porządkiem iloczynem w tym sensie jest $\min(a,b)$, a nie zwykłe mnożenie liczb.
\end{przyklad}

\begin{przyklad}[Dlaczego wybór morfizmów ma znaczenie?]\label{cat:ex:motywacja-morfizmy}
Rozważmy wszystkie zbiory. Jeśli jako morfizmy dopuszczamy wszystkie funkcje, każdy singleton $S$ jest obiektem, do którego z dowolnego zbioru $X$ prowadzi dokładnie jeden morfizm: funkcja przyjmująca jedyną możliwą wartość (dla $X=\varnothing$ jest to funkcja pusta).

Jeżeli dopuścimy tylko funkcje różnowartościowe, identyczności i złożenia nadal są dozwolone, lecz ze zbioru dwuelementowego do singletonu nie istnieje żaden morfizm. Co więcej, w tej kategorii nie istnieje obiekt, do którego z każdego zbioru prowadzi dokładnie jeden morfizm. Taki obiekt musiałby mieć dokładnie jeden element, ponieważ morfizmy z singletonu odpowiadają jego elementom. Tymczasem ze zbioru dwuelementowego nie ma do niego iniekcji. Pojęcie obiektu końcowego omówimy w rozdziale~\ref{sec:uniwersalne}.
\end{przyklad}

\begin{uwaga}[Jak czytać diagramy]
Zapis $A\xrightarrow{f}B$ oznacza morfizm $f$ z $A$ do $B$. Diagram jest \emph{przemienny}, gdy dowolne dwie skierowane drogi prowadzące od tego samego obiektu do tego samego obiektu końcowego wyznaczają ten sam morfizm. Przykładowo diagram
\begin{center}
\begin{tikzpicture}[>=stealth]
 \node (a) at (0,2) {$A$};
 \node (b) at (4,2) {$B$};
 \node (c) at (0,0) {$C$};
 \node (d) at (4,0) {$D$};
 \draw[->] (a) -- node[above] {$f$} (b);
 \draw[->] (a) -- node[left] {$g$} (c);
 \draw[->] (b) -- node[right] {$h$} (d);
 \draw[->] (c) -- node[below] {$k$} (d);
\end{tikzpicture}
\end{center}
jest przemienny wtedy i tylko wtedy, gdy $h\circ f=k\circ g$. Diagram jest więc graficznym zapisem równania. W $h\circ f$ najpierw wykonujemy $f$, a potem $h$. Sam rysunek kwadratu nie zapewnia przemienności; trzeba ją założyć lub wykazać.
\end{uwaga}


\section{Kategorie i ich pierwsze przykłady}\label{sec:kategorie}

\subsection{Definicja kategorii}

\begin{definicja}[Kategoria]\label{def:kategoria}
\emph{Kategorią} $\mathcal C$ nazywamy dane obejmujące:
\begin{enumerate}
\item klasę obiektów $\Ob(\mathcal C)$;
\item dla każdej uporządkowanej pary obiektów $A,B$ kolekcję morfizmów $\Hom_{\mathcal C}(A,B)$; w tym skrypcie zakładamy, że jest ona zbiorem; każdy morfizm traktujemy wraz z jego dziedziną $A$ i przeciwdziedziną $B$;
\item dla każdej trójki obiektów $A,B,C$ operację składania
$\Hom_{\mathcal C}(B,C)\times\Hom_{\mathcal C}(A,B)\to\Hom_{\mathcal C}(A,C)$, $(g,f)\mapsto g\circ f$;
\item dla każdego obiektu $A$ wyróżniony morfizm $\id_A\in\Hom_{\mathcal C}(A,A)$.
\end{enumerate}
Dane te muszą spełniać dwa aksjomaty:
\begin{enumerate}[label=(\roman*)]
\item \emph{Łączność}: jeśli $f:A\to B$, $g:B\to C$ i $h:C\to D$, to $h\circ(g\circ f)=(h\circ g)\circ f$.
\item \emph{Jednostki}: dla każdego $f:A\to B$ zachodzi $f\circ\id_A=f=\id_B\circ f$.
\end{enumerate}
Kategorię nazywamy \emph{lokalnie małą}, jeśli dla każdej pary obiektów jej kolekcja morfizmów jest zbiorem; to właśnie konwencja przyjęta dalej. Kategoria jest \emph{mała}, jeśli jej obiekty i wszystkie morfizmy tworzą zbiory. W kategorii lokalnie małej o zbiorze obiektów zbiór wszystkich morfizmów otrzymujemy jako sumę rozłączną zbiorów $\Hom_{\mathcal C}(A,B)$ indeksowaną parami obiektów. Kategorię, która nie jest mała, nazywamy dużą.
\end{definicja}

Morfizmy nie muszą być funkcjami pomiędzy zbiorami. Aksjomaty wymagają tylko możliwości ich składania i istnienia identyczności. W zapisie $g\circ f$ najpierw wykonujemy $f$, a następnie $g$.

\begin{lemat}[Jednoznaczność identyczności]
Jeśli $e,e':A\to A$ są morfizmami spełniającymi aksjomat jednostki dla wszystkich morfizmów o odpowiednich dziedzinach i przeciwdziedzinach, to $e=e'$.
\end{lemat}
\begin{proof}
Z faktu, że $e$ jest jednostką lewostronną, wynika $e\circ e'=e'$. Z faktu, że $e'$ jest jednostką prawostronną, wynika $e\circ e'=e$. Zatem $e=e'$.
\end{proof}


\begin{uwaga}[Kwestie teoriomnogościowe]
Obiekty kategorii nie muszą tworzyć zbioru: na przykład ``wszystkie zbiory'' tworzą klasę właściwą. Przyjmujemy standardową konwencję, że rozpatrywane kategorie są lokalnie małe, a słowo \emph{mały diagram} oznacza diagram indeksowany małą kategorią. Gdy chcemy tworzyć kategorię wszystkich funktorów pomiędzy dużymi kategoriami, trzeba dodatkowo kontrolować wielkość kolekcji transformacji naturalnych. Przy głównych zastosowaniach rozróżnienie to zaznaczymy jawnie.
\end{uwaga}


\subsection{Wielkość kategorii}

\begin{definicja}[Uniwersum Grothendiecka]\label{cat:def:uniwersum}
Uniwersum Grothendiecka $\mathcal U$ to zbiór przechodni domknięty na tworzenie par, zbiorów potęgowych oraz sum rodzin $(x_i)_{i\in I}$ takich, że $I\in\mathcal U$ i $x_i\in\mathcal U$ dla każdego $i\in I$. Przyjmujemy ponadto $\mathbb N\in\mathcal U$. Przechodniość oznacza, że $x\in y\in\mathcal U$ implikuje $x\in\mathcal U$. Zbiory należące do $\mathcal U$ nazywamy $\mathcal U$-małymi.
\end{definicja}

\begin{uwaga}[Relatywność pojęcia małości]\label{cat:rem:uniwersa}
Jeżeli $X,Y\in\mathcal U$, to $X\times Y\in\mathcal U$ (pary uporządkowane tworzymy za pomocą par nieuporządkowanych, a zbiór wszystkich takich par jest podzbiorem odpowiedniego iterowanego zbioru potęgowego). Każdy podzbiór $X\times Y$ również należy do $\mathcal U$: należy on do $\mathcal P(X\times Y)$, a ta należy do $\mathcal U$. W szczególności zbiór funkcji $X\to Y$, utożsamianych z ich wykresami, jest $\mathcal U$-mały, gdyż jest podzbiorem $\mathcal P(X\times Y)$. Natomiast kolekcja wszystkich $\mathcal U$-małych zbiorów to samo $\mathcal U$, które nie należy do siebie na mocy aksjomatu ufundowania. Kategoria tych zbiorów jest więc lokalnie $\mathcal U$-mała, ale nie jest $\mathcal U$-mała. Może być mała względem większego uniwersum zawierającego $\mathcal U$.

Istnienia dowolnie dużych uniwersów tego rodzaju nie przyjmujemy jako konsekwencji standardowych aksjomatów ZFC. Jeżeli posługujemy się uniwersami, odpowiednie założenie teoriomnogościowe wskazujemy jawnie; w pozostałym tekście wystarczy konwencja klas i zbiorów z definicji~\ref{def:kategoria}.
\end{uwaga}

\subsection{Kategorie znane z innych dziedzin}

\begin{przyklad}[Zbiory]
W kategorii $\Set$ obiektami są zbiory, morfizmami są funkcje, a składanie i identyczności są zwykłymi operacjami na funkcjach. Jest to przykład kategorii lokalnie małej, ale nie małej: nie istnieje zbiór wszystkich zbiorów.
\end{przyklad}

\begin{przyklad}[Grupy i monoidy]
W $\Grp$ obiektami są grupy, a morfizmami homomorfizmy grup. W $\Mon$ obiektami są monoidy (zbiory z łącznym działaniem i elementem neutralnym), a morfizmami homomorfizmy zachowujące działanie i element neutralny. Złożenie homomorfizmów i identyczność są homomorfizmami, dlatego otrzymujemy kategorie.
\end{przyklad}

\begin{przyklad}[Przestrzenie liniowe]
Ustalmy ciało $k$. W kategorii $\Vect_k$ obiektami są przestrzenie wektorowe nad $k$, a morfizmami odwzorowania $k$-liniowe. Kategoria $\Vect_k^{\mathrm{fd}}$ to jej pełna podkategoria przestrzeni skończenie wymiarowych: oznacza to, że rozpatrujemy tylko takie obiekty, lecz wszystkie odwzorowania liniowe pomiędzy nimi. W kategorii $\Ab$ obiektami są grupy abelowe, a morfizmami homomorfizmy.
\end{przyklad}

\begin{przyklad}[Topologia i pierścienie]
W $\Top$ obiektami są przestrzenie topologiczne, a morfizmami odwzorowania ciągłe. W $\Ring$ rozważamy pierścienie przemienne z jedynką i homomorfizmy zachowujące $1$. Ważne jest określenie konwencji: bez podania, jakie pierścienie i jakie homomorfizmy rozważamy, nazwa kategorii nie wyznacza całej struktury.
\end{przyklad}

\begin{przyklad}[Moduły i rozmaitości różniczkowe]\label{cat:ex:moduly-rozmaitosci}
Dla pierścienia $R$ z jedynką oznaczmy przez $\cat{Mod}_R^{\mathrm{l}}$ kategorię lewych unitarnych $R$-modułów i homomorfizmów $R$-liniowych. Nie zakładamy przemienności $R$. Jeśli $f:M\to N$ i $g:N\to P$ są liniowe, to dla $r,s\in R$ oraz $x,y\in M$ zachodzi
\[
(g\circ f)(rx+sy)=g(rf(x)+sf(y))
=r(g\circ f)(x)+s(g\circ f)(y).
\]
Identyczność również jest liniowa, więc składanie funkcji wyznacza kategorię. Jest to przykład oparty na konstrukcjach z rozdziału 16 skryptu \emph{Algebra abstrakcyjna}.

W kategorii $\cat{Man}$ obiektami są gładkie rozmaitości bez brzegu (Hausdorffa, z przeliczalną bazą topologii), a morfizmami są gładkie odwzorowania. Złożenie jest gładkie, ponieważ w mapach lokalnych jest złożeniem gładkich funkcji między otwartymi podzbiorami przestrzeni euklidesowych; identyczność w tej samej mapie lokalnej jest identycznością euklidesową. Korzystamy tu z pojęć w rozdziale 1 skryptu \emph{Geometria i topologia różniczkowa}.
\end{przyklad}

\begin{definicja}[Kategoria konkretna]\label{cat:def:konkretna}
Kategoria konkretna nad $\Set$ to kategoria $\mathcal C$ wyposażona w przyporządkowanie każdemu obiektowi $A$ zbioru $U(A)$ i każdemu morfizmowi $f:A\to B$ funkcji $U(f):U(A)\to U(B)$, zachowujące identyczności i składanie oraz rozróżniające morfizmy o ustalonych dziedzinie i przeciwdziedzinie: $U(f)=U(g)$ implikuje $f=g$. W języku rozdziału~\ref{sec:funktory} jest to wierny funktor $U:\mathcal C\to\Set$.
\end{definicja}

\begin{przyklad}
Kategorie grup, modułów, przestrzeni topologicznych i rozmaitości są konkretne dzięki przyporządkowaniu zbioru bazowego oraz tej samej funkcji bazowej. Konkretność nie wymaga, aby każda funkcja między zbiorami bazowymi zachowywała strukturę: np. nie każda funkcja między grupami jest homomorfizmem.
\end{przyklad}

\begin{przyklad}[Kategoria dyskretna]
Dla zbioru $S$ zdefiniujmy kategorię $\mathcal D_S$ o obiektach $S$, w której $\Hom(x,x)=\{\id_x\}$, a $\Hom(x,y)=\varnothing$ dla $x\neq y$. Morfizmami są wyłącznie identyczności. Taka kategoria przechowuje informację o elementach $S$, ale nie o relacjach między nimi.
\end{przyklad}

\begin{przyklad}[Zbiory uporządkowane]\label{przyklad:poset}
Niech $(P,\leq)$ będzie zbiorem częściowo uporządkowanym. Obiektami kategorii $\mathcal P$ są elementy $P$, a $\Hom_{\mathcal P}(x,y)$ jest zbiorem jednoelementowym, gdy $x\leq y$, i pustym w przeciwnym przypadku. Zwrotność zapewnia identyczności, przechodniość zapewnia składanie, a łączność wynika z tego, że między ustalonymi obiektami istnieje co najwyżej jeden morfizm.

Dla przykładu zbiór $\{1,2,3,6\}$ uporządkowany przez podzielność ma morfizm $2\to 6$, ale nie ma morfizmu $6\to 2$. Nie należy interpretować tych strzałek jako funkcji pomiędzy liczbami. Każdy \emph{praporządek} (relacja zwrotna i przechodnia, niekoniecznie antysymetryczna) także definiuje kategorię.
\end{przyklad}

\begin{przyklad}[Monoid jako kategoria z jednym obiektem]\label{przyklad:monoid-kat}
Niech $(M,\cdot,1)$ będzie monoidem. Tworzymy kategorię z jednym obiektem $*$, przyjmując $\Hom(*,*)=M$. Składanie definiujemy przez $g\circ f=g\cdot f$, a identycznością jest $1$. Aksjomaty kategorii to dokładnie łączność i istnienie jednostki w monoidzie.

W szczególności grupa $G$ jest kategorią z jednym obiektem, której wszystkie morfizmy są odwracalne. Tym samym teoria kategorii zawiera teorię monoidów i grup jako szczególne przypadki.
\end{przyklad}

\begin{definicja}[Grupoid]\label{cat:def:grupoid}
Grupoid to kategoria, w której każdy morfizm $f:A\to B$ ma odwrotny morfizm $g:B\to A$, czyli $g\circ f=\id_A$ oraz $f\circ g=\id_B$. Taki morfizm $g$ jest jedyny: dla drugiej odwrotności $g'$ mamy $g=g\circ\id_B=g\circ(f\circ g')=(g\circ f)\circ g'=\id_A\circ g'=g'$.
\end{definicja}

\begin{przyklad}[Grupoidy]\label{cat:ex:grupoid}
Kategoria dyskretna jest grupoidem. Grupa rozpatrywana jako kategoria z jednym obiektem również jest grupoidem. Dla dowolnego zbioru $S$ możemy zbudować grupoid par: obiektami są elementy $S$, a jedynym morfizmem $a\to b$ jest para $(a,b)$. Definiujemy $(b,c)\circ(a,b)=(a,c)$ oraz $\id_a=(a,a)$. Oba nawiasowania złożenia trzech strzałek dają tę samą parę z początkiem pierwszej i końcem ostatniej strzałki; identyczności nie zmieniają tych końców. Odwrotnością $(a,b)$ jest $(b,a)$. Inny ważny przykład stanowi grupoid podstawowy przestrzeni topologicznej, którego strzałkami są klasy homotopii dróg o ustalonych końcach (por. rozdział 4 skryptu \emph{Topologia algebraiczna}); jego formalną konstrukcję odkładamy do części o zastosowaniach.
\end{przyklad}

\begin{przyklad}[Macierze i standardowe współrzędne]\label{przyklad:mat}
W kategorii $\cat{Mat}_k$ obiektami są liczby $0,1,2,\ldots$, a morfizmami $n\to m$ macierze o $m$ wierszach i $n$ kolumnach z wyrazami w $k$. Składaniem jest mnożenie macierzy w tej samej kolejności co składanie odwzorowań, a identycznością macierz jednostkowa $I_n$ (dla $n=0$ pusta macierz $0\times0$).

Przykładowo macierz $\begin{pmatrix}1&2\\0&1\\3&0\end{pmatrix}$ jest morfizmem $2\to3$. Kategoria ta koduje wszystkie skończenie wymiarowe przestrzenie liniowe \emph{z wyborem wzorcowego modelu} $k^n$.
\end{przyklad}

\begin{przyklad}[Relacje zamiast funkcji]
W $\Rel$ obiektami są zbiory, a morfizm $R:A\to B$ to relacja $R\subseteq A\times B$. Dla $R:A\to B$ i $S:B\to C$ definiujemy $S\circ R=\{(a,c):\text{istnieje }b\in B\text{ z }(a,b)\in R,(b,c)\in S\}$. Identycznością na $A$ jest przekątna $\Delta_A=\{(a,a):a\in A\}$. Aby wykazać łączność, niech dodatkowo $T:C\to D$ będzie relacją. Dla $a\in A$ i $d\in D$ warunek $(a,d)\in T\circ(S\circ R)$ jest równoważny istnieniu $b\in B$ i $c\in C$ spełniających
\[
(a,b)\in R,\qquad (b,c)\in S,\qquad (c,d)\in T.
\]
Dokładnie ten sam warunek opisuje przynależność do $(T\circ S)\circ R$. Ponadto $(a,b)\in R\circ\Delta_A$ wtedy i tylko wtedy, gdy istnieje $a'\in A$ takie, że $a=a'$ i $(a',b)\in R$, czyli gdy $(a,b)\in R$. Analogicznie $\Delta_B\circ R=R$.

Na przykład dla $A=B=C=\{1,2\}$ oraz $R=\{(1,1),(1,2)\}$, $S=\{(1,2),(2,2)\}$ otrzymujemy $S\circ R=\{(1,2)\}$. Kategoria $\Rel$ pokazuje, że kategorie nie muszą być zbudowane z funkcji.
\end{przyklad}


\begin{przyklad}[Ścieżki w grafie skierowanym]\label{przyklad:sciezki}
Rozważmy graf z wierzchołkami $a,b,c$ i krawędziami $u:a\to b$, $v:b\to c$. Obiektami kategorii ścieżek są wierzchołki; morfizmami są skończone skierowane ścieżki, w tym ścieżki długości $0$ (identyczności). W szczególności istnieje morfizm $v\circ u:a\to c$, chociaż w wyjściowym grafie nie było osobnej krawędzi $a\to c$. Składanie polega na konkatenacji skończonych ciągów krawędzi: najpierw zapisujemy krawędzie pierwszej ścieżki, potem drugiej. Dwa nawiasowania konkatenacji trzech ciągów dają dokładnie ten sam ciąg, a dopisanie pustego ciągu w odpowiednim wierzchołku nie zmienia ścieżki. Otrzymujemy przykład \emph{kategorii swobodnej na grafie}. Pokazuje on, że sam graf skierowany nie jest jeszcze kategorią: trzeba uwzględnić również wszystkie wymagane złożenia.
\end{przyklad}

\begin{przyklad}[Przestrzenie z wyróżnionym punktem]
W kategorii $\Top_*$ obiektami są pary $(X,x_0)$ z przestrzenią topologiczną $X$ i wyróżnionym punktem $x_0\in X$. Morfizm $f:(X,x_0)\to(Y,y_0)$ to odwzorowanie ciągłe spełniające $f(x_0)=y_0$. Przykładowo funkcja $t\mapsto t^2$ daje morfizm $([0,1],0)\to([0,1],0)$, ale nie daje morfizmu $([0,1],1)\to([0,1],0)$. Ten sam wzór funkcji może więc być morfizmem w jednej kategorii, a w innej nie.
\end{przyklad}

\subsection{Podkategorie i kategoria przeciwna}

\begin{definicja}[Podkategoria]
\emph{Podkategorią} $\mathcal D$ kategorii $\mathcal C$ nazywamy kategorię, której obiekty są obiektami $\mathcal C$, zbiory morfizmów $\Hom_{\mathcal D}(A,B)$ są podzbiorami $\Hom_{\mathcal C}(A,B)$, a identyczności i składanie są odziedziczone z $\mathcal C$. Jest \emph{pełna}, jeśli dla dowolnych $A,B\in\Ob(\mathcal D)$ zachodzi $\Hom_{\mathcal D}(A,B)=\Hom_{\mathcal C}(A,B)$.
\end{definicja}

\begin{przyklad}
Przestrzenie skończenie wymiarowe tworzą pełną podkategorię $\Vect_k$. Z kolei kategoria mająca te same obiekty co $\Vect_k$, ale jedynie liniowe izomorfizmy jako morfizmy, nie jest pełną podkategorią $\Vect_k$: pomija np. odwzorowanie zerowe $k\to k$.
\end{przyklad}

\begin{definicja}[Kategoria przeciwna]\label{def:op}
Dla kategorii $\mathcal C$ \emph{kategorię przeciwną} $\mathcal C^{\op}$ otrzymujemy, pozostawiając te same obiekty i odwracając kierunki morfizmów: $\Hom_{\mathcal C^{\op}}(A,B)=\Hom_{\mathcal C}(B,A)$. Piszemy $f^{\op}:B\to A$ na oznaczenie strzałki odpowiadającej $f:A\to B$ w $\mathcal C$. Jeżeli $f:A\to B$ oraz $g:B\to C$, definiujemy $f^{\op}\circ_{\mathcal C^{\op}}g^{\op}=(g\circ_{\mathcal C}f)^{\op}$ oraz $\id_A^{\op}$ jako identyczność w kategorii przeciwnej. Dla dodatkowego $h:C\to D$ mamy
\[
(f^{\op}\circ g^{\op})\circ h^{\op}
=(h\circ(g\circ f))^{\op}
=((h\circ g)\circ f)^{\op}
=f^{\op}\circ(g^{\op}\circ h^{\op}).
\]
Obie tożsamości jednostkowe wynikają z $(f\circ\id_A)^{\op}=f^{\op}=(\id_B\circ f)^{\op}$. Zatem kategoria przeciwna spełnia aksjomaty kategorii.
\end{definicja}

\begin{uwaga}
Dwukrotne odwrócenie strzałek daje $(\mathcal C^{\op})^{\op}=\mathcal C$. Jeśli kategoria pochodzi ze zbioru uporządkowanego, jej kategoria przeciwna odpowiada porządkowi odwrotnemu. Pojęcie \emph{dualności} polega na tym, że każde twierdzenie wynikające wyłącznie z aksjomatów kategorii ma wersję z odwróconymi wszystkimi strzałkami.
\end{uwaga}


\section{Izomorfizmy, monomorfizmy i epimorfizmy}\label{sec:morfizmy}

\begin{definicja}[Izomorfizm]
Morfizm $f:A\to B$ w kategorii $\mathcal C$ jest \emph{izomorfizmem}, jeśli istnieje $g:B\to A$ takie, że $g\circ f=\id_A$ i $f\circ g=\id_B$. Morfizm $g$ nazywamy odwrotnością $f$ i oznaczamy $f^{-1}$. Piszemy $A\cong B$, jeśli istnieje izomorfizm między $A$ i $B$.
\end{definicja}

\begin{stwierdzenie}
Odwrotność izomorfizmu jest jednoznaczna. Złożenie dwóch izomorfizmów jest izomorfizmem, a odwrotnością złożenia jest złożenie odwrotności w odwrotnej kolejności.
\end{stwierdzenie}
\begin{proof}
Jeśli $g,h:B\to A$ są odwrotnościami $f:A\to B$, to $g=g\circ(f\circ h)=(g\circ f)\circ h=h$. Dla izomorfizmów $f:A\to B$, $u:B\to C$ mamy $(u\circ f)\circ(f^{-1}\circ u^{-1})=\id_C$ oraz $(f^{-1}\circ u^{-1})\circ(u\circ f)=\id_A$, korzystając z łączności.
\end{proof}

\begin{przyklad}[Co oznacza izomorfizm?]
W $\Set$ jest nim bijekcja, w $\Grp$ izomorfizm grup, w $\Vect_k$ odwracalne odwzorowanie liniowe, a w $\Top$ homeomorfizm. W kategorii pochodzącej z porządku częściowego jedynymi izomorfizmami są identyczności (bo $x\leq y$ i $y\leq x$ pociągają $x=y$). W kategorii praporządku różne obiekty mogą być izomorficzne.
\end{przyklad}

\begin{definicja}[Monomorfizm i epimorfizm]\label{def:mono-epi}
Morfizm $f:A\to B$ jest \emph{monomorfizmem}, jeśli dla każdego obiektu $X$ i wszystkich $u,v:X\to A$ zachodzi $f\circ u=f\circ v\Rightarrow u=v$. Morfizm $f:A\to B$ jest \emph{epimorfizmem}, jeśli dla każdego obiektu $Y$ i wszystkich $u,v:B\to Y$ zachodzi $u\circ f=v\circ f\Rightarrow u=v$.
\end{definicja}

Inaczej mówiąc, monomorfizmy można skracać z lewej strony równania, a epimorfizmy z prawej. Definicje te nie używają elementów ani obrazu odwzorowania.

\begin{twierdzenie}\label{tw:set-mono-epi}
Dla funkcji $f:A\to B$ w $\Set$ zachodzą równoważności: $f$ jest monomorfizmem wtedy i tylko wtedy, gdy jest różnowartościowa; $f$ jest epimorfizmem wtedy i tylko wtedy, gdy jest surjektywna.
\end{twierdzenie}
\begin{proof}
Jeżeli $f$ jest różnowartościowa, to równość $f\circ u=f\circ v$ daje $u(x)=v(x)$ dla każdego $x$. Odwrotnie, jeśli $f(a)=f(a')$, rozważmy funkcje $u,v:\{*\}\to A$ dane przez $u(*)=a$, $v(*)=a'$. Moniczność $f$ daje $u=v$, a więc $a=a'$.

Jeżeli $f$ jest surjektywna, to z $u\circ f=v\circ f$ wynika $u=v$ przez sprawdzenie na dowolnym $b=f(a)$. Jeśli $f$ nie jest surjektywna, wybierzmy $b_0\notin f(A)$ oraz funkcje $u,v:B\to\{0,1\}$, gdzie $u$ jest stale równa $0$, a $v(b_0)=1$ i $v(b)=0$ dla $b\neq b_0$. Mamy $u\neq v$, lecz $u\circ f=v\circ f$.
\end{proof}

\begin{przyklad}[Uwaga na kategorię]\label{przyklad:epi-pierscienie}
W kategorii $\Ring$ inkluzja $\mathbb Z\hookrightarrow\mathbb Q$ jest epimorfizmem, chociaż nie jest surjektywna. Istotnie, jeżeli $u,v:\mathbb Q\to R$ są homomorfizmami pierścieni z jedynką, które zgadzają się na $\mathbb Z$, to dla $n\in\mathbb Z\setminus\{0\}$ elementy $u(1/n)$ i $v(1/n)$ są obydwa odwrotnościami tego samego elementu $u(n)=v(n)$, a więc są równe. W konsekwencji $u(a/n)=v(a/n)$ dla wszystkich liczb wymiernych. Zatem $u=v$.

Przeciwnie, w kategorii pochodzącej z dowolnego porządku częściowego \emph{każdy istniejący} morfizm jest jednocześnie monomorfizmem i epimorfizmem, ponieważ każdy zbiór $\Hom(x,y)$ ma co najwyżej jeden element. Nie oznacza to, że każdy morfizm jest tam izomorfizmem.
\end{przyklad}

\begin{definicja}[Sekcja i retrakcja]
Dla morfizmu $f:A\to B$ morfizm $s:B\to A$ spełniający $f\circ s=\id_B$ nazywamy \emph{sekcją} $f$ (prawą odwrotnością $f$); sam $f$ jest wówczas \emph{rozszczepionym epimorfizmem}. Morfizm $r:B\to A$ spełniający $r\circ f=\id_A$ nazywamy \emph{retrakcją} $f$ (lewą odwrotnością $f$); sam $f$ jest wówczas \emph{rozszczepionym monomorfizmem}.
\end{definicja}

\begin{stwierdzenie}
Każdy rozszczepiony monomorfizm jest monomorfizmem, a każdy rozszczepiony epimorfizm jest epimorfizmem.
\end{stwierdzenie}
\begin{proof}
Jeżeli $r\circ f=\id_A$ i $f\circ u=f\circ v$, to $u=r\circ f\circ u=r\circ f\circ v=v$. Jeżeli $f\circ s=\id_B$ i $u\circ f=v\circ f$, to $u=u\circ f\circ s=v\circ f\circ s=v$.
\end{proof}

\begin{przyklad}
W $\Vect_k$ zanurzenie $i:k\to k^2$, $i(x)=(x,0)$, ma lewą odwrotność $p:k^2\to k$, $p(x,y)=x$, bo $p\circ i=\id_k$. Zatem $i$ jest rozszczepionym monomorfizmem, a $p$ rozszczepionym epimorfizmem. W $\Set$ funkcja $\varnothing\to\{*\}$ jest monomorfizmem, ale nie ma lewej odwrotności: nie istnieje funkcja $\{*\}\to\varnothing$.
\end{przyklad}

\begin{uwaga}
Morfizm będący jednocześnie monomorfizmem i epimorfizmem nazywamy \emph{bimorfizmem}. Nie musi on być izomorfizmem; świadczy o tym przykład $2\to6$ w kategorii porządku podzielności z przykładu~\ref{przyklad:poset}. Kategorie, w których każdy taki morfizm jest izomorfizmem, nazywa się \emph{zrównoważonymi} (ang. \emph{balanced}). Kategoria $\Set$ jest zrównoważona na mocy twierdzenia~\ref{tw:set-mono-epi}.
\end{uwaga}


\begin{przyklad}[Funktor zapominania nie musi zachowywać epimorfizmów]
Dla pierścieni przemiennych z jedynką funktor $U:\Ring\to\Set$ zapomina strukturę. Jak pokazano w przykładzie~\ref{przyklad:epi-pierscienie}, morfizm $\mathbb Z\hookrightarrow\mathbb Q$ jest epimorfizmem w $\Ring$. Jego obraz przez $U$ nie jest epimorfizmem w $\Set$, gdyż nie jest surjektywny. To konkretny kontrprzykład do błędnego przekonania, że każdy funktor zachowuje epimorfizmy.
\end{przyklad}

\begin{przyklad}[Monomorfizmy w kategorii grup]
Homomorfizm grup $f:G\to H$ jest monomorfizmem wtedy i tylko wtedy, gdy jest różnowartościowy. Jeden kierunek wynika ze skracania funkcji różnowartościowych. Jeśli $f(a)=f(b)$, weźmy dwa homomorfizmy $u,v:\mathbb Z\to G$ określone przez $u(1)=a$ i $v(1)=b$. Wówczas $f\circ u=f\circ v$, więc moniczność wymusza $u=v$, a stąd $a=b$. W ten sposób dla grup rolę jednoelementowego zbioru testowego z kategorii $\Set$ pełni grupa wolna na jednym generatorze.
\end{przyklad}

\subsection{Reguły skracania i odwrotności jednostronne}

\begin{stwierdzenie}\label{cat:prop:skracanie}
Złożenie monomorfizmów jest monomorfizmem, a złożenie epimorfizmów jest epimorfizmem. Jeśli $g\circ f$ jest monomorfizmem, to $f$ jest monomorfizmem; jeśli $g\circ f$ jest epimorfizmem, to $g$ jest epimorfizmem. Każdy izomorfizm jest bimorfizmem.
\end{stwierdzenie}
\begin{proof}
Niech $f:A\to B$, $g:B\to C$. Jeśli $g\circ f\circ u=g\circ f\circ v$, moniczność $g$ daje $f\circ u=f\circ v$, a moniczność $f$ daje $u=v$. Jeśli $u\circ g\circ f=v\circ g\circ f$, epiczność $f$ daje $u\circ g=v\circ g$, a epiczność $g$ daje $u=v$.
Gdy $f\circ u=f\circ v$, po złożeniu z $g$ z lewej strony i skróceniu monomorfizmu $g\circ f$ dostajemy $u=v$. Analogicznie z $u\circ g=v\circ g$ i epiczności $g\circ f$ otrzymujemy $u=v$. Izomorfizm ma obie odwrotności jednostronne, zatem jest zarazem rozszczepionym monomorfizmem i rozszczepionym epimorfizmem.
\end{proof}

\begin{stwierdzenie}\label{cat:prop:bimorfizm-rozszczepiony}
Monomorfizm mający prawą odwrotność jest izomorfizmem. Epimorfizm mający lewą odwrotność jest izomorfizmem.
\end{stwierdzenie}
\begin{proof}
Jeśli $f\circ s=\id_B$ i $f$ jest monomorfizmem, to $f\circ(s\circ f)=f=f\circ\id_A$. Po skróceniu dostajemy $s\circ f=\id_A$. Jeśli $r\circ f=\id_A$ i $f$ jest epimorfizmem, to $(f\circ r)\circ f=f=\id_B\circ f$, więc $f\circ r=\id_B$.
\end{proof}

\subsection{Przykłady z algebry i topologii}

\begin{twierdzenie}[Epimorfizmy grup]\label{cat:thm:epi-grp}
Epimorfizmy w $\Grp$ to dokładnie homomorfizmy surjektywne.
\end{twierdzenie}
\begin{proof}
Surjektywność wystarcza, ponieważ funkcje o tej samej dziedzinie, zgadzające się na całym obrazie surjekcji, są równe. Niech teraz $f:G\to H$ nie będzie surjekcją i niech $K=f(G)$ będzie właściwą podgrupą $H$. Rozważmy zbiór $X=H\times\{0,1\}$ i homomorfizm $u:H\to\operatorname{Sym}(X)$ określony przez $u(h)(x,i)=(hx,i)$. Zdefiniujmy permutację $\tau$ zbioru $X$, która zamienia $(x,0)$ z $(x,1)$ dla $x\in K$, a wszystkie pozostałe punkty pozostawia na miejscu. Niech $v(h)=\tau\circ u(h)\circ\tau^{-1}$. Jest to homomorfizm, ponieważ sprzężenie przez ustaloną permutację zachowuje składanie.

Dla $k\in K$ zachodzi $kx\in K$ wtedy i tylko wtedy, gdy $x\in K$. Stąd $u(k)$ komutuje z $\tau$, więc $v(k)=u(k)$ i $v\circ f=u\circ f$. Jeżeli $h\in H\setminus K$, to jednak $u(h)(1,0)=(h,0)$, podczas gdy $v(h)(1,0)=\tau(h,1)=(h,1)$. Zatem $u\ne v$ i $f$ nie jest epimorfizmem.
\end{proof}

\begin{twierdzenie}[Morfizmy modułów]\label{cat:thm:mono-epi-mod}
W kategorii lewych unitarnych modułów nad pierścieniem $R$ z jedynką monomorfizmy są dokładnie homomorfizmami różnowartościowymi, a epimorfizmy dokładnie homomorfizmami surjektywnymi. Kategoria ta jest zrównoważona.
\end{twierdzenie}
\begin{proof}
Dla $f:M\to N$ różnowartościowość i surjektywność implikują odpowiednie skracanie funkcji. Jeśli $f(x)=f(y)$, dwa homomorfizmy $a,b:R\to M$ dane przez $a(r)=rx$, $b(r)=ry$ mają równe złożenia z $f$. Moniczność daje $a=b$, więc $x=a(1)=b(1)=y$. Jeśli $f$ nie jest surjektywna, moduł $N/f(M)$ jest niezerowy. Rzut ilorazowy $q:N\to N/f(M)$ różni się od homomorfizmu zerowego, ale $q\circ f=0$. Zatem $f$ nie jest epimorfizmem. Bijektywny homomorfizm modułów ma liniową odwrotność: dla $n=f(x)$, $n'=f(y)$ zachodzi $f^{-1}(rn+sn')=rx+sy$. Wykorzystujemy konstrukcje ilorazów i homomorfizmów z rozdziału 16 skryptu \emph{Algebra abstrakcyjna}.
\end{proof}

\begin{przyklad}[Bimorfizm bez rozszczepienia]\label{cat:ex:mono-epi-top}
W $\Top$ monomorfizmy to dokładnie ciągłe funkcje różnowartościowe: do dowodu konieczności używamy dwóch odwzorowań z przestrzeni jednopunktowej. Epimorfizmy to dokładnie ciągłe surjekcje. Jeśli ciągła funkcja $f:X\to Y$ nie jest surjekcją, wybieramy $y_0\in Y\setminus f(X)$ i dwie funkcje do zbioru $\{0,1\}$ różniące się jedynie w $y_0$, tak jak w dowodzie twierdzenia~\ref{tw:set-mono-epi}. Nadajemy zbiorowi docelowemu topologię antydyskretną $\{\varnothing,\{0,1\}\}$, dzięki czemu obie funkcje są ciągłe. Dowodzi to braku epiczności.

Identyczność ze zbioru $\{0,1\}$ z topologią dyskretną do tego samego zbioru z topologią antydyskretną jest ciągłą bijekcją, ale jej odwrotność nie jest ciągła (przeciwobraz singletonu nie jest otwarty). Jest zatem bimorfizmem, który nie jest izomorfizmem. Na mocy stwierdzenia~\ref{cat:prop:bimorfizm-rozszczepiony} nie jest też rozszczepiony z żadnej strony.
\end{przyklad}

\begin{przyklad}[Gładka bijekcja nie musi być dyfeomorfizmem]
Odwzorowanie $f:\mathbb R\to\mathbb R$, $f(x)=x^3$, jest gładką bijekcją, więc można je skracać z obu stron w kategorii rozmaitości. Odwrotność nie jest różniczkowalna w zerze, gdyż dla $t>0$ iloraz różnicowy wynosi $t^{1/3}/t=t^{-2/3}$ i jest nieograniczony, gdy $t\to0^+$. Stąd $f$ nie jest izomorfizmem w tej kategorii. Pojęcie dyfeomorfizmu pochodzi z rozdziału 1 skryptu \emph{Geometria i topologia różniczkowa}.
\end{przyklad}

\begin{przyklad}[Lokalizacja jako epimorfizm]\label{cat:ex:epi-lokalizacja}
Uogólnijmy przykład~\ref{przyklad:epi-pierscienie}. Niech $R$ będzie pierścieniem przemiennym z jedynką, $S$ jego zbiorem multiplikatywnym, a $j:R\to S^{-1}R$ kanonicznym homomorfizmem. Własność uniwersalna lokalizacji (rozdział 12 skryptu \emph{Algebra abstrakcyjna}) mówi, że homomorfizm z $R$, który przekształca elementy $S$ w elementy odwracalne, ma dokładnie jedno przedłużenie wzdłuż $j$. Jeśli $u\circ j=v\circ j$, wspólne złożenie spełnia ten warunek, ponieważ $j(s)$ jest odwracalny dla każdego $s\in S$. Jednoznaczność daje $u=v$. Zatem $j$ jest epimorfizmem w $\Ring$ bez założenia, że $j$ jest różnowartościowy.
\end{przyklad}

\subsection{Monomorfizmy i epimorfizmy regularne}

\begin{definicja}[Regularność morfizmu]\label{cat:def:regularne}
Monomorfizm regularny to morfizm $m:A\to B$, dla którego istnieją $u,v:B\to C$ takie, że $u\circ m=v\circ m$, a każdy $t:X\to B$ spełniający $u\circ t=v\circ t$ ma dokładnie jedną faktoryzację $t=m\circ w$ z $w:X\to A$. Dualnie epimorfizm regularny $e:B\to A$ spełnia $e\circ u=e\circ v$ dla pewnych $u,v:C\to B$, a każdy $t:B\to X$ spełniający $t\circ u=t\circ v$ ma dokładnie jedną faktoryzację $t=w\circ e$. Są to konstrukcje uniwersalne omawiane w rozdziale~\ref{sec:equalizatory}.
\end{definicja}

\begin{stwierdzenie}\label{cat:prop:regularne-mono-epi}
Monomorfizm regularny jest monomorfizmem, a epimorfizm regularny jest epimorfizmem.
\end{stwierdzenie}
\begin{proof}
Jeśli $m\circ a=m\circ b=t$, to $t$ spełnia $u\circ t=v\circ t$ i zarówno $a$, jak i $b$ są jego faktoryzacjami. Jednoznaczność daje $a=b$. Jeśli $a\circ e=b\circ e=t$, to $t\circ u=t\circ v$ i obie faktoryzacje $a,b$ muszą być równe.
\end{proof}

\begin{przyklad}
W $\Set$ każda iniekcja $m:A\to B$ jest regularna: niech $u:B\to\{0,1\}$ będzie stale zerowa, a $v$ będzie zerowa dokładnie na $m(A)$ i jedynkowa poza nim. Funkcja $t:X\to B$ spełnia $u\circ t=v\circ t$ dokładnie wtedy, gdy $t(X)\subseteq m(A)$; wtedy $w(x)$ jest jedynym elementem $A$ o obrazie $t(x)$. Każda surjekcja $e:B\to A$ jest regularna: bierzemy relację $C=\{(b,b')\in B^2:e(b)=e(b')\}$ oraz jej dwa rzuty $u,v:C\to B$. Jeśli $t\circ u=t\circ v$, to $t$ jest stała na włóknach $e$. Wzór $w(a)=t(b)$ dla dowolnego $b$ z $e(b)=a$ jest poprawny i jednoznacznie określa $w$; $t=w\circ e$. Nie trzeba tutaj wybierać globalnego przekroju surjekcji.
\end{przyklad}


\section{Funktory: odwzorowania między kategoriami}\label{sec:funktory}

\subsection{Definicja i przykłady}

\begin{definicja}[Funktor kowariantny]\label{def:funktor}
Niech $\mathcal C,\mathcal D$ będą kategoriami. \emph{Funktor (kowariantny)} $F:\mathcal C\to\mathcal D$ przyporządkowuje każdemu obiektowi $A$ obiekt $F(A)$ i każdemu morfizmowi $f:A\to B$ morfizm $F(f):F(A)\to F(B)$ w taki sposób, że dla wszystkich odpowiednich morfizmów zachodzi $F(\id_A)=\id_{F(A)}$ oraz $F(g\circ f)=F(g)\circ F(f)$.
\end{definicja}

Funktor zachowuje \emph{sposób organizacji strzałek}. Nie musi zachowywać dosłownie wszystkich dodatkowych struktur, które możemy rozpoznać w wyjściowej kategorii.

\begin{przyklad}[Funktor zapominania]
Funktor $U:\Grp\to\Set$ przypisuje grupie $G$ jej zbiór elementów $U(G)$, a homomorfizmowi $\varphi:G\to H$ tę samą funkcję, traktowaną wyłącznie jako odwzorowanie zbiorów. Niczego nie zmienia w funkcji $\varphi$; jedynie \emph{zapomina}, że funkcja zachowuje mnożenie. Analogicznie mamy $U:\Vect_k\to\Set$ i $U:\Top\to\Set$.
\end{przyklad}

\begin{przyklad}[Fundamentalny przykład z algebry liniowej]
Funktor $F:\cat{Mat}_k\to\Vect_k^{\mathrm{fd}}$ definiujemy przez $F(n)=k^n$ i $F(A)(x)=Ax$. Zachowanie składania jest równością $(BA)x=B(Ax)$, a zachowanie identyczności wynika z $I_nx=x$. Każda skończenie wymiarowa przestrzeń wektorowa jest izomorficzna z pewnym $F(n)$, lecz nie jest dosłownie równa takiej przestrzeni bez wyboru bazy.
\end{przyklad}

\begin{przyklad}[Potęga symetryczna jako funktor]
Ustalmy $r\geq1$. Dla przestrzeni $V$ nad $k$ oznaczmy przez $\operatorname{Sym}^r(V)$ $r$-tą potęgę symetryczną, a dla $f:V\to W$ przez $\operatorname{Sym}^r(f)$ jedyne odwzorowanie liniowe określone na iloczynach symetrycznych wzorem $v_1\cdots v_r\mapsto f(v_1)\cdots f(v_r)$. Odwzorowanie to istnieje, ponieważ $r$-liniowe odwzorowanie $(v_1,\ldots,v_r)\mapsto f(v_1)\cdots f(v_r)$ jest symetryczne, więc faktoryzuje się przez potęgę symetryczną. Iloczyny symetryczne generują tę przestrzeń, a na każdym takim generatorze zarówno $\operatorname{Sym}^r(g\circ f)$, jak i $\operatorname{Sym}^r(g)\circ\operatorname{Sym}^r(f)$ dają $g(f(v_1))\cdots g(f(v_r))$. Dla identyczności każdy generator pozostaje niezmieniony. Zatem $\operatorname{Sym}^r:\Vect_k\to\Vect_k$ jest funktorem. To przykład konstrukcji algebraicznej działającej również na odwzorowaniach.
\end{przyklad}

\begin{przyklad}[Grupa działająca na zbiorze]\label{przyklad:dzialanie-funktor}
Traktujmy grupę $G$ jako kategorię z jednym obiektem $*$, jak w przykładzie~\ref{przyklad:monoid-kat}. Funktor $F:G\to\Set$ wybiera zbiór $X=F(*)$ i dla każdego $g\in G$ bijekcję $F(g):X\to X$ (bijekcję, ponieważ $F(g^{-1})$ jest odwrotnością $F(g)$). Równość $F(gh)=F(g)\circ F(h)$ oznacza dokładnie aksjomat lewego działania $g\cdot(h\cdot x)=(gh)\cdot x$. A zatem \emph{lewe działania grupy $G$ na zbiorach są tym samym co funktory $G\to\Set$}.
\end{przyklad}

\begin{przyklad}[Funktory między kategoriami z jednym obiektem]
Monoidy $M,N$ rozumiemy jako kategorie z jednym obiektem. Funktor $F:M\to N$ jest dokładnie homomorfizmem monoidów: funkcja $M\to N$ musi zachować jednostkę i iloczyn. Jeśli oba monoidy są grupami, otrzymujemy zwykłe homomorfizmy grup. Tak jeden i ten sam aksjomat funktora przybiera postać znajomego warunku algebraicznego.
\end{przyklad}


\begin{przyklad}[Zbiór składowych spójnych jako funktor]
Niech $c(X)$ oznacza zbiór składowych spójnych przestrzeni topologicznej $X$. Oznaczenie $\pi_0(X)$ rezerwujemy dla zbioru składowych łukowej spójności (klas punktów połączonych drogą), które nie zawsze pokrywają się ze składowymi spójnymi. Dla ciągłego $f:X\to Y$ obraz każdej składowej spójnej $K\subseteq X$ jest spójny, więc zawiera się w dokładnie jednej składowej $L\subseteq Y$. Możemy zatem zdefiniować $c(f)(K)=L$. Identyczność pozostawia każdą składową na miejscu. Jeżeli $g:Y\to Z$ jest ciągłe, a $f(K)\subseteq L=c(f)(K)$, to $g(f(K))\subseteq g(L)\subseteq c(g)(L)$, zatem jednoznaczność składowej zawierającej niepusty zbiór spójny daje $c(g\circ f)(K)=c(g)(c(f)(K))$. Otrzymujemy funktor $c:\Top\to\Set$. Dla sumy topologicznej dwóch niepustych przedziałów składowe są dwie; dla przedziału $[0,1]$ jest jedna. Funktor wydobywa niezmiennik topologiczny i pamięta działanie ciągłych odwzorowań na tym niezmienniku.
\end{przyklad}

\begin{przyklad}[Funktor określony przez diagram]
Niech $\mathcal J$ będzie kategorią z obiektami $0,1$ i jedną nietrywialną strzałką $a:0\to1$. Funktor $F:\mathcal J\to\Set$ jest dokładnie wyborem dwóch zbiorów $X=F(0)$, $Y=F(1)$ i funkcji $f=F(a):X\to Y$. Dla $X=\{1,2,3\}$, $Y=\{a,b\}$ i $f(1)=f(2)=a$, $f(3)=b$ otrzymujemy konkretny diagram, który sam jest funktorem. Zmiana kształtu kategorii indeksującej prowadzi do innych typów diagramów.
\end{przyklad}


\subsection{Funktory kontrawariantne}

\begin{definicja}[Funktor kontrawariantny]
\emph{Funktor kontrawariantny} z $\mathcal C$ do $\mathcal D$ jest funktorem kowariantnym $F:\mathcal C^{\op}\to\mathcal D$. W zapisie z morfizmami kategorii $\mathcal C$ oznacza to, że strzałce $f:A\to B$ przypisujemy strzałkę $F(f):F(B)\to F(A)$, przy czym $F(\id_A)=\id_{F(A)}$ oraz $F(g\circ f)=F(f)\circ F(g)$.
\end{definicja}

\begin{przyklad}[Przeciwobraz]
Dla zbioru $X$ niech $\mathcal P(X)$ oznacza zbiór jego podzbiorów. Funkcji $f:X\to Y$ przyporządkujmy $\mathcal P(f):\mathcal P(Y)\to\mathcal P(X)$, $B\mapsto f^{-1}(B)$. Równość $(g\circ f)^{-1}(C)=f^{-1}(g^{-1}(C))$ daje funktor kontrawariantny $\mathcal P:\Set^{\op}\to\Set$.

Istnieje też \emph{inny}, kowariantny funktor potęgowy $\mathcal P_+:\Set\to\Set$, który funkcji $f:X\to Y$ przypisuje $A\mapsto f(A)$. Warto nie mylić tych dwóch konstrukcji.
\end{przyklad}

\begin{przyklad}[Przestrzeń dualna]\label{przyklad:dualny}
Dla przestrzeni wektorowej $V$ niech $V^*=\Hom_k(V,k)$. Odwzorowaniu liniowemu $f:V\to W$ przypisujemy $f^*:W^*\to V^*$, $\lambda\mapsto\lambda\circ f$. Mamy $(g\circ f)^*=f^*\circ g^*$, a zatem $(-)^*:\Vect_k^{\op}\to\Vect_k$ jest funktorem kontrawariantnym. Dla skończenie wymiarowych $V$ jest on szczególnie dobrze związany z podwójną dualnością, którą omówimy w rozdziale~\ref{sec:naturalne}.
\end{przyklad}

\begin{stwierdzenie}\label{tw:funktory-iso}
Każdy funktor $F:\mathcal C\to\mathcal D$ przenosi izomorfizmy na izomorfizmy. Jeśli $f:A\to B$ jest izomorfizmem w $\mathcal C$, to $F(f)^{-1}=F(f^{-1})$.
\end{stwierdzenie}
\begin{proof}
Z $f^{-1}\circ f=\id_A$ otrzymujemy $F(f^{-1})\circ F(f)=F(\id_A)=\id_{F(A)}$. Podobnie $F(f)\circ F(f^{-1})=\id_{F(B)}$.
\end{proof}

\begin{uwaga}
Funktor nie musi przenosić dowolnych monomorfizmów na monomorfizmy ani dowolnych epimorfizmów na epimorfizmy. Powód jest subtelny: w definicji mono- lub epimorfizmu kwantyfikujemy po \emph{wszystkich} morfizmach w danej kategorii, a po przejściu do drugiej kategorii mogą pojawić się nowe morfizmy, których nie widzieliśmy w wyjściowej.
\end{uwaga}

\subsection{Wierność i pełność}

\begin{definicja}[Funktor wierny i pełny]\label{def:pelny}
Funktor $F:\mathcal C\to\mathcal D$ indukuje dla każdego $A,B\in\Ob(\mathcal C)$ funkcję $F_{A,B}:\Hom_{\mathcal C}(A,B)\to\Hom_{\mathcal D}(F(A),F(B))$, $f\mapsto F(f)$. Funktor jest \emph{wierny}, jeśli wszystkie $F_{A,B}$ są różnowartościowe; jest \emph{pełny}, jeśli wszystkie $F_{A,B}$ są surjektywne; jest \emph{w pełni wierny}, jeśli wszystkie są bijekcjami.
\end{definicja}

\begin{przyklad}
Funktor zapominania $U:\Grp\to\Set$ jest wierny: dwa homomorfizmy o tej samej funkcji na elementach są równe. Nie jest pełny, gdyż nie każda funkcja między zbiorami bazowymi grup jest homomorfizmem.
\end{przyklad}

\begin{przyklad}
Włączenie pełnej podkategorii $\Vect_k^{\mathrm{fd}}\hookrightarrow\Vect_k$ jest w pełni wierne. Funktor $\cat{Mat}_k\to\Vect_k^{\mathrm{fd}}$ z poprzedniego przykładu również jest w pełni wierny: każdemu odwzorowaniu liniowemu $k^n\to k^m$ odpowiada dokładnie jedna macierz o rozmiarach $m\times n$.
\end{przyklad}


\begin{stwierdzenie}[W pełni wierne funktory odzwierciedlają izomorfizmy]\label{tw:ff-reflect-iso}
Niech $F:\mathcal C\to\mathcal D$ będzie w pełni wierny. Jeśli $f:A\to B$ jest morfizmem w $\mathcal C$ i $F(f)$ jest izomorfizmem w $\mathcal D$, to $f$ jest izomorfizmem w $\mathcal C$.
\end{stwierdzenie}
\begin{proof}
Pełność daje $g:B\to A$ z $F(g)=F(f)^{-1}$. Z funktorialności wynika $F(g\circ f)=F(g)\circ F(f)=\id_{F(A)}=F(\id_A)$, więc wierność daje $g\circ f=\id_A$. Analogicznie $F(f\circ g)=F(f)\circ F(g)=\id_{F(B)}=F(\id_B)$, a wierność daje $f\circ g=\id_B$.
\end{proof}

\begin{przyklad}[Wierność bez pełności]
Funktor zapominania $U:\Grp\to\Set$ jest wierny, lecz niepełny. Na przykład funkcja $q:\mathbb Z\to\mathbb Z$, $q(n)=n^2$, jest funkcją zbiorów, ale nie jest homomorfizmem grup addytywnych, ponieważ $q(1+1)=4\neq 2=q(1)+q(1)$. Pełność wymaga reprezentowania \emph{każdej} funkcji między odpowiednimi zbiorami jako obrazu homomorfizmu, co w tym przykładzie zawodzi.
\end{przyklad}


\subsection{Funktory Hom}

\begin{definicja}[Funktory reprezentowane przez obiekt]\label{def:hom-funktory}
Niech $\mathcal C$ będzie kategorią lokalnie małą i $A\in\Ob(\mathcal C)$. Definiujemy funktor kowariantny $h^A:\mathcal C\to\Set$ przez $h^A(X)=\Hom_{\mathcal C}(A,X)$ oraz $h^A(f)(u)=f\circ u$ dla $f:X\to Y$ i $u:A\to X$. Definiujemy też funktor kontrawariantny $h_A:\mathcal C^{\op}\to\Set$ przez $h_A(X)=\Hom_{\mathcal C}(X,A)$ oraz $h_A(f)(v)=v\circ f$ dla $f:X\to Y$ i $v:Y\to A$.
\end{definicja}

\begin{proof}[Sprawdzenie aksjomatów funktora]
Dla $u:A\to X$ mamy $h^A(\id_X)(u)=\id_X\circ u=u$. Dla $f:X\to Y$, $g:Y\to Z$ zachodzi $h^A(g\circ f)(u)=(g\circ f)\circ u=g\circ(f\circ u)$. Dla $v:Z\to A$ analogicznie $h_A(g\circ f)(v)=v\circ(g\circ f)=(v\circ g)\circ f=h_A(f)(h_A(g)(v))$, a $h_A(\id_X)(w)=w\circ\id_X=w$ dla $w:X\to A$.
\end{proof}

\begin{przyklad}
W $\Set$ funktor $h^1=\Hom_{\Set}(\{*\},-)$ jest naturalnie izomorficzny z funktorem identycznościowym, ponieważ każda funkcja $\{*\}\to X$ jest wyznaczona przez obraz $*$. Natomiast $h^{\{0,1\}}(X)$ jest naturalnie utożsamiany ze zbiorem $X\times X$, ponieważ funkcję na zbiorze dwuelementowym opisują dwa jej obrazy.
\end{przyklad}

\subsection{Zachowywanie i odzwierciedlanie własności morfizmów}

\begin{stwierdzenie}\label{cat:prop:faithful-reflect}
Funktor wierny odzwierciedla monomorfizmy i epimorfizmy. Każdy funktor zachowuje rozszczepione monomorfizmy i rozszczepione epimorfizmy.
\end{stwierdzenie}
\begin{proof}
Niech $F$ będzie wierny i $F(f)$ monomorfizmem. Jeżeli $f\circ u=f\circ v$, to $F(f)\circ F(u)=F(f)\circ F(v)$, skąd $F(u)=F(v)$, a więc $u=v$. Jeżeli $F(f)$ jest epimorfizmem, z $u\circ f=v\circ f$ analogicznie otrzymujemy $F(u)\circ F(f)=F(v)\circ F(f)$, następnie $F(u)=F(v)$ i $u=v$. Jeżeli $r\circ f=\id_A$, to $F(r)\circ F(f)=\id_{F(A)}$; jeżeli $f\circ s=\id_B$, to $F(f)\circ F(s)=\id_{F(B)}$. To daje odpowiednie rozszczepienia w kategorii docelowej.
\end{proof}

\begin{przyklad}[Granice tych własności]\label{cat:ex:functor-limits}
Wierny funktor zapominania $\Top\to\Set$ nie odzwierciedla izomorfizmów: bijekcja ciągła z dwupunktowej przestrzeni dyskretnej do antydyskretnej jest kontrprzykładem z rozdziału~\ref{sec:morfizmy}. Funktor $\Ring\to\Set$ nie zachowuje epimorfizmów (przykład~\ref{przyklad:epi-pierscienie}). Funktor do kategorii z jednym obiektem i jedną strzałką jest pełny tylko wtedy, gdy każda para obiektów źródłowych ma przynajmniej jeden morfizm; dla kategorii z jednym obiektem odpowiadającej nietrywialnej grupie jest zatem pełny, ale niewierny. Pokazuje to niezależność pełności od wierności.
\end{przyklad}

\subsection{Kategorie iloczynowe i bifunktory}

\begin{definicja}[Kategoria iloczynowa i bifunktor]\label{cat:def:bifunctor}
Kategoria $\mathcal C\times\mathcal D$ ma obiekty $(A,B)$ i morfizmy $(f,g):(A,B)\to(A',B')$, gdzie $f:A\to A'$ i $g:B\to B'$. Identyczności i składanie definiujemy współrzędnie:
\[
(f',g')\circ(f,g)=(f'\circ f,g'\circ g),
\qquad \id_{(A,B)}=(\id_A,\id_B).
\]
Aksjomaty kategorii wynikają z odpowiednich równości w każdej współrzędnej. Bifunktor to funktor $F:\mathcal C\times\mathcal D\to\mathcal E$.
\end{definicja}

\begin{stwierdzenie}[Bifunktor Hom]\label{cat:prop:hom-bifunctor}
Dla kategorii lokalnie małej $\mathcal C$ istnieje bifunktor
\[
\Hom_{\mathcal C}:\mathcal C^{\op}\times\mathcal C\longrightarrow\Set.
\]
Na obiektach przyjmuje wartości $\Hom_{\mathcal C}(A,B)$, a dla $a:A'\to A$, $b:B\to B'$ i $u:A\to B$ działa przez $u\mapsto b\circ u\circ a$.
\end{stwierdzenie}
\begin{proof}
Obraz ma poprawny typ $A'\to B'$. Identyczności po obu stronach nie zmieniają $u$. Dla kolejnych strzałek $a':A''\to A'$ i $b':B'\to B''$ dwukrotne zastosowanie bifunktora daje $b'\circ(b\circ u\circ a)\circ a'=(b'\circ b)\circ u\circ(a\circ a')$, czyli dokładnie obraz złożenia w kategorii iloczynowej. Ustalając pierwszy lub drugi argument, odzyskujemy funktory z definicji~\ref{def:hom-funktory}.
\end{proof}

\begin{przyklad}[Iloczyn tensorowy jako bifunktor]\label{cat:ex:tensor-bifunctor}
Niech $R$ będzie pierścieniem przemiennym z jedynką. Korzystamy z własności uniwersalnej iloczynu tensorowego z rozdziału 18 skryptu \emph{Algebra abstrakcyjna}. Dla homomorfizmów $f:M\to M'$ i $g:N\to N'$ odwzorowanie $(x,y)\mapsto f(x)\otimes g(y)$ jest dwuliniowe, więc indukuje jedyny homomorfizm $f\otimes g:M\otimes_R N\to M'\otimes_R N'$. Dla kolejnych homomorfizmów $f',g'$ mamy
\[
((f'\otimes g')\circ(f\otimes g))(x\otimes y)
=(f'\circ f)(x)\otimes(g'\circ g)(y).
\]
Tensory proste generują iloczyn tensorowy jako moduł, zatem zachowana jest kompozycja; identyczności sprawdzamy tak samo na tensorach prostych. Otrzymujemy bifunktor kowariantny.

Funktor $M\mapsto M\otimes_{\mathbb Z}\mathbb Z/2\mathbb Z$ nie zachowuje monomorfizmów: iniekcję $2:\mathbb Z\to\mathbb Z$ przeprowadza na odwzorowanie zerowe niezerowego modułu $\mathbb Z/2\mathbb Z$ w siebie. Wykorzystujemy tutaj izomorfizm $\mathbb Z\otimes_{\mathbb Z}\mathbb Z/2\mathbb Z\cong\mathbb Z/2\mathbb Z$ oraz relację $2\otimes y=1\otimes 2y=0$.
\end{przyklad}

\subsection{Funktory w geometrii i topologii}

\begin{przyklad}[Spektrum jako funktor kontrawariantny]\label{cat:ex:spec-functor}
Dla pierścienia przemiennego $R$ z jedynką rozważmy przestrzeń $\operatorname{Spec} R$ ideałów pierwszych z topologią Zariskiego (rozdział 4 skryptu \emph{Geometria algebraiczna}). Homomorfizm $\varphi:R\to S$ zachowujący jedynkę wyznacza funkcję $\operatorname{Spec} S\to\operatorname{Spec} R$, $\mathfrak p\mapsto\varphi^{-1}(\mathfrak p)$. Przeciwobraz ideału pierwszego jest ideałem pierwszym: nie zawiera jedynki, a z $\varphi(ab)=\varphi(a)\varphi(b)\in\mathfrak p$ wynika $\varphi(a)\in\mathfrak p$ lub $\varphi(b)\in\mathfrak p$. Funkcja jest ciągła, gdyż przeciwobraz otwartego zbioru bazowego $D_R(r)=\{\mathfrak q:r\notin\mathfrak q\}$ to $D_S(\varphi(r))$. Dla $\psi:S\to T$ i $\mathfrak p\in\operatorname{Spec} T$ mamy
\[
(\psi\circ\varphi)^{-1}(\mathfrak p)=\varphi^{-1}(\psi^{-1}(\mathfrak p)).
\]
Identyczność działa przez identyczność. Otrzymujemy $\operatorname{Spec}:\Ring^{\op}\to\Top$. W tym miejscu rozpatrujemy jedynie przestrzeń topologiczną; konstrukcję morfizmu snopów strukturalnych i funktora do kategorii schematów afinicznych pozostawiamy do dalszej części.
\end{przyklad}

\begin{przyklad}[Homologia kompleksów]\label{cat:ex:homology-functor}
Dla kompleksu łańcuchowego $C_\bullet$ modułów nad ustalonym pierścieniem oznaczmy przez $H_n(C_\bullet)=\ker d_n/\operatorname{im}d_{n+1}$ jego homologię stopnia $n$. Morfizm kompleksów $f:C_\bullet\to D_\bullet$ spełnia $d_n^D f_n=f_{n-1}d_n^C$. Przeprowadza więc cykle na cykle, a brzegi na brzegi: jeśli $x=d_{n+1}^C y$, to $f_n(x)=d_{n+1}^D f_{n+1}(y)$. Wzór $H_n(f)([x])=[f_n(x)]$ jest dobrze określony na klasach ilorazowych. Zachowanie identyczności i składania wynika z $\id([x])=[x]$ oraz $H_n(g\circ f)([x])=[g_n(f_n(x))]$. Jest to konstrukcja z rozdziału 3 skryptu \emph{Topologia algebraiczna}, przedstawiona tutaj jako funktor kowariantny.
\end{przyklad}

\begin{przyklad}[Cofnięcie form różniczkowych i kohomologia de Rhama]\label{cat:ex:derham-functor}
Dla gładkiego $f:M\to N$ cofnięcie $f^*:\Omega^q(N)\to\Omega^q(M)$ określa wzór
\[
(f^*\omega)_x(v_1,\ldots,v_q)=
\omega_{f(x)}(\mathrm df_x(v_1),\ldots,\mathrm df_x(v_q)).
\]
Dla $g:N\to P$ reguła łańcuchowa $\mathrm d(g\circ f)_x=\mathrm dg_{f(x)}\circ\mathrm df_x$ daje $(g\circ f)^*=f^*\circ g^*$ przez podstawienie do powyższego wzoru. Cofnięcie identycznością pozostawia formę bez zmian. Korzystamy ponadto z udowodnionej w rozdziale 10 skryptu \emph{Geometria i topologia różniczkowa} zgodności $d(f^*\omega)=f^*(d\omega)$ z różniczką zewnętrzną. Wobec tego cofnięcie zachowuje formy zamknięte i dokładne oraz indukuje $H^q_{\mathrm{dR}}(f):H^q_{\mathrm{dR}}(N)\to H^q_{\mathrm{dR}}(M)$, $[\omega]\mapsto[f^*\omega]$. Równości dla cofnięć dają te same równości na klasach kohomologii, więc zarówno formy ustalonego stopnia, jak i kohomologia de Rhama wyznaczają funktory kontrawariantne do kategorii rzeczywistych przestrzeni wektorowych.
\end{przyklad}


\section{Transformacje naturalne}\label{sec:naturalne}

\providecommand{\atlasref}[2]{#2}

Dwa funktory mogą przyporządkowywać obiektom różne struktury, a mimo to być ze sobą systematycznie powiązane. Pojedyncza rodzina odwzorowań nie wystarczy: trzeba jeszcze wymagać zgodności z każdym morfizmem kategorii wyjściowej.

\begin{definicja}[Transformacja naturalna]\label{def:naturalna}
Dla funktorów $F,G:\mathcal C\to\mathcal D$ \emph{transformacja naturalna} $\eta:F\Rightarrow G$ jest rodziną morfizmów $\eta_A:F(A)\to G(A)$, po jednym dla każdego $A\in\Ob(\mathcal C)$, spełniającą dla każdego $f:A\to B$ warunek $G(f)\circ\eta_A=\eta_B\circ F(f)$. Innymi słowy, poniższy diagram jest przemienny:
\[
\begin{CD}
F(A) @>{\eta_A}>> G(A)\\
@V{F(f)}VV @VV{G(f)}V\\
F(B) @>>{\eta_B}> G(B).
\end{CD}
\]
Jeśli wszystkie $\eta_A$ są izomorfizmami, to $\eta$ nazywamy \emph{izomorfizmem naturalnym} i piszemy $F\cong G$.
\end{definicja}

\begin{stwierdzenie}[Funktory i transformacje naturalne między posetami]\label{tw:fun-posety}
Niech $(P,\le_P)$ i $(Q,\le_Q)$ będą zbiorami częściowo uporządkowanymi, traktowanymi jako kategorie. Funktory $F:P\to Q$ są dokładnie funkcjami monotonicznymi. Jeśli $F,G:P\to Q$ są monotoniczne, to istnieje transformacja naturalna $F\Rightarrow G$ wtedy i tylko wtedy, gdy $F(x)\le_Q G(x)$ dla każdego $x\in P$; jest ona wówczas jednoznaczna.
\end{stwierdzenie}
\begin{proof}
W $P$ istnieje morfizm $x\to y$ dokładnie wtedy, gdy $x\le_P y$. Funktor musi więc wysyłać każdą taką parę w parę $F(x)\le_Q F(y)$, co jest monotonicznością. Odwrotnie, funkcja monotoniczna definiuje morfizm $F(x)\to F(y)$ dla każdego $x\to y$; identyczności i złożenia są zachowane, gdyż odpowiednie zbiory morfizmów w $Q$ mają co najwyżej jeden element. Komponent transformacji naturalnej $\eta_x:F(x)\to G(x)$ istnieje dokładnie przy wskazanej nierówności. Dla $x\le_P y$ obie strony równania naturalności są wówczas morfizmami $F(x)\to G(y)$, a więc są równe.
\end{proof}

\begin{przyklad}[Zbiory uporządkowane przez inkluzję]
W kategorii pochodzącej ze zbioru częściowo uporządkowanego $(\mathcal P(\{1,2,3\}),\subseteq)$ istnieje strzałka $\{1\}\to\{1,2\}$, ale nie strzałka odwrotna. Funkcja $F(U)=U\cap\{1,2\}$ jest monotoniczna, zatem daje funktor. Ponadto $F(U)\subseteq U$ dla każdego $U$, co określa transformację naturalną $F\Rightarrow\id$. Nie jest ona izomorfizmem naturalnym, bo np. $F(\{3\})=\varnothing\neq\{3\}$.
\end{przyklad}

\begin{przyklad}[Naturalne odwzorowanie do przestrzeni podwójnie dualnej]\label{przyklad:podwojny-dual}
W $\Vect_k$ dla każdej przestrzeni $V$ definiujemy $\iota_V:V\to V^{**}$ przez $\iota_V(v)(\lambda)=\lambda(v)$ dla $v\in V$, $\lambda\in V^*$. Jest to odwzorowanie liniowe. Jeśli $f:V\to W$, to $f^{**}:V^{**}\to W^{**}$ spełnia
$f^{**}(\iota_V(v))(\mu)=\iota_V(v)(f^*(\mu))=\mu(f(v))=\iota_W(f(v))(\mu)$
dla wszystkich $v\in V$ i $\mu\in W^*$. Stąd $f^{**}\circ\iota_V=\iota_W\circ f$: rodzina $(\iota_V)$ jest transformacją naturalną $\id_{\Vect_k}\Rightarrow(-)^{**}$.

Odwzorowanie $\iota_V$ jest różnowartościowe: jeżeli $v\ne0$, rozszerzamy $\{v\}$ do bazy Hamela i definiujemy funkcjonał $\lambda$ przyjmujący wartość $1$ na $v$ i $0$ na pozostałych wektorach bazy. Wtedy $\iota_V(v)(\lambda)=1$. Używamy tutaj twierdzenia o istnieniu bazy (w ZFC). Jeżeli $\dim V=n<\infty$, z bazy $V$ otrzymujemy bazę dualną, więc $\dim V^{**}=\dim V^*=n$. Różnowartościowe odwzorowanie między przestrzeniami tej samej skończonej wymiarowości jest bijekcją. Na $\Vect_k^{\mathrm{fd}}$ otrzymujemy zatem izomorfizm naturalny.

Dla przykładu niech $V=k^{(\mathbb N)}$ będzie przestrzenią ciągów o skończonym nośniku z bazą $(e_j)_{j\in\mathbb N}$. Jej przestrzeń dualna utożsamia się z $k^{\mathbb N}$ przez wartości funkcjonału na $e_j$. W ilorazie $k^{\mathbb N}/k^{(\mathbb N)}$ klasa stałego ciągu $(1,1,\ldots)$ jest niezerowa, więc istnieje funkcjonał przyjmujący na niej wartość $1$ (ponownie przez rozszerzenie do bazy). Złożenie tego funkcjonału z rzutem ilorazowym daje $\Lambda\in V^{**}$, który zeruje wszystkie ciągi o skończonym nośniku, ale nie jest zerowy. Gdyby $\Lambda=\iota_V(v)$, podstawienie ciągu mającego jedynkę na miejscu $j$ i zera na pozostałych miejscach dawałoby $v_j=0$ dla każdego $j$. Wtedy $v=0$, sprzecznie z $\Lambda\ne0$. Zatem $\iota_V$ nie jest surjektywna w tym przykładzie.
\end{przyklad}

\begin{przyklad}[Naturalne przyporządkowanie singletonu]
Dla kowariantnego funktora potęgowego $\mathcal P_+$ z sekcji~\ref{sec:funktory} definiujemy $\eta_X:X\to\mathcal P(X)$ przez $\eta_X(x)=\{x\}$. Dla funkcji $f:X\to Y$ mamy $\mathcal P_+(f)(\eta_X(x))=\{f(x)\}=\eta_Y(f(x))$. Zatem $\eta:\id_{\Set}\Rightarrow\mathcal P_+$ jest naturalna.
\end{przyklad}

\begin{przyklad}[Naturalność i działania grup]
Niech $G$ będzie grupą, a $F,H:G\to\Set$ funktorami reprezentującymi lewe działania $G$ na zbiorach $X$ i $Y$ (przykład~\ref{przyklad:dzialanie-funktor}). Transformacja naturalna $\eta:F\Rightarrow H$ składa się z jednej funkcji $\eta_*:X\to Y$. Warunek naturalności dla morfizmu $g\in G$ ma postać $\eta_*(g\cdot x)=g\cdot\eta_*(x)$. Jest to dokładnie definicja odwzorowania $G$-ekwiwariantnego.
\end{przyklad}

\begin{twierdzenie}[Naturalne endomorfizmy identyczności w przestrzeniach liniowych]\label{tw:skalarnosc}
Niech $k$ będzie ciałem. Każda transformacja naturalna $\eta:\id_{\Vect_k}\Rightarrow\id_{\Vect_k}$ ma postać $\eta_V=\lambda\id_V$ dla pewnego $\lambda\in k$, niezależnego od $V$. Odwrotnie, każdy $\lambda\in k$ definiuje taką transformację naturalną.
\end{twierdzenie}
\begin{proof}
Ponieważ $\eta_k:k\to k$ jest liniowa, istnieje $\lambda=\eta_k(1)\in k$ takie, że $\eta_k(t)=\lambda t$. Dla dowolnego $V$ i $v\in V$ rozważmy liniową funkcję $f_v:k\to V$, $t\mapsto tv$. Naturalność daje $\eta_V(f_v(1))=f_v(\eta_k(1))$, więc $\eta_V(v)=\lambda v$. Odwrotnie, dla dowolnego liniowego $f:V\to W$ mamy $f(\lambda v)=\lambda f(v)$, co jest warunkiem naturalności.
\end{proof}

\begin{uwaga}[Naturalny nie znaczy dowolnie wybrany]
Każda skończenie wymiarowa przestrzeń wektorowa ma bazę, ale \emph{wybór bazy} zazwyczaj nie jest naturalny: różne automorfizmy tej samej przestrzeni nie muszą zachowywać wybranej bazy. Natomiast odwzorowanie $V\to V^{**}$ z przykładu~\ref{przyklad:podwojny-dual} nie zależy od wyboru bazy. Naturalność formalizuje to rozróżnienie.
\end{uwaga}

\begin{stwierdzenie}[Składanie pionowe transformacji]\label{cat:prop:vertical}
Jeśli $F,G,H:\mathcal C\to\mathcal D$ i $\eta:F\Rightarrow G$, $\theta:G\Rightarrow H$ są transformacjami naturalnymi, to rodzina $(\theta\circ\eta)_A=\theta_A\circ\eta_A$ jest transformacją naturalną $F\Rightarrow H$. Dla małych $\mathcal C$ i lokalnie małych $\mathcal D$ funktory $\mathcal C\to\mathcal D$ z transformacjami naturalnymi tworzą kategorię $[\mathcal C,\mathcal D]$.
\end{stwierdzenie}
\begin{proof}
Dla $f:A\to B$ mamy $H(f)\theta_A\eta_A=\theta_B G(f)\eta_A=\theta_B\eta_B F(f)$, gdzie pominięto znaki składania. Łączność wynika z łączności w $\mathcal D$ w każdej składowej, a identycznością funktora $F$ jest rodzina $(\id_{F(A)})_A$. Jest ona naturalna, gdyż $F(f)\circ\id_{F(A)}=\id_{F(B)}\circ F(f)$; złożenie z nią nie zmienia żadnej składowej. Jeżeli $\mathcal C$ jest mała, zbiór transformacji $F\Rightarrow G$ jest podzbiorem produktu zbiorów $\prod_{A\in\Ob(\mathcal C)}\Hom_{\mathcal D}(F(A),G(A))$. Produkt jest zbiorem, ponieważ indeksujemy go zbiorem obiektów i $\mathcal D$ jest lokalnie mała. Kategoria funktorów jest więc lokalnie mała. Nie twierdzimy, że jest mała, jeśli $\mathcal D$ jest duża.
\end{proof}


\begin{twierdzenie}[Jedyna naturalna operacja jednoargumentowa na zbiorach]\label{tw:only-id-set}
Każda transformacja naturalna $\eta:\id_{\Set}\Rightarrow\id_{\Set}$ jest identycznościowa.
\end{twierdzenie}
\begin{proof}
Dla jednoelementowego zbioru $1=\{*\}$ istnieje tylko jedna funkcja $1\to1$, zatem $\eta_1=\id_1$. Niech $X$ będzie dowolnym zbiorem, a $x\in X$. Dla funkcji $f_x:1\to X$, $f_x(*)=x$, naturalność daje $\eta_X(x)=\eta_X(f_x(*))=f_x(\eta_1(*))=x$. Ponieważ równość zachodzi dla każdego $x\in X$, otrzymujemy $\eta_X=\id_X$ (również gdy $X=\varnothing$, gdzie funkcja $X\to X$ jest jednoznaczna).
\end{proof}

\begin{przyklad}[Dlaczego wybór bazy nie jest naturalny]
Dla $V=\mathbb R^2$ wybierzmy bazę $(v_1,v_2)$. Istnieje automorfizm liniowy $T:V\to V$ spełniający $T(v_1)=v_1+v_2$ i $T(v_2)=v_2$. Wybór uporządkowanej bazy zgodny z \emph{każdym} izomorfizmem wymagałby, aby automorfizm $T$ zachowywał oba wybrane wektory, co jest fałszem. Przeciwieństwem jest naturalne odwzorowanie $V\to V^{**}$: dla niego zgodność z dowolnym $T$ została wykazana w przykładzie~\ref{przyklad:podwojny-dual}.
\end{przyklad}

\begin{stwierdzenie}[Kiedy transformacja naturalna jest odwracalna]\label{tw:nat-iso-components}
Niech $F,G:\mathcal C\to\mathcal D$ będą funktorami i $\eta:F\Rightarrow G$ transformacją naturalną. Istnieje transformacja naturalna $\theta:G\Rightarrow F$ z $\theta\circ\eta=\id_F$ i $\eta\circ\theta=\id_G$ wtedy i tylko wtedy, gdy każda składowa $\eta_A:F(A)\to G(A)$ jest izomorfizmem w $\mathcal D$.
\end{stwierdzenie}
\begin{proof}
Jeśli $\theta$ jest odwrotnością naturalną, to $\theta_A\eta_A=\id_{F(A)}$ i $\eta_A\theta_A=\id_{G(A)}$ dla każdego $A$, więc $\eta_A$ jest izomorfizmem. Odwrotnie, załóżmy odwracalność wszystkich $\eta_A$ i zdefiniujmy $\theta_A=\eta_A^{-1}$. Dla $f:A\to B$ naturalność $\eta$ to $G(f)\eta_A=\eta_BF(f)$. Składając z lewej strony z $\eta_B^{-1}$, a z prawej z $\eta_A^{-1}$, otrzymujemy $F(f)\eta_A^{-1}=\eta_B^{-1}G(f)$, czyli naturalność $\theta$. Równości $\theta\eta=\id_F$ i $\eta\theta=\id_G$ zachodzą składowa po składowej.
\end{proof}

\subsection{Składanie poziome i prawo wymiany}

\begin{definicja}[Składanie poziome]\label{cat:def:horizontal}
Niech $F,G:\mathcal C\to\mathcal D$, $H,K:\mathcal D\to\mathcal E$ oraz $\alpha:F\Rightarrow G$, $\beta:H\Rightarrow K$. Składanie pionowe zapisujemy dalej jako $\circ$, a składanie poziome jako $\ast$. Definiujemy
\[
(\beta\ast\alpha)_A=\beta_{G(A)}\circ H(\alpha_A)
=K(\alpha_A)\circ\beta_{F(A)}:HF(A)\to KG(A).
\]
Zapis $HF$ oznacza złożenie funktorów $H\circ F$. Równość dwóch wzorów wynika z naturalności $\beta$ względem morfizmu $\alpha_A$ w $\mathcal D$. Rodzinę $\beta\ast\alpha:HF\Rightarrow KG$ nazywamy złożeniem poziomym.
\end{definicja}

\begin{stwierdzenie}\label{cat:prop:horizontal}
Złożenie poziome jest transformacją naturalną. W szczególności można składać transformację naturalną z funktorem z lewej lub z prawej strony.
\end{stwierdzenie}
\begin{proof}
Dla $f:A\to B$ korzystamy kolejno z naturalności $\beta$ względem $G(f)$ i naturalności $\alpha$:
\begin{align*}
KG(f)\circ(\beta\ast\alpha)_A
&=K(G(f))\circ\beta_{G(A)}\circ H(\alpha_A)\\
&=\beta_{G(B)}\circ H(G(f))\circ H(\alpha_A)\\
&=\beta_{G(B)}\circ H(G(f)\circ\alpha_A)\\
&=\beta_{G(B)}\circ H(\alpha_B\circ F(f))\\
&=(\beta\ast\alpha)_B\circ HF(f).
\end{align*}
Jeżeli $\beta=\id_H$, otrzymujemy transformację $H\alpha:HF\Rightarrow HG$ o składowych $H(\alpha_A)$. Jeżeli $\alpha=\id_F$, otrzymujemy $\beta F:HF\Rightarrow KF$ o składowych $\beta_{F(A)}$.
\end{proof}

\begin{twierdzenie}[Prawo wymiany]\label{cat:thm:interchange}
Niech $F_0,F_1,F_2:\mathcal C\to\mathcal D$, $G_0,G_1,G_2:\mathcal D\to\mathcal E$ oraz
\[
\alpha:F_0\Rightarrow F_1,\quad \alpha':F_1\Rightarrow F_2,
\qquad
\beta:G_0\Rightarrow G_1,\quad \beta':G_1\Rightarrow G_2.
\]
Wtedy jako transformacje $G_0F_0\Rightarrow G_2F_2$ mamy
\[
(\beta'\circ\beta)\ast(\alpha'\circ\alpha)
=(\beta'\ast\alpha')\circ(\beta\ast\alpha).
\]
\end{twierdzenie}
\begin{proof}
W obiekcie $A$ prawa strona jest równa
\[
\beta'_{F_2(A)}\circ G_1(\alpha'_A)\circ
\beta_{F_1(A)}\circ G_0(\alpha_A).
\]
Naturalność $\beta$ względem $\alpha'_A:F_1(A)\to F_2(A)$ daje
$G_1(\alpha'_A)\circ\beta_{F_1(A)}
=\beta_{F_2(A)}\circ G_0(\alpha'_A)$.
Po podstawieniu i wykorzystaniu funktorialności $G_0$ otrzymujemy
\[
(\beta'_{F_2(A)}\circ\beta_{F_2(A)})\circ
G_0(\alpha'_A\circ\alpha_A),
\]
czyli składową lewej strony. Równość wszystkich składowych dowodzi tezy.
\end{proof}

\begin{stwierdzenie}[Łączność i jednostki składania poziomego]\label{cat:prop:horizontal-assoc}
Dla dodatkowych funktorów $L,M:\mathcal E\to\mathcal K$ i transformacji $\gamma:L\Rightarrow M$ zachodzi
$\gamma\ast(\beta\ast\alpha)=(\gamma\ast\beta)\ast\alpha$.
Składanie poziome z transformacją identycznościową funktora identycznościowego odpowiedniej kategorii nie zmienia transformacji.
\end{stwierdzenie}
\begin{proof}
W obiekcie $A$ lewa strona ma składową
\[
\gamma_{KG(A)}\circ L(\beta_{G(A)}\circ H(\alpha_A))
=\gamma_{KG(A)}\circ L(\beta_{G(A)})\circ LH(\alpha_A).
\]
Prawa strona ma tę samą składową z definicji składania poziomego. Dla jednostki z lewej strony wzór daje $\id_{G(A)}\circ\alpha_A=\alpha_A$. Dla jednostki z prawej strony zastosowany do $\alpha\ast\id_{\id_{\mathcal C}}$ daje $\alpha_A\circ F(\id_A)=\alpha_A$.
\end{proof}

\subsection{Naturalność konstrukcji algebraicznych i geometrycznych}

\begin{przyklad}[Naturalny izomorfizm tensorowy]\label{cat:ex:tensor-unit-natural}
Dla pierścienia przemiennego $R$ z jedynką odwzorowanie $\ell_N:R\otimes_R N\to N$, $r\otimes n\mapsto rn$, jest izomorfizmem o odwrotności $n\mapsto1\otimes n$. Istnienie i wzajemną odwrotność wykazano w rozdziale 18 skryptu \emph{Algebra abstrakcyjna}. Dla homomorfizmu $f:N\to P$ mamy
\[
f(\ell_N(r\otimes n))=f(rn)=rf(n)=\ell_P((\id_R\otimes f)(r\otimes n)).
\]
Równość na tensorach prostych wystarcza, bo generują one moduł. Otrzymujemy izomorfizm naturalny $R\otimes_R(-)\cong\id$ w kategorii $R$-modułów.
\end{przyklad}

\begin{przyklad}[Naturalność odwzorowania łączącego]\label{cat:ex:connecting-natural}
Rozważmy kategorię krótkich ciągów dokładnych kompleksów łańcuchowych $0\to A_\bullet\xrightarrow{i}B_\bullet\xrightarrow{p}C_\bullet\to0$ nad ustalonym pierścieniem. Morfizmy są trójkami odwzorowań łańcuchowych $(f_A,f_B,f_C)$ zgodnymi z $i,p$, czyli $f_Bi=i'f_A$ oraz $p'f_B=f_Cp$. Składamy je współrzędnie. Funktory przypisujące ciągowi $H_n(C_\bullet)$ oraz $H_{n-1}(A_\bullet)$ działają na morfizmach przez $H_n(f_C)$ oraz $H_{n-1}(f_A)$.

Odwzorowania łączące $\delta_n:H_n(C_\bullet)\to H_{n-1}(A_\bullet)$, skonstruowane w \atlasref{TOP-03::uwaga:les-naturalnosc}{TOP-03: naturalność odwzorowania łączącego}, tworzą transformację naturalną tych funktorów. Rzeczywiście, dla cyklu $c\in C_n$ wybieramy $b\in B_n$ z $p(b)=c$. Ponieważ $p(d_Bb)=d_Cc=0$, istnieje jedyne $a\in A_{n-1}$ z $i(a)=d_Bb$. Jest cyklem, gdyż $i(d_Aa)=d_B^2b=0$ i $i$ jest różnowartościowe. Definicja odwzorowania łączącego daje $\delta_n[c]=[a]$.

Dla drugiego ciągu $f_B(b)$ jest podniesieniem $f_C(c)$, gdyż $p'f_B(b)=f_Cp(b)$. Ponadto
\[
d_{B'}f_B(b)=f_B(d_Bb)=f_Bi(a)=i'f_A(a).
\]
A zatem $\delta'_n(H_n(f_C)[c])=[f_A(a)]=H_{n-1}(f_A)(\delta_n[c])$. Jest to warunek naturalności. Poprawność samej konstrukcji $\delta_n$ na klasach ilorazowych pochodzi z przywołanego rozdziału.
\end{przyklad}

\begin{przyklad}[Izomorfizm naturalny de Rhama]\label{cat:ex:derham-natural}
Twierdzenie de Rhama (udowodnione w rozdziale 10 skryptu \emph{Geometria i topologia różniczkowa}) daje izomorfizmy
\[
I_M:H^q_{\mathrm{dR}}(M)\longrightarrow H^q_{\mathrm{sing}}(M;\mathbb R).
\]
Sprawdźmy sens ich naturalności. Na kołańcuchach określonych na gładkich sympleksach odwzorowanie całkowania ma postać $I_M(\omega)(\sigma)=\int_{\Delta^q}\sigma^*\omega$. Jeżeli $f:M\to N$ jest gładkie, to dla formy $\omega$ na $N$ i gładkiego sympleksu $\sigma:\Delta^q\to M$ mamy
\[
I_M(f^*\omega)(\sigma)=\int_{\Delta^q}(f\circ\sigma)^*\omega
=I_N(\omega)(f\circ\sigma).
\]
Prawa strona jest wartością cofnięcia kołańcucha $I_N(\omega)$ na $\sigma$. Twierdzenie Stokesa zapewnia zgodność całkowania z różniczkami kompleksów. Ograniczenie kołańcuchów singularnych do gładkich sympleksów także komutuje z cofnięciami; w przywołanym dowodzie indukuje ono izomorfizm na kohomologii. Jego odwrotność jest naturalna na mocy stwierdzenia~\ref{tw:nat-iso-components}. W konsekwencji $I_M\circ H^q_{\mathrm{dR}}(f)=H^q_{\mathrm{sing}}(f)\circ I_N$, czyli otrzymujemy izomorfizm dwóch funktorów kontrawariantnych.
\end{przyklad}

\begin{przyklad}[Zapowiedź porównania GAGA]\label{cat:ex:gaga-natural}
Niech $X$ będzie schematem właściwym skończonego typu nad $\mathbb C$. Twierdzenie Serre'a GAGA, przywołane jako wynik zewnętrzny w rozdziale 9 skryptu \emph{Geometria algebraiczna}, daje dla ustalonego $q\ge0$ naturalne izomorfizmy
\[
H^q(X,\mathcal F)\longrightarrow H^q(X^{\mathrm{an}},\mathcal F^{\mathrm{an}})
\]
dla koherentnych snopów algebraicznych $\mathcal F$ na $X$. Naturalność dotyczy morfizmów takich snopów: porównanie kohomologii komutuje z odwzorowaniami indukowanymi przez morfizm $\mathcal F\to\mathcal G$ i jego analityfikację. Jest to zapowiedź zastosowania transformacji naturalnych do snopów; nie dowodzimy tu twierdzenia GAGA ani nie usuwamy założenia właściwości.
\end{przyklad}


\section{Obiekty uniwersalne: początkowe, końcowe, zerowe}\label{sec:uniwersalne}

\begin{definicja}[Obiekt początkowy i końcowy]\label{def:poczatkowy}
Obiekt $I$ kategorii $\mathcal C$ jest \emph{początkowy}, jeśli dla każdego $A\in\Ob(\mathcal C)$ istnieje dokładnie jeden morfizm $I\to A$. Obiekt $T$ jest \emph{końcowy}, jeśli dla każdego $A\in\Ob(\mathcal C)$ istnieje dokładnie jeden morfizm $A\to T$. Obiekt będący jednocześnie początkowym i końcowym nazywamy \emph{obiektem zerowym}.
\end{definicja}

\begin{przyklad}[Zbiory, grupy, przestrzenie]
W $\Set$ zbiór pusty $\varnothing$ jest początkowy, ponieważ istnieje dokładnie jedna funkcja $\varnothing\to A$ dla każdego zbioru $A$. Dowolny singleton $\{*\}$ jest końcowy: jedyna możliwa funkcja przypisuje każdemu elementowi $A$ punkt $*$ (również dla $A=\varnothing$). W $\Grp$ grupa trywialna $\{e\}$ jest jednocześnie początkowa i końcowa, czyli zerowa: homomorfizm z niej musi wysłać $e$ na element neutralny, a do niej istnieje tylko funkcja stała, która jest homomorfizmem. W $\Vect_k$ analogicznie obiektem zerowym jest przestrzeń $\{0\}$, ponieważ odwzorowania liniowe z niej i do niej są jednoznacznie wyznaczone. W $\Top$ obiektem początkowym jest przestrzeń pusta, a końcowym dowolna przestrzeń jednoelementowa: jedyne funkcje o odpowiednich typach są ciągłe, gdyż przeciwobrazy zbiorów otwartych są puste lub równe całej dziedzinie.
\end{przyklad}

\begin{przyklad}[Minimum i maksimum]
W kategorii częściowego porządku $(P,\leq)$ obiekt początkowy to najmniejszy element $P$, jeśli istnieje; obiekt końcowy to największy element. Na przykład w $(\mathcal P(S),\subseteq)$ są to odpowiednio $\varnothing$ i $S$.
\end{przyklad}

\begin{twierdzenie}[Jednoznaczność z dokładnością do jedynego izomorfizmu]\label{tw:unique-universal}
Jeśli $I,J$ są obiektami początkowymi tej samej kategorii, to istnieje dokładnie jeden izomorfizm $I\to J$. Analogicznie dla dwóch obiektów końcowych.
\end{twierdzenie}
\begin{proof}
Z początkowości $I$ istnieje jedyny $f:I\to J$, a z początkowości $J$ jedyny $g:J\to I$. Złożenie $g\circ f:I\to I$ musi być równe $\id_I$, ponieważ istnieje dokładnie jeden morfizm $I\to I$; podobnie $f\circ g=\id_J$. Jedyność izomorfizmu wynika z jedyności dowolnego morfizmu $I\to J$. Jeżeli $T,U$ są końcowe, końcowość $U$ daje jedyny morfizm $a:T\to U$, a końcowość $T$ jedyny $b:U\to T$. Ponieważ $b\circ a$ i $\id_T$ są morfizmami do $T$, są równe; analogicznie $a\circ b=\id_U$. Jedyność izomorfizmu wynika z końcowości $U$. W szczególności dwa obiekty zerowe są połączone jedynym izomorfizmem.
\end{proof}

\begin{uwaga}
Wyrażenie \emph{jednoznaczny z dokładnością do jedynego izomorfizmu} jest mocniejsze od zwykłego \emph{jednoznaczny z dokładnością do izomorfizmu}: nie tylko obiekty muszą być izomorficzne, lecz także odpowiedni izomorfizm jest wymuszony przez daną własność uniwersalną. Ta idea wróci przy iloczynach, granicach i sprzężeniach.
\end{uwaga}

\subsection{Dalsze przykłady i ograniczenia}

\begin{przyklad}[Pierścienie i moduły]\label{cat:ex:initial-ring}
W kategorii pierścieni przemiennych z jedynką i homomorfizmów zachowujących jedynkę obiektem początkowym jest $\mathbb Z$. Jedyny homomorfizm do pierścienia $R$ ma postać $n\mapsto n1_R$: addytywność i zachowanie jedynki wymuszają ten wzór, który zachowuje także mnożenie, gdyż $(n1_R)(m1_R)=(nm)1_R$. Przyjmujemy konwencję z rozdziału 9 skryptu \emph{Algebra abstrakcyjna}, dopuszczającą pierścień zerowy, czyli $0=1$. Jest on obiektem końcowym: jedyna funkcja $R\to\{0\}$ zachowuje wszystkie działania i jedynkę, która w przeciwdziedzinie jest zerem. Nie istnieje obiekt zerowy, gdyż $\mathbb Z$ nie jest izomorficzny z pierścieniem zerowym.

Gdyby wykluczyć pierścień zerowy, nadal istniałby obiekt początkowy $\mathbb Z$, lecz nie istniałby obiekt końcowy. Istotnie, homomorfizmy z $\mathbb F_2$ i $\mathbb F_3$ do hipotetycznego obiektu końcowego $T$ wymuszałyby jednocześnie $2\cdot1_T=0$ i $3\cdot1_T=0$, a po odjęciu $1_T=0$, sprzecznie z przyjętą konwencją.

Dla ustalonego pierścienia $R$ z jedynką kategoria lewych unitarnych $R$-modułów ma obiekt zerowy $\{0\}$. Każdy homomorfizm z modułu zerowego lub do niego jest zerowy, a taki homomorfizm zawsze istnieje. Moduł ten jest również modułem wolnym na pustym zbiorze generatorów (por. rozdział 16 skryptu \emph{Algebra abstrakcyjna}).
\end{przyklad}

\begin{przyklad}[Przestrzenie z wyróżnionym punktem i kategoria relacji]\label{cat:ex:pointed-zero}
W kategorii $\Top_*$ przestrzeni z wyróżnionym punktem singleton z jego jedynym punktem wyróżnionym jest zerowy. Jedyna dopuszczalna strzałka ze singletonu do $(X,x_0)$ wybiera $x_0$ i jest ciągła; jedyna strzałka w przeciwną stronę jest stała i zachowuje wyróżniony punkt. Zbiór pusty nie jest obiektem tej kategorii.

W kategorii $\Rel$ relacji z rozdziału~\ref{sec:kategorie} zbiór pusty jest zerowy, ponieważ zarówno $\varnothing\times A$, jak i $A\times\varnothing$ są puste i każdy ma dokładnie jeden podzbiór. Widzimy, że ten sam zbiór może być tylko początkowy w $\Set$, ale zerowy w $\Rel$, zależnie od tego, jakie morfizmy dopuszczamy.
\end{przyklad}

\begin{przyklad}[Kategorie bez obiektu początkowego lub końcowego]
W kategorii odpowiadającej grupie nietrywialnej z jednym obiektem jedyny obiekt nie jest ani początkowy, ani końcowy: ma więcej niż jeden endomorfizm. W kategorii zbiorów i funkcji różnowartościowych zbiór pusty jest początkowy, ale obiekt końcowy nie istnieje. Gdyby $T$ był końcowy, iniekcje ze singletonu do $T$ odpowiadałyby jego elementom; zatem $T$ miałby dokładnie jeden element. Nie istnieje jednak iniekcja ze zbioru dwuelementowego do singletonu.
\end{przyklad}

\begin{stwierdzenie}[Własności obiektów początkowych i końcowych]\label{cat:prop:initial-terminal-arrows}
Każdy morfizm \emph{wychodzący} z obiektu końcowego jest monomorfizmem rozszczepionym, a każdy morfizm \emph{wchodzący} do obiektu początkowego jest epimorfizmem rozszczepionym.
\end{stwierdzenie}
\begin{proof}
Niech $T$ będzie końcowy i $f:T\to A$. Jedyny morfizm $r:A\to T$ spełnia $r\circ f=\id_T$, ponieważ $T$ ma dokładnie jeden endomorfizm. Jeśli $I$ jest początkowy i $g:A\to I$, jedyny morfizm $s:I\to A$ spełnia $g\circ s=\id_I$. Otrzymujemy odpowiednie rozszczepienia z rozdziału~\ref{sec:morfizmy}. Dla kontrastu morfizm $\varnothing\to\{*\}$ w $\Set$ nie jest epimorfizmem, mimo że jego przeciwdziedzina jest obiektem końcowym.
\end{proof}

\begin{definicja}[Ścisły obiekt początkowy]\label{cat:def:strict-initial}
Obiekt początkowy $I$ jest \emph{ścisły}, jeżeli każdy morfizm $A\to I$ jest izomorfizmem.
\end{definicja}

\begin{przyklad}
W $\Set$ i $\Top$ obiekt pusty jest ścisły: istnienie funkcji $A\to\varnothing$ wymusza pustą dziedzinę, a jedyne odwzorowanie między pustymi obiektami jest izomorfizmem. W $\Grp$ obiekt początkowy nie jest ścisły, ponieważ homomorfizm z dowolnej nietrywialnej grupy do grupy trywialnej nie jest izomorfizmem.
\end{przyklad}

\subsection{Morfizmy zerowe}

\begin{definicja}[Morfizm zerowy]\label{cat:def:zero-morphism}
Jeżeli kategoria ma obiekt zerowy $O$, to dla dowolnych obiektów $A,B$ definiujemy morfizm zerowy $0_{A,B}:A\to B$ jako złożenie
\[
\begin{CD}
A @>>> O @>>> B.
\end{CD}
\]
Pierwsza strzałka jest jedyna dzięki końcowości $O$, a druga dzięki jego początkowości.
\end{definicja}

\begin{stwierdzenie}[Niezależność wyboru i pochłanianie przez zero]\label{cat:prop:zero-absorption}
Morfizm zerowy nie zależy od wyboru obiektu zerowego. Dla $u:X\to A$ oraz $v:B\to Y$ zachodzą
\[
0_{A,B}\circ u=0_{X,B},\qquad
v\circ0_{A,B}=0_{A,Y}.
\]
Ponadto obiekt $A$ jest zerowy wtedy i tylko wtedy, gdy $\id_A=0_{A,A}$.
\end{stwierdzenie}
\begin{proof}
Oznaczmy jedyne strzałki $A\to O$ i $O\to B$ przez $p_A$ i $j_B$. Jeżeli $O'$ jest drugim obiektem zerowym, jedyny izomorfizm $h:O\to O'$ spełnia $p'_A=h\circ p_A$ i $j'_B\circ h=j_B$, z jedyności strzałek do $O'$ i z $O$. Zatem $j'_B\circ p'_A=j_B\circ p_A$.

Ponieważ $p_A\circ u$ jest strzałką $X\to O$, mamy $p_A\circ u=p_X$. Podobnie $v\circ j_B=j_Y$. Stąd
\[
(j_B\circ p_A)\circ u=j_B\circ p_X,
\qquad v\circ(j_B\circ p_A)=j_Y\circ p_A.
\]
Jeżeli $A$ jest zerowy, ma jedyny endomorfizm, więc $\id_A=0_{A,A}$. Odwrotnie, przy tej równości dla każdego $f:A\to B$ mamy $f=f\circ\id_A=f\circ0_{A,A}=0_{A,B}$. Dla każdego $g:B\to A$ analogicznie $g=\id_A\circ g=0_{B,A}$. Istnienie morfizmów zerowych zapewnia, że oba zbiory morfizmów są jednoelementowe, więc $A$ jest zerowy.
\end{proof}

\begin{przyklad}[Postać morfizmów zerowych]
W kategorii grup $0_{G,H}$ jest homomorfizmem wysyłającym każdy element $G$ na element neutralny $H$; dla modułów jest zwykłym odwzorowaniem zerowym. W $\Top_*$ jest odwzorowaniem stałym o wartości w punkcie wyróżnionym przeciwdziedziny, a w $\Rel$ relacją pustą. Każda z tych postaci wynika ze złożenia przez odpowiedni obiekt zerowy.
\end{przyklad}

\begin{stwierdzenie}[Funktory zachowujące obiekt zerowy]\label{cat:prop:zero-functor}
Jeżeli $F:\mathcal C\to\mathcal D$ jest funktorem między kategoriami mającymi obiekty zerowe i $F$ przeprowadza obiekt zerowy $\mathcal C$ na obiekt zerowy $\mathcal D$, to $F(0_{A,B})=0_{F(A),F(B)}$.
\end{stwierdzenie}
\begin{proof}
Z funktorialności $F(j_B\circ p_A)=F(j_B)\circ F(p_A)$. Prawa strona faktoryzuje się przez obiekt zerowy $F(O)$, a obie strzałki do i z niego są jedyne. Niezależność wyboru obiektu zerowego daje tezę.
\end{proof}

\subsection{Jednoznaczność jako schemat własności uniwersalnej}

\begin{uwaga}\label{cat:rem:universal-coherence}
Początkowość i końcowość są własnościami obiektów, ale w dalszych konstrukcjach uniwersalnych trzeba pamiętać także o dodatkowych danych (np. rzutach iloczynu z rozdziału~\ref{sec:produkty}). Jedyny izomorfizm jest wtedy jedyny wśród morfizmów zgodnych z tymi danymi, niekoniecznie wśród wszystkich morfizmów między obiektami bazowymi.

Jeżeli $I,J,K$ są początkowe, a $u_{I,J}:I\to J$ oznacza jedyny morfizm, to $u_{J,K}\circ u_{I,J}=u_{I,K}$, gdyż obie strony są morfizmami z $I$ do $K$. Podobna równość zachodzi dla obiektów końcowych. Własność uniwersalna wymusza zatem nie tylko istnienie porównań, lecz także ich zgodność przy składaniu.
\end{uwaga}


\providecommand{\atlasref}[2]{#2}
\section{Iloczyny, koprodukty i konstrukcje uniwersalne}\label{sec:produkty}

\subsection{Iloczyn}

\begin{definicja}[Iloczyn dwóch obiektów]\label{def:produkt}
Niech $A,B$ będą obiektami kategorii $\mathcal C$. \emph{Iloczynem} $A$ i $B$ nazywamy obiekt $P$ wraz z morfizmami $p_A:P\to A$ i $p_B:P\to B$ o własności: dla każdego obiektu $X$ i każdej pary morfizmów $f:X\to A$, $g:X\to B$ istnieje dokładnie jeden morfizm $\langle f,g\rangle:X\to P$ taki, że $p_A\circ\langle f,g\rangle=f$ i $p_B\circ\langle f,g\rangle=g$.
\end{definicja}

Własność uniwersalną iloczynu przedstawia jeden diagram z dwoma przemiennymi trójkątami; przerywana strzałka $\langle f,g\rangle$ jest jednoznacznie wyznaczona:
\begin{center}
\begin{tikzpicture}[>=stealth]
\node (x) at (0,3) {$X$};
\node (p) at (0,1.5) {$P$};
\node (a) at (-2.5,0) {$A$};
\node (b) at (2.5,0) {$B$};
\draw[->,dashed] (x) -- node[right] {$\langle f,g\rangle$} (p);
\draw[->] (x) -- node[left] {$f$} (a);
\draw[->] (x) -- node[right] {$g$} (b);
\draw[->] (p) -- node[below left] {$p_A$} (a);
\draw[->] (p) -- node[below right] {$p_B$} (b);
\end{tikzpicture}
\end{center}
Przyjmujemy zwykle zapis $P=A\times B$, ale trzeba pamiętać, że chodzi o obiekt z określonymi rzutami, a nie tylko samą klasę izomorfizmu obiektu.

\begin{przyklad}[Zbiory]
W $\Set$ iloczynem jest zbiór par $A\times B$. Unikalna funkcja z definicji ma postać $\langle f,g\rangle(x)=(f(x),g(x))$. W szczególności funkcję do iloczynu można podać, podając osobno obie jej współrzędne.
\end{przyklad}

\begin{przyklad}[Grupy i topologia]
W $\Grp$ iloczynem jest $A\times B$ z mnożeniem $(a,b)(a',b')=(aa',bb')$. W $\Top$ jest nim kartezjański iloczyn wyposażony w topologię produktową. W obu przypadkach uniwersalny morfizm jest odwzorowaniem $h:x\mapsto(f(x),g(x))$. Dla homomorfizmów grup $h(xx')=(f(x)f(x'),g(x)g(x'))=h(x)h(x')$; dla odwzorowań ciągłych przeciwobraz bazowego otwartego prostokąta $U\times V$ wynosi $f^{-1}(U)\cap g^{-1}(V)$, więc jest otwarty. W obu przypadkach rzuty są morfizmami, a ich wartości wyznaczają obie współrzędne $h$, co dowodzi jednoznaczności.
\end{przyklad}

\begin{przyklad}[Iloczyn w przestrzeniach liniowych]
W $\Vect_k$ iloczynem dwóch przestrzeni $V,W$ jest przestrzeń $V\oplus W$ z rzutami na współrzędne. Dla \emph{skończonej} liczby przestrzeni suma prosta i iloczyn prosty są kanonicznie izomorficzne; przy nieskończonej rodzinie należy je rozróżnić.
\end{przyklad}

\begin{przyklad}[Iloczyn w porządku]
W kategorii zbioru częściowo uporządkowanego $P$ iloczyn obiektów $a,b$ istnieje wtedy i tylko wtedy, gdy istnieje ich największe ograniczenie dolne (kres dolny) $a\wedge b$. Rzuty istnieją, bo $a\wedge b\leq a,b$. Każde $x$ spełniające $x\leq a,b$ spełnia $x\leq a\wedge b$. Zatem logiczne \emph{``i''} odpowiada kategorii iloczynu.
\end{przyklad}

\begin{twierdzenie}[Jednoznaczność iloczynu]\label{tw:produkt-unique}
Jeśli $(P,p_A,p_B)$ i $(Q,q_A,q_B)$ są dwoma iloczynami tych samych obiektów $A,B$, to istnieje dokładnie jeden izomorfizm $u:P\to Q$ spełniający $q_A\circ u=p_A$ i $q_B\circ u=p_B$.
\end{twierdzenie}
\begin{proof}
Z własności $Q$ wynika istnienie jedynego $u:P\to Q$ o wymaganych równościach. Z własności $P$ wynika analogicznie istnienie jedynego $v:Q\to P$ zgodnego z rzutami. Teraz $p_A\circ(v\circ u)=q_A\circ u=p_A$ i tak samo dla $p_B$. Ponieważ $\id_P$ również spełnia te dwa równania, jedyność w definicji iloczynu daje $v\circ u=\id_P$. Symetrycznie $u\circ v=\id_Q$.
\end{proof}

\subsection{Koprodukt}

\begin{definicja}[Koprodukt dwóch obiektów]\label{def:koprodukt}
Niech $A,B\in\Ob(\mathcal C)$. \emph{Koprodukt} $A,B$ to obiekt $Q$ z morfizmami $i_A:A\to Q$, $i_B:B\to Q$ taki, że dla każdego $X$ oraz morfizmów $f:A\to X$, $g:B\to X$ istnieje dokładnie jeden $[f,g]:Q\to X$ spełniający $[f,g]\circ i_A=f$ i $[f,g]\circ i_B=g$.
\end{definicja}

Koprodukt jest pojęciem dualnym do iloczynu: definicję otrzymujemy przez odwrócenie wszystkich strzałek, czyli przejście do kategorii przeciwnej. Stosując twierdzenie~\ref{tw:produkt-unique} w $\mathcal C^{\op}$, otrzymujemy jedyny izomorfizm między dwoma koproduktami zgodny z ich morfizmami kanonicznymi.

\begin{przyklad}[Suma rozłączna zbiorów]
W $\Set$ koproduktem jest $A\sqcup B=(A\times\{0\})\cup(B\times\{1\})$, a $i_A(a)=(a,0)$, $i_B(b)=(b,1)$. Dla $f:A\to X$ i $g:B\to X$ jedyna funkcja $[f,g]$ jest określona na obu rozłącznych częściach przez $[f,g](a,0)=f(a)$, $[f,g](b,1)=g(b)$. Nawet gdy $A\cap B\neq\varnothing$, koprodukt zachowuje informacje, z której kopii pochodzi element.
\end{przyklad}

\begin{przyklad}[Koprodukt w grupach]
W $\Grp$ koproduktem jest \emph{iloczyn wolny} $G*H$. Korzystamy z konstrukcji grup wolnych i prezentacji z \atlasref{AA-07::aa:tw:free-exists}{AA-07: istnienie grupy wolnej}. Niech $F$ będzie grupą wolną na rozłącznych kopiach zbiorów $G$ i $H$, a $N$ normalnym domknięciem relacji
\[
[g]_G[g']_G[gg']_G^{-1},\qquad
[h]_H[h']_H[hh']_H^{-1}
\]
dla wszystkich $g,g'\in G$, $h,h'\in H$. Definiujemy $G*H=F/N$. Relacje sprawiają, że odwzorowania $i_G:g\mapsto[g]_GN$ i $i_H:h\mapsto[h]_HN$ zachowują mnożenie, a w szczególności jedynkę (idempotent w grupie jest elementem neutralnym). Dla homomorfizmów $f:G\to K$, $g:H\to K$ własność grupy wolnej daje jedyny homomorfizm $F\to K$ wysyłający oznaczone generatory na ich obrazy przez $f$ i $g$. Każda relacja trafia na jedynkę, więc jądro tego homomorfizmu zawiera $N$ i otrzymujemy homomorfizm $F/N\to K$. Jest jedyny, ponieważ obrazy oznaczonych generatorów generują $F/N$. Wybierając $K=G$, $f=\id_G$ oraz trywialny $g:H\to G$, otrzymujemy lewą odwrotność $i_G$; analogicznie $i_H$ ma lewą odwrotność. Oba morfizmy są więc zanurzeniami grup.

Iloczyn prosty nie realizuje tej własności z kanonicznymi zanurzeniami: weźmy $G=H=\mathbb Z/2\mathbb Z$ i homomorfizmy do grupy permutacji $S_3$, wysyłające generatory odpowiednio na $(12)$ i $(23)$. W $G\times H$ generatory różnych czynników komutują, ale ich zadane obrazy nie komutują. Nie istnieje zatem homomorfizm rozszerzający tę parę.
\end{przyklad}

\begin{przyklad}[Suma prosta przestrzeni]
W $\Vect_k$ koproduktem $V,W$ jest także $V\oplus W$, ale tym razem z zanurzeniami $i_V(v)=(v,0)$, $i_W(w)=(0,w)$. Jedyny morfizm $[f,g]:V\oplus W\to X$ ma wzór $[f,g](v,w)=f(v)+g(w)$. Iloczyn i koprodukt mogą więc mieć ten sam obiekt nośny, lecz są określone przez \emph{inne strzałki i inne własności uniwersalne}.
\end{przyklad}

\begin{przyklad}[Koprodukt w porządku]
Dla kategorii zbioru częściowo uporządkowanego $P$ koprodukt $a,b$ jest ich kresem górnym $a\vee b$, o ile istnieje. Zestawienie przykładów pokazuje dwie kategoryczne wersje słów \emph{``i''} (iloczyn) oraz \emph{``lub''} (koprodukt) w teorii porządku.
\end{przyklad}


\begin{przyklad}[Suma prosta a iloczyn tensorowy]
Niech $V=k^2$ i $W=k^3$. Ich produkt w $\Vect_k$ ma wymiar $2+3=5$; jest nim $V\oplus W$, wyposażony w dwa rzuty. Iloczyn tensorowy $V\otimes_k W$ ma wymiar $2\cdot3=6$, więc nie może być produktem w tej kategorii. Tensor jest opisywany przez \emph{inną} własność uniwersalną: dla każdej przestrzeni $Z$ odwzorowania liniowe $V\otimes_k W\to Z$ są w naturalnej bijekcji z odwzorowaniami dwuliniowymi $V\times W\to Z$. Sama nazwa ``iloczyn'' nie przesądza zatem o konstrukcji kategorialnej.
\end{przyklad}

\begin{przyklad}[Puste produkty i puste koprodukty]
Przyjmując definicję dla rodziny indeksowanej zbiorem (podaną poniżej), iloczyn pustej rodziny obiektów w kategorii $\mathcal C$, jeśli istnieje, to obiekt końcowy: dla każdego $X$ ma istnieć dokładnie jedna strzałka $X\to P$, bo nie ma żadnych rzutów, których zgodność trzeba sprawdzać. Dualnie koprodukt pustej rodziny jest obiektem początkowym. W $\Set$ odpowiadają im odpowiednio singleton i zbiór pusty. To ważny przypadek brzegowy definicji granicy i kogranicy.
\end{przyklad}

\begin{przyklad}[Biprodukt dwóch przestrzeni liniowych]
W $\Vect_k$ dla $V\oplus W$ występują zanurzenia $i_V(v)=(v,0)$, $i_W(w)=(0,w)$ i rzuty $p_V(v,w)=v$, $p_W(v,w)=w$. Spełniają one $p_Vi_V=\id_V$, $p_Wi_W=\id_W$, $p_Vi_W=0$, $p_Wi_V=0$ oraz $i_Vp_V+i_Wp_W=\id_{V\oplus W}$. Ostatnią równość sprawdzamy na $(v,w)$: lewa strona daje $(v,0)+(0,w)=(v,w)$. Te wzory uzasadniają nazwę \emph{biprodukt}: jeden obiekt realizuje jednocześnie produkt i koprodukt, przy ściśle zgodnych zanurzeniach i rzutach.
\end{przyklad}


\subsection{Nie należy mylić istnienia z konstrukcją}

\begin{uwaga}
Nie w każdej kategorii istnieją iloczyny i koprodukty. Dla zbioru częściowo uporządkowanego, w którym dwa elementy nie mają kresu dolnego, odpowiedni iloczyn nie istnieje. Również wtedy, gdy w danej kategorii znamy wygodną konstrukcję na elementach, \emph{własność uniwersalna} jest właściwym narzędziem sprawdzania, czy ta konstrukcja rzeczywiście jest produktem lub koproduktem.
\end{uwaga}

\begin{stwierdzenie}[Zachowanie własności przez izomorfizm]
Jeżeli $(P,p_A,p_B)$ jest iloczynem $A,B$ i $u:Q\to P$ jest izomorfizmem, to $(Q,p_A\circ u,p_B\circ u)$ również jest iloczynem $A,B$. Analogiczne stwierdzenie zachodzi dla koproduktów: jeśli $(P,i_A,i_B)$ jest koproduktem, a $u:P\to Q$ izomorfizmem, to $(Q,u\circ i_A,u\circ i_B)$ jest koproduktem.
\end{stwierdzenie}
\begin{proof}
Dla danych $f:X\to A$, $g:X\to B$ wyznaczmy jedyne $h:X\to P$ zgodne z rzutami. Wtedy $u^{-1}\circ h:X\to Q$ jest zgodny z nowymi rzutami. Jeśli $t:X\to Q$ także jest zgodny, to $u\circ t=h$ na mocy jedyności w $P$, czyli $t=u^{-1}\circ h$. Dla koproduktu, przy danych $f:A\to X$, $g:B\to X$, jedyny morfizm $h:P\to X$ zgodny z $i_A,i_B$ daje morfizm $h\circ u^{-1}:Q\to X$. Każdy inny $t:Q\to X$ zgodny z $u\circ i_A,u\circ i_B$ spełnia $t\circ u=h$, więc $t=h\circ u^{-1}$.
\end{proof}

\subsection{Rodziny dowolnej mocy}

\begin{definicja}[Iloczyn i koprodukt rodziny]\label{cat:def:indexed-products}
Niech $(A_i)_{i\in I}$ będzie rodziną obiektów indeksowaną zbiorem $I$. Iloczyn to obiekt $P$ z rzutami $p_i:P\to A_i$ taki, że każda rodzina $f_i:X\to A_i$ wyznacza jedyny $f:X\to P$ z $p_i\circ f=f_i$ dla wszystkich $i$. Koprodukt to obiekt $Q$ z morfizmami $j_i:A_i\to Q$ taki, że każda rodzina $g_i:A_i\to Y$ wyznacza jedyny $g:Q\to Y$ z $g\circ j_i=g_i$ dla wszystkich $i$. Oznaczamy je odpowiednio przez $\prod_{i\in I}A_i$ i $\coprod_{i\in I}A_i$.
\end{definicja}

W kategorii lokalnie małej własności te odpowiadają bijekcjom
\[
\Hom(X,P)\longrightarrow\prod_{i\in I}\Hom(X,A_i),
\quad f\longmapsto(p_i\circ f)_i,
\]
\[
\Hom(Q,Y)\longrightarrow\prod_{i\in I}\Hom(A_i,Y),
\quad g\longmapsto(g\circ j_i)_i.
\]
Surjektywność jest warunkiem istnienia, a różnowartościowość warunkiem jednoznaczności. Bijekcje są naturalne odpowiednio w $X$ (kontrawariantnie) i $Y$ (kowariantnie): dla $u:X'\to X$ mamy $p_i\circ(f\circ u)=(p_i\circ f)\circ u$, a dla $v:Y\to Y'$ mamy $(v\circ g)\circ j_i=v\circ(g\circ j_i)$. Przy $I=\varnothing$ iloczyn zbiorów morfizmów po prawej stronie jest singletonem, co odzyskuje rozdział~\ref{sec:uniwersalne}.

\begin{przyklad}[Zbiory, grupy i przestrzenie topologiczne]\label{cat:ex:indexed-set-top}
W $\Set$ iloczyn jest zbiorem wszystkich rodzin $(a_i)_{i\in I}$ z $a_i\in A_i$, a koprodukt sumą rozłączną $\{(i,a):i\in I,\ a\in A_i\}$. Dla rodziny funkcji do czynników jedyny morfizm do iloczynu ma wartość $(f_i(x))_i$; do opisu funkcji z koproduktu wystarczy wzór $g(i,a)=g_i(a)$. Dla $I=\varnothing$ są to odpowiednio singleton pustych rodzin i zbiór pusty. Nie zakładamy, że iloczyn niepustych zbiorów jest niepusty bez odpowiedniej postaci aksjomatu wyboru.

W $\Grp$ iloczyn dowolnej rodziny grup ma działanie współrzędne; wzór na morfizm do iloczynu zachowuje mnożenie w każdej współrzędnej.

W $\Top$ na iloczynie zbiorów bierzemy topologię generowaną przez podbazę $p_i^{-1}(U)$, gdzie $U$ jest otwarty w $A_i$ (\atlasref{GT-04::thm:topog-iloczyn-uniwersalnosc}{GT-04: własność uniwersalna iloczynu}). Funkcja $f:X\to\prod_i A_i$ jest ciągła wtedy i tylko wtedy, gdy wszystkie $p_i\circ f$ są ciągłe: w jedną stronę przez składanie, w drugą przez przeciwobrazy podbazy, $f^{-1}(p_i^{-1}(U))=(p_i\circ f)^{-1}(U)$. Na sumie rozłącznej bierzemy topologię, w której $W$ jest otwarty dokładnie wtedy, gdy każde $j_i^{-1}(W)$ jest otwarte. Warunek ten definiuje topologię, gdyż przeciwobrazy zachowują dowolne sumy i skończone przecięcia. Morfizmy $j_i$ są ciągłe, a dla złożonej z $g_i$ funkcji $g$ mamy $j_i^{-1}(g^{-1}(U))=g_i^{-1}(U)$, co dowodzi ciągłości $g$ dokładnie wtedy, gdy wszystkie $g_i$ są ciągłe.
\end{przyklad}

\begin{przyklad}[Iloczyn prosty i suma prosta modułów]\label{cat:ex:module-indexed}
Dla rodziny lewych unitarnych $R$-modułów $(M_i)_{i\in I}$ korzystamy z konstrukcji z \atlasref{AA-16::aa:def:modules-direct-sum}{AA-16: suma prosta i iloczyn prosty modułów}. Iloczyn prosty zawiera wszystkie rodziny elementów, a suma prosta jest podmodułem
\[
\bigoplus_{i\in I}M_i
=\{(m_i)_{i\in I}:\{i:m_i\ne0\}\text{ jest skończony}\}.
\]
Iloczyn prosty jest iloczynem kategorialnym: $x\mapsto(f_i(x))_i$ jest liniowe współrzędnie, a rzuty wymuszają ten wzór. Suma prosta jest koproduktem: dla homomorfizmów $g_i:M_i\to Y$ definiujemy
\[
g((m_i)_i)=\sum_{i\in I}g_i(m_i).
\]
Suma jest skończona, bo $m_i=0$ poza skończonym zbiorem. Addytywność sprawdzamy na skończonej sumie nośników dwóch rodzin, a jednorodność wynika z liniowości wszystkich $g_i$. Ponieważ każdy element sumy prostej jest skończoną sumą elementów z obrazów kanonicznych zanurzeń $j_i$, wzór jest wymuszony warunkami $g\circ j_i=g_i$. Ten sam argument dla grup abelowych wykorzystuje konstrukcję z \atlasref{AA-08::aa:def:abelian-directsum}{AA-08: suma prosta grup abelowych}.

Dla skończonego $I$ oba moduły są równe przy tej realizacji. Dla $I=\mathbb N$ i $M_i=\mathbb Z$ stała rodzina jedynek należy do iloczynu, ale nie do sumy prostej. W szczególności kanoniczna inkluzja sumy prostej w iloczyn nie jest izomorfizmem. Dla $M_i=\mathbb Z/2\mathbb Z$ suma prosta jest przeliczalna (suma zbiorów rodzin o nośniku zawartym w $\{1,\ldots,n\}$), a iloczyn nieprzeliczalny (argument przekątniowy Cantora). W tym przypadku nie istnieje nawet dowolny izomorfizm tych grup.
\end{przyklad}

\subsection{Funktorialność i kanoniczne porównania}

\begin{stwierdzenie}[Iloczyn i koprodukt jako bifunktory]\label{cat:prop:product-functor}
Jeżeli dla każdej pary obiektów kategorii $\mathcal C$ ustalono iloczyn, to przyporządkowanie $(A,B)\mapsto A\times B$ rozszerza się do bifunktora $\mathcal C\times\mathcal C\to\mathcal C$. Analogicznie ustalony koprodukt jest bifunktorem.
\end{stwierdzenie}
\begin{proof}
Dla $f:A\to A'$ i $g:B\to B'$ definiujemy $f\times g$ jako jedyny morfizm $A\times B\to A'\times B'$ z
\[
p_{A'}\circ(f\times g)=f\circ p_A,\qquad
p_{B'}\circ(f\times g)=g\circ p_B.
\]
Rzuty z iloczynu rozróżniają morfizmy do niego. Dla kolejnej pary $f',g'$ złożenie $(f'\times g')\circ(f\times g)$ ma pierwszą współrzędną $f'\circ f\circ p_A$ i drugą $g'\circ g\circ p_B$, tak samo jak $(f'\circ f)\times(g'\circ g)$. Identyczność ma wymagane rzuty, więc $\id_A\times\id_B=\id_{A\times B}$.

Dla koproduktu definiujemy $f\sqcup g$ warunkami $(f\sqcup g)\circ j_A=j_{A'}\circ f$ i $(f\sqcup g)\circ j_B=j_{B'}\circ g$. Po złożeniu kolejnych dwóch takich morfizmów ich ograniczenia wzdłuż $j_A,j_B$ wynoszą $j_{A''}\circ f'\circ f$ i $j_{B''}\circ g'\circ g$. Jednoznaczność dla koproduktu daje zachowanie składania; identyczność spełnia oba warunki dla pary identyczności.
\end{proof}

\begin{stwierdzenie}[Łączność, symetria i jednostka]\label{cat:prop:product-coherence}
W kategorii z ustalonymi iloczynami binarnymi i obiektem końcowym $T$ istnieją kanoniczne izomorfizmy
\[
(A\times B)\times C\cong A\times(B\times C),\quad
A\times B\cong B\times A,\quad A\times T\cong A.
\]
Są naturalne w występujących obiektach. Analogiczne izomorfizmy dla koproduktów otrzymujemy przy obiekcie początkowym.
\end{stwierdzenie}
\begin{proof}
Oba nawiasowania iloczynu trzech obiektów mają trzy rzuty (złożone z odpowiednich rzutów binarnych). Dla dowolnych $f:X\to A$, $g:X\to B$, $h:X\to C$ jedyne morfizmy zgodne z trzema rzutami to odpowiednio $\langle\langle f,g\rangle,h\rangle$ i $\langle f,\langle g,h\rangle\rangle$. Własność iloczynu daje więc w obu kierunkach jedyne porównania zgodne z trzema rzutami. Ich złożenia zachowują wszystkie rzuty, więc są identycznościami. Zamiana kolejności czynników daje porównanie $\langle p_B,p_A\rangle$; zamiana powtórna zachowuje oba rzuty. W przypadku $A\times T$ rzut $p_A$ ma odwrotność $\langle\id_A,!_A\rangle$, gdzie $!_A:A\to T$ jest jedynym morfizmem. Pierwsze złożenie jest identycznością na $A$, a drugie zachowuje rzut na $A$ oraz jedyny morfizm do $T$, więc jest identycznością na $A\times T$.

Dla naturalności każdego porównania składamy obie strony odpowiedniego kwadratu z rzutami na pojedyncze czynniki docelowe. Z definicji morfizmów iloczynowych obie strony dają w każdej współrzędnej ten sam morfizm czynnikowy złożony z odpowiednim rzutem źródłowym. Jednoznaczność iloczynu daje przemienność kwadratu; dla jednostki wystarczy rzut na $A$. W kategorii przeciwnej ten sam argument daje tezy dla koproduktów.
\end{proof}

\begin{stwierdzenie}[Porównanie koproduktu z iloczynem w kategorii z zerami]\label{cat:prop:coprod-product}
Jeżeli $\mathcal C$ ma obiekt zerowy oraz iloczyn $A\times B$ i koprodukt $A\sqcup B$, istnieje kanoniczny morfizm $c:A\sqcup B\to A\times B$ wyznaczony przez
\[
c\circ j_A=\langle\id_A,0_{A,B}\rangle,\qquad
c\circ j_B=\langle0_{B,A},\id_B\rangle.
\]
Nie musi on być izomorfizmem.
\end{stwierdzenie}
\begin{proof}
Oba morfizmy po prawej stronie istnieją przez własność iloczynu i definicję zer z rozdziału~\ref{sec:uniwersalne}. Własność koproduktu daje jedyny $c$. W kategorii modułów dla dwóch czynników jest to identyczność przy realizacji przez pary. W kategorii grup jest to morfizm $G*H\to G\times H$ wysyłający czynniki na odpowiednie podgrupy. Dla $G=H=\mathbb Z/2\mathbb Z$ nie jest iniekcją: komutator generatorów przechodzi na jedynkę, ale sam nie jest jedynką, ponieważ homomorfizm do $S_3$ skonstruowany wcześniej wysyła go na niejednostkowy komutator $(12)(23)(12)^{-1}(23)^{-1}$. Zatem $c$ nie jest izomorfizmem.
\end{proof}

W kategoriach addytywnych wrócimy do warunku, kiedy ten morfizm jest izomorfizmem (biprodukt). Sama obecność obiektu zerowego nie wystarcza.


\section{Equalizatory, koequalizatory i kwadraty kartezjańskie}\label{sec:equalizatory}

\providecommand{\atlasref}[2]{#2}

W konstrukcji iloczynu porównujemy dwie strzałki o wspólnej dziedzinie. Czasem potrzebujemy z kolei \emph{wymusić równość} dwóch strzałek lub połączyć dwa obiekty nad wspólną bazą. Te potrzeby prowadzą do equalizatorów i kwadratów kartezjańskich.

\subsection{Equalizator i koequalizator}

\begin{definicja}[Equalizator]\label{def:equalizer}
Niech $f,g:A\to B$ będą dwoma morfizmami w kategorii $\mathcal C$. \emph{Equalizatorem} $f$ i $g$ nazywamy morfizm $e:E\to A$ taki, że $f\circ e=g\circ e$, oraz taki, że dla każdego $u:X\to A$ z $f\circ u=g\circ u$ istnieje dokładnie jeden $\bar u:X\to E$ spełniający $e\circ\bar u=u$.
\end{definicja}
Dla $fu=gu$ własność uniwersalną przedstawia diagram:
\begin{center}
\begin{tikzpicture}[>=stealth,every node/.style={font=\small}]
 \node (x) at (0,2) {$X$};
 \node (e) at (0,0) {$E$};
 \node (a) at (3.4,0) {$A$};
 \node (b) at (7,0) {$B$};
 \draw[->] (x) -- node[above right] {$u$} (a);
 \draw[->,dashed] (x) -- node[left] {$\bar u$} (e);
 \draw[{Hooks[right]}-Stealth] (e) -- node[below] {$e$} (a);
 \draw[->] (a) to[bend left=18] node[above] {$f$} (b);
 \draw[->] (a) to[bend right=18] node[below] {$g$} (b);
\end{tikzpicture}
\end{center}

\begin{przyklad}[Equalizator w zbiorach]
Dla funkcji $f,g:A\to B$ przyjmijmy $E=\{a\in A:f(a)=g(a)\}$, a za $e:E\hookrightarrow A$ weźmy włączenie. Jeśli $u:X\to A$ spełnia $fu=gu$, to $u(x)\in E$ dla każdego $x$, zatem $u$ ma dokładnie jedno rozłożenie przez $e$. Na przykład dla $f,g:\mathbb R\to\mathbb R$ danych przez $f(x)=x^2$ i $g(x)=x$ otrzymujemy $E=\{0,1\}$.
\end{przyklad}

\begin{przyklad}[Jądro jako equalizator]
W $\Vect_k$ jądro odwzorowania $f:V\to W$ jest equalizatorem pary $f,0:V\to W$, gdzie $0$ jest odwzorowaniem zerowym. Istotnie, $f|_{\Ker f}=0$, a każde $u:X\to V$ z $fu=0$ przyjmuje wartości w $\Ker f$. To pierwsza wskazówka, że jądro jest konstrukcją kategorialną, a nie wyłącznie zbiorem rozwiązań równania.
\end{przyklad}

\begin{definicja}[Koequalizator]\label{def:coeq}
Dla morfizmów $f,g:A\to B$ \emph{koequalizatorem} nazywamy morfizm $q:B\to Q$ taki, że $q\circ f=q\circ g$ oraz dla każdego $v:B\to X$ spełniającego $v\circ f=v\circ g$ istnieje dokładnie jeden $\bar v:Q\to X$ o własności $\bar v\circ q=v$.
\end{definicja}
Dla $vf=vg$ diagram dualny ma postać:
\begin{center}
\begin{tikzpicture}[>=stealth,every node/.style={font=\small}]
 \node (a) at (0,0) {$A$};
 \node (b) at (3.6,0) {$B$};
 \node (q) at (7,0) {$Q$};
 \node (x) at (7,2) {$X$};
 \draw[->] (a) to[bend left=18] node[above] {$f$} (b);
 \draw[->] (a) to[bend right=18] node[below] {$g$} (b);
 \draw[->>] (b) -- node[below] {$q$} (q);
 \draw[->] (b) -- node[above left] {$v$} (x);
 \draw[->,dashed] (q) -- node[right] {$\bar v$} (x);
\end{tikzpicture}
\end{center}

\begin{przyklad}[Iloraz zbioru]
W $\Set$ koequalizator można zbudować, biorąc najmniejszą relację równoważności $\sim$ na $B$, dla której $f(a)\sim g(a)$ dla każdego $a\in A$, oraz rzut $q:B\to B/{\sim}$. Relacja $\sim$ istnieje jako przecięcie wszystkich relacji równoważności zawierających wskazane pary. Dla funkcji $v:B\to X$ spełniającej $vf=vg$ relacja $b\equiv b'$ określona przez $v(b)=v(b')$ jest relacją równoważności zawierającą te pary. Zatem $\sim$ zawiera się w $\equiv$, a wzór $\bar v([b])=v(b)$ jest dobrze określony. Surjektywność $q$ wymusza jednoznaczność.

Jeśli $A=\{*\}$, $B=\{0,1,2\}$, $f(*)=0$ i $g(*)=1$, to $Q$ ma dwie klasy: $\{0,1\}$ i $\{2\}$. Koequalizator \emph{skleja} obrazy wskazane przez parę morfizmów.
\end{przyklad}

\begin{przyklad}[Kojądro]
W $\Vect_k$ \emph{kojądro} odwzorowania $f:V\to W$ jest koequalizatorem pary $f,0:V\to W$. Ma postać rzutu $q:W\to W/\im f$. Jeśli $v:W\to X$ jest liniowa i $vf=0$, to $\im f\subseteq\Ker v$, więc $v$ przechodzi jednoznacznie przez $q$.
\end{przyklad}

\begin{stwierdzenie}
Każdy equalizator jest monomorfizmem. Każdy koequalizator jest epimorfizmem.
\end{stwierdzenie}
\begin{proof}
Niech $e:E\to A$ będzie equalizatorem $f,g:A\to B$. Jeśli $eu=ev$ dla $u,v:X\to E$, to wspólny morfizm $w=eu=ev$ spełnia $fw=gw$, więc na mocy jedyności rozkładu przez $e$ mamy $u=v$. Jeżeli $q:B\to Q$ jest koequalizatorem, a $u,v:Q\to X$ spełniają $u\circ q=v\circ q$, to wspólny morfizm $w=u\circ q$ spełnia $w\circ f=w\circ g$. Jednoznaczność faktoryzacji przez $q$ daje $u=v$, więc $q$ jest epimorfizmem.
\end{proof}


\begin{przyklad}[Equalizator dwóch homomorfizmów]
W kategorii $\Ab$ rozważmy $f,g:\mathbb Z\to\mathbb Z/6\mathbb Z$ określone przez $f(n)=[2n]_6$ oraz $g(n)=[5n]_6$. Warunek $f(n)=g(n)$ jest równoważny $3n\equiv0\pmod 6$, a więc $n$ jest parzyste. Equalizatorem jest zanurzenie $2\mathbb Z\hookrightarrow\mathbb Z$. Jeśli homomorfizm $u:H\to\mathbb Z$ spełnia $fu=gu$, to jego obraz jest zawarty w $2\mathbb Z$; dokładnie stąd wynika istnienie i jedyność wymaganego rozkładu. To przykład, w którym definicja uniwersalna daje równanie arytmetyczne.
\end{przyklad}

\begin{przyklad}[Koequalizator w grupach]
Niech $f,g:A\to B$ będą homomorfizmami grup. Oznaczmy przez $N$ najmniejszą podgrupę normalną $B$ zawierającą elementy $f(a)g(a)^{-1}$ dla wszystkich $a\in A$. Wówczas projekcja $q:B\to B/N$ jest koequalizatorem $f,g$. Rzeczywiście $qf=qg$. Jeśli $v:B\to K$ spełnia $vf=vg$, to $f(a)g(a)^{-1}\in\Ker v$, a normalność jądra daje $N\subseteq\Ker v$; powstaje jedyny homomorfizm $\bar v:B/N\to K$ z $\bar vq=v$. W kategorii grup nie wystarcza więc zwykła relacja równoważności zbiorów: utożsamienia muszą być zgodne z mnożeniem.
\end{przyklad}


\subsection{Pullback: para spełniająca warunek zgodności}

\begin{definicja}[Kwadrat kartezjański, pullback]\label{def:pullback}
Niech $f:A\to C$ i $g:B\to C$ będą morfizmami kategorii $\mathcal C$. \emph{Pullbackiem} (iloczynem włóknistym) tej pary jest obiekt $P$ z morfizmami $p_A:P\to A$, $p_B:P\to B$ spełniającymi $f\circ p_A=g\circ p_B$, takimi że dla każdego $X$ i pary $u:X\to A$, $v:X\to B$ z $f\circ u=g\circ v$ istnieje dokładnie jeden $h:X\to P$ spełniający $p_A\circ h=u$, $p_B\circ h=v$. Diagram
\begin{center}
\begin{tikzpicture}[>=stealth]
 \node (a) at (0,2) {$P$};
 \node (b) at (4,2) {$B$};
 \node (c) at (0,0) {$A$};
 \node (d) at (4,0) {$C$};
 \draw[->] (a) -- node[above] {$p_B$} (b);
 \draw[->] (a) -- node[left] {$p_A$} (c);
 \draw[->] (b) -- node[right] {$g$} (d);
 \draw[->] (c) -- node[below] {$f$} (d);
\end{tikzpicture}
\end{center}
nazywamy wówczas \emph{kwadratem kartezjańskim}.
\end{definicja}

\begin{przyklad}[Iloczyn włóknisty zbiorów]\label{przyklad:włokna}
W $\Set$ konstrukcja ma bardzo konkretną postać:
\[
A\times_C B:=\{(a,b)\in A\times B:f(a)=g(b)\}.
\]
Rzutami są $p_A(a,b)=a$, $p_B(a,b)=b$. Jeśli $u:X\to A$ i $v:X\to B$ spełniają $fu=gv$, to jedyna możliwa funkcja $h$ jest dana przez $h(x)=(u(x),v(x))$. Warunek zgodności gwarantuje, że wartości $h$ rzeczywiście należą do $A\times_C B$.
\end{przyklad}

\begin{przyklad}[Włókno funkcji]
Dla $f:A\to C$ i ustalonego $c\in C$ weźmy $B=\{*\}$ i $g(*)=c$. Wówczas $A\times_C\{*\}$ można utożsamić ze zbiorem $f^{-1}(\{c\})$. A zatem \emph{włókno nad punktem} jest szczególnym przypadkiem pullbacku. Ten sam pomysł działa dla wiązek w topologii i geometrii różniczkowej.
\end{przyklad}

\begin{przyklad}[Ograniczenie rodziny do podzbioru]
Niech $p:E\to B$ będzie ciągłą funkcją, a $i:U\hookrightarrow B$ włączeniem podprzestrzeni. Pullback w $\Top$ jest $E\times_B U\cong p^{-1}(U)$ z topologią podprzestrzeni, wyposażony w odwzorowanie do $U$. Istotnie, rzut $(e,u)\mapsto e$ ma odwrotność $e\mapsto(e,p(e))$ na $p^{-1}(U)$; obie funkcje są ciągłe, bo ograniczenie $p:p^{-1}(U)\to U$ jest ciągłe dla topologii podprzestrzeni. Własność uniwersalna wynika z ciągłości parowania i ograniczenia przeciwdziedziny do podprzestrzeni. W języku wiązek oznacza to \emph{ograniczenie obiektu nad $B$ do $U$}.
\end{przyklad}

\begin{przyklad}[Wspólny fragment dwóch odwzorowań]
Niech $A=B=\mathbb R$, $C=\mathbb R$ oraz $f(a)=a^2$, $g(b)=b^2$. Pullback w $\Set$ jest zbiorem $\{(a,b):a^2=b^2\}$, czyli sumą prostych $b=a$ oraz $b=-a$ w $\mathbb R^2$. Nie jest nim zwykły iloczyn $\mathbb R^2$: warunek wspólnego obrazu usuwa większość par.
\end{przyklad}


\begin{przyklad}[Pullback jako konkretne równanie liniowe]
Pracujmy w $\Vect_{\mathbb R}$. Niech $f:\mathbb R^2\to\mathbb R$ będzie dane przez $f(x,y)=x+y$, a $g:\mathbb R\to\mathbb R$ przez $g(t)=2t$. Iloczyn włóknisty to podprzestrzeń
\[
P=\{(x,y,t)\in\mathbb R^3:x+y=2t\}=\{(x,2t-x,t):x,t\in\mathbb R\}.
\]
W szczególności $\dim P=2$. Jeżeli $u:Z\to\mathbb R^2$ i $v:Z\to\mathbb R$ są liniowe i spełniają $fu=gv$, to odwzorowanie $z\mapsto(u(z),v(z))$ jest liniowe, przyjmuje wartości w $P$ i jest jedynym odwzorowaniem zgodnym z rzutami. Obliczyliśmy więc zarówno przestrzeń, jak i jej własność uniwersalną.
\end{przyklad}

\begin{przyklad}[Przeciwobrazy jako pullbacki]
Dla funkcji $f:X\to Y$ oraz włączenia $i:B\hookrightarrow Y$ produkt włóknisty $X\times_Y B$ jest naturalnie izomorficzny ze zbiorem $f^{-1}(B)$. Ustalmy $f:\mathbb Z\to\mathbb Z$, $f(n)=n^2$, i $B=\{0,1,4\}$. Wówczas $f^{-1}(B)=\{-2,-1,0,1,2\}$, a rzut na $X$ identyfikuje produkt włóknisty z tym podzbiorem. Ten sposób rozumowania działa również dla podprzestrzeni topologicznych oraz dla przeciwobrazów podrozmaitości, gdy odpowiednie struktury istnieją.
\end{przyklad}


\subsection{Pushout: sklejanie po wspólnej części}

\begin{definicja}[Pushout]\label{def:pushout}
Niech $f:C\to A$ i $g:C\to B$ będą morfizmami kategorii $\mathcal C$. \emph{Pushoutem} pary $(f,g)$ nazywamy obiekt $Q$ z morfizmami $i_A:A\to Q$, $i_B:B\to Q$ spełniającymi $i_A\circ f=i_B\circ g$, takimi że dla każdego $X$ oraz $u:A\to X$, $v:B\to X$ z $u\circ f=v\circ g$ istnieje dokładnie jeden $h:Q\to X$ z $h\circ i_A=u$, $h\circ i_B=v$. Odpowiedni kwadrat
\begin{center}
\begin{tikzpicture}[>=stealth]
 \node (a) at (0,2) {$C$};
 \node (b) at (4,2) {$B$};
 \node (c) at (0,0) {$A$};
 \node (d) at (4,0) {$Q$};
 \draw[->] (a) -- node[above] {$g$} (b);
 \draw[->] (a) -- node[left] {$f$} (c);
 \draw[->] (b) -- node[right] {$i_B$} (d);
 \draw[->] (c) -- node[below] {$i_A$} (d);
\end{tikzpicture}
\end{center}
nazywamy \emph{kokartezjańskim}.
\end{definicja}

\begin{przyklad}[Sklejanie zbiorów]
W $\Set$ pushout $A\sqcup_C B$ ma postać $(A\sqcup B)/{\sim}$, gdzie $\sim$ jest najmniejszą relacją równoważności identyfikującą $j_A(f(c))$ z $j_B(g(c))$ dla każdego $c\in C$, a $j_A,j_B$ są zanurzeniami do sumy rozłącznej. Jeżeli $q:A\sqcup B\to Q$ jest rzutem ilorazowym, morfizmami pushoutu są $i_A=q\circ j_A$ i $i_B=q\circ j_B$. Dla zgodnej pary $u,v$ morfizm $[u,v]$ z sumy rozłącznej identyfikuje wszystkie wskazane pary, więc faktoryzuje się jednoznacznie przez $q$ według własności koequalizatora. Otrzymujemy wymagany morfizm $h:Q\to X$.
\end{przyklad}

\begin{przyklad}[Sklejanie przedziałów w topologii]
Weźmy dwie kopie przedziału $[0,1]$ i jednopunktowy zbiór $C=\{*\}$, który odwzorowujemy w punkt $1$ pierwszego przedziału i $0$ drugiego. Pushout w $\Top$ jest przestrzenią ilorazową otrzymaną przez sklejenie końców; jest homeomorficzny z $[0,2]$. Szczegółowy dowód homeomorfizmu podajemy poniżej. Wersja z identyfikacją obu końców \emph{jednego} przedziału prowadzi z kolei do okręgu $S^1$ i może być opisana przez odpowiedni koequalizator.
\end{przyklad}

\begin{uwaga}
Pullback to granica diagramu $A\to C\leftarrow B$, a pushout to kogranica diagramu $A\leftarrow C\to B$. Związek ten uzasadnimy ogólnie w sekcji~\ref{sec:granice}. Nazwy angielskie są często używane również w polskich tekstach, więc warto je rozpoznawać.
\end{uwaga}


\begin{przyklad}[Pushout dwóch homomorfizmów liczb całkowitych]
W kategorii $\Ab$ weźmy $C=A=B=\mathbb Z$ i morfizmy $f(n)=2n$, $g(n)=3n$. Pushout jest ilorazem
\[
Q=(\mathbb Z\oplus\mathbb Z)/\langle(2,-3)\rangle.
\]
Mapy kanoniczne są określone przez $i_A(a)=[(a,0)]$ oraz $i_B(b)=[(0,b)]$. Ponieważ $2i_A(1)=3i_B(1)$, diagram komutuje. Każda para homomorfizmów $u,v:\mathbb Z\to H$ do grupy abelowej $H$ z $u(2n)=v(3n)$ definiuje homomorfizm $\widehat h:\mathbb Z\oplus\mathbb Z\to H$ dany przez $(a,b)\mapsto u(a)+v(b)$; ma on w jądrze $(2,-3)$, więc jednoznacznie przechodzi na $Q$. Ponadto $Q\cong\mathbb Z$: funkcja $(a,b)\mapsto3a+2b$ jest surjektywna, ponieważ przyjmuje wartość $1$ na $(1,-1)$. Jeśli $3a+2b=0$, to $a$ jest parzyste, więc $a=2k$ i $b=-3k$. Odwrotnie każda para $(2k,-3k)$ należy do jądra. Pierwsze twierdzenie o izomorfizmie daje żądany izomorfizm. Zwykłe sklejanie zbiorów nie pozwoliłoby uzyskać tej własności addytywnej.
\end{przyklad}

\begin{przyklad}[Sklejenie dwóch odcinków i jego topologia]
Dla dwóch kopii $[0,1]$ ze sklejanymi końcami $(1,\mathrm{lewy})\sim(0,\mathrm{prawy})$ otrzymujemy przestrzeń $Q=([0,1]\sqcup[0,1])/{\sim}$ z topologią ilorazową. Funkcja określona na sumie rozłącznej wzorami $t\mapsto t$ na pierwszej kopii i $t\mapsto1+t$ na drugiej jest ciągła i stała na klasach relacji. Indukuje bijekcję $Q\to[0,2]$. Jest to homeomorfizm, gdyż $Q$ jest zwarty, a $[0,2]$ jest Hausdorffa. Ta konstrukcja realizuje pushout w $\Top$, a nie tylko w $\Set$, ponieważ topologia ilorazowa zapewnia ciągłość każdego odwzorowania określonego zgodnie na obu odcinkach.
\end{przyklad}

\subsection{Konstrukcje z iloczynów i equalizatorów}

\begin{stwierdzenie}\label{cat:prop:pullback-equalizer}
Jeżeli istnieje iloczyn $A\times B$ oraz equalizator $e:P\to A\times B$ pary $f\circ p_A,g\circ p_B:A\times B\to C$, to $P$ z rzutami $p_A\circ e,p_B\circ e$ jest iloczynem włóknistym $f,g$. Dualnie pushout $f:C\to A$, $g:C\to B$ można skonstruować jako koequalizator pary $j_A\circ f,j_B\circ g:C\to A\sqcup B$.
\end{stwierdzenie}
\begin{proof}
Równość po złożeniu z $e$ daje przemienność kwadratu. Dla zgodnej pary $u:X\to A$, $v:X\to B$ jedyny morfizm $\langle u,v\rangle:X\to A\times B$ spełnia $(f\circ p_A)\circ\langle u,v\rangle=(g\circ p_B)\circ\langle u,v\rangle$. Faktoryzacja przez $e$ daje morfizm do $P$ o wymaganych rzutach. Każdy inny taki morfizm po złożeniu z $e$ ma te same rzuty, więc daje $\langle u,v\rangle$ i jest równy pierwszemu przez jednoznaczność equalizatora. Dowód dla pushoutu otrzymujemy w kategorii przeciwnej, w której koprodukt jest iloczynem, a koequalizator equalizatorem.
\end{proof}

\begin{przyklad}[Moduły]\label{cat:ex:module-pull-push}
Dla lewych $R$-modułów i homomorfizmów $f:A\to C$, $g:B\to C$ pullback jest jądrem homomorfizmu $A\oplus B\to C$, $(a,b)\mapsto f(a)-g(b)$, z ograniczeniami rzutów. Dla $f:C\to A$, $g:C\to B$ pushout jest modułem
\[
(A\oplus B)/\operatorname{im}(f,-g)
\]
z morfizmami $a\mapsto[(a,0)]$, $b\mapsto[(0,b)]$. W obu przypadkach stwierdzenie~\ref{cat:prop:pullback-equalizer} stosujemy do equalizatora lub koequalizatora, które są odpowiednio jądrem różnicy i kojądrem różnicy. Istnienie tych modułów i własność ilorazu pochodzą z \atlasref{AA-16::aa:tw:modules-first-isomorphism}{AA-16: jądro, obraz i iloraz modułów}. Dla pushoutu relacja $[(f(c),0)]=[(0,g(c))]$ zachodzi, gdyż różnica tych par należy do podmodułu, przez który dzielimy.
\end{przyklad}

\begin{przyklad}[Topologia]\label{cat:ex:top-pull-push}
W $\Top$ equalizator otrzymujemy jako podprzestrzeń zbioru zgodności $E=\{a:f(a)=g(a)\}$; odwzorowanie o obrazie w $E$ jest ciągłe do $E$ wtedy i tylko wtedy, gdy jest ciągłe do $A$. Koequalizator jest ilorazem zbioru z topologią ilorazową, dla której $\bar v$ jest ciągłe wtedy i tylko wtedy, gdy $\bar v\circ q$ jest ciągłe (\atlasref{GT-04::thm:topog-iloraz-uniwersalnosc}{GT-04: własność uniwersalna ilorazu}). Stąd, na mocy stwierdzenia~\ref{cat:prop:pullback-equalizer}, pullback ma topologię podprzestrzeni iloczynu, a pushout topologię ilorazową sumy rozłącznej. Nie zakładamy, że przestrzenie ilorazowe są Hausdorffa.
\end{przyklad}

\subsection{Sklejanie kwadratów i stabilność morfizmów}

\begin{twierdzenie}[Lemat o sklejaniu kwadratów kartezjańskich]\label{cat:thm:pullback-pasting}
Niech poniższy diagram będzie przemienny i niech prawy kwadrat będzie kartezjański.
\begin{center}
\begin{tikzpicture}[>=stealth]
 \node (p) at (0,2) {$P$};
 \node (q) at (3.4,2) {$Q$};
 \node (r) at (6.8,2) {$R$};
 \node (x) at (0,0) {$X$};
 \node (y) at (3.4,0) {$Y$};
 \node (z) at (6.8,0) {$Z$};
 \draw[->] (p) -- node[above] {$a$} (q);
 \draw[->] (q) -- node[above] {$b$} (r);
 \draw[->] (p) -- node[left] {$p$} (x);
 \draw[->] (q) -- node[right] {$q$} (y);
 \draw[->] (r) -- node[right] {$r$} (z);
 \draw[->] (x) -- node[below] {$u$} (y);
 \draw[->] (y) -- node[below] {$v$} (z);
\end{tikzpicture}
\end{center}
Wówczas lewy kwadrat jest kartezjański wtedy i tylko wtedy, gdy zewnętrzny prostokąt jest kartezjański.
\end{twierdzenie}
\begin{proof}
Załóżmy najpierw kartezjańskość lewego kwadratu. Dla $s:W\to X$, $t:W\to R$ spełniających $v\circ u\circ s=r\circ t$, prawy kwadrat daje jedyny $k:W\to Q$ z $q\circ k=u\circ s$ i $b\circ k=t$. Lewy kwadrat daje jedyny $h:W\to P$ z $p\circ h=s$ i $a\circ h=k$. Wtedy $b\circ a\circ h=t$. Jeżeli $h'$ spełnia warunki zewnętrznego prostokąta, to $a\circ h'$ spełnia warunki na $k$, zatem $a\circ h'=k$; jedyność w lewym kwadracie daje $h'=h$.

Odwrotnie załóżmy kartezjańskość prostokąta. Dla $s:W\to X$, $k:W\to Q$ z $u\circ s=q\circ k$ para $s,b\circ k$ jest zgodna względem $v\circ u,r$. Prostokąt daje jedyny $h:W\to P$ z $p\circ h=s$ i $b\circ a\circ h=b\circ k$. Mamy $q\circ a\circ h=u\circ p\circ h=u\circ s=q\circ k$. Jednoznaczność w prawym kwadracie daje $a\circ h=k$. Każde inne rozwiązanie dla lewego kwadratu spełnia również warunki prostokąta, więc jest równe $h$.
\end{proof}

\begin{uwaga}
W kategorii przeciwnej i po odwróceniu kolejności kolumn otrzymujemy wersję dualną: jeżeli lewy kwadrat jest kokartezjański, to prawy jest kokartezjański wtedy i tylko wtedy, gdy kokartezjański jest zewnętrzny prostokąt.
\end{uwaga}

\begin{stwierdzenie}[Stabilność monomorfizmów i epimorfizmów]\label{cat:prop:base-change-mono}
W kwadracie kartezjańskim z definicji~\ref{def:pullback}, jeżeli $g:B\to C$ jest monomorfizmem, to $p_A:P\to A$ jest monomorfizmem. Dualnie w kwadracie kokartezjańskim, jeżeli $g:C\to B$ jest epimorfizmem, to $i_A:A\to Q$ jest epimorfizmem.
\end{stwierdzenie}
\begin{proof}
Niech $s,t:W\to P$ spełniają $p_A\circ s=p_A\circ t$. Z przemienności wynika
$g\circ p_B\circ s=f\circ p_A\circ s=f\circ p_A\circ t=g\circ p_B\circ t$.
Moniczność $g$ daje $p_B\circ s=p_B\circ t$. Jednoznaczność w iloczynie włóknistym daje $s=t$. W wersji dualnej dla $s,t:Q\to W$ z $s\circ i_A=t\circ i_A$ mamy $s\circ i_B\circ g=t\circ i_B\circ g$. Epiczność $g$ daje $s\circ i_B=t\circ i_B$, a własność pushoutu daje $s=t$.
\end{proof}

\begin{stwierdzenie}[Cofnięcie izomorfizmu]\label{cat:prop:base-change-iso}
Jeżeli w tym samym kwadracie kartezjańskim $g$ jest izomorfizmem, to $p_A$ jest izomorfizmem.
\end{stwierdzenie}
\begin{proof}
Para $\id_A,g^{-1}\circ f$ jest zgodna, więc daje morfizm $s:A\to P$ z $p_A\circ s=\id_A$ oraz $p_B\circ s=g^{-1}\circ f$. Morfizm $p_A$ jest monomorfizmem na mocy stwierdzenia~\ref{cat:prop:base-change-mono}. Z $p_A\circ s\circ p_A=p_A$ wynika $s\circ p_A=\id_P$, więc $s=p_A^{-1}$. Dualnie wypchnięcie izomorfizmu w kwadracie kokartezjańskim jest izomorfizmem.
\end{proof}

\subsection{Iloczyny włókniste rozmaitości}

\begin{przyklad}[Transwersalność i istnienie iloczynu włóknistego]\label{cat:ex:manifold-pullback}
Dla gładkich rozmaitości bez brzegu $M,N,S$ oraz gładkich odwzorowań $f:M\to S$, $g:N\to S$ załóżmy, że dla każdej pary $f(x)=g(y)=z$ zachodzi
\[
\mathrm df_x(T_xM)+\mathrm dg_y(T_yN)=T_zS.
\]
\atlasref{DG-15::thm:transwersalne-przeciecie}{DG-15: twierdzenie o transwersalnym iloczynie włóknistym} daje na $P=\{(x,y):f(x)=g(y)\}$ strukturę podrozmaitości zanurzonej w $M\times N$, o wymiarze $\dim M+\dim N-\dim S$ (jeśli $P$ jest niepuste). Jest to także pullback w kategorii gładkich rozmaitości: zgodna para gładkich odwzorowań daje gładkie odwzorowanie do $M\times N$ o obrazie zawartym w $P$. Jest ono gładkie jako odwzorowanie do $P$, ponieważ w lokalnych mapach dostosowanych do podrozmaitości składowe normalne są zerowe, a pozostałe składowe gładkie. Jednoznaczność wynika z rzutów.

Bez transwersalności pullback może nie istnieć. Dla $f=g:\mathbb R\to\mathbb R$, $f(t)=t^2$, przypuśćmy, że istnieje pullback $P$ z rzutami $p,q:P\to\mathbb R$. Z własności uniwersalnej dla punktu istnieje jedyny $o\in P$ o $p(o)=q(o)=0$. Zgodne pary $(t,t)$ i $(t,-t)$ dają gładkie krzywe $\gamma_+,\gamma_-:\mathbb R\to P$ przechodzące przez $o$. Dla $F=(p,q):P\to\mathbb R^2$ mamy
\[
\mathrm dF_o(\gamma'_+(0))=(1,1),\qquad
\mathrm dF_o(\gamma'_-(0))=(1,-1).
\]
Te wektory generują $\mathbb R^2$, więc $\mathrm dF_o$ jest surjektywna. Z twierdzenia o submersji obraz pewnego otoczenia $o$ zawiera otoczenie $(0,0)$ w $\mathbb R^2$. Jednak z przemienności $p^2=q^2$ obraz $F$ zawiera się w sumie dwóch prostych $y=x$ i $y=-x$, która nie zawiera żadnego otoczenia $(0,0)$. Sprzeczność. Sam zbiór zgodności nie zapewnia więc istnienia iloczynu włóknistego w kategorii rozmaitości.
\end{przyklad}


\section{Granice i kogranice: wspólna teoria konstrukcji}\label{sec:granice}

\providecommand{\atlasref}[2]{#2}

\subsection{Diagram jako funktor}

\begin{definicja}[Diagram, stożek, granica]\label{def:limit}
Niech $\mathcal J$ będzie małą kategorią (\emph{kategorią indeksującą}), $\mathcal C$ kategorią i $D:\mathcal J\to\mathcal C$ funktorem, zwanym \emph{diagramem kształtu $\mathcal J$}. \emph{Stożek} nad $D$ o wierzchołku $X\in\Ob(\mathcal C)$ jest rodziną morfizmów $p_j:X\to D(j)$, po jednym dla każdego obiektu $j$ kategorii $\mathcal J$, taką że dla każdego $\alpha:j\to j'$ w $\mathcal J$ zachodzi $D(\alpha)\circ p_j=p_{j'}$.

\emph{Granicą} diagramu $D$ nazywamy stożek $(L,(\ell_j)_j)$ taki, że dla każdego stożka $(X,(p_j)_j)$ istnieje dokładnie jeden morfizm $u:X\to L$ spełniający $\ell_j\circ u=p_j$ dla wszystkich $j$. Oznaczamy $L=\varprojlim_{j\in\mathcal J}D(j)$ lub $L=\lim D$.
\end{definicja}
Dla każdej strzałki $\alpha:j\to j'$ diagram stożków ma postać poniżej. Przerywany morfizm $u$ jest ten sam dla całego diagramu $D$:
\begin{center}
\begin{tikzpicture}[>=stealth,every node/.style={font=\small}]
 \node (x) at (0,3.6) {$X$};
 \node (l) at (0,1.9) {$L$};
 \node (j) at (-2.8,0) {$D(j)$};
 \node (k) at (2.8,0) {$D(j')$};
 \draw[->,dashed] (x) -- node[right] {$u$} (l);
 \draw[->] (x) -- node[left] {$p_j$} (j);
 \draw[->] (x) -- node[right] {$p_{j'}$} (k);
 \draw[->] (l) -- node[above left] {$\ell_j$} (j);
 \draw[->] (l) -- node[above right] {$\ell_{j'}$} (k);
 \draw[->] (j) -- node[below] {$D(\alpha)$} (k);
\end{tikzpicture}
\end{center}

Przekładając definicję: diagram to układ obiektów połączonych strzałkami; stożek to jeden obiekt z kompatybilnymi strzałkami do całego układu; granica jest \emph{najbardziej uniwersalnym} takim obiektem, ponieważ przez nią w jeden sposób przechodzą wszystkie inne stożki.

\begin{przyklad}[Iloczyn jako granica]
Niech $\mathcal J$ będzie dyskretną kategorią o dwóch obiektach, a $D(0)=A$, $D(1)=B$. Nie ma żadnych strzałek zgodności między różnymi obiektami. Stożek nad $D$ to dowolne dwa morfizmy $X\to A$ i $X\to B$, a uniwersalny stożek jest właśnie iloczynem $A\times B$ z definicji~\ref{def:produkt}. Dla rodziny indeksowanej zbiorem $I$ otrzymujemy w ten sposób iloczyn $\prod_{i\in I}A_i$.
\end{przyklad}

\begin{przyklad}[Equalizator jako granica]
Niech $\mathcal J$ ma dwa obiekty $0,1$ oraz dwa różne równoległe morfizmy $\alpha,\beta:0\to1$, oprócz identyczności. Wybór diagramu $D$ jest równoważny wyborowi pary $f,g:A\to B$. Stożek o wierzchołku $X$ ma strzałki $p_0:X\to A$ i $p_1:X\to B$ z $fp_0=p_1=gp_0$. Wszystkie dane poza $p_0$ są wymuszone, a warunek istnienia stożka to $fp_0=gp_0$. Granica jest więc equalizatorem $f$ i $g$.
\end{przyklad}

\begin{przyklad}[Granice w porządkach]
Dla kategorii posetu $(P,\leq)$ stożek o wierzchołku $x$ nad diagramem $D$ istnieje dokładnie wtedy, gdy $x\leq D(j)$ dla wszystkich $j$. Granica, jeśli istnieje, jest kresem dolnym zbioru wartości diagramu. Relacje między obiektami diagramu nie wnoszą tu dodatkowych równań, ponieważ odpowiednie hom-zbiory są co najwyżej jednoelementowe.
\end{przyklad}

\begin{twierdzenie}[Jednoznaczność granicy]\label{tw:limit-unique}
Jeśli dwa stożki $(L,(\ell_j)_j)$ i $(L',(\ell'_j)_j)$ są granicami tego samego diagramu $D$, to istnieje dokładnie jeden izomorfizm $u:L\to L'$ spełniający $\ell'_j\circ u=\ell_j$ dla wszystkich $j$.
\end{twierdzenie}
\begin{proof}
Z uniwersalności $(L',(\ell'_j)_j)$ zastosowanej do stożka o wierzchołku $L$ otrzymujemy jedyną strzałkę $u:L\to L'$ z $\ell'_ju=\ell_j$. Z uniwersalności $(L,(\ell_j)_j)$ zastosowanej do stożka o wierzchołku $L'$ otrzymujemy $v:L'\to L$ z $\ell_jv=\ell'_j$. Dla każdego $j$ mamy
\[
\ell_j(vu)=(\ell_jv)u=\ell'_ju=\ell_j=\ell_j\id_L.
\]
Zatem $vu$ i $\id_L$ faktoryzują ten sam stożek przez $L$. Jednoznaczność w definicji granicy daje $vu=\id_L$. Analogicznie
$\ell'_j(uv)=\ell_jv=\ell'_j=\ell'_j\id_{L'}$, skąd $uv=\id_{L'}$. Otrzymaliśmy izomorfizm, jedyny spośród morfizmów spełniających wymagane równania.
\end{proof}


\begin{twierdzenie}[Konstrukcja wszystkich małych granic w $\Set$]\label{tw:limits-set}
Niech $\mathcal J$ będzie małą kategorią i $D:\mathcal J\to\Set$ funktorem. Wówczas granica diagramu $D$ istnieje. Można ją zrealizować jako zbiór
\[
L=\left\{(x_j)_{j\in\Ob(\mathcal J)}\in\prod_{j\in\Ob(\mathcal J)}D(j):
D(\alpha)(x_j)=x_{j'}\ \text{dla każdego }\alpha:j\to j'\right\},
\]
z rzutami $\ell_j:L\to D(j)$, $\ell_j((x_i)_i)=x_j$.
\end{twierdzenie}
\begin{proof}
Rodzina rzutów jest stożkiem, ponieważ dla $\alpha:j\to j'$ oraz $x=(x_i)_i\in L$ z definicji mamy $D(\alpha)(\ell_j(x))=D(\alpha)(x_j)=x_{j'}=\ell_{j'}(x)$. Niech $X$ będzie dowolnym zbiorem ze stożkiem $p_j:X\to D(j)$. Zdefiniujmy $u:X\to\prod_jD(j)$ przez $u(x)=(p_j(x))_j$. Zgodność stożka daje $D(\alpha)(p_j(x))=p_{j'}(x)$, więc $u(x)\in L$. Dla każdego $j$ zachodzi $\ell_ju=p_j$. Jeśli $v:X\to L$ spełnia te same równości, to dla każdego $x\in X$ i każdego $j$ współrzędne $v(x)$ i $u(x)$ są równe. Zatem $u=v$. Otrzymaliśmy własność uniwersalną granicy, również gdy $\mathcal J$ jest pusta.
\end{proof}

\begin{przyklad}[Granica nad pustym diagramem i nad ciągiem stałym]
Jeżeli $\mathcal J$ nie ma obiektów, to iloczyn pustej rodziny zbiorów jest singletonem, a dodatkowe warunki nie występują. Granica pustego diagramu w $\Set$ jest więc jednoelementowa. Jeśli natomiast $X_n=\{0,1\}$ dla wszystkich $n$, a funkcje $p_n:X_{n+1}\to X_n$ są identycznościami, to granica odwrotna ma dokładnie dwa elementy: ciągi $(0,0,\ldots)$ oraz $(1,1,\ldots)$. Gdy wszystkie $p_n$ są stałe równe $0$, jedynym zgodnym ciągiem jest $(0,0,\ldots)$.
\end{przyklad}

\begin{przyklad}[Granice skończone w kategoriach przestrzeni liniowych]
Dla diagramu $D:\mathcal J\to\Vect_k$ o skończenie wielu obiektach i morfizmach definiujemy mapę liniową
\[
A:\prod_jD(j)\longrightarrow\prod_{\alpha:j\to j'}D(j'),
\qquad
A((x_j)_j)_\alpha=D(\alpha)x_j-x_{j'}.
\]
Jej jądro jest podprzestrzenią zgodnych krotek. Rzuty ograniczone do jądra są liniowe. Dla stożka liniowych map $p_j:X\to D(j)$ funkcja $x\mapsto(p_j(x))_j$ jest liniowa współrzędnie i przyjmuje wartości w $\Ker A$ dzięki równaniom stożka. Jednoznaczność wynika z jednoznaczności współrzędnych. Zatem $\Ker A$ wraz z rzutami jest granicą w $\Vect_k$. Konstrukcja z twierdzenia~\ref{tw:limits-set} nie jest przypadkowym trikiem zbiorowym: w algebrze liniowej sprowadza granice do układów równań liniowych.
\end{przyklad}


\subsection{Kogranica}

\begin{definicja}[Kostożek i kogranica]\label{def:colimit}
Dla diagramu $D:\mathcal J\to\mathcal C$ \emph{kostożek} o wierzchołku $X$ to rodzina $i_j:D(j)\to X$ spełniająca $i_{j'}\circ D(\alpha)=i_j$ dla każdego $\alpha:j\to j'$. \emph{Kogranicą} jest kostożek $(Q,(q_j)_j)$ taki, że każdy kostożek $(X,(i_j)_j)$ wyznacza dokładnie jeden $u:Q\to X$ spełniający $u\circ q_j=i_j$ dla wszystkich $j$. Oznaczamy $Q=\varinjlim_{j\in\mathcal J}D(j)$ lub $Q=\colim D$.
\end{definicja}
Dualnie dla każdego $\alpha:j\to j'$ otrzymujemy diagram kostożków z jednym wspólnym morfizmem $u$:
\begin{center}
\begin{tikzpicture}[>=stealth,every node/.style={font=\small}]
 \node (j) at (-2.8,3.6) {$D(j)$};
 \node (k) at (2.8,3.6) {$D(j')$};
 \node (q) at (0,1.7) {$Q$};
 \node (x) at (0,0) {$X$};
 \draw[->] (j) -- node[above] {$D(\alpha)$} (k);
 \draw[->] (j) -- node[below right] {$q_j$} (q);
 \draw[->] (k) -- node[below left] {$q_{j'}$} (q);
 \draw[->] (j) -- node[left] {$i_j$} (x);
 \draw[->] (k) -- node[right] {$i_{j'}$} (x);
 \draw[->,dashed] (q) -- node[right] {$u$} (x);
\end{tikzpicture}
\end{center}

\begin{przyklad}[Koprodukt, koequalizator i pushout]
Koprodukt to kogranica diagramu dyskretnego. Koequalizator to kogranica diagramu dwóch równoległych strzałek. Pushout to kogranica diagramu $A\leftarrow C\to B$. Na odwrót, produkt, equalizator i pullback są odpowiednimi granicami. Formalne definicje tych sześciu pojęć to po prostu różne kształty tej samej własności uniwersalnej.
\end{przyklad}

\begin{przyklad}[Granica odwrotna ciągu zbiorów]\label{przyklad:inverse-limit}
Niech $X_0,X_1,X_2,\ldots$ będą zbiorami, a $p_n:X_{n+1}\to X_n$ funkcjami dla $n\geq0$. Granica odwrotna tego diagramu w $\Set$ jest zbiorem
\[
\varprojlim_{n\geq0}X_n=\{(x_0,x_1,\ldots)\in\prod_{n\geq0}X_n:p_n(x_{n+1})=x_n\text{ dla każdego }n\geq0\}.
\]
Zatem granica wybiera \emph{zgodne ciągi} elementów. Przykładowo dla $X_n=\mathbb Z/2^{n+1}\mathbb Z$ i redukcji modulo $2^{n+1}$ otrzymujemy zbiór $2$-adycznych liczb całkowitych $\mathbb Z_2$ (z odpowiednią strukturą pierścieniową).
\end{przyklad}

\begin{przyklad}[Kogranica rosnącego ciągu zbiorów]
Jeśli $X_0\subseteq X_1\subseteq\cdots$ i morfizmami są włączenia, to kogranicą w $\Set$ jest zwykła suma zbiorów $\bigcup_{n\geq0}X_n$. Jeżeli zamiast włączeń występują dowolne funkcje $X_n\to X_{n+1}$, odpowiednią kogranicę tworzymy z sumy rozłącznej wszystkich $X_n$, identyfikując element $x\in X_n$ z jego obrazem w $X_{n+1}$ i domykając te identyfikacje do relacji równoważności.
\end{przyklad}

\begin{przyklad}[Granice w rachunku liniowym]
W $\Vect_k$ granice skończonych diagramów można skonstruować za pomocą skończonych iloczynów i jąder, a kogranice za pomocą skończonych koproduktów i kojąder. Przykładowo dla $f:V\to W$ jądro $\Ker f$ to equalizator $f$ i zera; przestrzeń ilorazowa $W/\im f$ to ich koequalizator. Jest to początek języka kategorii addytywnych i abelowych.
\end{przyklad}

\begin{uwaga}[Ograniczenia wielkości]
Zakładamy, że $\mathcal J$ jest mała, więc wszystkie obiekty i strzałki diagramu są indeksowane zbiorami. W przeciwnym razie formuły typu $\prod_{j\in\Ob(\mathcal J)}D(j)$ mogą prowadzić do problemów teoriomnogościowych. Istnienie wszystkich \emph{małych} granic (odpowiednio kogranic) w kategorii nazywamy \emph{zupełnością} (odpowiednio \emph{kozupełnością}).
\end{uwaga}


\begin{twierdzenie}[Konstrukcja wszystkich małych kogranic w $\Set$]\label{tw:colimits-set}
Niech $\mathcal J$ będzie małą kategorią i $D:\mathcal J\to\Set$ diagramem. Zbudujmy sumę rozłączną $S=\coprod_{j\in\Ob(\mathcal J)}D(j)$, której elementy zapisujemy jako $(j,x)$ z $x\in D(j)$. Niech $\sim$ będzie najmniejszą relacją równoważności zawierającą pary
\[
(j,x)\sim(j',D(\alpha)(x))\quad \text{dla każdego }\alpha:j\to j'\text{ i }x\in D(j).
\]
Wówczas $Q=S/{\sim}$ wraz z mapami $q_j(x)=[(j,x)]$ jest kogranicą diagramu $D$.
\end{twierdzenie}
\begin{proof}
Dla $\alpha:j\to j'$ z definicji relacji mamy $q_{j'}(D(\alpha)(x))=[(j',D(\alpha)(x))]=[(j,x)]=q_j(x)$, więc $q_j$ tworzą kostożek. Niech $i_j:D(j)\to X$ będzie dowolnym kostożkiem. Funkcję $h:S\to X$ definiujemy na każdej składowej przez $h(j,x)=i_j(x)$. Warunek $i_{j'}D(\alpha)=i_j$ daje $h(j',D(\alpha)x)=h(j,x)$. Relacja $s\mathrel{R}t$ wtedy i tylko wtedy, gdy $h(s)=h(t)$, jest relacją równoważności zawierającą wszystkie generatory $\sim$. Z minimalności $\sim$ wynika $\sim\,\subseteq R$, czyli $h$ jest stała na klasach $\sim$. Zatem istnieje jednoznaczna funkcja $\bar h:Q\to X$, $\bar h([(j,x)])=i_j(x)$, i spełnia ona $\bar hq_j=i_j$. Jedyność wynika z tego, że każdy element $Q$ jest klasą pewnego $(j,x)$; wartości na obrazach wszystkich $q_j$ wyznaczają więc całą funkcję. Dla pustego $\mathcal J$ otrzymujemy $Q=\varnothing$.
\end{proof}

\begin{przyklad}[Kogranica diagramu, który scala dwa punkty]
Niech $X_n=\{0,1\}$ dla $n\geq0$, a wszystkie mapy $f_n:X_n\to X_{n+1}$ będą stałe i równe $0$. W kogranicy elementy $(n,0)$ oraz $(n,1)$ są utożsamione, gdyż oba przechodzą do $(n+1,0)$. Wszystkie elementy całej sumy rozłącznej należą do jednej klasy, więc kogranica jest singletonem. Gdy zamiast tego wszystkie $f_n$ są identycznościami, pozostają dokładnie dwie klasy, odpowiadające $0$ i $1$.
\end{przyklad}

\begin{przyklad}[Liczby $2$-adyczne jako zgodne reszty]
W przykładzie~\ref{przyklad:inverse-limit} weźmy układ reszt modulo $2,4,8,16,\ldots$. Ciąg $([1]_2,[3]_4,[7]_8,[15]_{16},\ldots)$ jest zgodny, ponieważ $2^{n+2}-1$ redukuje się modulo $2^{n+1}$ do $2^{n+1}-1$. Reprezentuje on liczbę $-1$ w pierścieniu $\mathbb Z_2$. Natomiast ciąg $([0]_2,[1]_4,\ldots)$ nie jest zgodny: $[1]_4$ po redukcji modulo $2$ daje $[1]_2$, a nie $[0]_2$. Ten prosty rachunek wyjaśnia, jakie dane dopuszcza granica odwrotna.
\end{przyklad}

\subsection{Istnienie granic i kogranic w ogólnych kategoriach}

\begin{twierdzenie}[Iloczyny i equalizatory wystarczają]\label{cat:thm:limits-products}
Jeżeli kategoria $\mathcal C$ ma wszystkie małe iloczyny i equalizatory, to ma wszystkie małe granice. Wystarczy istnienie skończonych iloczynów (w tym pustego) i equalizatorów, aby istniały granice wszystkich skończonych kategorii indeksujących. Dualnie, małe koprodukty i koequalizatory zapewniają wszystkie małe kogranice; wersja skończona również zachodzi.
\end{twierdzenie}
\begin{proof}
Dla diagramu $D:\mathcal J\to\mathcal C$ tworzymy
$P=\prod_jD(j)$ oraz $R=\prod_{\alpha:j\to j'}D(j')$.
Niech $s,t:P\to R$ będą wyznaczone przez składowe
$\pi_\alpha s=D(\alpha)\pi_j$ oraz $\pi_\alpha t=\pi_{j'}$.
Jeśli $e:L\to P$ jest equalizatorem $s,t$, to $\pi_je$ tworzą stożek. Dowolny stożek $(p_j:X\to D(j))$ wyznacza jedyną mapę $p:X\to P$. Równości stożka są równoważne $sp=tp$, więc istnieje jedyna $u:X\to L$ z $eu=p$. Jedyność mapy do iloczynu zapewnia, że warunki $\pi_jeu=p_j$ są równoważne $eu=p$. To dowodzi własności granicy. Dla kogranicy bierzemy $S=\coprod_jD(j)$ i $U=\coprod_{\alpha:j\to j'}D(j)$. Oznaczmy włączenia przez $i_j:D(j)\to S$ oraz $k_\alpha:D(j)\to U$. Definiujemy $r,t:U\to S$ przez $rk_\alpha=i_j$ oraz $tk_\alpha=i_{j'}D(\alpha)$. Jeśli $q:S\to Q$ jest ich koequalizatorem, to $qi_j=qrk_\alpha=qtk_\alpha=qi_{j'}D(\alpha)$, więc $qi_j$ tworzą kostożek. Dowolny kostożek $(h_j:D(j)\to X)$ wyznacza jedyną mapę $h:S\to X$ z $hi_j=h_j$. Zgodność kostożka daje $hrk_\alpha=htk_\alpha$ dla wszystkich $\alpha$, a z własności koproduktu $U$ wynika $hr=ht$. Koequalizator daje jedyną mapę $v:Q\to X$ z $vq=h$. Warunek ten jest równoważny $vqi_j=h_j$ dla wszystkich $j$, z jednoznaczności mapy z koproduktu $S$. To dowodzi istnienia i jednoznaczności faktoryzacji również dla pustego diagramu.
\end{proof}

\begin{uwaga}
Odwrotność jest natychmiastowa: iloczyny i equalizatory są szczególnymi granicami. Granica pustego diagramu to obiekt końcowy, a kogranica pustego diagramu to obiekt początkowy, o ile istnieją.
\end{uwaga}

\subsection{Granice punktowe i iterowane}

\begin{twierdzenie}[Granice i kogranice funktorów]\label{cat:thm:pointwise-limits}
Niech $\mathcal A$ będzie małą kategorią, $\mathcal J$ małą kategorią indeksującą, a $F:\mathcal J\to[\mathcal A,\mathcal C]$ diagramem funktorów, gdzie $\mathcal C$ jest lokalnie mała. Jeśli dla każdego $a$ istnieje granica diagramu $j\mapsto F(j)(a)$, to granica $F$ istnieje w kategorii funktorów i jest obliczana punktowo. Dualnie dla kogranic.
\end{twierdzenie}
\begin{proof}
Wybierzmy granice $L(a)$ z rzutami $\ell_{j,a}$. Dla $f:a\to b$ oraz $\alpha:j\to k$ naturalność $F(\alpha)$ i zgodność stożka dają
\[
F(\alpha)_bF(j)(f)\ell_{j,a}
=F(k)(f)F(\alpha)_a\ell_{j,a}
=F(k)(f)\ell_{k,a}.
\]
Zatem strzałki $F(j)(f)\ell_{j,a}$ tworzą stożek nad diagramem w $b$. Jedyna mapa $L(f):L(a)\to L(b)$ spełnia
$\ell_{j,b}L(f)=F(j)(f)\ell_{j,a}$.
Dla identyczności mamy $\ell_{j,a}L(\id_a)=\ell_{j,a}$, więc $L(\id_a)=\id_{L(a)}$. Dla $g:b\to c$ zachodzi
\[
\ell_{j,c}L(g)L(f)=F(j)(g)F(j)(f)\ell_{j,a}
=F(j)(gf)\ell_{j,a}=\ell_{j,c}L(gf).
\]
Jednoznaczność w granicy $L(c)$ daje $L(g)L(f)=L(gf)$. Równania definiujące $L(f)$ są naturalnością rzutów $\ell_j:L\to F(j)$. Stożek $(p_j:H\to F(j))_j$ wyznacza punktowo jedyne mapy $u_a:H(a)\to L(a)$ z $\ell_{j,a}u_a=p_{j,a}$. Dla $f:a\to b$ obliczamy
\[
\ell_{j,b}L(f)u_a=F(j)(f)p_{j,a}
=p_{j,b}H(f)=\ell_{j,b}u_bH(f).
\]
Z jednoznaczności w $L(b)$ otrzymujemy $L(f)u_a=u_bH(f)$. Rodzina $u_a$ jest więc transformacją naturalną i jedynym morfizmem stożków. Dla kogranic stosujemy argument do diagramu $\mathcal J^{\op}\to[\mathcal A^{\op},\mathcal C^{\op}]$, otrzymanego przez odwrócenie strzałek i transformacji naturalnych. Odwrócenie wyniku daje dokładnie kostożek i jego własność uniwersalną w $[\mathcal A,\mathcal C]$.
\end{proof}

\begin{twierdzenie}[Iterowanie granic tego samego rodzaju]\label{cat:thm:iterated-limits}
Dla diagramu $D:\mathcal I\times\mathcal J\to\mathcal C$ w kategorii zupełnej, przy małych $\mathcal I,\mathcal J$, istnieją kanoniczne izomorfizmy
\[
\lim_i\lim_jD(i,j)\cong\lim_{(i,j)}D(i,j)
\cong\lim_j\lim_iD(i,j).
\]
W kategorii kozupełnej zachodzi dualne twierdzenie dla dwóch kogranic.
\end{twierdzenie}
\begin{proof}
Dla każdego $i$ wybierzmy granicę $L_i=\lim_jD(i,j)$ i rzuty $r_{i,j}:L_i\to D(i,j)$. Jeśli $\alpha:i\to i'$ oraz $\beta:j\to j'$, to funktorialność $D$ daje
\[
D(\id_{i'},\beta)D(\alpha,\id_j)
=D(\alpha,\id_{j'})D(\id_i,\beta).
\]
Po złożeniu z $r_{i,j}$ prawa strona jest równa $D(\alpha,\id_{j'})r_{i,j'}$. Rodzina $D(\alpha,\id_j)r_{i,j}$ jest więc stożkiem, co wyznacza jedyną strzałkę $L(\alpha):L_i\to L_{i'}$ taką, że
\[
r_{i',j}L(\alpha)=D(\alpha,\id_j)r_{i,j}.
\]
Dla identyczności równania te spełnia identyczność $L_i$, a dla złożenia dwóch strzałek spełnia je złożenie odpowiednich map $L(\alpha)$. Jednoznaczność zapewnia zatem funktorialność $L$.

Niech $M=\lim_iL_i$ z rzutami $m_i:M\to L_i$. Rodzina $r_{i,j}m_i$ jest zgodna w kierunku $j$ dzięki stożkom $r_{i,j}$, a w kierunku $i$, ponieważ
\[
D(\alpha,\id_j)r_{i,j}m_i
=r_{i',j}L(\alpha)m_i=r_{i',j}m_{i'}.
\]
Każda strzałka $(\alpha,\beta)$ jest złożeniem strzałek w tych dwóch kierunkach; rodzina jest więc stożkiem nad $D$.

Niech teraz $p_{i,j}:X\to D(i,j)$ będzie dowolnym stożkiem. Dla ustalonego $i$ istnieje jedyna $v_i:X\to L_i$ z $r_{i,j}v_i=p_{i,j}$. Zgodność w kierunku $i$ daje
\[
r_{i',j}L(\alpha)v_i
=D(\alpha,\id_j)p_{i,j}=p_{i',j}=r_{i',j}v_{i'}.
\]
Jednoznaczność w granicy $L_{i'}$ daje $L(\alpha)v_i=v_{i'}$. Istnieje zatem jedyna $u:X\to M$ z $m_iu=v_i$; spełnia ona $r_{i,j}m_iu=p_{i,j}$. Jeśli $w:X\to M$ spełnia te same równania, najpierw jednoznaczność w $L_i$ daje $m_iw=v_i$ dla każdego $i$, a następnie jednoznaczność w $M$ daje $w=u$. Jest to własność granicy $D$.

Po zamianie $\mathcal I$ i $\mathcal J$ dostajemy drugi stożek graniczny, więc twierdzenie~\ref{tw:limit-unique} daje oba izomorfizmy zgodne z rzutami. Dla kogranic stosujemy udowodniony wynik do $D^{\op}:\mathcal I^{\op}\times\mathcal J^{\op}\to\mathcal C^{\op}$: granice w kategorii przeciwnej są dokładnie kogranicami w wyjściowej, z odwróconymi morfizmami stożków.
\end{proof}

\begin{przyklad}[Punktowe granice reprezentowalnych funktorów]
Jeśli $L$ jest granicą diagramu $D$ w lokalnie małej kategorii, to dla każdego $X$ zbiór $\Hom(X,L)$ jest naturalnie izomorficzny z $\lim_j\Hom(X,D(j))$: bijekcja przypisuje $u$ rodzinę $(\ell_ju)_j$, a naturalność w $X$ wynika z prekompozycji. Dualnie $\Hom(Q,X)\cong\lim_j\Hom(D(j),X)$ dla kogranicy $Q$, gdzie po prawej indeksujemy przez $\mathcal J^{\op}$.
\end{przyklad}

\subsection{Kogranice filtrowane}

\begin{definicja}[Kategoria filtrowana]\label{cat:def:filtered}
Mała kategoria $\mathcal J$ jest filtrowana, jeśli jest niepusta, każde dwa obiekty mają strzałki do wspólnego obiektu, a każda para równoległych strzałek $f,g:i\to j$ ma strzałkę $h:j\to k$ z $hf=hg$. Przykładem jest niepusty porządek skierowany, w szczególności $\mathbb N$ z naturalnym porządkiem.
\end{definicja}

\begin{lemat}[Równość reprezentantów]\label{cat:lem:filtered-equality}
Dla diagramu zbiorów $D$ o małej filtrowanej kategorii indeksującej klasy $[(i,x)]$ i $[(j,y)]$ w kogranicy są równe wtedy i tylko wtedy, gdy istnieją strzałki $i\to k$, $j\to k$, po których obrazy $x,y$ są równe.
\end{lemat}
\begin{proof}
Najpierw pokażemy, że dla skończonej listy obiektów i strzałek między nimi można znaleźć obiekt $k$ i mapy $b_i:i\to k$ zgodne z każdą z wymienionych strzałek. Dla listy pustej wybieramy dowolny obiekt (niepustość). Dla listy niepustej stosujemy kolejno drugi aksjomat filtrowania, składając dotychczasowe mapy z mapą do nowego wspólnego obiektu. Otrzymujemy mapy $b_i:i\to k$, początkowo bez warunków zgodności.

Dla strzałki $\alpha:i\to j$ porównujemy równoległe mapy $b_i,b_j\alpha:i\to k$. Trzeci aksjomat daje $h:k\to k'$ z $hb_i=hb_j\alpha$. Zastępujemy wszystkie $b_l$ przez $hb_l$. Równania już uzyskane pozostają prawdziwe po złożeniu z $h$. Po skończenie wielu takich krokach wszystkie wymagane równania są spełnione.

Najmniejsza relacja równoważności z twierdzenia~\ref{tw:colimits-set} składa się dokładnie z par połączonych skończonym ciągiem identyfikacji generujących, ich odwróceń lub równości. Istotnie, relacja opisana przez takie ciągi jest zwrotna (ciąg pusty), symetryczna (odwrócenie ciągu) i przechodnia (konkatenacja ciągów), a każda relacja równoważności zawierająca generatory musi ją zawierać.

Jeżeli $[(i,x)]=[(j,y)]$, wybierzmy taki ciąg. Zastosujmy pierwszą część dowodu do wszystkich obiektów i strzałek użytych w tym ciągu. Każda identyfikacja $(l,z)\sim(l',D(\alpha)z)$ ma po przeniesieniu do $k$ równe obrazy, gdyż
\[
D(b_{l'})D(\alpha)z=D(b_{l'}\alpha)z=D(b_l)z.
\]
Równość zachodzi także dla identyfikacji odwróconych. Przechodniość równości daje $D(b_i)x=D(b_j)y$.

Odwrotnie, jeżeli dla $b:i\to k$ oraz $c:j\to k$ zachodzi $D(b)x=D(c)y$, to $(i,x)\sim(k,D(b)x)=(k,D(c)y)\sim(j,y)$.
\end{proof}

\begin{twierdzenie}[Kogranice filtrowane a granice skończone w $\Set$]\label{cat:thm:filtered-finite}
W kategorii zbiorów kogranice filtrowane komutują z granicami skończonych diagramów: dla małej filtrowanej $\mathcal J$, skończonej $\mathcal I$ i $D:\mathcal J\times\mathcal I\to\Set$ kanoniczna mapa
\[
\colim_j\lim_i D(j,i)\longrightarrow
\lim_i\colim_j D(j,i)
\]
jest bijekcją.
\end{twierdzenie}
\begin{proof}
Dla każdego $j$ oznaczmy przez $L_j=\lim_iD(j,i)$ zbiór zgodnych krotek. Morfizm $j\to k$ działa współrzędnie na krotkach, ponieważ $D$ jest funktorem na kategorii produktowej. Kanoniczna funkcja
\[
\Phi:\colim_j L_j\longrightarrow\lim_i\colim_jD(j,i)
\]
przypisuje klasie $[(j,(x_i)_i)]$ krotkę $([(j,x_i)])_i$. Jest dobrze określona: identyfikacja w kogranicy po lewej daje identyfikacje współrzędnych po prawej, a zgodność krotki w $L_j$ daje zgodność jej klas.

\emph{Surjektywność.} Weźmy zgodną rodzinę klas $(c_i)_i$ po prawej. Wybierzmy reprezentanty $c_i=[(j_i,x_i)]$. Obiektów $i$ jest skończenie wiele, zatem istnieje wspólny etap $j$ oraz strzałki $j_i\to j$. Po przeniesieniu reprezentantów mamy $c_i=[(j,y_i)]$. Jeżeli $\beta:i\to i'$ jest strzałką w $\mathcal I$, zgodność klas oznacza
\[
[(j,D(j,\beta)y_i)]=[(j,y_{i'})].
\]
Lemat~\ref{cat:lem:filtered-equality} daje dwie strzałki z $j$ do pewnego $k_\beta$, po których obrazy tych elementów są równe. Wyrównujemy te dwie strzałki trzecim aksjomatem filtrowania, otrzymując jedną strzałkę $a_\beta:j\to k'_\beta$, po której obrazy są równe. Ponieważ strzałek $\beta$ jest skończenie wiele, stosujemy konstrukcję z dowodu lematu do układu wszystkich $a_\beta$, otrzymując jedną strzałkę $a:j\to k$ faktoryzującą każdą z nich. Jeśli strzałek nie ma, bierzemy dowolny wspólny etap (można wziąć $j$). Definiujemy $z_i=D(a,i)y_i$. Dla każdej $\beta:i\to i'$ funktorialność daje
\[
D(k,\beta)z_i=D(a,i')D(j,\beta)y_i=D(a,i')y_{i'}=z_{i'}.
\]
Stąd $(z_i)_i\in L_k$ i $\Phi([(k,(z_i)_i)])=(c_i)_i$.

\emph{Injektywność.} Niech obrazy klas reprezentowanych przez $(x_i)_i\in L_j$ oraz $(y_i)_i\in L_{j'}$ będą równe. Dla każdej współrzędnej lemat daje etap $k_i$ oraz strzałki $b_i:j\to k_i$, $c_i:j'\to k_i$, po których obrazy $x_i,y_i$ są równe. Skończony układ tych strzałek ma zgodny kostożek do pewnego $k$ (konstrukcja z dowodu lematu). Otrzymujemy więc dwie strzałki $b:j\to k$, $c:j'\to k$ faktoryzujące odpowiednio wszystkie $b_i,c_i$. Obrazy krotek są równe we wszystkich współrzędnych, czyli są równe w $L_k$. Lemat zastosowany do diagramu $j\mapsto L_j$ daje równość klas po lewej.

Dla pustej $\mathcal I$ każdy $L_j$ jest singletonem, podobnie granica po prawej. Kogranica stałego singletonu jest singletonem, bo $\mathcal J$ jest niepusta, a każde dwa obiekty mają wspólny etap. Tym samym uwzględniono również granicę pustą.
\end{proof}

\begin{przyklad}[Dlaczego skończoność jest istotna]\label{cat:ex:filtered-infinite}
Niech $A_n=\{0,\ldots,n\}$ z włączeniami i niech indeks współrzędnych przebiega $\mathbb N$. Mapa
$\colim_n A_n^{\mathbb N}\to(\colim_n A_n)^{\mathbb N}=\mathbb N^{\mathbb N}$
jest włączeniem ciągów ograniczonych w zbiór wszystkich ciągów. Ciąg $a_k=k$ nie należy do obrazu. Kogranica filtrowana nie musi więc komutować z nieskończonym iloczynem.
\end{przyklad}

\subsection{Granice odwrotne w topologii i algebrze}

\begin{przyklad}[Solenoid]
Dla liczby całkowitej $p\ge2$ weźmy okręgi $S^1$ i mapy $z\mapsto z^p$. Ich granica odwrotna w $\Top$ jest podprzestrzenią produktu $(S^1)^{\mathbb N}$ złożoną z ciągów spełniających $z_n=z_{n+1}^p$. Rzuty są ciągłe, a mapa do produktu jest ciągła dokładnie wtedy, gdy wszystkie jej składowe są ciągłe; to dowodzi własności uniwersalnej z topologią podprzestrzeni. Tak otrzymuje się solenoid, omawiany szerzej w \atlasref{TOP-10::thm:solenoid-v4}{TOP-10: homologia i kohomologia Čecha solenoidu}.
\end{przyklad}

\begin{przyklad}[Granica odwrotna nie zachowuje epimorfizmów]\label{cat:ex:inverse-not-exact}
Dla $n\ge1$ rozważmy zgodne krótkie ciągi dokładne grup abelowych
$0\to2^n\mathbb Z\to\mathbb Z\to\mathbb Z/2^n\mathbb Z\to0$.
Mapami przejścia są włączenia po lewej, identyczności pośrodku i redukcje po prawej. Granice tych trzech układów wynoszą odpowiednio $0$, $\mathbb Z$, $\mathbb Z_2$. Mapa $\mathbb Z\to\mathbb Z_2$ nie jest surjektywna: wybór dowolnego nieskończonego ciągu cyfr binarnych daje zgodne reszty, różne ciągi dają różne elementy, więc $\mathbb Z_2$ jest nieprzeliczalny, a $\mathbb Z$ przeliczalny. Granice w $\Ab$ obliczamy jako podgrupy produktu zgodnych krotek, zatem przykład pokazuje brak prawostronnej dokładności granicy odwrotnej. Jest to adaptacja \atlasref{TOP-10::prz:cech-lim-niedokladnosc}{TOP-10: przykładu niedokładności granicy odwrotnej}.
\end{przyklad}

\begin{uwaga}[Pierwsza przeszkoda: $\lim^1$]\label{cat:rem:lim-one}
Dla przeliczalnego układu odwrotnego grup abelowych $(A_n,u_n)$ kładziemy $T((a_n)_n)=(u_n(a_{n+1}))_n$. Wtedy
\[
\lim_n A_n=\Ker(1-T),\qquad
\lim\nolimits^1_n A_n=\operatorname{coker}(1-T)
\]
na produkcie $\prod_n A_n$. Pierwsza równość to równania zgodności; druga jest standardowym opisem pierwszego pochodnego funktora granicy dla takich układów (AT-29). Jeśli wszystkie $u_n$ są surjektywne, to $1-T$ jest surjektywne: dla zadanego $(b_n)$ wybieramy $a_0=0$, a następnie $a_{n+1}$ z $u_n(a_{n+1})=a_n-b_n$. Stąd $\lim^1 A_n=0$. Ogólnie wystarcza warunek Mittaga--Lefflera, czyli stabilizacja obrazów $A_m\to A_n$ dla każdego ustalonego $n$; dowód podano w \atlasref{TOP-29::thm:xviii-ml}{TOP-29: warunek Mittaga--Lefflera dla układów przeliczalnych}. Nie należy utożsamiać zwykłej granicy z granicą homotopijną ani automatycznie przenosić tej formuły na dowolne kategorie indeksujące.
\end{uwaga}

\section{Własności uniwersalne jako reprezentowanie funktorów}\label{sec:reprezentacja}

Definicje z poprzednich sekcji mają jeszcze jedną wspólną cechę. Dla dowolnego obiektu testowego $X$ opisują wszystkie strzałki $X\to U$ albo $U\to X$ przez prostsze dane; kierunek strzałek rozstrzyga o wariancji reprezentowanego funktora. Taki sposób myślenia prowadzi do pojęcia funktora reprezentowalnego.

\begin{definicja}[Reprezentowalność]\label{def:representable}
Niech $\mathcal C$ będzie kategorią lokalnie małą i $F:\mathcal C\to\Set$ funktorem. Mówimy, że $F$ jest \emph{reprezentowalny}, jeśli istnieje obiekt $A\in\Ob(\mathcal C)$ oraz izomorfizm naturalny $\Hom_{\mathcal C}(A,-)\cong F$. Wtedy $A$ \emph{reprezentuje} $F$. Dla funktora kontrawariantnego $G:\mathcal C^{\op}\to\Set$ reprezentowalność oznacza istnienie $A$ oraz izomorfizmu naturalnego $\Hom_{\mathcal C}(-,A)\cong G$.
\end{definicja}

\begin{przyklad}[Iloczyn jako reprezentowanie]
Ustalmy obiekty $A,B$ i rozważmy kontrawariantny funktor $G(X)=\Hom_{\mathcal C}(X,A)\times\Hom_{\mathcal C}(X,B)$, z działaniem na morfizmach przez składanie z prawej strony. Jeśli iloczyn $A\times B$ istnieje, to uniwersalna własność daje bijekcje
\[
\Hom_{\mathcal C}(X,A\times B)\longrightarrow\Hom_{\mathcal C}(X,A)\times\Hom_{\mathcal C}(X,B),\qquad h\longmapsto(p_Ah,p_Bh).
\]
Są one naturalne względem $X$: dla $u:Y\to X$ obie drogi prowadzą od $h$ do $(p_Ahu,p_Bhu)$. Zatem iloczyn reprezentuje funktor $G$.
\end{przyklad}

\begin{przyklad}[Koprodukt jako reprezentowanie]
Niech $Q=A\sqcup B$ będzie koproduktem z włączeniami $i_A:A\to Q$ i $i_B:B\to Q$. Rozważmy kowariantny funktor $F(X)=\Hom_{\mathcal C}(A,X)\times\Hom_{\mathcal C}(B,X)$. Dla $g:X\to Y$ określamy $F(g)(a,b)=(ga,gb)$; prawa funktora wynikają z łączności składania. Funkcja
\[
\Phi_X:\Hom_{\mathcal C}(Q,X)\to F(X),\qquad
h\mapsto(hi_A,hi_B)
\]
jest surjektywna, gdyż dla każdej pary $(a,b)$ własność koproduktu daje morfizm $h$ z $hi_A=a$, $hi_B=b$. Ta sama własność daje jego jednoznaczność, czyli injektywność $\Phi_X$. Dla $g:X\to Y$ mamy
$\Phi_Y(gh)=(ghi_A,ghi_B)=F(g)(\Phi_X(h))$.
Jest to naturalność; zatem $Q$ reprezentuje $F$. Różnica kierunku argumentu $X$ w porównaniu z iloczynem jest istotna.
\end{przyklad}

\begin{twierdzenie}[Reprezentowalność i element uniwersalny]\label{tw:univ-element}
Niech $\mathcal C$ będzie lokalnie mała, a $F:\mathcal C\to\Set$ funktorem. Następujące warunki są równoważne:
\begin{enumerate}[label=(\roman*)]
\item istnieje obiekt $A\in\Ob(\mathcal C)$ i izomorfizm naturalny $\Phi:\Hom_{\mathcal C}(A,-)\Rightarrow F$;
\item istnieje para $(A,u)$, gdzie $A\in\Ob(\mathcal C)$ oraz $u\in F(A)$, taka że dla każdego $X\in\Ob(\mathcal C)$ i każdego $x\in F(X)$ istnieje dokładnie jeden morfizm $f:A\to X$ spełniający $F(f)(u)=x$.
\end{enumerate}
W przypadku (i) można wziąć $u=\Phi_A(\id_A)$.
\end{twierdzenie}
\begin{proof}
Załóżmy (i) i ustalmy $X$. Dla $f:A\to X$ funktor $\Hom(A,-)$ wysyła $\id_A$ na $f\circ\id_A=f$. Naturalność $\Phi$ oznacza
\[
\Phi_X\circ\Hom(A,f)=F(f)\circ\Phi_A.
\]
Obliczając obie strony na $\id_A$, otrzymujemy $\Phi_X(f)=F(f)(\Phi_A(\id_A))=F(f)(u)$. Funkcja $\Phi_X$ jest bijekcją, więc dla każdego $x\in F(X)$ istnieje dokładnie jeden $f$ taki, że $F(f)(u)=x$. Otrzymujemy (ii).

Załóżmy (ii). Dla każdego $X$ definiujemy $\Phi_X:\Hom(A,X)\to F(X)$ przez $\Phi_X(f)=F(f)(u)$. Dla dowolnego $x\in F(X)$ istnienie strzałki $f$ z $F(f)(u)=x$ daje surjektywność. Jeżeli $\Phi_X(f)=\Phi_X(f')=x$, jednoznaczność takiej strzałki daje $f=f'$, czyli injektywność. Dla strzałki $g:X\to Y$ i $f:A\to X$ zachodzi $\Phi_Y(gf)=F(gf)(u)=F(g)(F(f)(u))=F(g)(\Phi_X(f))$, czyli rodzina $\Phi_X$ jest naturalna. Ponieważ każda $\Phi_X$ jest bijekcją, równanie naturalności po złożeniu z $\Phi_Y^{-1}$ z lewej i $\Phi_X^{-1}$ z prawej daje $\Hom(A,g)\Phi_X^{-1}=\Phi_Y^{-1}F(g)$. Rodzina odwrotności jest zatem naturalna. Mamy izomorfizm naturalny wymagany w (i).
\end{proof}



\begin{przyklad}[Reprezentowanie funktora zapominania dla zbiorów punktowanych]
W kategorii $\Set_*$ zbiorów z wyróżnionym punktem i funkcji zachowujących ten punkt niech $D=\{*,a\}$, gdzie $*$ jest punktem wyróżnionym. Funktor zapominania $U:\Set_*\to\Set$ jest reprezentowany przez $D$: odwzorowanie punktowane $f:D\to(X,x_0)$ jest jednoznacznie wyznaczone przez wartość $f(a)\in X$, ponieważ musi spełniać $f(*)=x_0$. Odwrotnie, dla każdego $x\in X$ funkcja $f(*)=x_0$, $f(a)=x$ jest punktowana. Jeśli $h:(X,x_0)\to(Y,y_0)$ jest morfizmem, to $(hf)(a)=h(f(a))$, co dowodzi naturalności bijekcji $\Hom_{\Set_*}(D,(X,x_0))\cong X$. Elementem uniwersalnym jest $a\in U(D)$, nie zaś punkt wyróżniony $*$.
\end{przyklad}

\subsection{Elementy uniwersalne i jednoznaczność reprezentacji}

\begin{twierdzenie}[Kontrawariantny element uniwersalny]\label{cat:thm:contra-univ-element}
Niech $G:\mathcal C^{\op}\to\Set$ będzie funktorem, gdzie $\mathcal C$ jest lokalnie mała. Dla obiektu $A$ i elementu $u\in G(A)$ istnieje izomorfizm naturalny $\Psi:\Hom(-,A)\Rightarrow G$ z $\Psi_A(\id_A)=u$ wtedy i tylko wtedy, gdy dla każdego $X$ i $x\in G(X)$ istnieje dokładnie jeden morfizm $f:X\to A$ taki, że $G(f)(u)=x$. Odpowiadający izomorfizm naturalny ma składowe $f\mapsto G(f)(u)$.
\end{twierdzenie}
\begin{proof}
Załóżmy, że $\Psi:\Hom(-,A)\Rightarrow G$ jest izomorfizmem naturalnym z $\Psi_A(\id_A)=u$. Dla $f:X\to A$ naturalność względem $f$, rozpatrywanego jako strzałka $A\to X$ w $\mathcal C^{\op}$, daje
\[
\Psi_X(f)=\Psi_X(\id_A\circ f)=G(f)(\Psi_A(\id_A))=G(f)(u).
\]
Bijektywność $\Psi_X$ daje istnienie i jednoznaczność $f$.
Odwrotnie, załóżmy własność uniwersalną. Funkcja
$\Psi_X:\Hom(X,A)\to G(X)$ zadana przez $\Psi_X(f)=G(f)(u)$
jest surjektywna przez istnienie morfizmu $f$ dla każdego $x$, a injektywna przez jego jednoznaczność. Ponadto $\Psi_A(\id_A)=G(\id_A)(u)=u$. Dla $h:Y\to X$ kontrawariantność daje
\[
\Psi_Y(fh)=G(fh)(u)=G(h)G(f)(u)=G(h)(\Psi_X(f)).
\]
To jest równanie naturalności, a składowe odwrotne są naturalne przez odwrócenie bijekcji, jak w dowodzie twierdzenia~\ref{tw:univ-element}.
\end{proof}

\begin{twierdzenie}[Jednoznaczność pary reprezentującej]\label{cat:thm:representing-unique}
Jeżeli $(A,u)$ oraz $(B,v)$ są parami uniwersalnymi dla $F:\mathcal C\to\Set$, to istnieje dokładnie jeden izomorfizm $a:A\to B$ spełniający $F(a)(u)=v$. Dla $G:\mathcal C^{\op}\to\Set$ istnieje dokładnie jeden izomorfizm $a:A\to B$ z $G(a)(v)=u$.
\end{twierdzenie}
\begin{proof}
W przypadku kowariantnym uniwersalność $(A,u)$ daje jedyny $a:A\to B$ z $F(a)(u)=v$, a uniwersalność $(B,v)$ daje jedyny $b:B\to A$ z $F(b)(v)=u$. Wówczas
\[
F(ba)(u)=F(b)(v)=u=F(\id_A)(u).
\]
Jednoznaczność dla $(A,u)$ wymusza $ba=\id_A$. Tak samo
$F(ab)(v)=F(a)(u)=v=F(\id_B)(v)$, zatem $ab=\id_B$.
W przypadku kontrawariantnym zastosujmy uniwersalność $(B,v)$ do elementu $u\in G(A)$: istnieje dokładnie jedna strzałka $a:A\to B$ z $G(a)(v)=u$. Z uniwersalności $(A,u)$ dla $v\in G(B)$ dostajemy jedyną strzałkę $b:B\to A$ z $G(b)(u)=v$. Kontrawariantność daje
\[
G(ba)(u)=G(a)G(b)(u)=G(a)(v)=u,
\qquad
G(ab)(v)=G(b)G(a)(v)=G(b)(u)=v.
\]
Identyczności $\id_A$ i $\id_B$ spełniają te same równania co odpowiednio $ba$ i $ab$. Uniwersalność obu par daje $ba=\id_A$ i $ab=\id_B$. Zatem $a$ jest izomorfizmem, a jego jednoznaczność jako morfizmu zgodnego z elementami uniwersalnymi została wykazana już na początku.
\end{proof}

\begin{uwaga}
Jednoznaczny jest izomorfizm zachowujący wybrane elementy uniwersalne, nie dowolny izomorfizm obiektów reprezentujących. Na przykład $\mathbb Z$ reprezentuje funktor zapominania $\Ab\to\Set$ za pomocą elementu $1$ albo $-1$: dla $x$ w dowolnej grupie abelowej jedynym homomorfizmem wysyłającym $1$ na $x$ jest $n\mapsto nx$, a jedynym homomorfizmem wysyłającym $-1$ na $x$ jest $n\mapsto -nx$. Izomorfizm zgodny z tymi dwoma elementami to mnożenie przez $-1$, natomiast identyczność nie jest z nimi zgodna.
\end{uwaga}

\subsection{Kategoria elementów}

\begin{definicja}[Kategoria elementów funktora]\label{cat:def:elements}
Niech $\mathcal C$ będzie lokalnie mała. Dla $F:\mathcal C\to\Set$ obiektami kategorii elementów $\int F$ są pary $(X,x)$ z $x\in F(X)$. Morfizmem $(X,x)\to(Y,y)$ jest strzałka $f:X\to Y$ z $F(f)(x)=y$. Dla $G:\mathcal C^{\op}\to\Set$ przyjmujemy tę samą orientację strzałek w $\mathcal C$, lecz warunek ma postać $G(f)(y)=x$. Kategoria elementów może być duża.
\end{definicja}

Istotnie, identyczność spełnia warunek, ponieważ $F(\id_X)=\id_{F(X)}$. Jeżeli $F(f)(x)=y$ i $F(g)(y)=z$, to $F(gf)(x)=F(g)(y)=z$. W wersji kontrawariantnej $G(f)(y)=x$, $G(g)(z)=y$ dają $G(gf)(z)=G(f)G(g)(z)=x$. Łączność oraz prawa identyczności dziedziczymy z $\mathcal C$.

\begin{stwierdzenie}[Reprezentowalność przez obiekt początkowy lub końcowy]\label{cat:prop:elements-representable}
Dla lokalnie małej kategorii $\mathcal C$ funktor kowariantny $F:\mathcal C\to\Set$ jest reprezentowalny dokładnie wtedy, gdy $\int F$ ma obiekt początkowy. Funktor kontrawariantny $G:\mathcal C^{\op}\to\Set$ jest reprezentowalny dokładnie wtedy, gdy $\int G$ ma obiekt końcowy.
\end{stwierdzenie}
\begin{proof}
Strzałka z $(A,u)$ do $(X,x)$ w $\int F$ to dokładnie morfizm $f:A\to X$ spełniający $F(f)(u)=x$. Istnienie dokładnie jednej takiej strzałki dla każdego $(X,x)$ jest równoważne warunkowi twierdzenia~\ref{tw:univ-element}. W wersji kontrawariantnej strzałka z $(X,x)$ do $(A,u)$ ma postać $f:X\to A$ z $G(f)(u)=x$. Stosujemy twierdzenie~\ref{cat:thm:contra-univ-element}.
\end{proof}

\subsection{Przykłady algebraiczne i geometryczne}

\begin{przyklad}[Moduł wolny]\label{cat:ex:free-module-rep}
Niech $R$ będzie pierścieniem z jedynką, $S$ zbiorem, a $\mathcal C$ kategorią lewych $R$-modułów. Funktor
$F(M)=\operatorname{Map}(S,M)$ działa na homomorfizmie $h:M\to N$ przez $F(h)(u)=h\circ u$. Reprezentuje go moduł wolny $R^{(S)}$ z elementem uniwersalnym $e:S\to R^{(S)}$, $s\mapsto e_s$.

Dla funkcji $u:S\to M$ określamy $\bar u:R^{(S)}\to M$ przez
\[
\bar u\left(\sum_{s\in S}r_se_s\right)=\sum_{s\in S}r_su(s).
\]
Rodzina współczynników $(r_s)_s$ ma skończony nośnik, więc obie sumy zawierają jedynie skończenie wiele niezerowych składników. Jednoznaczność współrzędnych zapewnia dobre określenie. Jeśli $x=\sum_s r_se_s$ i $y=\sum_s t_se_s$, to
\[
\bar u(x+y)=\sum_s(r_s+t_s)u(s)
=\sum_s r_su(s)+\sum_s t_su(s)=\bar u(x)+\bar u(y).
\]
Sumujemy nad sumą skończonych nośników obu rodzin. Dla $a\in R$ mamy
$\bar u(ax)=\sum_s(ar_s)u(s)=a\sum_s r_su(s)=a\bar u(x)$.
Kolejność czynników pozostaje zachowana również dla nieprzemiennego $R$. Dla $S=\varnothing$ otrzymujemy moduł zerowy i jedyny homomorfizm z niego do $M$. Mamy $\bar u e=u$. Jeżeli $w$ jest innym homomorfizmem z $we=u$, to
$w(\sum_s r_se_s)=\sum_s r_sw(e_s)=\sum_s r_su(s)$, więc $w=\bar u$. Twierdzenie~\ref{tw:univ-element} daje reprezentowalność; naturalność można też sprawdzić przez $(h\bar u)e=hu$. Konstrukcja modułu wolnego jest opisana w tomie algebry abstrakcyjnej (AA-16).
\end{przyklad}

\begin{przyklad}[Lokalizacja pierścienia]\label{cat:ex:localization-rep}
Ustalmy przemienny pierścień $R$ z jedynką oraz zbiór multiplikatywny $S$ zawierający $1$, niezawierający $0$. W kategorii przemiennych pierścieni z jedynką i homomorfizmów zachowujących jedynkę rozważmy funktor
\[
F(T)=\{f:R\to T:\ f\text{ jest homomorfizmem},\
 f(S)\subseteq T^\times\}.
\]
Na strzałce $h:T\to T'$ kładziemy $F(h)(f)=hf$. Jest to dobrze określone, gdyż $h(t)h(t^{-1})=1$ dla jednostki $t$.

Elementem uniwersalnym jest $\iota:R\to S^{-1}R$, $a\mapsto a/1$. Dla $f\in F(T)$ formuła
$\bar f(a/s)=f(a)f(s)^{-1}$ daje jedyną faktoryzację przez $\iota$. Rzeczywiście, z równości $a/s=b/t$ wynika $u(at-bs)=0$ dla pewnego $u\in S$. Po zastosowaniu $f$ i pomnożeniu przez $f(u)^{-1}f(s)^{-1}f(t)^{-1}$ otrzymujemy $f(a)f(s)^{-1}=f(b)f(t)^{-1}$. Wzory dodawania i mnożenia ułamków oraz przemienność $T$ dają
\[
\begin{aligned}
\bar f(a/s+b/t)
&=f(at+bs)f(st)^{-1}\\
&=(f(a)f(t)+f(b)f(s))f(s)^{-1}f(t)^{-1}\\
&=f(a)f(s)^{-1}+f(b)f(t)^{-1},\\
\bar f((a/s)(b/t))
&=f(ab)f(st)^{-1}\\
&=f(a)f(s)^{-1}f(b)f(t)^{-1}.
\end{aligned}
\]
Ponadto $\bar f(1)=1$, $\bar f\iota=f$. Każdy homomorfizm spełniający ostatnią równość musi wysyłać $(a/1)(s/1)^{-1}$ na $f(a)f(s)^{-1}$. Zatem $(S^{-1}R,\iota)$ reprezentuje $F$. Dla homomorfizmu $h:T\to T'$ mamy $(h\bar f)\iota=hf$. Jednoznaczność przedłużenia daje $h\bar f=\overline{hf}$, co wyraża naturalność w pierścieniu docelowym. Pełną konstrukcję lokalizacji podano w AA-12 i w tomie geometrii algebraicznej (ALG-04).
\end{przyklad}

\begin{przyklad}[Iloczyn tensorowy]\label{cat:ex:tensor-rep}
Niech $M$ będzie prawym, a $N$ lewym modułem nad pierścieniem $R$ z jedynką. Dla grupy abelowej $A$ niech $B(A)$ będzie zbiorem funkcji $b:M\times N\to A$ addytywnych w obu zmiennych i spełniających $b(mr,n)=b(m,rn)$. Dla homomorfizmu $h:A\to A'$ funkcja $h\circ b$ nadal jest addytywna i zbilansowana: na przykład
$(hb)(mr,n)=h(b(m,rn))=(hb)(m,rn)$.
Identyczności i złożenia zachowuje się przez składanie funkcji. Zatem $B:\Ab\to\Set$ jest funktorem.

Z konstrukcji iloczynu tensorowego (AA-18) każda taka funkcja $b$ ma jedyną addytywną faktoryzację $\bar b:M\otimes_R N\to A$ z $\bar b(m\otimes n)=b(m,n)$. Jest to dokładnie warunek twierdzenia~\ref{tw:univ-element} dla pary $(M\otimes_R N,\tau)$, gdzie $\tau(m,n)=m\otimes n$. Naturalność bijekcji
\[
\Hom_{\Ab}(M\otimes_R N,A)\cong B(A)
\]
wynika z $(h\bar b)\tau=h(\bar b\tau)=hb$. Nie zakładamy, że dla nieprzemiennego $R$ iloczyn tensorowy dowolnej pary takich modułów ma kanoniczną strukturę $R$-modułu. Dla przemiennego $R$ wersja $R$-dwuliniowa reprezentowana jest w kategorii $R$-modułów (AA-18).
\end{przyklad}

\begin{przyklad}[Różniczki Kählera]\label{cat:ex:kahler-rep}
Dla przemiennych pierścieni z jedynką $k,A$ i ustalonego homomorfizmu $k\to A$ rozważmy funktor $A$-modułów
$M\mapsto\operatorname{Der}_k(A,M)$. Jego elementy to addytywne funkcje $\delta:A\to M$ z $\delta(ab)=a\delta(b)+b\delta(a)$, zerujące się na obrazie $k$. Homomorfizm $A$-modułów $h:M\to N$ działa przez $\delta\mapsto h\delta$: liniowość $h$ zapewnia regułę Leibniza dla $h\delta$, a addytywność i zerowanie się na obrazie $k$ są zachowane.

Przypomnijmy konstrukcję uniwersalnej derywacji, aby uzasadnić istnienie obiektu reprezentującego także dla dowolnej $k$-algebry przemiennej. Niech $E$ będzie wolnym $A$-modułem na symbolach $[a]$, $a\in A$, a $L$ podmodułem generowanym przez
\[
[a+b]-[a]-[b],\quad [ab]-a[b]-b[a],
\quad [\iota(c)]
\]
dla $a,b\in A$, $c\in k$, gdzie $\iota:k\to A$ jest homomorfizmem strukturalnym. Niech $\Omega^1_{A/k}=E/L$ i $d(a)=[a]+L$. Każdy generator $L$ znika w ilorazie, zatem $d$ jest addytywna, spełnia regułę Leibniza i zeruje się na obrazie $k$.

Dla derywacji $\delta:A\to M$ własność wolnego modułu daje jedyny homomorfizm $A$-liniowy $v:E\to M$ z $v([a])=\delta(a)$. Obrazy trzech rodzajów generatorów $L$ są odpowiednio
$\delta(a+b)-\delta(a)-\delta(b)$, $\delta(ab)-a\delta(b)-b\delta(a)$ oraz $\delta(\iota(c))$, czyli zero. Z liniowości $v$ wynika $L\subseteq\Ker v$. Dlatego $\varphi(e+L)=v(e)$ jest dobrze określonym homomorfizmem $A$-liniowym, a $\varphi d=\delta$. Jeżeli $\varphi'$ ma tę samą własność, to $\varphi'$ i $\varphi$ są równe na wszystkich $[a]+L$, które generują iloraz, więc są równe na całym module. Szczególne prezentacje różniczek Kählera omawiano w ALG-04. Zatem
\[
\Hom_A(\Omega^1_{A/k},M)\cong\operatorname{Der}_k(A,M).
\]
Naturalność względem $h$ wyraża równość $(h\varphi)d=h\delta$.
To przykład reprezentowania funktora, a nie utożsamienia wszystkich derywacji z jednym odwzorowaniem $d$.
\end{przyklad}

\subsection{Kryterium braku reprezentowalności}

\begin{stwierdzenie}[Reprezentowalność a obiekt końcowy]\label{cat:prop:rep-terminal}
Jeżeli lokalnie mała kategoria $\mathcal C$ ma obiekt końcowy $1$, to każdy kowariantny funktor reprezentowalny $F:\mathcal C\to\Set$ spełnia $|F(1)|=1$.
\end{stwierdzenie}
\begin{proof}
Dla obiektu reprezentującego $A$ mamy bijekcję $\Hom(A,1)\cong F(1)$. Z definicji obiektu końcowego istnieje dokładnie jedna strzałka $A\to1$, więc oba zbiory mają jeden element.
\end{proof}

\begin{przyklad}[Funktor niereprezentowalny]\label{cat:ex:nonrepresentable}
Funktor zbioru potęgowego $\mathcal P:\Set\to\Set$ działający na funkcjach przez obraz prosty nie jest reprezentowalny. Istotnie, dla funkcji $f:X\to Y$, $g:Y\to Z$ i podzbioru $B\subseteq X$ mamy $(gf)[B]=g[f[B]]$ oraz $\id_X[B]=B$, co sprawdza prawa funktora. Natomiast $\mathcal P(\{*\})=\{\varnothing,\{*\}\}$ ma dwa elementy, co przeczy stwierdzeniu~\ref{cat:prop:rep-terminal}. Nie dotyczy to funktora kontrawariantnego zbioru potęgowego działającego przez przeciwobraz: ten jest reprezentowany przez $\{0,1\}$, ponieważ podzbiór $B\subseteq X$ odpowiada funkcji charakterystycznej $\chi_B:X\to\{0,1\}$. Odwrotność tego przyporządkowania wysyła funkcję $v:X\to\{0,1\}$ na $v^{-1}(\{1\})$: rzeczywiście $\chi_B^{-1}(\{1\})=B$ i $\chi_{v^{-1}(\{1\})}=v$. Dla $h:Y\to X$ i $y\in Y$ obie wartości $\chi_{h^{-1}(B)}(y)$ oraz $\chi_B(h(y))$ są równe $1$ dokładnie wtedy, gdy $h(y)\in B$. Stąd $\chi_{h^{-1}(B)}=\chi_B\circ h$, co daje naturalność.
\end{przyklad}

\section{Lemat Yonedy: obiekt poznajemy przez wszystkie morfizmy}\label{sec:yoneda}

Lemat Yonedy jest jednym z centralnych twierdzeń teorii kategorii. Jego treść można zapamiętać tak: jeśli znamy wszystkie morfizmy wychodzące z obiektu $A$ i wiemy, jak zachowują się przy składaniu, to opisaliśmy ten obiekt w sposób wystarczający, by odzyskać każdą naturalną operację na tych morfizmach.

\begin{twierdzenie}[Lemat Yonedy, postać kowariantna]\label{tw:yoneda}
Niech $\mathcal C$ będzie kategorią lokalnie małą, $A\in\Ob(\mathcal C)$ oraz $F:\mathcal C\to\Set$ funktorem. Istnieje kanoniczna bijekcja
\[
\Nat\bigl(\Hom_{\mathcal C}(A,-),F\bigr)\ \cong\ F(A),
\]
która transformacji naturalnej $\alpha$ przypisuje element $\alpha_A(\id_A)$. Odwrotnie, element $a\in F(A)$ wyznacza transformację naturalną $\alpha^a$, określoną dla $X\in\Ob(\mathcal C)$ i $f:A\to X$ wzorem $\alpha^a_X(f)=F(f)(a)$.
\end{twierdzenie}
\begin{proof}
Najpierw sprawdzamy, że wzór definiuje transformację naturalną. Dla $u:X\to Y$ i $f:A\to X$ mamy $F(u)(\alpha^a_X(f))=F(u)(F(f)(a))=F(u\circ f)(a)=\alpha^a_Y(u\circ f)$, co jest dokładnie warunkiem naturalności.

Załóżmy następnie, że $\alpha:\Hom(A,-)\Rightarrow F$ jest dowolna. Naturalność względem $f:A\to X$ daje
\[
\alpha_X(f)=\alpha_X(f\circ\id_A)=F(f)(\alpha_A(\id_A)).
\]
Zatem $\alpha$ jest jednoznacznie wyznaczona przez $a=\alpha_A(\id_A)$. Odwrotnie, dla transformacji $\alpha^a$ mamy $\alpha^a_A(\id_A)=F(\id_A)(a)=a$. Obie konstrukcje są więc wzajemnie odwrotne.
\end{proof}

\begin{twierdzenie}[Lemat Yonedy, postać kontrawariantna]\label{cat:thm:yoneda-contra}
Dla lokalnie małej kategorii $\mathcal C$, obiektu $A$ oraz $G:\mathcal C^{\op}\to\Set$ istnieje kanoniczna bijekcja
\[
\Nat(\Hom_{\mathcal C}(-,A),G)\cong G(A).
\]
Transformacji $\alpha$ odpowiada $\alpha_A(\id_A)$, a elementowi $a\in G(A)$ odpowiada rodzina $\alpha^a_X(f)=G(f)(a)$ dla $f:X\to A$.
\end{twierdzenie}
\begin{proof}
Dla $v:Y\to X$ i $f:X\to A$ mamy
\[
G(v)(\alpha^a_X(f))=G(v)G(f)(a)
=G(fv)(a)=\alpha^a_Y(fv).
\]
To jest naturalność, ponieważ $\Hom(-,A)$ działa na $v$ przez składanie z prawej strony. Dla dowolnej transformacji $\alpha$ naturalność względem $f:X\to A$ zastosowana do $\id_A$ daje
$\alpha_X(f)=G(f)(\alpha_A(\id_A))$. Zatem po odczytaniu elementu $\alpha_A(\id_A)$ odzyskujemy całą transformację. W drugą stronę $\alpha^a_A(\id_A)=G(\id_A)(a)=a$. Obie konstrukcje są wzajemnie odwrotne.
\end{proof}

\begin{twierdzenie}[Naturalność bijekcji Yonedy]\label{cat:thm:yoneda-natural}
W postaci kowariantnej bijekcja $\Theta_{A,F}:\Nat(\Hom(A,-),F)\to F(A)$ jest naturalna względem obiektu $A$ i funktora $F$. W postaci kontrawariantnej jest naturalna względem $A$ jako obiektu $\mathcal C^{\op}$ oraz względem $G$.
\end{twierdzenie}
\begin{proof}
Piszmy $k_A=\Hom(A,-)$. Dla $r:A\to B$ prekompozycja z $r$ definiuje transformację $k_r:k_B\Rightarrow k_A$ o składowych $f\mapsto fr$. Naturalność tej transformacji wynika z $(gf)r=g(fr)$. Zatem $\alpha:k_A\Rightarrow F$ przechodzi na $\alpha k_r:k_B\Rightarrow F$. Obliczamy
\[
\Theta_{B,F}(\alpha k_r)=\alpha_B(r)
=F(r)(\alpha_A(\id_A))=F(r)(\Theta_{A,F}(\alpha)).
\]
Środkowa równość wynika z naturalności $\alpha$. Dla transformacji $\eta:F\Rightarrow F'$ mamy
\[
\Theta_{A,F'}(\eta\alpha)=\eta_A(\alpha_A(\id_A))
=\eta_A(\Theta_{A,F}(\alpha)).
\]
Te dwa równania dają naturalność w obu zmiennych; ich zgodność dla jednoczesnej zmiany wynika z naturalności $\eta$ względem $r$.

W wersji kontrawariantnej oznaczmy $h_A=\Hom(-,A)$. Strzałka $r:A\to B$ wyznacza $h_r:h_A\Rightarrow h_B$, $f\mapsto rf$. Dla $\beta:h_B\Rightarrow G$ otrzymujemy
\[
\Theta_{A,G}(\beta h_r)=\beta_A(r)
=G(r)(\beta_B(\id_B))=G(r)(\Theta_{B,G}(\beta)).
\]
Dla $\eta:G\Rightarrow G'$ mamy analogicznie
$\Theta_{A,G'}(\eta\beta)=\eta_A(\Theta_{A,G}(\beta))$.
Kierunek pierwszego równania odpowiada dokładnie strzałce $B\to A$ w $\mathcal C^{\op}$.
\end{proof}

\begin{uwaga}[Wielkość rodziny transformacji]
W kategorii lokalnie małej, która nie musi być mała, rodzina składowych transformacji może być indeksowana klasą właściwą. Przy zwykłej konwencji klasowej (lub odpowiedniej hierarchii uniwersów) lemat Yonedy identyfikuje takie transformacje z elementami zbioru $F(A)$ albo $G(A)$, co daje ich małą parametryzację. Nie wynika stąd lokalna małość dowolnej kategorii funktorów o dużej dziedzinie. W twierdzeniu o włożeniu Yonedy przyjmujemy małość $\mathcal C$.
\end{uwaga}


\begin{przyklad}[Wszystkie naturalne wybory wartości funkcji]
Ustalmy zbiór $A$ i rozważmy funktory na $\Set$. Lemat Yonedy daje $\Nat(\Hom_{\Set}(A,-),\id_{\Set})\cong A$. Rzeczywiście, każdy element $a\in A$ definiuje naturalną rodzinę \emph{ewaluacji} $\ev_a^X:\Hom(A,X)\to X$, $f\mapsto f(a)$. Yoneda mówi, że \emph{każda} naturalna rodzina funkcji o tej postaci pochodzi z wyboru jednego $a\in A$.
\end{przyklad}

\begin{przyklad}[Przekształcenia na zbiorach elementów]
W kategorii $\Grp$ niech $A=\mathbb Z$. Homomorfizm $f:\mathbb Z\to G$ jest jednoznacznie zadany przez $f(1)\in G$, zatem funktor $\Hom_{\Grp}(\mathbb Z,-)$ jest naturalnie izomorficzny z funktorem zapominania $U:\Grp\to\Set$. Sformułowanie to wyjaśnia kategorialnie, dlaczego $\mathbb Z$ jest \emph{grupą wolną na jednym generatorze}: morfizmy wychodzące z $\mathbb Z$ są tym samym co pojedyncze elementy grupy.
\end{przyklad}


\begin{przyklad}[Yoneda klasyfikuje operacje na parach i trójkach]\label{przyklad:yoneda-tuple}
W $\Set$ funktor $X\mapsto X^n$ (wysyłający $f:X\to Y$ na $(x_1,\ldots,x_n)\mapsto(f(x_1),\ldots,f(x_n))$) możemy rozumieć jako $\Hom_{\Set}(\{1,\ldots,n\},X)$, gdzie $n\geq1$. Dla $m,n\geq1$ lemat Yonedy daje bijekcję
\[
\Nat\bigl((- )^n,(- )^m\bigr)\cong\Hom_{\Set}(\{1,\ldots,m\},\{1,\ldots,n\}).
\]
Każda naturalna operacja zamieniająca $n$-kę w $m$-kę polega więc wyłącznie na wybieraniu współrzędnych wejścia (z możliwymi powtórzeniami). Na przykład dla $n=2$, $m=3$ jest dokładnie $2^3=8$ takich operacji, m.in. $(x,y)\mapsto(x,y,x)$ oraz $(x,y)\mapsto(y,y,y)$. Nie istnieje naturalne ``dodawanie'' elementów dowolnego zbioru: wymagałoby dodatkowej struktury.
\end{przyklad}

\begin{przyklad}[Wszystkie naturalne mapy $X\to\mathcal P(X)$]
Niech $\mathcal P_+:\Set\to\Set$ oznacza kowariantny funktor potęgowy, działający na funkcjach przez obrazy podzbiorów. Ponieważ $\Hom_{\Set}(1,-)\cong\id_{\Set}$, z lematu Yonedy wynika
$\Nat(\id_{\Set},\mathcal P_+)\cong\mathcal P_+(1)$. Jednoelementowy zbiór ma dokładnie dwa podzbiory. Stąd istnieją dokładnie dwie naturalne rodziny funkcji $X\to\mathcal P(X)$: $x\mapsto\varnothing$ oraz $x\mapsto\{x\}$. Ich naturalność sprawdzamy przez równości $f(\varnothing)=\varnothing$ i $f(\{x\})=\{f(x)\}$. Wniosek jest znacznie silniejszy niż samo wskazanie dwóch przykładów: innych rodzin naturalnych nie ma.
\end{przyklad}

\begin{przyklad}[Yoneda w najmniejszej nietrywialnej kategorii]
Rozważmy kategorię posetu $0<1$. Jej funktory $h_0=\Hom(-,0)$ i $h_1=\Hom(-,1)$ spełniają
$h_0(0)=\{*\}$, $h_0(1)=\varnothing$, $h_1(0)=\{*\}$, $h_1(1)=\{*\}$. Istnieje dokładnie jedna transformacja naturalna $h_0\Rightarrow h_1$, ale żadna $h_1\Rightarrow h_0$. Tak funktory Hom pamiętają, że w wyjściowej kategorii istnieje morfizm $0\to1$, lecz nie $1\to0$.
\end{przyklad}


\begin{twierdzenie}[Włożenie Yonedy]\label{tw:yoneda-embedding}
Niech $\mathcal C$ będzie małą kategorią. Przyporządkowanie $A\mapsto h_A=\Hom_{\mathcal C}(-,A)$, $f:A\to B\mapsto h_f:h_A\Rightarrow h_B$, gdzie $(h_f)_X(u)=f\circ u$, definiuje w pełni wierny funktor $y:\mathcal C\to[\mathcal C^{\op},\Set]$.
\end{twierdzenie}
\begin{proof}
Przyporządkowanie $h_f$ jest naturalne, ponieważ dla $v:Y\to X$ w $\mathcal C$ mamy $(f\circ u)\circ v=f\circ(u\circ v)$. Dla $u:X\to A$ mamy $(h_{\id_A})_X(u)=u$. Dla $f:A\to B$, $g:B\to C$ otrzymujemy $(h_{gf})_X(u)=(gf)u=g(fu)=((h_g)_X\circ(h_f)_X)(u)$. Równość wszystkich składowych daje $h_{gf}=h_gh_f$. Lemat Yonedy w postaci kontrawariantnej daje dla każdych $A,B$ bijekcję $\Nat(h_A,h_B)\cong h_B(A)=\Hom_{\mathcal C}(A,B)$. Jej odwrotność jest dokładnie funkcją $f\mapsto h_f$. Stąd pełność i wierność.
\end{proof}

\begin{wniosek}[Obiekt jest wyznaczony przez swoje morfizmy]
Dla małej kategorii $\mathcal C$ i jej obiektów $A,B$ zachodzi $h_A\cong h_B$ wtedy i tylko wtedy, gdy $A\cong B$.
\end{wniosek}
\begin{proof}
Izomorfizm $A\to B$ daje izomorfizm naturalny $h_A\to h_B$. Odwrotnie, izomorfizm naturalny i jego naturalna odwrotność odpowiadają przez w pełni wierny funktor Yonedy morfizmom $f:A\to B$ i $g:B\to A$, z $h_f=\alpha$ oraz $h_g=\alpha^{-1}$. Wtedy $h_{gf}=h_gh_f=\id_{h_A}=h_{\id_A}$ i $h_{fg}=\id_{h_B}=h_{\id_B}$. Wierność daje $gf=\id_A$ oraz $fg=\id_B$, więc $f$ i $g$ są wzajemnie odwrotne.
\end{proof}

\begin{uwaga}[Na czym polega siła Yonedy?]
Twierdzenie nie stwierdza, że obiekt można po prostu utożsamić z \emph{zbiorem} jego morfizmów. Istotna jest cała rodzina zbiorów $\Hom(X,A)$ dla różnych $X$ i ich zgodność ze składaniem. Inaczej sam zbiór morfizmów, bez informacji o dziedzinach i składaniu, byłby zbyt ubogim opisem.
\end{uwaga}

\begin{stwierdzenie}[Kowariantny odpowiednik włożenia Yonedy]\label{cat:prop:yoneda-cov-embedding}
Dla małej kategorii $\mathcal C$ funktor $\mathcal C^{\op}\to[\mathcal C,\Set]$ zadany przez $A\mapsto k_A=\Hom(A,-)$ oraz $r:A\to B$ w $\mathcal C$ przez $k_r:k_B\Rightarrow k_A$, $(k_r)_X(f)=fr$, jest w pełni wierny.
\end{stwierdzenie}
\begin{proof}
Naturalność $k_r$ sprawdzono w dowodzie twierdzenia~\ref{cat:thm:yoneda-natural}. Dla $r:A\to B$, $s:B\to C$ mamy $(k_r)_X((k_s)_X(f))=(fs)r=f(sr)=(k_{sr})_X(f)$ dla $f:C\to X$, więc $k_{sr}=k_rk_s$. Identyczności zachowują się, ponieważ $f\id_A=f$. Lemat Yonedy daje
$\Nat(k_B,k_A)\cong k_A(B)=\Hom(A,B)$; odwrotność tej bijekcji wysyła $r$ na $k_r$. Jest to dokładnie bijekcja na zbiorach morfizmów w kategorii przeciwnej, co dowodzi pełności i wierności.
\end{proof}

\begin{przyklad}[Naturalne operacje na grupach]\label{cat:ex:yoneda-group-words}
Niech $U:\Grp\to\Set$ będzie funktorem zapominania i $n\geq1$. Grupa wolna $F_n$ na generatorach $s_1,\ldots,s_n$ (konstrukcja i własność uniwersalna w AA-07) reprezentuje $U^n$: homomorfizm $f:F_n\to G$ odpowiada krotce $(f(s_1),\ldots,f(s_n))$. Własność uniwersalna daje istnienie i jednoznaczność $f$, a dla homomorfizmu $g:G\to H$ krotka odpowiadająca $gf$ jest obrazem krotki $f$ przez $U(g)^n$. Jest to więc naturalna bijekcja.

Lemat Yonedy zastosowany do $U$ daje
\[
\Nat(U^n,U)\cong U(F_n).
\]
Dokładniej, element $w\in F_n$ wyznacza operację $t_G(x_1,\ldots,x_n)=f(w)$, gdzie $f(s_i)=x_i$. Homomorfizmy zachowują mnożenie, odwrotności i jedynkę, więc jest to podstawienie elementów $x_i$ do słowa $w$. Dla $g:G\to H$ własność uniwersalna zapewnia, że $gf$ jest homomorfizmem odpowiadającym krotce $(g(x_i))_i$; stąd naturalność. Odwrotnie, dowolna naturalna operacja $t$ ma wartość $w=t_{F_n}(s_1,\ldots,s_n)$ i naturalność względem $f$ wymusza $t_G(x_1,\ldots,x_n)=f(w)$. Dwa różne elementy $F_n$ dają różne operacje już na krotce generatorów. W przeciwieństwie do kategorii zbiorów pojawia się tu mnożenie, np. $xy$ i $yx$ są różnymi operacjami naturalnymi, gdyż $s_1s_2\ne s_2s_1$ w $F_2$.
\end{przyklad}

\section{Obiekty wykładnicze i przestrzenie odwzorowań}\label{sec:exponentials}

Iloczyn pozwala zbudować parę elementów. Inną naturalną konstrukcją jest zbiór \emph{wszystkich odwzorowań} pomiędzy dwoma zbiorami. Poniżej zobaczymy, że również ona ma własność uniwersalną, lecz wymaga już posiadania iloczynów.

\begin{definicja}[Obiekt wykładniczy]\label{def:exponential}
Niech $\mathcal C$ będzie lokalnie małą kategorią mającą wybrane iloczyny dwóch obiektów i niech $A,B\in\Ob(\mathcal C)$. \emph{Obiektem wykładniczym} $B^A$ nazywamy obiekt z morfizmem ewaluacji $\ev:B^A\times A\to B$ o własności: dla każdego $X\in\Ob(\mathcal C)$ i morfizmu $f:X\times A\to B$ istnieje dokładnie jeden morfizm $\widetilde f:X\to B^A$ taki, że $\ev\circ(\widetilde f\times\id_A)=f$. Operacja $(u,v)\mapsto u\times v$ na morfizmach iloczynów jest określona jednoznacznie przez warunki zgodności z rzutami.
\end{definicja}

\subsection{Transpozycja i jej naturalność}

\begin{twierdzenie}[Jednoznaczność obiektu wykładniczego]\label{cat:thm:exp-unique}
Jeżeli $(E,e)$ i $(E',e')$ są obiektami wykładniczymi dla pary $(A,B)$, to istnieje dokładnie jeden izomorfizm $u:E\to E'$ taki, że $e'(u\times\id_A)=e$.
\end{twierdzenie}
\begin{proof}
Własność uniwersalna $(E',e')$ zastosowana do $e:E\times A\to B$ daje jedyny morfizm $u:E\to E'$ z żądaną własnością. Analogicznie otrzymujemy $v:E'\to E$ z $e(v\times\id_A)=e'$. Ponieważ iloczyn morfizmów zachowuje złożenia, mamy
\[
e((vu)\times\id_A)
=e(v\times\id_A)(u\times\id_A)=e.
\]
Również $\id_E$ spełnia to równanie. Jednoznaczność transpozycji $e$ daje $vu=\id_E$. Zamieniając role obu par, otrzymujemy $uv=\id_{E'}$. Zatem $u$ jest izomorfizmem, a jego jednoznaczność wynika z pierwszego zastosowania własności uniwersalnej.
\end{proof}

\begin{twierdzenie}[Naturalna bijekcja wykładnicza]\label{cat:thm:exp-naturality}
Dla ustalonego $A$ i wybranych obiektów wykładniczych $B^A$ bijekcje
\[
\Hom_{\mathcal C}(X\times A,B)
\cong\Hom_{\mathcal C}(X,B^A),\qquad
f\longmapsto\widetilde f
\]
są naturalne w $X$ i $B$. Jeżeli obiekty $B^A$ istnieją dla wszystkich $B$, to przyporządkowanie $B\mapsto B^A$ jest funktorem i $(-)\times A\dashv(-)^A$.
\end{twierdzenie}
\begin{proof}
Oznaczmy ewaluację dla $B$ przez $e_B$. Odwzorowanie odwrotne wysyła $h:X\to B^A$ na $e_B(h\times\id_A)$. Definicja transpozycji daje $e_B(\widetilde f\times\id_A)=f$, natomiast jednoznaczność transpozycji daje $\widetilde{e_B(h\times\id_A)}=h$.

Dla $r:Y\to X$ obliczamy
\[
e_B((\widetilde f r)\times\id_A)
=e_B(\widetilde f\times\id_A)(r\times\id_A)
=f(r\times\id_A).
\]
Stąd $\widetilde{f(r\times\id_A)}=\widetilde f r$. Jest to naturalność w zmiennej kontrawariantnej $X$.

Dla $g:B\to C$ definiujemy $g^A:B^A\to C^A$ jako transpozycję $ge_B$. Zatem $e_C(g^A\times\id_A)=ge_B$. Dla $f:X\times A\to B$ otrzymujemy
\[
e_C((g^A\widetilde f)\times\id_A)
=ge_B(\widetilde f\times\id_A)=gf,
\]
więc $\widetilde{gf}=g^A\widetilde f$. Jest to naturalność w zmiennej $B$. Obie reguły łącznie dają
$\widetilde{gf(r\times\id_A)}=g^A\widetilde f r$.

Ponieważ $e_B(\id_{B^A}\times\id_A)=e_B$, jednoznaczność daje $(\id_B)^A=\id_{B^A}$. Dla $g:B\to C$ i $k:C\to D$ mamy
\[
e_D((k^Ag^A)\times\id_A)
=ke_C(g^A\times\id_A)=kge_B.
\]
Stąd $(kg)^A=k^Ag^A$, co dowodzi funktorialności. Otrzymana naturalna bijekcja jest dokładnie warunkiem sprzężenia; ogólną teorię sprzężeń rozwiniemy w sekcji~\ref{sec:adjunctions}.
\end{proof}

Operację $f\mapsto\widetilde f$ nazywamy transpozycją wykładniczą (ang. \emph{currying}, od nazwiska Haskella Curry'ego).

\begin{stwierdzenie}[Zmiana wykładnika]\label{cat:prop:exp-argument}
Jeśli istnieją $B^A$ i $B^{A'}$, to każda strzałka $a:A'\to A$ wyznacza morfizm $B^a:B^A\to B^{A'}$ przez równanie
\[
e_{A',B}(B^a\times\id_{A'})
=e_{A,B}(\id_{B^A}\times a).
\]
Przyporządkowanie jest kontrawariantne względem $A$ i zgodne z kowariantnym działaniem na $B$.
\end{stwierdzenie}
\begin{proof}
Prawa strona jest morfizmem $B^A\times A'\to B$, więc własność wykładnicza daje istnienie i jednoznaczność $B^a$. Dla $a':A''\to A'$ po zastosowaniu ewaluacji do $B^{a'}B^a$ i użyciu funktorialności iloczynu otrzymujemy
\[
\begin{aligned}
e_{A'',B}((B^{a'}B^a)\times\id_{A''})
&=e_{A',B}(\id\times a')(B^a\times\id_{A''})\\
&=e_{A',B}(B^a\times\id_{A'})(\id\times a')\\
&=e_{A,B}(\id\times(aa')).
\end{aligned}
\]
Jednoznaczność daje $B^{aa'}=B^{a'}B^a$; dla identyczności daje $B^{\id_A}=\id_{B^A}$. Dla $g:B\to C$ obie strzałki $C^ag^A$ i $g^{A'}B^a$ z $B^A$ do $C^{A'}$, po pomnożeniu przez $\id_{A'}$ i złożeniu z $e_{A',C}$, dają $ge_{A,B}(\id\times a)$. Są więc równe. Gdy istnieją wszystkie wykładniki, mamy bifunktor $\mathcal C^{\op}\times\mathcal C\to\mathcal C$, $(A,B)\mapsto B^A$.
\end{proof}

\subsection{Przykłady i ograniczenia}
\begin{przyklad}[Zbiór wszystkich funkcji]\label{przyklad:exponential-set}
W $\Set$ obiekt $B^A$ jest po prostu zbiorem wszystkich funkcji $A\to B$. Ewaluacja ma wzór $\ev(\varphi,a)=\varphi(a)$. Dla $f:X\times A\to B$ określamy $\widetilde f(x)(a)=f(x,a)$. Odwrotnie, każda $h:X\to B^A$ daje funkcję $(x,a)\mapsto h(x)(a)$. Dla $f$ złożenie obu operacji daje w punkcie $(x,a)$ wartość $\widetilde f(x)(a)=f(x,a)$. Dla $h$ transpozycja funkcji $(x,a)\mapsto h(x)(a)$ w punkcie $x$ daje funkcję równą $h(x)$ na każdym $a\in A$. Zatem obie operacje są wzajemnie odwrotne, co dowodzi istnienia i jednoznaczności.

Jeśli $A=\{0,1\}$, to $B^A\cong B\times B$: funkcja na zbiorze dwuelementowym jest opisana przez parę wartości. Jeśli $A=\varnothing$, to $B^A$ jest singletonem, nawet gdy $B=\varnothing$, bo istnieje dokładnie jedna funkcja $\varnothing\to B$.
\end{przyklad}

\begin{przyklad}[Implikacja logiczna jako obiekt wykładniczy]
Niech $S$ będzie zbiorem i rozważmy poset $(\mathcal P(S),\subseteq)$ jako kategorię. Iloczynem obiektów jest przecięcie, a obiektem wykładniczym $B^A$ jest podzbiór $(S\setminus A)\cup B$. Istotnie, dla każdego $X\subseteq S$ mamy równoważność
$X\cap A\subseteq B\ \Longleftrightarrow\ X\subseteq(S\setminus A)\cup B$.
Jeżeli $X\cap A\subseteq B$ i $x\in X$, to albo $x\notin A$, albo $x\in B$. Odwrotnie, jeśli $X\subseteq(S\setminus A)\cup B$ i $x\in X\cap A$, to $x$ nie należy do $S\setminus A$, więc należy do $B$. Ewaluacja jest inkluzją $((S\setminus A)\cup B)\cap A\subseteq B$. Jest to dokładnie bijekcja zbiorów morfizmów potrzebna w definicji~\ref{def:exponential}, ponieważ w posetach każdy niepusty zbiór morfizmów jest jednoelementowy. W algebrach Heytinga takie $B^A$ interpretuje się jako implikację $A\Rightarrow B$.
\end{przyklad}

\begin{przyklad}[Dlaczego $\Grp$ nie ma wszystkich obiektów wykładniczych]
Przypuśćmy, że w $\Grp$ istnieje obiekt wykładniczy $\mathbb Z^{\mathbb Z}$. Ponieważ grupa trywialna $1$ jest końcowa i $1\times\mathbb Z\cong\mathbb Z$, własność uniwersalna wymagałaby bijekcji
$\Hom_{\Grp}(\mathbb Z,\mathbb Z)\cong\Hom_{\Grp}(1,\mathbb Z^{\mathbb Z})$.
Jednak lewy zbiór jest nieskończony: homomorfizm $\mathbb Z\to\mathbb Z$ może wysłać $1$ na dowolną liczbę całkowitą. Prawy zbiór jest jednoelementowy, bo z grupy trywialnej do dowolnej grupy istnieje tylko jeden homomorfizm. Sprzeczność. Obiekt wykładniczy nie istnieje więc dla tej pary.
\end{przyklad}

\begin{definicja}[Kategoria kartezjańsko domknięta]
Kategorię nazywamy \emph{kartezjańsko domkniętą}, jeśli ma obiekt końcowy, iloczyny wszystkich par obiektów i obiekty wykładnicze $B^A$ dla wszystkich par $A,B$. Kategoria $\Set$ jest kartezjańsko domknięta, ale $\Grp$ nie jest.
\end{definicja}

\begin{przyklad}[Przeszkoda w kategoriach z obiektem zerowym]\label{cat:ex:exp-zero}
Jeśli $\mathcal C$ jest lokalnie mała, ma obiekt zerowy $0$ i iloczyny, to istnienie $B^A$ wymusza, że $\Hom(A,B)$ jest jednoelementowy. Istotnie, $0$ jest końcowy, więc $0\times A\cong A$. Z własności wykładniczej
\[
\Hom(A,B)\cong\Hom(0\times A,B)\cong\Hom(0,B^A).
\]
Ostatni zbiór jest jednoelementowy, ponieważ $0$ jest też początkowy. W szczególności w kategorii przestrzeni wektorowych nad ciałem $k$ z iloczynem kartezjańskim nie istnieje $k^k$: identyczność i odwzorowanie zerowe są różnymi endomorfizmami $k$. To nie przeczy istnieniu sprzężenia tensor--Hom (AA-18), gdzie zamiast iloczynu kartezjańskiego występuje iloczyn tensorowy.
\end{przyklad}

\subsection{Prawo wykładnicze w topologii}

\begin{twierdzenie}[Obiekt wykładniczy dla lokalnie zwartej przestrzeni Hausdorffa]\label{cat:thm:exp-top}
Niech $A$ będzie lokalnie zwartą przestrzenią Hausdorffa, a $B$ dowolną przestrzenią topologiczną. Zbiór $C(A,B)$ odwzorowań ciągłych z topologią zwarto-otwartą, oznaczany $C_k(A,B)$, jest obiektem wykładniczym $B^A$ w kategorii $\Top$. Ewaluacja to $e(h,a)=h(a)$.
\end{twierdzenie}
\begin{proof}
Przypomnijmy aparat z TO-18: podbazę topologii zwarto-otwartej tworzą zbiory
\[
[K,U]=\{h\in C(A,B):h(K)\subseteq U\},
\]
gdzie $K\subseteq A$ jest zwarty, a $U\subseteq B$ otwarty. Najpierw sprawdzamy ciągłość ewaluacji. Dla $(h,a)$ i otwartego $U$ zawierającego $h(a)$, lokalna zwartość i własność Hausdorffa dają otwarte $W$ oraz zwarty $K$ z $a\in W\subseteq K\subseteq h^{-1}(U)$ (lemat o zwartych domknięciach w TO-14, użyty w TO-18). Zatem $[K,U]\times W$ jest otwartym otoczeniem $(h,a)$, a dla $(g,x)$ w tym otoczeniu mamy $g(x)\in U$. To dowodzi ciągłości $e$.

Niech $f:X\times A\to B$ będzie ciągłe. Każdy przekrój $f_x(a)=f(x,a)$ jest ciągły jako złożenie z odwzorowaniem $a\mapsto(x,a)$. Definiujemy $\widetilde f(x)=f_x$. Jeśli $x_0$ należy do przeciwobrazu $[K,U]$, to $\{x_0\}\times K\subseteq f^{-1}(U)$. Dla każdego $a\in K$ wybieramy otwarte $V_a\ni x_0$, $W_a\ni a$ z $V_a\times W_a\subseteq f^{-1}(U)$. Ze zwartości $K$ wybieramy skończone pokrycie przez $W_{a_1},\ldots,W_{a_m}$. Otwarte otoczenie $V=\bigcap_{i=1}^m V_{a_i}$ spełnia $V\times K\subseteq f^{-1}(U)$. Dla $K=\varnothing$ bierzemy $V=X$. Stąd $V\subseteq\widetilde f^{-1}([K,U])$, więc wszystkie przeciwobrazy zbiorów podbazowych są otwarte. To dowodzi ciągłości transpozycji.

Odwrotnie, dla ciągłego $h:X\to C_k(A,B)$ funkcja $e(h\times\id_A)$ jest ciągła jako złożenie ciągłych odwzorowań. Obie operacje są wzajemnie odwrotne: równości sprawdzamy punktowo, tak jak w przykładzie~\ref{przyklad:exponential-set}. Otrzymujemy zatem własność uniwersalną w $\Top$, a nie tylko bijekcję zbiorów wszystkich funkcji.
\end{proof}

\begin{uwaga}
Hipoteza lokalnej zwartości Hausdorffa jest wystarczająca, nie twierdzimy, że jest konieczna w dowolnej kategorii przestrzeni. Sama konstrukcja topologii zwarto-otwartej ma sens bez tych hipotez, ale ciągłość wspólnej ewaluacji może zawodzić. Przykład dla przestrzeni liczb wymiernych omówiono w TO-18. Nie wnioskujemy stąd, że $\Top$ jest kartezjańsko domknięta.
\end{uwaga}

\subsection{Rachunek obiektów wykładniczych}

\begin{stwierdzenie}[Podstawowe izomorfizmy]\label{cat:prop:exp-laws}
W lokalnie małej kategorii kartezjańsko domkniętej, dla obiektu końcowego $1$, zachodzą kanoniczne izomorfizmy
\[
B^1\cong B,\qquad 1^A\cong 1,\qquad
(B\times C)^A\cong B^A\times C^A,\qquad
(B^A)^D\cong B^{D\times A}.
\]
\end{stwierdzenie}
\begin{proof}
W każdym przypadku pokażemy naturalną w $X$ bijekcję zbiorów morfizmów z $X$ do obu stron. Dla pierwszego izomorfizmu mamy
$\Hom(X,B^1)\cong\Hom(X\times1,B)\cong\Hom(X,B)$, gdyż rzut $X\times1\to X$ jest naturalnym izomorfizmem, którego odwrotnością jest para identyczności i jedynej strzałki do $1$.
Dla drugiego $\Hom(X,1^A)\cong\Hom(X\times A,1)$ jest singletonem dla każdego $X$, więc $1^A$ jest końcowy.
Dla trzeciego korzystamy z kolejnych własności uniwersalnych:
\[
\begin{aligned}
\Hom(X,(B\times C)^A)
&\cong\Hom(X\times A,B\times C)\\
&\cong\Hom(X\times A,B)\times\Hom(X\times A,C)\\
&\cong\Hom(X,B^A)\times\Hom(X,C^A)\\
&\cong\Hom(X,B^A\times C^A).
\end{aligned}
\]
Środkowa bijekcja jest iloczynem bijekcji wykładniczych; naturalność każdej wynika z twierdzenia~\ref{cat:thm:exp-naturality} albo własności iloczynu.
Wreszcie
\[
\begin{aligned}
\Hom(X,(B^A)^D)
&\cong\Hom(X\times D,B^A)\\
&\cong\Hom((X\times D)\times A,B)\\
&\cong\Hom(X\times(D\times A),B)\\
&\cong\Hom(X,B^{D\times A}).
\end{aligned}
\]
Trzecia bijekcja jest indukowana przez kanoniczny izomorfizm łączności iloczynu, naturalny w $X$ (sekcja~\ref{sec:produkty}). Pozostałe są naturalne na mocy twierdzenia~\ref{cat:thm:exp-naturality}. Zatem złożone bijekcje są izomorfizmami naturalnymi funktorów reprezentowalnych. Wniosek o jednoznaczności obiektów reprezentujących (sekcja~\ref{sec:yoneda}) daje żądane izomorfizmy.
\end{proof}

\section{Obiekty nad ustaloną bazą i kategorie przecinkowe}\label{sec:slice}

W geometrii często nie rozważa się po prostu przestrzeni $X$, lecz przestrzeń $X$ \emph{nad bazą} $B$, czyli odwzorowanie $p:X\to B$. Na przykład wiązka wektorowa ma projekcję na rozmaitość bazową. Naturalną kategorią dla takich danych jest kategoria obiektów nad $B$.

\begin{definicja}[Kategoria obiektów nad ustaloną bazą, slice category]\label{def:slice}
Niech $\mathcal C$ będzie kategorią i $B\in\Ob(\mathcal C)$. \emph{Kategorią obiektów nad $B$}, oznaczaną $\mathcal C/B$, nazywamy kategorię, której obiektami są morfizmy $p:X\to B$ w $\mathcal C$, a morfizm z $(X,p)$ do $(Y,q)$ jest morfizmem $u:X\to Y$ w $\mathcal C$ spełniającym $q\circ u=p$. Składanie i identyczności dziedziczymy z $\mathcal C$.
\end{definicja}

Sprawdźmy aksjomaty: $p\circ\id_X=p$, a jeśli $q\circ u=p$ i $r\circ v=q$, to
\[
r\circ(v\circ u)=(r\circ v)\circ u=q\circ u=p.
\]
Złożenie pozostaje morfizmem nad $B$. Łączność i prawa identyczności wynikają z aksjomatów $\mathcal C$. Funktor zapominający do $\mathcal C$ wysyła $(X,p)$ na $X$, a $u$ na $u$; zachowuje złożenia i identyczności oraz jest wierny, ponieważ na zbiorach morfizmów działa przez inkluzje.

\begin{stwierdzenie}[Obiekt początkowy i końcowy]\label{cat:prop:slice-initial-terminal}
Obiektem końcowym $\mathcal C/B$ jest $(B,\id_B)$. Jeżeli $\mathcal C$ ma obiekt początkowy $I$, to jedyna strzałka $i_B:I\to B$ jest obiektem początkowym $\mathcal C/B$.
\end{stwierdzenie}
\begin{proof}
Morfizm z $(X,p)$ do $(B,\id_B)$ musi spełniać $\id_B\circ u=p$, czyli $u=p$. Ta strzałka istnieje i jest jedyna. Z początkowości $I$ istnieje dokładnie jedna strzałka $i_X:I\to X$. Ponieważ $p\circ i_X$ i $i_B$ mają źródło $I$ i cel $B$, są równe. Zatem $i_X$ jest morfizmem nad $B$; każdy inny musiałby być inną strzałką $I\to X$, co jest niemożliwe.
\end{proof}

\begin{przyklad}
W $\Set/B$ obiektem początkowym jest $\varnothing\to B$, a końcowym $\id_B$. Jeśli $B=\varnothing$, istnienie funkcji $X\to B$ wymusza $X=\varnothing$.
\end{przyklad}

Obiekt kategorii $\mathcal C/B$ to więc cała strzałka, a nie samo $X$. Morfizm wyznacza przemienny diagram
\[
\begin{CD}
X @>{u}>> Y\\
@V{p}VV @VV{q}V\\
B @>>{\id_B}> B.
\end{CD}
\]

\begin{przyklad}[Zbiory nad zbiorem bazowym]
W $\Set/B$ obiektami są funkcje $p:X\to B$. Możemy rozumieć je jako \emph{rodziny zbiorów} $(X_b)_{b\in B}$, gdzie $X_b=p^{-1}(\{b\})$. Morfizm $u:(X,p)\to(Y,q)$ zachowuje włókna: z $p(x)=b$ wynika $q(u(x))=b$, więc $u$ wysyła $X_b$ do $Y_b$. Odwrotnie, dla rodziny funkcji $u_b:X_b\to Y_b$ określamy $u(x)=u_{p(x)}(x)$. Jest to funkcja dobrze określona, ponieważ $x\in X_{p(x)}$, a jej wartość należy do $Y_{p(x)}$, więc $q(u(x))=p(x)$. Ograniczenia $u$ do włókien są zadanymi $u_b$; jednoznaczność wynika z tego, że każdy $x$ należy do dokładnie jednego włókna.

Każda rodzina $(S_b)_{b\in B}$ powstaje, z dokładnością do kanonicznych bijekcji włókien, z projekcji
\[
\{(b,s):b\in B,\ s\in S_b\}\longrightarrow B,\qquad (b,s)\longmapsto b.
\]
Dla pierwotnego $p:X\to B$ odwzorowania $x\mapsto(p(x),x)$ i $(b,x)\mapsto x$ są wzajemnie odwrotne nad $B$. Są naturalne, bo dla $u$ nad $B$ zachodzi $(q(u(x)),u(x))=(p(x),u(x))$. Ograniczenie do włókien zachowuje identyczności i złożenia punktowo; konstrukcja sumy rozłącznej działa na morfizmach przez $(b,s)\mapsto(b,u_b(s))$, więc również zachowuje identyczności i złożenia. Bijekcja włókna z $S_b$ jest naturalna, gdyż $(b,s)\mapsto s$ komutuje z każdą funkcją $u_b$. Otrzymujemy równoważność kategorii $\Set/B$ i kategorii rodzin zbiorów indeksowanych przez $B$.
\end{przyklad}

\begin{stwierdzenie}[Pullback jako iloczyn nad bazą]\label{cat:prop:slice-product}
Niech $p:X\to B$ i $q:Y\to B$ będą obiektami kategorii $\mathcal C/B$. Ich iloczyn istnieje wtedy i tylko wtedy, gdy istnieje pullback $X\times_B Y$ w $\mathcal C$ (definicja~\ref{def:pullback}). Wówczas wraz ze wspólną projekcją $r=p\circ\pi_X=q\circ\pi_Y$ jest on ich iloczynem nad $B$.
\end{stwierdzenie}
\begin{proof}
Morfizmy rzutów $\pi_X:X\times_B Y\to X$ i $\pi_Y:X\times_B Y\to Y$ są morfizmami w $\mathcal C/B$, ponieważ $p\pi_X=r=q\pi_Y$. Niech $(Z,t)$ będzie obiektem nad $B$ z morfizmami $u:Z\to X$, $v:Z\to Y$ nad $B$. Wówczas $pu=t=qv$, więc z uniwersalności pullbacku otrzymujemy dokładnie jeden $h:Z\to X\times_B Y$ o własnościach $\pi_Xh=u$, $\pi_Yh=v$. Zachodzi $rh=p\pi_Xh=pu=t$, a zatem $h$ jest także morfizmem nad $B$. Każdy morfizm nad $B$ spełniający te same warunki jest równy $h$ z jednoznaczności pullbacku.

Odwrotnie, niech $(P,r)$ z rzutami $\pi_X,\pi_Y$ będzie iloczynem nad $B$. Rzuty spełniają $p\circ\pi_X=r=q\circ\pi_Y$. Dla dowolnej pary strzałek $u:Z\to X$, $v:Z\to Y$ o własności $p\circ u=q\circ v$ oznaczmy wspólną strzałkę przez $t$. Wówczas $u,v$ są morfizmami nad $B$ z $(Z,t)$, więc istnieje jedyny morfizm $h:(Z,t)\to(P,r)$ o zadanych złożeniach z rzutami. Jeśli dowolna strzałka $k:Z\to P$ w $\mathcal C$ ma te same złożenia, to $r\circ k=p\circ\pi_X\circ k=p\circ u=t$. Jest więc morfizmem nad $B$ i $k=h$. To dowodzi własności uniwersalnej pullbacku.
\end{proof}

\begin{przyklad}[Włókna iloczynu nad bazą]
Dla $p:X\to B$ i $q:Y\to B$ w $\Set$ mamy $(X\times_B Y)_b\cong X_b\times Y_b$. Para $(x,y)$ należy do włókna nad $b$ wtedy i tylko wtedy, gdy $p(x)=b=q(y)$. Ta równość wyjaśnia, dlaczego pullback nazywa się \emph{iloczynem włóknistym}.
\end{przyklad}

\begin{definicja}[Ogólna kategoria przecinkowa]\label{cat:def:comma}
Jeżeli $F:\mathcal A\to\mathcal C$ i $G:\mathcal D\to\mathcal C$ są funktorami, to \emph{kategoria przecinkowa} $(F\downarrow G)$ ma obiekty $(A,D,u)$, gdzie $A\in\Ob(\mathcal A)$, $D\in\Ob(\mathcal D)$ oraz $u:F(A)\to G(D)$. Morfizm $(a,d):(A,D,u)\to(A',D',u')$ składa się ze strzałek $a:A\to A'$ i $d:D\to D'$ spełniających $G(d)\circ u=u'\circ F(a)$. Identyczności i złożenia określamy współrzędnie.
\end{definicja}

Para $(\id_A,\id_D)$ spełnia warunek przemienności, gdyż funktory zachowują identyczności. Jeżeli dodatkowo $(a',d'):(A',D',u')\to(A'',D'',u'')$ jest morfizmem, to
\[
\begin{aligned}
G(d'\circ d)\circ u
&=G(d')\circ G(d)\circ u\\
&=G(d')\circ u'\circ F(a)\\
&=u''\circ F(a')\circ F(a)\\
&=u''\circ F(a'\circ a).
\end{aligned}
\]
Zatem złożenie jest morfizmem. Łączność i prawa identyczności zachodzą w obu współrzędnych. Rzuty na $\mathcal A$ i $\mathcal D$ są funktorami, ponieważ identyczności i złożenia określono współrzędnie.

\begin{przyklad}[Obiekty nad bazą i pod bazą]
Niech $\mathbf 1$ będzie kategorią z jednym obiektem i jego identycznością, a $\Delta_B:\mathbf 1\to\mathcal C$ funktorem wskazującym $B$. Kategoria $(\id_{\mathcal C}\downarrow\Delta_B)$ jest kanonicznie izomorficzna z $\mathcal C/B$: obiektem jest strzałka $p:X\to B$, a warunek na morfizm redukuje się do $q\circ u=p$.

Kategoria $(\Delta_B\downarrow\id_{\mathcal C})$ opisuje obiekty pod $B$ i jest oznaczana $B/\mathcal C$. Jej obiektami są strzałki $p:B\to X$, a morfizm do $q:B\to Y$ jest strzałką $u:X\to Y$ z $u\circ p=q$. Obiektem początkowym jest $\id_B$, ponieważ jedyną strzałką spełniającą $u\circ\id_B=p$ jest $p$. W obu identyfikacjach usunięcie jedynego obiektu i morfizmu kategorii $\mathbf 1$ ma odwrotność polegającą na ich dołączeniu; obie operacje zachowują składanie i identyczności.
\end{przyklad}

Kategorie przecinkowe są przydatne przy opisie obiektów uniwersalnych i granic. W geometrii nie należy zakładać istnienia dowolnych iloczynów włóknistych rozmaitości; zob. przykład~\ref{cat:ex:manifold-pullback}.



\section{Równoważności kategorii}\label{sec:equivalence}

W algebrze odróżniamy równość od izomorfizmu: przestrzeń $V$ może być izomorficzna z $k^n$, choć nie jest tym samym zbiorem. Podobne zjawisko występuje na poziomie całych kategorii. Równość kategorii jest zwykle zbyt rygorystyczna, dlatego wprowadzamy równoważność.

\begin{definicja}[Równoważność kategorii]\label{def:equivalence}
Funktor $F:\mathcal C\to\mathcal D$ jest \emph{równoważnością kategorii}, jeśli istnieje funktor $G:\mathcal D\to\mathcal C$ oraz izomorfizmy naturalne $G\circ F\cong\id_{\mathcal C}$ i $F\circ G\cong\id_{\mathcal D}$. Mówimy wtedy, że $\mathcal C$ i $\mathcal D$ są \emph{równoważne}, i piszemy $\mathcal C\simeq\mathcal D$. Funktor $G$ jest \emph{quasi-odwrotnością} $F$.
\end{definicja}

Nie wymagamy, aby $G\circ F$ oraz $F\circ G$ były \emph{równe} odpowiednim identycznościom. Wystarcza, że można je z identycznościami porównać przez izomorfizmy naturalne.

\begin{definicja}[Istotna surjektywność na obiektach]
Funktor $F:\mathcal C\to\mathcal D$ jest \emph{istotnie surjektywny na obiektach}, jeśli dla każdego obiektu $Y\in\Ob(\mathcal D)$ istnieje $X\in\Ob(\mathcal C)$ oraz izomorfizm $F(X)\cong Y$ w $\mathcal D$.
\end{definicja}

\begin{twierdzenie}[Kryterium równoważności]\label{tw:equivalence}
Niech $\mathcal C$ i $\mathcal D$ będą małymi kategoriami. Funktor $F:\mathcal C\to\mathcal D$ jest równoważnością wtedy i tylko wtedy, gdy jest w pełni wierny i istotnie surjektywny na obiektach.
\end{twierdzenie}
\begin{proof}
Załóżmy najpierw, że $F$ jest równoważnością. Wybierzmy quasi-odwrotność $G$ i izomorfizmy naturalne $\alpha:GF\Rightarrow\id_{\mathcal C}$ oraz $\beta:FG\Rightarrow\id_{\mathcal D}$. Dla każdego $Y\in\Ob(\mathcal D)$ morfizm $\beta_Y:FG(Y)\to Y$ jest izomorfizmem, co dowodzi istotnej surjektywności. Jeśli $f,g:A\to B$ i $F(f)=F(g)$, to $GF(f)=GF(g)$. Naturalność $\alpha$ daje $\alpha_BGF(f)=f\alpha_A$ oraz $\alpha_BGF(g)=g\alpha_A$; ponieważ $\alpha_A$ jest odwracalny, wynika $f=g$. To dowodzi wierności.

Sprawdźmy pełność bez nieuzasadnionego zakładania, że $F(\alpha_X)=\beta_{F(X)}$. Dla $X\in\Ob(\mathcal C)$ zdefiniujmy automorfizm $c_X:F(X)\to F(X)$ przez $c_X=F(\alpha_X)\circ\beta_{F(X)}^{-1}$. Niech $u:F(A)\to F(B)$ będzie dowolnym morfizmem. Połóżmy $u'=c_B^{-1}\circ u\circ c_A$ i $f=\alpha_B\circ G(u')\circ\alpha_A^{-1}:A\to B$. Naturalność $\beta$ daje
$FG(u')=\beta_{F(B)}^{-1}\circ u'\circ\beta_{F(A)}$, zatem
\[
F(f)=F(\alpha_B)\circ FG(u')\circ F(\alpha_A)^{-1}
=c_B\circ u'\circ c_A^{-1}=u.
\]
Każdy $u$ ma więc przeciwobraz przez $F_{A,B}$; funktor $F$ jest pełny.

Na odwrót, załóżmy pełną wierność i istotną surjektywność $F$. Dla każdego $Y\in\Ob(\mathcal D)$ wybierzmy obiekt $G(Y)\in\Ob(\mathcal C)$ i izomorfizm $\varepsilon_Y:F(G(Y))\to Y$. Dla każdego $Y$ zbiór dopuszczalnych par $(X,\varepsilon)$ jest niepusty dzięki istotnej surjektywności. Jest zbiorem, ponieważ obie kategorie są małe. Rodzina tych zbiorów jest indeksowana zbiorem $\Ob(\mathcal D)$; aksjomat wyboru daje jednoczesny wybór par $(G(Y),\varepsilon_Y)$. Dla strzałki $u:Y\to Y'$ pełność i wierność dają dokładnie jeden morfizm $G(u):G(Y)\to G(Y')$ spełniający
\[
F(G(u))=\varepsilon_{Y'}^{-1}\circ u\circ\varepsilon_Y.
\]
Dla identyczności mamy $F(G(\id_Y))=\varepsilon_Y^{-1}\circ\varepsilon_Y=\id_{F(G(Y))}=F(\id_{G(Y)})$, więc wierność daje $G(\id_Y)=\id_{G(Y)}$. Jeśli $u:Y\to Y'$ i $v:Y'\to Y''$, to
\[
\begin{aligned}
F(G(v)\circ G(u))
&=(\varepsilon_{Y''}^{-1}\circ v\circ\varepsilon_{Y'})
  \circ(\varepsilon_{Y'}^{-1}\circ u\circ\varepsilon_Y)\\
&=\varepsilon_{Y''}^{-1}\circ(v\circ u)\circ\varepsilon_Y\\
&=F(G(v\circ u)).
\end{aligned}
\]
Obie strzałki pod funktorem $F$ należą do $\Hom(G(Y),G(Y''))$, więc z wierności $G(v\circ u)=G(v)\circ G(u)$. Tak powstaje funktor $G:\mathcal D\to\mathcal C$. Mnożąc definiującą równość z lewej przez $\varepsilon_{Y'}$, otrzymujemy $\varepsilon_{Y'}\circ FG(u)=u\circ\varepsilon_Y$. Jest to naturalność $\varepsilon:FG\Rightarrow\id_{\mathcal D}$, a każda składowa jest izomorfizmem z wyboru.

Dla $A\in\Ob(\mathcal C)$ pełna wierność wyznacza dokładnie jeden $\eta_A:A\to GF(A)$ z $F(\eta_A)=\varepsilon_{F(A)}^{-1}$. Z pełności istnieje także $\theta_A:GF(A)\to A$ z $F(\theta_A)=\varepsilon_{F(A)}$. Zatem
\[
\begin{aligned}
F(\theta_A\circ\eta_A)&=\varepsilon_{F(A)}\circ\varepsilon_{F(A)}^{-1}=F(\id_A),\\
F(\eta_A\circ\theta_A)&=\varepsilon_{F(A)}^{-1}\circ\varepsilon_{F(A)}=F(\id_{GF(A)}).
\end{aligned}
\]
Wierność daje $\theta_A\circ\eta_A=\id_A$ i $\eta_A\circ\theta_A=\id_{GF(A)}$. Tym samym $\eta_A$ jest izomorfizmem o odwrotności $\theta_A$. Dla $f:A\to A'$ naturalność $\varepsilon$ daje
$FGF(f)\circ\varepsilon_{F(A)}^{-1}=\varepsilon_{F(A')}^{-1}\circ F(f)$. Stąd
\[
\begin{aligned}
F(GF(f)\circ\eta_A)
&=FGF(f)\circ\varepsilon_{F(A)}^{-1}\\
&=\varepsilon_{F(A')}^{-1}\circ F(f)\\
&=F(\eta_{A'}\circ f).
\end{aligned}
\]
Z wierności $GF(f)\circ\eta_A=\eta_{A'}\circ f$. Zatem $\eta:\id_{\mathcal C}\Rightarrow GF$ jest izomorfizmem naturalnym, a $F$ ma quasi-odwrotność.
\end{proof}


\begin{uwaga}[Zakres założenia o małości]
Małość kategorii w twierdzeniu~\ref{tw:equivalence} jest wygodnym warunkiem uzasadniającym jednoczesny wybór reprezentantów. W praktyce kryterium stosuje się również do kategorii lokalnie małych, przy odpowiedniej konwencji dotyczącej wyborów dla dużych klas. Kategoria wszystkich skończenie wymiarowych przestrzeni wektorowych nie jest dosłownie mała; można jednak pracować z jej małą podkategorią reprezentantów $k^0,k^1,k^2,\ldots$ lub przeprowadzić konstrukcję w ustalonym uniwersum zbiorów. W przykładzie macierzy rozumiemy zatem równoważność w ustalonym uniwersum (z wyborem w większym uniwersum) albo jako równoważność z małym modelem. Bez konwencji umożliwiającej wybór klasowy nie wnioskujemy automatycznie o istnieniu globalnego quasi-odwrotnego funktora na właściwej klasie obiektów.
\end{uwaga}


\begin{przyklad}[Macierze i przestrzenie skończenie wymiarowe]\label{przyklad:mat-equivalence}
Funktor $F:\cat{Mat}_k\to\Vect_k^{\mathrm{fd}}$ z przykładu~\ref{przyklad:mat} jest w pełni wierny, bo każdemu odwzorowaniu liniowemu $k^n\to k^m$ odpowiada dokładnie jedna macierz. Jest też istotnie surjektywny, ponieważ każda skończenie wymiarowa przestrzeń $V$ jest izomorficzna z $k^{\dim V}$. Zatem
\[
\cat{Mat}_k\simeq\Vect_k^{\mathrm{fd}}.
\]
Kategorie te nie są jednak dosłownie równe: w pierwszej obiektami są liczby naturalne, a w drugiej przestrzenie wektorowe (zwykle rozpatrywane bez uprzedniego wyboru baz). To modelowa sytuacja, w której równoważność jest właściwszym pojęciem niż równość.
\end{przyklad}

\begin{przyklad}[Szkielet kategorii]
\emph{Szkieletem} kategorii $\mathcal C$ nazywamy pełną podkategorię zawierającą dokładnie po jednym reprezentancie każdej klasy izomorfizmu obiektów (o ile taki wybór można przeprowadzić). Dla $\Vect_k^{\mathrm{fd}}$ można wybrać reprezentantów $k^0,k^1,k^2,\ldots$, co daje kategorię równoważną $\cat{Mat}_k$. Każda mała kategoria jest równoważna swojemu szkieletowi, lecz szkielet nie musi być dosłownie jednoznaczny. W przykładzie przestrzeni wektorowych, dla pełnej formalności teoriomnogościowej, można pracować w ustalonym uniwersum zbiorów albo zastąpić rozważaną kategorię małym modelem zawierającym reprezentantów wszystkich jej klas izomorfizmu.
\end{przyklad}

\begin{stwierdzenie}[Szkielet małej kategorii]\label{cat:prop:skeleton}
Przy aksjomacie wyboru każda mała kategoria ma szkielet. Inkluzja szkieletu jest równoważnością, a dowolne dwa szkielety tej samej małej kategorii są izomorficzne jako kategorie.
\end{stwierdzenie}
\begin{proof}
Relacja izomorfizmu na zbiorze obiektów jest równoważnością: identyczność daje zwrotność, odwrotność izomorfizmu symetrię, a złożenie izomorfizmów przechodniość. Aksjomat wyboru pozwala wybrać jeden obiekt z każdej niepustej klasy. Pełna podkategoria $\mathcal S$ na wybranych obiektach jest szkieletem. Jej inkluzja jest w pełni wierna z definicji pełnej podkategorii oraz istotnie surjektywna dzięki wyborowi reprezentantów. Stosujemy twierdzenie~\ref{tw:equivalence}.

Dla dwóch szkieletów $\mathcal S,\mathcal T$ tej samej kategorii złożenie inkluzji $\mathcal S\to\mathcal C$ z quasi-odwrotnością inkluzji $\mathcal T\to\mathcal C$ daje funktor $H:\mathcal S\to\mathcal T$ w pełni wierny i istotnie surjektywny: złożenie bijekcji na zbiorach morfizmów jest bijekcją, a istotna surjektywność złożenia wynika z tego, że funktory zachowują izomorfizmy. Quasi-odwrotność inkluzji jest sama równoważnością na mocy definicji~\ref{def:equivalence}.

W kategorii $\mathcal T$ obiekty izomorficzne są równe, więc istotna surjektywność $H$ daje rzeczywistą surjektywność na obiektach. Jeżeli $H(S)=H(S')$, to pełna wierność odzwierciedla izomorfizmy: podnosimy identyczność w obu kierunkach, a ich złożenia są identycznościami dzięki wierności (por. stwierdzenie~\ref{tw:ff-reflect-iso}). Stąd $S\cong S'$, a ponieważ $\mathcal S$ jest szkieletem, mamy $S=S'$. Funktor $H$ jest więc bijekcją na obiektach i na każdym zbiorze morfizmów. Definiujemy jego ścisłą odwrotność przez odwrotność bijekcji na obiektach i jedyne przeciwobrazy morfizmów; zachowanie identyczności i złożeń wynika z ich zachowania przez $H$ oraz jednoznaczności przeciwobrazów. Oba złożenia są identycznościami na obiektach i morfizmach, czyli $H$ jest izomorfizmem kategorii.
\end{proof}

\begin{uwaga}[Kiedy używać jakiego pojęcia?]
Izomorfizm $A\cong B$ porównuje \emph{obiekty tej samej kategorii}. Izomorfizm naturalny $F\cong G$ porównuje \emph{funktory z tymi samymi dziedziną i przeciwdziedziną}. Równoważność $\mathcal C\simeq\mathcal D$ porównuje \emph{całe kategorie}. Są to trzy różne poziomy pojęcia ``to samo z dokładnością do odpowiedniej struktury''.
\end{uwaga}


\section{Sprzężenia funktorów}\label{sec:adjunctions}

Sprzężenie opisuje sytuację, w której konstrukcja nowego obiektu jest \emph{najswobodniejsza} lub \emph{najbardziej uniwersalna} względem zapomnienia pewnej struktury. To jeden z najczęściej występujących wzorców we współczesnej algebrze i topologii.

\subsection{Definicja przez naturalne bijekcje}

\begin{definicja}[Para funktorów sprzężonych]\label{def:adjunction}
Niech $\mathcal C,\mathcal D$ będą lokalnie małymi kategoriami, a $F:\mathcal C\to\mathcal D$ i $G:\mathcal D\to\mathcal C$ funktorami. Mówimy, że $F$ jest \emph{lewym sprzężeniem} $G$, a $G$ \emph{prawym sprzężeniem} $F$, jeśli dla wszystkich $A\in\Ob(\mathcal C)$ i $B\in\Ob(\mathcal D)$ mamy bijekcje
\[
\Phi_{A,B}:\Hom_{\mathcal D}(F(A),B)\ \overset{\cong}{\longrightarrow}\ \Hom_{\mathcal C}(A,G(B)),
\]
naturalne jednocześnie względem $A$ i $B$. Piszemy $F\dashv G$.
\end{definicja}

Naturalność względem obu zmiennych oznacza konkretnie, że dla $u:A'\to A$, $v:B\to B'$ i $h:F(A)\to B$ zachodzą
$\Phi_{A',B}(h\circ F(u))=\Phi_{A,B}(h)\circ u$ oraz $\Phi_{A,B'}(v\circ h)=G(v)\circ\Phi_{A,B}(h)$.

\begin{przyklad}[Wolna przestrzeń wektorowa]\label{przyklad:free-vect}
Dla zbioru $S$ niech $k^{(S)}$ oznacza przestrzeń wektorową funkcji $a:S\to k$ o \emph{skończonym nośniku} (tj. zbiór $\{s:a(s)\neq0\}$ jest skończony). Dla $s\in S$ oznaczmy przez $e_s$ funkcję przyjmującą wartość $1$ w $s$ i $0$ poza $s$. Każdy wektor ma jednoznaczny zapis $\sum_{s\in S}a_s e_s$ z tylko skończenie wieloma niezerowymi współczynnikami. Definiujemy funktor $F:\Set\to\Vect_k$ przez $F(S)=k^{(S)}$ i $F(f)(e_s)=e_{f(s)}$ dla funkcji $f:S\to T$, a następnie rozszerzamy liniowo.

Jeśli $U:\Vect_k\to\Set$ jest funktorem zapominania, to każda funkcja $\varphi:S\to U(V)$ ma dokładnie jedno przedłużenie liniowe $\widetilde\varphi:k^{(S)}\to V$, dane przez $\widetilde\varphi(\sum a_s e_s)=\sum a_s\varphi(s)$. Odwrotnie, przekształceniu liniowemu $L:k^{(S)}\to V$ odpowiada funkcja $s\mapsto L(e_s)$. Otrzymujemy naturalną bijekcję
\[
\Hom_{\Vect_k}(k^{(S)},V)\cong\Hom_{\Set}(S,U(V)),
\]
czyli $F\dashv U$. Wolna przestrzeń jest tak skonstruowana, aby \emph{każda} funkcja ze zbioru generatorów jednoznacznie przedłużała się liniowo.
\end{przyklad}

\begin{przyklad}[Uniwersalny element wolnej przestrzeni]
Dla ustalonego zbioru $S$ i funktora $G:\Vect_k\to\Set$ określonego przez $G(V)=\Hom_{\Set}(S,U(V))$ reprezentującym obiektem jest $k^{(S)}$. Element uniwersalny $u\in G(k^{(S)})$ to funkcja $s\mapsto e_s$. Dla dowolnej $\varphi:S\to U(V)$ istnieje dokładnie jedno odwzorowanie liniowe $L:k^{(S)}\to V$, dla którego $G(L)(u)=\varphi$, czyli $L(e_s)=\varphi(s)$ dla wszystkich $s$. Jest to element uniwersalny odpowiadający sprzężeniu wolnej przestrzeni z funktorem zapominania.
\end{przyklad}

\begin{przyklad}[Grupa wolna]
Dla zbioru $S$ istnieje grupa wolna $F_{\Grp}(S)$ na generatorach $S$ (elementami są zredukowane słowa w symbolach $s$ oraz $s^{-1}$). Każdą funkcję $S\to U(G)$ do zbioru elementów grupy $G$ przedłużamy jednoznacznie do homomorfizmu $F_{\Grp}(S)\to G$. Otrzymujemy sprzężenie $F_{\Grp}:\Set\rightleftarrows\Grp:U$. Dla $S=\{*\}$ grupa wolna jest izomorficzna z $\mathbb Z$, co zgadza się z przykładem w sekcji~\ref{sec:yoneda}.
\end{przyklad}

\begin{przyklad}[Dwie konstrukcje topologiczne]
Niech $U:\Top\to\Set$ zapomina topologię. Dla zbioru $S$ oznaczmy przez $D(S)$ zbiór $S$ z topologią dyskretną (wszystkie podzbiory są otwarte), a przez $I(S)$ zbiór $S$ z topologią antydyskretną (otwarte są tylko $\varnothing$ i $S$). Każda funkcja $D(S)\to X$ jest ciągła, a każda funkcja $X\to I(S)$ jest ciągła. Dlatego mamy dwa sprzężenia
\[
D\dashv U\dashv I.
\]
Rzeczywiście $\Hom_{\Top}(D(S),X)\cong\Hom_{\Set}(S,U(X))$ oraz $\Hom_{\Set}(U(X),S)\cong\Hom_{\Top}(X,I(S))$. Jeden funktor zapominania może zatem mieć zarówno lewe, jak i prawe sprzężenie.
\end{przyklad}

\begin{przyklad}[Przekątna, iloczyn i koprodukt]\label{przyklad:diag-adj}
W $\Set$ niech $\Delta:\Set\to\Set\times\Set$ będzie funktorem przekątnym, $\Delta(X)=(X,X)$. Niech $\times:\Set\times\Set\to\Set$ oznacza iloczyn dwóch zbiorów, a $\sqcup:\Set\times\Set\to\Set$ ich sumę rozłączną. Uniwersalne własności iloczynu i koproduktu można zapisać jako dwa sprzężenia
\[
\sqcup\ \dashv\ \Delta\ \dashv\ \times.
\]
Pierwsze wynika z $\Hom(A\sqcup B,X)\cong\Hom(A,X)\times\Hom(B,X)$, a drugie z $\Hom(X,A\times B)\cong\Hom(X,A)\times\Hom(X,B)$.
\end{przyklad}


\begin{przyklad}[Sprzężenie funkcji monotonicznych na liczbach całkowitych]\label{przyklad:adj-poset}
Traktujmy $(\mathbb Z,\leq)$ jako kategorię. Funkcje $F,G:\mathbb Z\to\mathbb Z$ dane przez $F(n)=2n$ i $G(m)=\lfloor m/2\rfloor$ są monotoniczne, a więc są funktorami. Dla dowolnych liczb całkowitych $n,m$ zachodzi
$F(n)\leq m\ \Longleftrightarrow\ 2n\leq m\ \Longleftrightarrow\ n\leq\lfloor m/2\rfloor\ \Longleftrightarrow\ n\leq G(m)$.
Każda ze stron oznacza niepustość odpowiedniego zbioru morfizmów, a skoro te zbiory są puste lub jednoelementowe, dostajemy bijekcję naturalną $\Hom(F(n),m)\cong\Hom(n,G(m))$. Zatem $F\dashv G$. Na przykład dla $n=3,m=7$ nierówności $6\leq7$ i $3\leq3$ są prawdziwe; dla $n=4,m=7$ obie są fałszywe.
\end{przyklad}

\begin{przyklad}[Obraz, przeciwobraz i kwantyfikator ``dla każdego'']\label{przyklad:adj-inverse-image}
Niech $f:X\to Y$ będzie funkcją. Na posetach podzbiorów mamy trzy monotoniczne funkcje:
\begin{itemize}
\item $\exists_f:\mathcal P(X)\to\mathcal P(Y)$, $\exists_f(A)=f(A)$;
\item $f^{-1}:\mathcal P(Y)\to\mathcal P(X)$, $f^{-1}(B)=\{x:f(x)\in B\}$;
\item $\forall_f:\mathcal P(X)\to\mathcal P(Y)$, $\forall_f(A)=\{y\in Y:f^{-1}(\{y\})\subseteq A\}$.
\end{itemize}
Dla $A\subseteq X$, $B\subseteq Y$ zachodzą równoważności
\[
f(A)\subseteq B\ \Longleftrightarrow\ A\subseteq f^{-1}(B),\qquad
f^{-1}(B)\subseteq A\ \Longleftrightarrow\ B\subseteq\forall_f(A).
\]
Są to dokładnie własności dwóch sprzężeń $\exists_f\dashv f^{-1}\dashv\forall_f$. Pierwszą równoważność sprawdzamy bezpośrednio: $f(x)\in B$ dla wszystkich $x\in A$ wtedy i tylko wtedy, gdy $x\in f^{-1}(B)$ dla wszystkich $x\in A$. Druga mówi, że każde włókno nad $y\in B$ jest zawarte w $A$ wtedy i tylko wtedy, gdy $y\in\forall_f(A)$.

Przykładowo dla $X=\{a,b,c\}$, $Y=\{0,1\}$, $f(a)=0$, $f(b)=f(c)=1$ i $A=\{a,b\}$ otrzymujemy $\exists_f(A)=\{0,1\}$, lecz $\forall_f(A)=\{0\}$. Warto uwzględnić także puste włókna: jeśli $y\notin f(X)$, to $y\notin\exists_f(A)$, lecz $y\in\forall_f(A)$ dla każdego $A\subseteq X$, gdyż $\varnothing\subseteq A$. Obie strony sprzężeń mają więc wyraźnie różne zadania: ``istnieje element włókna'' i ``wszystkie elementy włókna''.
\end{przyklad}

\begin{przyklad}[Naturalność sprzężenia wolnej przestrzeni]
Dla zbioru $S$, przestrzeni $V$, funkcji $\varphi:S\to U(V)$ odpowiadające jej liniowe przedłużenie spełnia $\widetilde\varphi(e_s)=\varphi(s)$. Sprawdźmy obie zmienne. Jeśli $a:S'\to S$, to $(\widetilde\varphi\circ F(a))(e_{s'})=\widetilde\varphi(e_{a(s')})=\varphi(a(s'))$, czyli liniowy morfizm odpowiadający $\varphi\circ a$. Jeśli $b:V\to W$ jest liniowe, to $(b\circ\widetilde\varphi)(e_s)=b(\varphi(s))$, a zatem odpowiada funkcji $U(b)\circ\varphi$. Wykazaliśmy naturalność bijekcji z przykładu~\ref{przyklad:free-vect} względem $S$ i $V$, a nie tylko bijektywność dla każdej ustalonej pary.
\end{przyklad}


\subsection{Jednostka, kojednostka i tożsamości trójkątne}

\begin{definicja}[Jednostka i kojednostka sprzężenia]\label{def:unit-counit}
Dla sprzężenia $F:\mathcal C\rightleftarrows\mathcal D:G$ z bijekcjami $\Phi_{A,B}$ definiujemy \emph{jednostkę} $\eta:\id_{\mathcal C}\Rightarrow GF$ przez $\eta_A:=\Phi_{A,F(A)}(\id_{F(A)})$. Definiujemy \emph{kojednostkę} $\varepsilon:FG\Rightarrow\id_{\mathcal D}$ przez $\varepsilon_B:=\Phi_{G(B),B}^{-1}(\id_{G(B)})$.
\end{definicja}

W sprzężeniach konstrukcji swobodnej z funktorem zapominania jednostka opisuje odwzorowanie generatorów do obiektu swobodnego, a kojednostka jego realizację w obiekcie z daną strukturą. W ogólnym sprzężeniu jednostka i kojednostka nie muszą być odpowiednio monomorfizmem i epimorfizmem. W przykładzie wolnej przestrzeni wektorowej $\eta_S(s)=e_s$, natomiast
$\varepsilon_V:k^{(U(V))}\to V$ spełnia $\varepsilon_V(e_v)=v$ i jest liniowa.

\begin{twierdzenie}[Tożsamości trójkątne i odtworzenie sprzężenia]\label{tw:triangles}
Dla każdego sprzężenia $F\dashv G$ jednostka $\eta$ i kojednostka $\varepsilon$ są transformacjami naturalnymi spełniającymi dla wszystkich $A\in\Ob(\mathcal C)$ i $B\in\Ob(\mathcal D)$ równości
\[
\varepsilon_{F(A)}\circ F(\eta_A)=\id_{F(A)},\qquad G(\varepsilon_B)\circ\eta_{G(B)}=\id_{G(B)}.
\]
Odwrotnie, jeśli dane są funktory $F,G$ i transformacje naturalne $\eta:\id_{\mathcal C}\Rightarrow GF$, $\varepsilon:FG\Rightarrow\id_{\mathcal D}$ spełniające te równości, to definiują sprzężenie $F\dashv G$.
\end{twierdzenie}
\begin{proof}
Najpierw wychodzimy od naturalnych bijekcji $\Phi_{A,B}$ z definicji sprzężenia. Zdefiniowane w sekcji jednostka $\eta$ i kojednostka $\varepsilon$ są naturalne: dla $u:A\to A'$ naturalność $\Phi$ w pierwszej i drugiej zmiennej daje
$GF(u)\eta_A=\Phi_{A,F(A')}(F(u))=\eta_{A'}u$. Dla $v:B\to B'$ naturalność bijekcji daje
\[
\begin{aligned}
\Phi_{G(B),B'}(v\circ\varepsilon_B)
&=G(v)\circ\Phi_{G(B),B}(\varepsilon_B)=G(v),\\
\Phi_{G(B),B'}(\varepsilon_{B'}\circ FG(v))
&=\Phi_{G(B'),B'}(\varepsilon_{B'})\circ G(v)=G(v).
\end{aligned}
\]
Injektywność $\Phi_{G(B),B'}$ daje $v\circ\varepsilon_B=\varepsilon_{B'}\circ FG(v)$, czyli naturalność kojednostki.

Z naturalności $\Phi$ w drugiej zmiennej, zastosowanej do $h:F(A)\to B$ i $\id_{F(A)}$, wynika ważny wzór
$\Phi_{A,B}(h)=G(h)\eta_A$.
Z naturalności $\Phi^{-1}$ w pierwszej zmiennej, zastosowanej do $u:A\to G(B)$ i $\id_{G(B)}$, wynika wzór odwrotny
$\Phi_{A,B}^{-1}(u)=\varepsilon_BF(u)$.
Zastosujmy drugi wzór przy $B=F(A)$ i $u=\eta_A$. Ponieważ $\eta_A=\Phi_{A,F(A)}(\id_{F(A)})$, mamy
\[
\varepsilon_{F(A)}F(\eta_A)
=\Phi_{A,F(A)}^{-1}(\eta_A)=\id_{F(A)}.
\]
Z kolei pierwszy wzór przy $A=G(B)$ i $h=\varepsilon_B$ daje
$G(\varepsilon_B)\eta_{G(B)}=\Phi_{G(B),B}(\varepsilon_B)=\id_{G(B)}$.
To dowodzi obu tożsamości trójkątnych, bez pomijania kroku prowadzącego od naturalności do wzorów na morfizmy sprzężone.

Odwrotnie, niech dane będą naturalne transformacje $\eta:\id_{\mathcal C}\Rightarrow GF$ i $\varepsilon:FG\Rightarrow\id_{\mathcal D}$ spełniające tożsamości trójkątne. Definiujemy dla $h:F(A)\to B$ oraz $u:A\to G(B)$ morfizmy
\[
\Phi_{A,B}(h)=G(h)\eta_A,\qquad
\Psi_{A,B}(u)=\varepsilon_BF(u).
\]
Mają one odpowiednie dziedziny i przeciwdziedziny. Naturalność $\varepsilon$ dla $h$ daje $\varepsilon_BFG(h)=h\varepsilon_{F(A)}$, stąd
$\Psi(\Phi(h))=\varepsilon_BFG(h)F(\eta_A)=h\varepsilon_{F(A)}F(\eta_A)=h$.
Naturalność $\eta$ dla $u$ daje $GF(u)\eta_A=\eta_{G(B)}u$, zatem
$\Phi(\Psi(u))=G(\varepsilon_B)GF(u)\eta_A=G(\varepsilon_B)\eta_{G(B)}u=u$.
Funkcje $\Phi_{A,B}$ i $\Psi_{A,B}$ są więc wzajemnie odwrotne. Naturalność $\Phi$ względem $u_0:A'\to A$ wynika z rachunku
$\Phi_{A',B}(hF(u_0))=G(h)GF(u_0)\eta_{A'}=G(h)\eta_Au_0=\Phi_{A,B}(h)u_0$.
Jej naturalność względem $v:B\to B'$ wynika z $\Phi_{A,B'}(vh)=G(v)G(h)\eta_A=G(v)\Phi_{A,B}(h)$. Tak otrzymujemy wymagane sprzężenie.
\end{proof}

\begin{przyklad}[Sprawdzenie pierwszego trójkąta dla wolnej przestrzeni]
Dla $S\in\Ob(\Set)$ i generatora $e_s\in F(S)$ mamy $F(\eta_S)(e_s)=e_{e_s}\in F(U(F(S)))$, po czym $\varepsilon_{F(S)}(e_{e_s})=e_s$. Obie strony $\varepsilon_{F(S)}\circ F(\eta_S)$ i $\id_{F(S)}$ są liniowe i zgadzają się na bazie $(e_s)_{s\in S}$, zatem są równe. To konkretna realizacja pierwszej tożsamości trójkątnej. Dla drugiej, przy $V\in\Ob(\Vect_k)$ i $v\in U(V)$ mamy $\eta_{U(V)}(v)=e_v$, a następnie $U(\varepsilon_V)(e_v)=v$. Złożenie $U(\varepsilon_V)\circ\eta_{U(V)}$ jest więc identycznością na $U(V)$.
\end{przyklad}


\begin{przyklad}[Oba trójkąty dla sprzężenia dyskretnego]
W sprzężeniu $D:\Set\rightleftarrows\Top:U$ z przykładu w sekcji~\ref{sec:adjunctions} mamy $UD=\id_{\Set}$ dosłownie, więc $\eta_S=\id_S$. Dla przestrzeni $X$ kojednostka $\varepsilon_X:D(U(X))\to X$ jest identycznością jako funkcja między zbiorami, a ciągłość wynika z dyskretności dziedziny. Pierwszy trójkąt mówi, że złożenie $D(S)\xrightarrow{D(\eta_S)}DUD(S)\xrightarrow{\varepsilon_{D(S)}}D(S)$ jest identycznością. Drugi mówi, że $U(X)\xrightarrow{\eta_{U(X)}}UDU(X)\xrightarrow{U(\varepsilon_X)}U(X)$ również jest identycznością. Oba wnioski są natychmiastowe po podstawieniu definicji, ale pokazują różnicę między identycznością funkcji a homeomorfizmem: $\varepsilon_X$ nie musi być homeomorfizmem.
\end{przyklad}


\subsection{Sprzężenia a zachowanie konstrukcji}

\begin{twierdzenie}[Lewe sprzężenia zachowują kogranice]\label{tw:adjoints-preserve}
Niech $F:\mathcal C\to\mathcal D$ i $G:\mathcal D\to\mathcal C$ będą funktorami ze sprzężeniem $F\dashv G$. Niech $J$ będzie małą kategorią i $D:J\to\mathcal C$ diagramem mającym kogranicę $(Q,(q_j)_j)$. Wówczas $(F(Q),(F(q_j))_j)$ jest kogranicą diagramu $F\circ D$ w $\mathcal D$. Dualnie, prawy sprzężony $G$ zachowuje wszystkie istniejące małe granice.
\end{twierdzenie}
\begin{proof}
Rodzina $F(q_j)$ jest kostożkiem, ponieważ dla $\alpha:j\to j'$ zachodzi $F(q_{j'})\circ F(D(\alpha))=F(q_{j'}\circ D(\alpha))=F(q_j)$.

Niech $(b_j:F(D(j))\to Y)_j$ będzie dowolnym kostożkiem. Definiujemy $a_j=\Phi_{D(j),Y}(b_j)$. Naturalność w pierwszej zmiennej daje
\[
a_{j'}\circ D(\alpha)
=\Phi_{D(j),Y}(b_{j'}\circ F(D(\alpha)))=a_j.
\]
Z własności uniwersalnej $Q$ istnieje dokładnie jeden $a:Q\to G(Y)$ taki, że $a\circ q_j=a_j$ dla każdego $j$. Niech $b=\Phi_{Q,Y}^{-1}(a)$. Ponownie z naturalności
\[
\Phi_{D(j),Y}(b\circ F(q_j))=a\circ q_j=a_j
=\Phi_{D(j),Y}(b_j).
\]
Injektywność daje $b\circ F(q_j)=b_j$. Jeśli $b'$ spełnia te same warunki, to $\Phi_{Q,Y}(b')\circ q_j=a_j$ dla wszystkich $j$, więc $\Phi_{Q,Y}(b')=a$ z jednoznaczności kogranicy, a następnie $b'=b$. Obejmuje to również pusty diagram.

Dla granic niech $E:J\to\mathcal D$ ma granicę $(L,(\ell_j)_j)$. Rodzina $G(\ell_j)$ jest stożkiem z funktorialności $G$. Dla dowolnego stożka $(c_j:X\to G(E(j)))_j$ definiujemy $d_j=\Phi_{X,E(j)}^{-1}(c_j)$. Dla $\alpha:j\to j'$ mamy
\[
\Phi_{X,E(j')}(E(\alpha)\circ d_j)
=G(E(\alpha))\circ c_j=c_{j'},
\]
więc $E(\alpha)\circ d_j=d_{j'}$. Istnieje zatem dokładnie jeden $d:F(X)\to L$ z $\ell_j\circ d=d_j$. Dla $c=\Phi_{X,L}(d)$ naturalność daje
\[
G(\ell_j)\circ c=\Phi_{X,E(j)}(\ell_j\circ d)=c_j.
\]
Jeśli $c'$ również spełnia te warunki, to $d'=\Phi_{X,L}^{-1}(c')$ ma $\ell_j\circ d'=d_j$, zatem $d'=d$ i $c'=c$. To dowodzi zachowania granicy przez $G$, również dla pustego diagramu.
\end{proof}

\begin{stwierdzenie}[Pełna wierność a jednostka i kojednostka]\label{cat:prop:adj-full-faithful}
W sprzężeniu $F\dashv G$ funktor $F$ jest w pełni wierny wtedy i tylko wtedy, gdy $\eta$ jest izomorfizmem naturalnym. Funktor $G$ jest w pełni wierny wtedy i tylko wtedy, gdy $\varepsilon$ jest izomorfizmem naturalnym.
\end{stwierdzenie}
\begin{proof}
Dla $A,A'\in\Ob(\mathcal C)$ złożenie odwzorowania na morfizmach $F$ z bijekcją sprzężenia wysyła $f:A\to A'$ na
\[
\Phi_{A,F(A')}(F(f))=GF(f)\circ\eta_A=\eta_{A'}\circ f.
\]
Jeżeli $\eta$ jest izomorfizmem, to ostatnie przyporządkowanie jest bijekcją, której odwrotność to składanie z $\eta_{A'}^{-1}$. Stąd $F$ jest w pełni wierny.

Odwrotnie, z pełnej wierności $F$ wynika, że dla każdego $A$ składanie z $\eta_{A'}$ jest bijekcją. Przy $A=GF(A')$ wybieramy jedyny $r:GF(A')\to A'$ taki, że $\eta_{A'}\circ r=\id_{GF(A')}$. Ponieważ
$\eta_{A'}\circ(r\circ\eta_{A'})=\eta_{A'}=\eta_{A'}\circ\id_{A'}$,
injektywność tej bijekcji przy $A=A'$ daje $r\circ\eta_{A'}=\id_{A'}$. Wszystkie składowe $\eta$ są izomorfizmami.

Dla $B,B'\in\Ob(\mathcal D)$ złożenie odwzorowania na morfizmach $G$ z odwrotną bijekcją sprzężenia wysyła $g:B\to B'$ na
\[
\Phi_{G(B),B'}^{-1}(G(g))=\varepsilon_{B'}\circ FG(g)=g\circ\varepsilon_B.
\]
Gdy $\varepsilon$ jest izomorfizmem, składanie z nią jest bijekcją, więc $G$ jest w pełni wierny. Jeżeli $G$ jest w pełni wierny, dla każdego $B'$ składanie z $\varepsilon_B$ jest bijekcją. Przy $B'=FG(B)$ otrzymujemy $t:B\to FG(B)$ z $t\circ\varepsilon_B=\id_{FG(B)}$. Ponieważ $(\varepsilon_B\circ t)\circ\varepsilon_B=\varepsilon_B=\id_B\circ\varepsilon_B$, injektywność przy $B'=B$ daje $\varepsilon_B\circ t=\id_B$. Zatem także $\varepsilon$ jest izomorfizmem naturalnym.
\end{proof}

\begin{przyklad}[Dlaczego wolna przestrzeń zachowuje koprodukty]
Dla zbiorów $S,T$ mamy kanoniczny izomorfizm $k^{(S\sqcup T)}\cong k^{(S)}\oplus k^{(T)}$. Można go sprawdzić bezpośrednio: funkcja o skończonym nośniku na sumie rozłącznej to para funkcji o skończonych nośnikach na jej składnikach. Twierdzenie~\ref{tw:adjoints-preserve} przewiduje ten izomorfizm bez rachunku na współrzędnych, ponieważ funktor wolnej przestrzeni jest lewym sprzężeniem i musi zachowywać koprodukty.
\end{przyklad}

\begin{uwaga}[Lewe i prawe są asymetryczne]
Lewe sprzężenie zachowuje kogranice, a prawe granice. Nie można bez dodatkowych założeń zamienić tych słów. Na przykład wolna przestrzeń $F:\Set\to\Vect_k$ nie zachowuje wszystkich iloczynów: dla zbiorów $S$ i $T$ o odpowiednio $2$ i $3$ elementach przestrzeń $F(S\times T)$ ma wymiar $6$, podczas gdy iloczyn $F(S)\times F(T)$ w $\Vect_k$ ma wymiar $2+3=5$.
\end{uwaga}


\begin{przyklad}[Prawy sprzężony zachowuje iloczyny]
Niech $U:\Grp\to\Set$ będzie funktorem zapominania. Ponieważ $U$ jest prawym sprzężeniem do funktora grupy wolnej, zachowuje granice. Dla grup $G,H$ dostajemy zatem kanoniczną bijekcję $U(G\times H)\cong U(G)\times U(H)$. Konkretna grupa $G\times H$ ma elementy $(g,h)$, więc bijekcja jest tu oczywista. Znaczenie twierdzenia polega jednak na tym, że działa ono także dla bardziej skomplikowanych granic, np. equalizatorów dowolnych homomorfizmów grup.
\end{przyklad}



\section{Monady: jak opisać konstrukcję swobodną na jednym poziomie}\label{sec:monads}

Monady pojawiają się naturalnie w teorii algebry, semantyce języków programowania i teorii obliczeń. Na poziomie podstawowym można je rozumieć jako funktory wyposażone w operację \emph{włożenia pojedynczego elementu} i operację \emph{spłaszczenia dwóch warstw tej samej konstrukcji}.

\begin{definicja}[Monada]\label{def:monad}
\emph{Monadą} na kategorii $\mathcal C$ jest trójka $(T,\eta,\mu)$, gdzie $T:\mathcal C\to\mathcal C$ jest funktorem, $\eta:\id_{\mathcal C}\Rightarrow T$ i $\mu:T\circ T\Rightarrow T$ są transformacjami naturalnymi, a dla każdego obiektu $A$ zachodzą równości
\[
\mu_A\circ T(\mu_A)=\mu_A\circ\mu_{T(A)},\qquad\mu_A\circ T(\eta_A)=\id_{T(A)}=\mu_A\circ\eta_{T(A)}.
\]
Pierwsza równość to \emph{łączność mnożenia}, a dwie pozostałe to prawa jednostki. Piszemy $T^2=T\circ T$ i $T^3=T\circ T\circ T$. Naturalność oznacza, że dla $f:A\to B$ zachodzą
\[
T(f)\circ\eta_A=\eta_B\circ f,\qquad
T(f)\circ\mu_A=\mu_B\circ T^2(f).
\]
Łączność jest przemiennością kwadratu, a prawa jednostki przemiennością dwóch trójkątów w diagramach:
\[
\begin{CD}
T^3(A) @>{T(\mu_A)}>> T^2(A)\\
@V{\mu_{T(A)}}VV @VV{\mu_A}V\\
T^2(A) @>>{\mu_A}> T(A).
\end{CD}
\qquad
\begin{CD}
T(A) @>{T(\eta_A)}>> T^2(A) @<{\eta_{T(A)}}<< T(A)\\
@V{\id}VV @VV{\mu_A}V @VV{\id}V\\
T(A) @= T(A) @= T(A).
\end{CD}
\]
\end{definicja}

\begin{przyklad}[Monada skończonych list]\label{przyklad:list-monad}
Dla zbioru $X$ niech $T(X)$ oznacza zbiór wszystkich skończonych list elementów $X$, włącznie z listą pustą. Funkcja $f:X\to Y$ działa na listach przez zastosowanie $f$ do każdego elementu. Definiujemy $\eta_X(x)=[x]$ oraz $\mu_X([[x_1,x_2],[x_3],[],[x_4]])=[x_1,x_2,x_3,x_4]$. W ogólności $\mu_X$ skleja listę list w jedną listę. Funktorialność $T$ wynika z punktowego działania funkcji na listach. Dla $f:X\to Y$ mamy $T(f)([x])=[f(x)]$, co daje naturalność $\eta$. Zastosowanie $f$ do konkatenacji list daje konkatenację list otrzymanych przez zastosowanie $f$ do każdej z nich; zatem $T(f)\circ\mu_X=\mu_Y\circ T^2(f)$. Dla listy list list obie strony równania łączności wypisują elementy $x_{ijk}$ w kolejności leksykograficznej indeksów $(i,j,k)$, niezależnie od tego, który poziom konkatenacji usuwamy najpierw. Dotyczy to także pustych list na każdym poziomie. Równania jednostek mówią, że spłaszczenie $[[x_1,\ldots,x_n]]$ i spłaszczenie $[[x_1],\ldots,[x_n]]$ daje $[x_1,\ldots,x_n]$.
\end{przyklad}

\begin{przyklad}[Monada z wolnego monoidu]
Skończone listy elementów zbioru $X$ tworzą \emph{monoid wolny} na $X$ z działaniem konkatenacji i pustą listą jako jednostką. Niech $F:\Set\to\Mon$ będzie funktorem wolnego monoidu, a $U:\Mon\to\Set$ funktorem zapominania. Dla funkcji $f:X\to U(M)$ jej przedłużenie na listy ma postać $[x_1,\ldots,x_n]\mapsto f(x_1)\cdots f(x_n)$, z pustym iloczynem równym jednostce $M$. Zachowuje konkatenację, a każdy homomorfizm przedłużający $f$ musi mieć tę postać. Składanie z funkcją na generatorach lub homomorfizmem monoidów zachowuje tę własność przedłużenia, co dowodzi naturalności bijekcji. Mamy sprzężenie $F\dashv U$, a złożenie $T=UF$ jest dokładnie funktorem list z poprzedniego przykładu. Jednostka sprzężenia jest operacją $x\mapsto[x]$, a jej kojednostka indukuje spłaszczanie list list.
\end{przyklad}


\begin{przyklad}[Monada podzbiorów]\label{przyklad:powerset-monad}
Na kategorii $\Set$ weźmy kowariantny funktor $T=\mathcal P_+$, który funkcji $f:X\to Y$ przypisuje obraz bezpośredni $A\mapsto f(A)$. Definiujemy $\eta_X(x)=\{x\}$ i $\mu_X:\mathcal P(\mathcal P(X))\to\mathcal P(X)$ przez $\mu_X(\mathscr A)=\bigcup_{A\in\mathscr A}A$. Symbol $\mathcal P_+$ oznacza tutaj funktor kowariantny (nie ograniczamy się do podzbiorów niepustych). Naturalność $\eta$ wynika z $f(\{x\})=\{f(x)\}$. Naturalność $\mu$ wynika z tożsamości $f(\bigcup_{A\in\mathscr A}A)=\bigcup_{A\in\mathscr A}f(A)$: obraz sumy jest sumą obrazów. Dla $\mathfrak A\subseteq\mathcal P(\mathcal P(X))$ element $x$ należy do obu stron aksjomatu łączności wtedy i tylko wtedy, gdy istnieją $\mathscr A\in\mathfrak A$ oraz $A\in\mathscr A$ takie, że $x\in A$. Dowodzi to łączności również dla pustych rodzin; aksjomaty jednostki wynikają z $\bigcup\{A\}=A$ i $\bigcup_{x\in A}\{x\}=A$. Dla $X=\{1,2,3\}$ i $\mathscr A=\{\{1,2\},\{2,3\},\varnothing\}$ dostajemy $\mu_X(\mathscr A)=\{1,2,3\}$. Jest to inny przykład monady niż listy: \emph{kolejność i powtórzenia} nie mają tu znaczenia.
\end{przyklad}

\begin{przyklad}[Monada ``możliwego braku wyniku'']
Dla zbioru $X$ zdefiniujmy $T(X)=X\sqcup\{\bot\}$; symbol $\bot$ oznacza nowy element, który nie należy do $X$ (formalnie stosujemy oznakowaną sumę rozłączną, więc kolejne symbole braku wyniku są rozróżnialne). Dla funkcji $f:X\to Y$ określamy $T(f)(x)=f(x)$, gdy $x\in X$, oraz $T(f)(\bot)=\bot$. Jednostka $\eta_X:X\to T(X)$ jest włączeniem. Zbiór $T(T(X))$ zawiera zwykłe wyniki oraz dwa rodzaje symbolu braku wyniku (wewnętrzny i zewnętrzny); mnożenie $\mu_X:T(T(X))\to T(X)$ zachowuje zwykłe wyniki i wysyła oba symbole na $\bot$. Spłaszczenie trzech warstw zawsze daje pierwotny wynik, jeśli taki istnieje, albo jeden symbol $\bot$, niezależnie od kolejności spłaszczania. Oba prawa jednostki zachodzą, ponieważ każda z dwóch map $T(\eta_X)$ i $\eta_{T(X)}$ dodaje jedynie warstwę oznakowania, którą $\mu_X$ usuwa, a istniejący brak wyniku pozostaje brakiem wyniku. Funktorialność i naturalność obu transformacji sprawdzamy na zwykłych wynikach i symbolach braku wyniku: na zwykłym wyniku obie strony odpowiedniego równania stosują tę samą funkcję, a symbol braku wyniku pozostaje symbolem braku wyniku. Otrzymujemy monadę przydatną przy opisie obliczeń częściowych.
\end{przyklad}

\begin{przyklad}[Konkatenacja list i składanie Kleisliego]
Niech $X=\{1,2\}$, $Y=\{a,b\}$, $Z=\{0,1\}$. Przyjmijmy funkcje $f:X\to T(Y)$ i $g:Y\to T(Z)$ dane przez
$f(1)=[a,b]$, $f(2)=[b]$, $g(a)=[0]$, $g(b)=[1,1]$.
Ich składanie w kategorii Kleisliego (definicja~\ref{cat:def:kleisli}) monady list ma wzór $(g\star f)(x)=\mu_Z(T(g)(f(x)))$. Dla $x=1$ otrzymujemy kolejno $f(1)=[a,b]$, następnie $T(g)([a,b])=[[0],[1,1]]$, a na końcu $(g\star f)(1)=[0,1,1]$. Dla $x=2$ wynikiem jest $[1,1]$. Nie jest to zwykłe składanie funkcji: wartością pośrednią jest lista, a nie pojedynczy element $Y$.
\end{przyklad}


\begin{twierdzenie}[Monada z każdego sprzężenia]\label{cat:thm:adjunction-monad}
Niech $F:\mathcal C\to\mathcal D$ i $G:\mathcal D\to\mathcal C$ spełniają $F\dashv G$, z jednostką $\eta:\id_{\mathcal C}\Rightarrow GF$ i kojednostką $\varepsilon:FG\Rightarrow\id_{\mathcal D}$. Wówczas $T=GF$ z jednostką $\eta$ i mnożeniem $\mu_A:=G(\varepsilon_{F(A)}):GFGF(A)\to GF(A)$ jest monadą na $\mathcal C$.
\end{twierdzenie}
\begin{proof}
Dla $f:A\to B$ naturalność kojednostki (twierdzenie~\ref{tw:triangles}) względem $F(f)$ daje $F(f)\circ\varepsilon_{F(A)}=\varepsilon_{F(B)}\circ FG(F(f))$. Stosując $G$, otrzymujemy
\[
T(f)\circ\mu_A=\mu_B\circ T^2(f).
\]
Naturalność $\eta$ jest częścią danych sprzężenia. Najpierw dwie tożsamości jednostkowe z definicji~\ref{def:monad} wynikają z tożsamości trójkątnych:
$\mu_A\circ T(\eta_A)=G(\varepsilon_{F(A)}\circ F(\eta_A))=\id_{GF(A)}$ oraz
$\mu_A\circ\eta_{T(A)}=G(\varepsilon_{F(A)})\circ\eta_{GF(A)}=\id_{GF(A)}$.

Pozostaje łączność. Naturalność kojednostki $\varepsilon$ dla morfizmu $\varepsilon_{F(A)}:FGF(A)\to F(A)$ daje
$\varepsilon_{F(A)}\circ FG(\varepsilon_{F(A)})=\varepsilon_{F(A)}\circ\varepsilon_{FGF(A)}$. Po zastosowaniu $G$ lewa strona jest równa $\mu_A\circ T(\mu_A)$, a prawa $\mu_A\circ\mu_{T(A)}$. Otrzymujemy trzeci aksjomat monady.
\end{proof}

\subsection{Kategoria Kleisliego}

\begin{definicja}[Kategoria Kleisliego]\label{cat:def:kleisli}
Dla monady $(T,\eta,\mu)$ na $\mathcal C$ kategorię Kleisliego $\mathcal C_T$ określamy następująco. Obiekty są obiektami $\mathcal C$, a morfizmy $A\to B$ to morfizmy $f:A\to T(B)$ w $\mathcal C$. Składanie morfizmów $f:A\to T(B)$ i $g:B\to T(C)$ jest dane wzorem
\[
g\star f=\mu_C\circ T(g)\circ f:A\to T(C),
\]
a identycznością obiektu $A$ jest $\eta_A$.
\end{definicja}

\begin{stwierdzenie}
Powyższe dane rzeczywiście określają kategorię.
\end{stwierdzenie}
\begin{proof}
Dla $f:A\to T(B)$ naturalność jednostki i prawa jednostki monady dają
\[
\eta_B\star f=\mu_B\circ T(\eta_B)\circ f=f,
\qquad
f\star\eta_A=\mu_B\circ\eta_{T(B)}\circ f=f.
\]
Niech $h:C\to T(D)$ będzie trzecim morfizmem. Rozpisując składanie i stosując naturalność $\mu$ dla $h:C\to T(D)$, dostajemy
\[
\begin{aligned}
h\star(g\star f)
&=\mu_D\circ T(h)\circ\mu_C\circ T(g)\circ f\\
&=\mu_D\circ\mu_{T(D)}\circ T^2(h)\circ T(g)\circ f\\
&=\mu_D\circ T(\mu_D)\circ T^2(h)\circ T(g)\circ f\\
&=\mu_D\circ T(\mu_D\circ T(h)\circ g)\circ f\\
&=(h\star g)\star f.
\end{aligned}
\]
Trzecia równość jest aksjomatem łączności monady, czwarta wynika z funktorialności $T$.
\end{proof}

\begin{stwierdzenie}[Sprzężenie Kleisliego]\label{cat:prop:kleisli-adj}
Jeżeli $\mathcal C$ jest lokalnie mała, to istnieje sprzężenie $J:\mathcal C\rightleftarrows\mathcal C_T:K$, którego monadą jest $(T,\eta,\mu)$. Lokalna małość $\mathcal C_T$ wynika z jej definicji.
\end{stwierdzenie}
\begin{proof}
Definiujemy $J(A)=A$, $J(f)=\eta_B\circ f$ dla $f:A\to B$, oraz $K(A)=T(A)$, $K(g)=\mu_B\circ T(g)$ dla $g:A\to T(B)$. Funktor $J$ zachowuje identyczności z definicji, a dla $f:A\to B$ i $g:B\to C$ mamy
\[
J(g)\star J(f)
=\mu_C\circ T(\eta_C)\circ T(g)\circ\eta_B\circ f
=\eta_C\circ g\circ f=J(g\circ f).
\]
Użyliśmy prawa jednostki i naturalności $\eta$. Funktor $K$ zachowuje identyczności, gdyż $\mu_A\circ T(\eta_A)=\id_{T(A)}$. Dla morfizmów Kleisliego $f:A\to T(B)$, $g:B\to T(C)$ mamy
\[
\begin{aligned}
K(g\star f)&=\mu_C\circ T(\mu_C)\circ T^2(g)\circ T(f)\\
&=\mu_C\circ\mu_{T(C)}\circ T^2(g)\circ T(f)\\
&=\mu_C\circ T(g)\circ\mu_B\circ T(f)=K(g)\circ K(f).
\end{aligned}
\]
Kolejne kroki wykorzystują łączność monady i naturalność $\mu$.

Z definicji $\Hom_{\mathcal C_T}(J(A),B)=\Hom_{\mathcal C}(A,K(B))$. Ta identyczność zbiorów jest naturalna: dla $u:A'\to A$ i morfizmu $v:B\to T(B')$ odpowiednie operacje na $g:A\to T(B)$ są $g\star J(u)=g\circ u$ (naturalność $\eta$ i prawo jednostki) oraz $v\star g=K(v)\circ g$. Jednostką sprzężenia jest więc $\eta_A$. Kojednostka w $B$ jest reprezentowana przez $\id_{T(B)}:T(B)\to T(B)$, wobec czego zastosowanie $K$ daje $\mu_B$. Ponadto $KJ(f)=\mu_B\circ T(\eta_B)\circ T(f)=T(f)$, a na obiektach $KJ=T$. Zatem indukowana monada jest dokładnie monadą wyjściową.
\end{proof}

\subsection{Kategoria Eilenberga--Moore'a}

\begin{definicja}[Algebra monady]\label{cat:def:em-algebra}
Algebrą monady $(T,\eta,\mu)$ jest para $(A,a)$, gdzie $a:T(A)\to A$ spełnia
\[
a\circ\eta_A=\id_A,\qquad a\circ T(a)=a\circ\mu_A.
\]
Morfizm algebr $f:(A,a)\to(B,b)$ to morfizm $f:A\to B$ taki, że $f\circ a=b\circ T(f)$. Kategorię tych algebr oznaczamy $\mathcal C^T$.
\end{definicja}
Rzeczywiście identyczność jest morfizmem algebr, a dla dwóch morfizmów $f:(A,a)\to(B,b)$ i $g:(B,b)\to(C,c)$ mamy
\[
(g\circ f)\circ a=g\circ b\circ T(f)
=c\circ T(g)\circ T(f)=c\circ T(g\circ f).
\]
Składanie i identyczności dziedziczą zatem aksjomaty kategorii $\mathcal C$.

\begin{stwierdzenie}[Sprzężenie Eilenberga--Moore'a]\label{cat:prop:em-adj}
Jeżeli $\mathcal C$ jest lokalnie mała, to również $\mathcal C^T$ jest lokalnie mała, ponieważ jej zbiory morfizmów są podzbiorami zbiorów morfizmów w $\mathcal C$. Funktor zapominający $U_T:\mathcal C^T\to\mathcal C$ ma lewe sprzężenie $F_T$, określone przez $F_T(A)=(T(A),\mu_A)$ i $F_T(f)=T(f)$. Indukowana monada jest równa $(T,\eta,\mu)$.
\end{stwierdzenie}
\begin{proof}
Aksjomaty algebry $(T(A),\mu_A)$ to prawo jednostki i łączność monady. Naturalność $\mu$ daje warunek morfizmu algebr dla $T(f)$, a funktorialność wynika z funktorialności $T$.

Dla $f:A\to B$ i algebry $(B,b)$ określamy $\widehat f=b\circ T(f):T(A)\to B$. Jest to morfizm algebr, gdyż
\[
\begin{aligned}
\widehat f\circ\mu_A
&=b\circ\mu_B\circ T^2(f)\\
&=b\circ T(b)\circ T^2(f)=b\circ T(\widehat f).
\end{aligned}
\]
Naturalność $\eta$ i prawo algebry dają $\widehat f\circ\eta_A=b\circ\eta_B\circ f=f$. Jeśli $h:(T(A),\mu_A)\to(B,b)$ jest morfizmem algebr, to
\[
b\circ T(h\circ\eta_A)=h\circ\mu_A\circ T(\eta_A)=h.
\]
Otrzymujemy więc wzajemnie odwrotne bijekcje $h\mapsto h\circ\eta_A$ i $f\mapsto\widehat f$. Naturalność pierwszej względem $u:A'\to A$ wynika z $h\circ T(u)\circ\eta_{A'}=h\circ\eta_A\circ u$, a względem morfizmu algebr $v:(B,b)\to(B',b')$ z łączności składania $(v\circ h)\circ\eta_A=v\circ(h\circ\eta_A)$. Jednostką sprzężenia jest $\eta$, a kojednostka w $(B,b)$ ma morfizm bazowy $b$ (odpowiada $\id_B$). Ponieważ $U_TF_T=T$ oraz $U_T(\varepsilon_{F_T(A)})=\mu_A$, indukowana monada jest wyjściową.
\end{proof}

\begin{przyklad}[Algebry monady list]\label{cat:ex:list-algebras}
Algebry monady list w $\Set$ odpowiadają dokładnie monoidom. Jeśli $a:T(M)\to M$ jest algebrą, połóżmy $e=a([])$ i $x\cdot y=a([x,y])$. Z $a\circ\eta_M=\id_M$ mamy $a([x])=x$. Aksjomat $a\circ T(a)=a\circ\mu_M$ dla $[[],[x]]$ i $[[x],[]]$ daje $e\cdot x=x=x\cdot e$. Dla $[[x,y],[z]]$ i $[[x],[y,z]]$ daje $(x\cdot y)\cdot z=a([x,y,z])=x\cdot(y\cdot z)$. Dla listy niepustej, stosując aksjomat do $[[x_1,\ldots,x_{n-1}],[x_n]]$, otrzymujemy $a([x_1,\ldots,x_n])=a([x_1,\ldots,x_{n-1}])\cdot x_n$. Indukcja daje iloczyn wszystkich elementów, włącznie z pustym iloczynem $e$.

Odwrotnie, w monoidzie funkcja mnożąca elementy listy w ich kolejności spełnia oba aksjomaty algebry: iloczyn listy jednoelementowej daje jej element, a iloczyn iloczynów kolejnych bloków jest iloczynem ich konkatenacji z łączności i praw jednostki monoidu. Morfizm algebr zachowuje $e$ i mnożenie przez zastosowanie warunku morfizmu do list pustych i dwuelementowych. Odwrotnie, homomorfizm monoidów zachowuje iloczyny wszystkich skończonych list przez indukcję, więc jest morfizmem algebr. Obie konstrukcje są wzajemnie odwrotne na obiektach i morfizmach.
\end{przyklad}

\section{Powiązania z innymi działami matematyki}\label{sec:powiazania}

\subsection{Algebra: od jąder do kategorii abelowych}

Kategoria $\Vect_k$ ma obiekt zerowy, jądra i kojądra wszystkich morfizmów, a każdy zbiór morfizmów $\Hom(V,W)$ jest grupą abelową względem dodawania odwzorowań liniowych. Ponadto składanie jest addytywne w każdej zmiennej:
$(f_1+f_2)\circ g=f_1\circ g+f_2\circ g$ oraz $f\circ(g_1+g_2)=f\circ g_1+f\circ g_2$.

\begin{definicja}[Kategorie addytywne i abelowe]\label{cat:def:abelian-category}
Kategoria preaddytywna to kategoria lokalnie mała, której zbiory morfizmów są grupami abelowymi, a składanie jest addytywne w obu zmiennych. Kategoria addytywna jest kategorią preaddytywną z obiektem zerowym i skończonymi biproduktami (obiektami będącymi jednocześnie iloczynami i koproduktami, z kanonicznymi projekcjami i włączeniami). Dokładniej dla biproduktu $P=A\oplus B$ wymagamy, aby projekcje $p_A,p_B$ i włączenia $i_A,i_B$ spełniały
\[
p_Ai_A=\id_A,\quad p_Bi_B=\id_B,\quad
p_Ai_B=0,\quad p_Bi_A=0,\quad
i_Ap_A+i_Bp_B=\id_P.
\]
Równości te zapewniają obie własności uniwersalne: morfizm do $P$ o składowych $u,v$ to $i_Au+i_Bv$, a morfizm z $P$ o ograniczeniach $r,s$ to $rp_A+sp_B$; jednoznaczność w obu przypadkach wynika z ostatniej równości. Kategoria abelowa jest kategorią addytywną, w której każdy morfizm ma jądro i kojądro, każdy monomorfizm jest jądrem pewnego morfizmu, a każdy epimorfizm jest kojądrem pewnego morfizmu.
\end{definicja}
Jądro $f:A\to B$ jest equalizatorem $f$ i morfizmu zerowego $A\to B$, a kojądro jest ich koequalizatorem. Obraz definiujemy jako jądro kojądra, traktowane jako podobiekt przeciwdziedziny. Ciąg $A\xrightarrow{f}B\xrightarrow{g}C$ w kategorii abelowej jest dokładny w $B$, gdy obraz $f$ i jądro $g$ są tym samym podobiektem $B$ (tj. ich morfizmy do $B$ są zgodne poprzez izomorfizm dziedzin). W $\Vect_k$ i $\Ab$ są to odpowiednio zwykłe jądro, rzut na iloraz przez obraz oraz włączenie obrazu.

\begin{stwierdzenie}\label{cat:prop:abelian-examples}
Kategorie $\Vect_k$ i $\Ab$ są abelowe.
\end{stwierdzenie}
\begin{proof}
W obu kategoriach dodawanie morfizmów jest punktowe i składanie jest dwuliniowe względem grup abelowych morfizmów. Obiekt zerowy to przestrzeń lub grupa trywialna. Dla dwóch obiektów $A,B$ suma prosta $A\oplus B$ jest iloczynem: morfizmy $u:X\to A$, $v:X\to B$ wyznaczają jedyny $x\mapsto(u(x),v(x))$. Jest też koproduktem: morfizmy $r:A\to Y$, $s:B\to Y$ wyznaczają jedyny $(a,b)\mapsto r(a)+s(b)$, ponieważ $(a,b)=(a,0)+(0,b)$. Projekcje i włączenia współrzędnych spełniają pięć podanych równości, co sprawdzamy na dowolnej parze $(a,b)$. Obiekt zerowy obsługuje pusty biprodukt, a iteracja konstrukcji daje wszystkie skończone biprodukty.

Jądro $f$ jest włączeniem $\Ker f\to A$: morfizm $h$ spełniający $fh=0$ ma wartości w $\Ker f$, więc faktoryzuje się jednoznacznie przez to włączenie. Kojądro jest rzutem $q:B\to B/\im f$: dla $g:B\to Z$ z $gf=0$ wzór $\bar g(b+\im f)=g(b)$ jest dobrze określony, bo $g$ zeruje się na $\im f$; daje jedyną faktoryzację przez surjekcję $q$.

Monomorfizmy w obu kategoriach są dokładnie morfizmami injektywnymi: injekcja daje skracanie z lewej, a przy niezerowym elemencie jądra dwa różne morfizmy z $k$ (odpowiednio $\mathbb Z$), wysyłające generator na ten element i na zero, mają to samo złożenie. Epimorfizmy są dokładnie surjekcjami: surjekcja daje skracanie z prawej; gdy obraz jest właściwy, rzut na niezerowy iloraz przez obraz i morfizm zerowy są różne, lecz ich złożenia z $f$ są równe. Monomorfizm identyfikuje dziedzinę izomorficznie z obrazem, więc jest jądrem rzutu na iloraz przez obraz. Dla epimorfizmu $f$ odwzorowanie $A/\Ker f\to B$, $a+\Ker f\mapsto f(a)$, jest dobrze określoną bijekcją liniową lub izomorfizmem grup: injektywność wynika z definicji jądra, surjektywność z surjektywności $f$. Zatem $f$ jest kojądrem włączenia jądra. Spełniono wszystkie warunki definicji.
\end{proof}

\begin{przyklad}[Ciąg dokładny związany z morfizmem]\label{cat:ex:exact-sequence}
Dla liniowego $f:V\to W$ mamy ciąg $0\to\Ker f\xrightarrow{i}V\xrightarrow{f}W\xrightarrow{q}W/\im f\to0$. Jest dokładny: $\Ker i=0$, $\im i=\Ker f$, $\Ker q=\im f$, a $q$ jest surjektywny. W tej postaci nie jest to krótki ciąg dokładny (krótki ciąg ma trzy obiekty między skrajnymi zerami). Dwa związane z nim krótkie ciągi dokładne to
\[
0\to\Ker f\to V\to\im f\to0,\qquad
0\to\im f\to W\to W/\im f\to0.
\]
W pierwszym morfizm $V\to\im f$ jest zawężeniem przeciwdziedziny $f$, w drugim $\im f\to W$ jest włączeniem. Dokładność wynika z tych samych równości jąder i obrazów. Por. skrypt \emph{Topologia algebraiczna}, rozdział o algebrze homologicznej.
\end{przyklad}

\subsection{Topologia: funktorialność niezmienników}

Z punktu widzenia teorii kategorii wielu niezmienników topologicznych nie należy traktować jako pojedynczych liczb lub grup, lecz jako funktory. Przykładowo grupa podstawowa przyporządkowuje \emph{punktowanej} przestrzeni topologicznej $(X,x_0)$ grupę $\pi_1(X,x_0)$. Każde ciągłe odwzorowanie punktowane $f:(X,x_0)\to(Y,y_0)$ indukuje homomorfizm $\pi_1(f):\pi_1(X,x_0)\to\pi_1(Y,y_0)$ przez $[\gamma]\mapsto[f\circ\gamma]$. Jeżeli pętle są homotopijne przy ustalonych końcach, złożenie homotopii z $f$ zachowuje tę własność, więc przyporządkowanie klas jest dobrze określone. Złożenie $f$ z konkatenacją pętli jest konkatenacją ich obrazów, stąd homomorfizm. Identyczności i składanie zachowuje, ponieważ $g\circ(f\circ\gamma)=(g\circ f)\circ\gamma$. Jest więc funktorem z kategorii punktowanych przestrzeni topologicznych do $\Grp$.

\begin{przyklad}[Homologia jako funktor]
Dla każdego $n\geq0$ homologia singularna o współczynnikach całkowitych jest funktorem $H_n:\Top\to\Ab$. Na wolnych grupach abelowych $C_n(X)$ generowanych przez sympleksy $\sigma:\Delta^n\to X$ definiujemy $C_n(f)(\sigma)=f\circ\sigma$ i rozszerzamy addytywnie. Dla $n\geq1$ i włączeń ścian $\delta_i:\Delta^{n-1}\to\Delta^n$ mamy
\[
\partial C_n(f)(\sigma)=\sum_i(-1)^i f\circ\sigma\circ\delta_i
=C_{n-1}(f)(\partial\sigma).
\]
Dla $n=0$ różniczka jest zerowa. Zatem cykle przechodzą na cykle, a brzegi na brzegi, więc wzór $H_n(f)([c])=[C_n(f)(c)]$ jest dobrze określony: zmiana reprezentanta o brzeg zmienia obraz o brzeg. Na generatorach zachodzą $C_n(\id_X)=\id_{C_n(X)}$ i $C_n(g\circ f)=C_n(g)C_n(f)$, zatem te same równości zachodzą na ilorazach homologii. Por. skrypt \emph{Topologia algebraiczna}, rozdział o homologii przestrzeni. To dlatego homeomorficzne przestrzenie mają izomorficzne grupy homologii: stosujemy ogólne twierdzenie~\ref{tw:funktory-iso} o zachowywaniu izomorfizmów przez funktory.
\end{przyklad}

\subsection{Analiza: przestrzenie funkcji jako funktory Hom}

Jeśli rozważamy kategorię odpowiednio dobranych przestrzeni analitycznych i dozwolonych odwzorowań, to zapis $\Hom(X,Y)$ oznacza \emph{przestrzeń odwzorowań zgodnych ze strukturą}, np. ciągłych lub holomorficznych. Własności uniwersalne mogą wtedy opisywać obiekty nie przez ich punkty, lecz przez wszystkie rodziny odwzorowań do nich i z nich. Zapis $\Hom(X,Y)$ sam w sobie oznacza zbiór morfizmów; dodatkowa topologia lub struktura liniowa wymaga osobnego określenia. To jest podstawowy punkt widzenia geometrii algebraicznej, teorii schematów i geometrii analitycznej.

\begin{przyklad}[Iloraz przestrzeni unormowanej jako konstrukcja uniwersalna]\label{cat:ex:normed-quotient}
Niech $X$ będzie przestrzenią unormowaną nad $k=\mathbb R$ lub $\mathbb C$, a $M$ jej domkniętą podprzestrzenią liniową. W kategorii przestrzeni unormowanych i ograniczonych odwzorowań liniowych rzut $q:X\to X/M$ z normą ilorazową $\|x+M\|=\inf_{m\in M}\|x+m\|$ jest kojądrem włączenia $M\to X$ (względem morfizmu zerowego). Istotnie, dla ograniczonego $f:X\to Y$ zerującego się na $M$ wzór $\bar f(x+M)=f(x)$ jest dobrze określony i liniowy. Dla każdego $m\in M$ mamy $\|\bar f(x+M)\|=\|f(x+m)\|\leq\|f\|\|x+m\|$, więc po wzięciu infimum $\|\bar f\|\leq\|f\|$. Ponieważ $f=\bar f q$ i $\|q\|\leq1$, również $\|f\|\leq\|\bar f\|$. Jednoznaczność wynika z surjektywności $q$. Por. skrypt \emph{Analiza funkcjonalna}, rozdział o konstrukcjach ilorazowych.
\end{przyklad}

\begin{przyklad}[Kategoria przestrzeni unormowanych nie jest abelowa]\label{cat:ex:normed-nonabelian}
Niech $c_{00}$ oznacza przestrzeń ciągów o skończenie wielu niezerowych wyrazach, z normą $\ell^2$. Włączenie $j:c_{00}\to\ell^2$ jest ograniczone i injektywne, więc jest monomorfizmem. Jest także epimorfizmem: jeśli ograniczone odwzorowania liniowe $g,h:\ell^2\to Y$ spełniają $gj=hj$, to dla każdego $x\in\ell^2$ jego obcięcia $x^{(N)}\in c_{00}$ spełniają
\[
\|x-x^{(N)}\|_2^2=\sum_{n>N}|x_n|^2\longrightarrow0.
\]
Z ciągłości $g,h$ otrzymujemy $g(x)=\lim_N g(x^{(N)})=\lim_N h(x^{(N)})=h(x)$. Jednak $j$ nie jest surjekcją: ciąg $(2^{-n})_{n\geq1}$ należy do $\ell^2$ i nie należy do $c_{00}$. Nie jest zatem izomorfizmem.

Gdyby ta kategoria była abelowa, $j$ jako monomorfizm byłby jądrem pewnego morfizmu $r$. Mielibyśmy $rj=0$, więc z epiczności $j$ wynikałoby $r=0$. Jądro morfizmu zerowego o dziedzinie $\ell^2$ jest identycznością $\ell^2$; z jednoznaczności jąder z dokładnością do izomorfizmu wynikałoby, że $j$ jest izomorfizmem. Sprzeczność. Przykład pokazuje, dlaczego algebraiczne własności jąder i kojąder nie przenoszą się automatycznie na kategorie z warunkiem ciągłości.
\end{przyklad}

\subsection{Geometria: presnopy i snopy}

\begin{przyklad}[Presnop funkcji ciągłych]\label{cat:ex:continuous-sheaf}
Dla przestrzeni topologicznej $X$ kategoria $\mathcal O(X)$ ma zbiory otwarte jako obiekty i inkluzje jako morfizmy. Presnop zbiorów jest funktorem $\mathcal O(X)^{\op}\to\Set$. Przykładem jest $\mathcal F(U)=C(U,\mathbb R)$, z ograniczeniami $r_{U,V}:\mathcal F(U)\to\mathcal F(V)$ dla $V\subseteq U$. Zachodzą $r_{U,U}=\id$ i $r_{V,W}r_{U,V}=r_{U,W}$ dla $W\subseteq V\subseteq U$.

Presnop jest snopem, jeśli dla każdego pokrycia otwartego $U=\bigcup_i U_i$ rodzina $s_i\in\mathcal F(U_i)$ zgodna po ograniczeniu do $U_i\cap U_j$ ma dokładnie jedno sklejenie w $\mathcal F(U)$. Dla funkcji ciągłych definiujemy $s(x)=s_i(x)$, gdy $x\in U_i$; zgodność daje niezależność od wyboru $i$. Dla otwartego $O\subseteq\mathbb R$ mamy $s^{-1}(O)=\bigcup_i s_i^{-1}(O)$, otwarty w $U$, bo każdy $U_i$ jest otwarty w $U$. Zatem $s$ jest ciągła. Jednoznaczność wynika z pokrycia, również dla $U=\varnothing$ (istnieje dokładnie jedna funkcja pusta). Tak warunek lokalnego sklejania uzupełnia samą funktorialność ograniczeń. Por. skrypt \emph{Geometria algebraiczna}, rozdział o presnopach i snopach.
\end{przyklad}

\begin{uwaga}[Dalsza ścieżka]
Po opanowaniu podstaw warto kolejno poznać kategorie addytywne i abelowe, kategorie monoidalne, sprzężenia i monady w większej ogólności, algebry nad monadą, a następnie kategorie pochodne lub kategorie wyższe. Nie są one wymagane do zrozumienia definicji funktora, transformacji naturalnej czy lematu Yonedy, lecz rozwijają te same idee.
\end{uwaga}


\section{Przykłady przekrojowe: jedno zadanie, wiele kategorii}\label{sec:przekrojowe}

Zestawmy teraz w jednym miejscu kilka konstrukcji, żeby wyraźniej zobaczyć, co jest wspólne, a co zależy od kategorii.

\begin{przyklad}[Co jest iloczynem dwóch obiektów?]\label{cat:ex:cross-products}
\leavevmode\par\noindent
\begin{itemize}
\item W $\Set$: $A\times B$ i zwykłe rzuty.
\item W $\Grp$: $G\times H$ z mnożeniem współrzędnym.
\item W $\Top$: $X\times Y$ z topologią produktową.
\item W $\Vect_k$: $V\oplus W$ (dla dwóch czynników).
\item W posetach: $a\wedge b$ (kres dolny), jeśli istnieje.
\end{itemize}
We wszystkich przypadkach testem jest \emph{ta sama} własność uniwersalna z definicji~\ref{def:produkt}. Dla morfizmów $f:Z\to A$, $g:Z\to B$ do czynników w $\Set$ jedynym morfizmem o zadanych składowych jest $z\mapsto(f(z),g(z))$. W $\Grp$ jest on homomorfizmem, bo mnożenie jest współrzędne; w $\Vect_k$ jest liniowy z tego samego powodu dla dodawania i mnożenia przez skalar. W $\Top$ przeciwobraz bazowego zbioru otwartego $U\times V$ wynosi $f^{-1}(U)\cap g^{-1}(V)$, więc odwzorowanie jest ciągłe. Jednoznaczność we wszystkich czterech przypadkach wynika z jednoznaczności współrzędnych. W kategorii zbioru częściowo uporządkowanego warunek istnienia morfizmów $z\to a$ i $z\to b$ znaczy $z\leq a$ i $z\leq b$; istnienie jedynego morfizmu $z\to a\wedge b$ jest dokładnie własnością największego ograniczenia dolnego.
\end{przyklad}

\begin{przyklad}[Co jest koproduktem dwóch obiektów?]\label{cat:ex:cross-coproducts}
\leavevmode\par\noindent
\begin{itemize}
\item W $\Set$: suma rozłączna $A\sqcup B$.
\item W $\Grp$: iloczyn wolny $G*H$.
\item W $\Top$: suma rozłączna z topologią sumy topologicznej.
\item W $\Vect_k$: $V\oplus W$.
\item W posetach: $a\vee b$ (kres górny), jeśli istnieje.
\end{itemize}
Dla morfizmów $f:A\to Z$, $g:B\to Z$ w $\Set$ określamy jedyny morfizm na sumie rozłącznej, stosując $f$ na pierwszym składniku i $g$ na drugim. W $\Top$ przeciwobraz zbioru otwartego jest otwarty na obu składnikach, czyli otwarty w topologii sumy. W $\Vect_k$ jedynym liniowym morfizmem o zadanych ograniczeniach jest $(a,b)\mapsto f(a)+g(b)$, ponieważ $(a,b)=(a,0)+(0,b)$. Dla posetu istnienie morfizmów $a\to z$ i $b\to z$ jest równoważne $a\vee b\leq z$; jednoznaczność wynika z definicji kategorii posetu.

W $\Grp$ dla iloczynu wolnego $G*H$ przedłużenie homomorfizmów $f:G\to K$, $g:H\to K$ otrzymujemy, zastępując kolejne litery z $G$ lub $H$ ich obrazami i mnożąc w $K$. Redukcje słów (usuwanie jednostek i mnożenie sąsiednich liter z tego samego czynnika) zachowują wartość, więc ewaluacja jest homomorfizmem; jednoznaczność wynika z generowania przez oba czynniki. Dla porównania homomorfizm $G\times H\to K$ o tych samych ograniczeniach do czynników istnieje wtedy i tylko wtedy, gdy $f(x)g(y)=g(y)f(x)$ dla wszystkich $x\in G$, $y\in H$. Konieczność wynika z przemienności $(x,e)$ i $(e,y)$ w iloczynie; przy tym warunku wzór $(x,y)\mapsto f(x)g(y)$ zachowuje mnożenie i jest jedynym możliwym przedłużeniem. Zatem iloczyn i koprodukt nie mają tej samej własności w $\Grp$.
\end{przyklad}

\begin{przyklad}[Co oznacza strzałka?]
W $\Set$ strzałką jest funkcja; w $\Grp$ homomorfizm; w $\Rel$ relacja; w kategorii zbioru częściowo uporządkowanego jedyny morfizm $x\to y$ istnieje dokładnie wtedy, gdy $x\leq y$ (nie rozróżniamy dowodów nierówności); w kategorii monoidu z jednym obiektem strzałką jest element monoidu. Zatem nie ma sensu pytać, czy morfizm kategorii jest \emph{zawsze ciągły} lub \emph{zawsze różnowartościowy}: właściwe pytanie zależy od wybranej kategorii.
\end{przyklad}

\begin{przyklad}[Rozwiązywanie problemu przez własność uniwersalną]
Załóżmy, że chcemy określić homomorfizm $h:G*H\to K$. W podejściu elementarnym musielibyśmy opisać, jak $h$ działa na zredukowanych słowach, i sprawdzić zgodność z mnożeniem. W podejściu kategorialnym wystarczy podać dwa homomorfizmy $f:G\to K$ i $g:H\to K$: własność koproduktu gwarantuje, że istnieje dokładnie jeden $h$ z $h\circ i_G=f$, $h\circ i_H=g$. To jest praktyczna oszczędność pracy, a nie jedynie zmiana notacji.
\end{przyklad}

\begin{przyklad}[Funktor może ujawniać strukturę ukrytą]
Każda grupa $G$ działa na sobie przez lewe mnożenie. Traktując $G$ jako kategorię z jednym obiektem, ze składaniem $g\circ h=gh$, otrzymujemy funktor do $\Set$, który przypisuje jedynemu obiektowi zbiór $G$, a morfizmowi $g\in G$ permutację $L_g(x)=gx$. Mamy $L_e=\id_G$ i $L_g\circ L_h=L_{gh}$, a odwrotnością $L_g$ jest $L_{g^{-1}}$. Ta obserwacja prowadzi do twierdzenia Cayleya: odwzorowanie $g\mapsto L_g$ jest homomorfizmem $G\to\operatorname{Sym}(G)$, a jest różnowartościowe, ponieważ $L_g(e)=g$.
\end{przyklad}


\begin{przyklad}[Dlaczego własność uniwersalna daje jednoznaczność konstrukcji]
Załóżmy, że obiekt $P$ z rzutami $p_A,p_B$ ma uniwersalną własność iloczynu, podobnie jak $Q$ z rzutami $q_A,q_B$. Nie wybieramy żadnych reprezentantów elementów. Uniwersalność $Q$ daje jedyny $u:P\to Q$ z $q_Au=p_A$, $q_Bu=p_B$, a uniwersalność $P$ daje jedyny $v:Q\to P$ z $p_Av=q_A$, $p_Bv=q_B$. Wtedy $p_Avu=q_Au=p_A$, $p_Bvu=q_Bu=p_B$, więc $vu=\id_P$ z jednoznaczności morfizmu do $P$. Podobnie $q_Auv=p_Av=q_A$ i $q_Buv=p_Bv=q_B$, stąd $uv=\id_Q$. Izomorfizm $u$ jest jedyny wśród morfizmów zachowujących rzuty. Ten dowód jest typowym wzorcem: \emph{aby porównać dwa obiekty uniwersalne, zastosuj ich uniwersalności do siebie nawzajem}.
\end{przyklad}

\begin{przyklad}[Operacje na zbiorach a operacje na przestrzeniach liniowych]
Dla zbioru $X$ nie ma naturalnego działania $X\times X\to X$ wyznaczonego samą strukturą zbioru poza wyborem pierwszej albo drugiej współrzędnej. W istocie lemat Yonedy z przykładu~\ref{przyklad:yoneda-tuple} daje dokładnie dwa naturalne odwzorowania $X^2\to X$: $(x,y)\mapsto x$ i $(x,y)\mapsto y$. Natomiast dla przestrzeni wektorowej $V$ istnieje naturalne dodawanie $V\oplus V\to V$, $(v,w)\mapsto v+w$, ponieważ każdy liniowy $f:V\to W$ spełnia $f(v+w)=f(v)+f(w)$. \emph{Naturalność zawsze zależy od kategorii i jej morfizmów}, a nie tylko od wzoru na elementach.
\end{przyklad}


\begin{przyklad}[Iloczyn i koprodukt w kategorii relacji]\label{cat:ex:rel-biproduct}
W $\Rel$ morfizm $R:A\to B$ jest podzbiorem $A\times B$, a składanie jest składaniem relacji. Suma rozłączna $P=A\sqcup B$ jest tu zarówno iloczynem, jak i koproduktem. Niech $i_A,i_B$ będą relacjami będącymi grafami włączeń, a $p_A,p_B$ relacjami odwrotnymi do tych grafów (nie funkcjami na całym $P$).

Dla relacji $R:Z\to A$, $S:Z\to B$ definiujemy $H:Z\to P$ przez
\[
z\,H\,i_A(a)\ \Longleftrightarrow\ z\,R\,a,\qquad
z\,H\,i_B(b)\ \Longleftrightarrow\ z\,S\,b,
\]
gdzie $i_A(a),i_B(b)$ oznaczają oznakowane elementy sumy. Z definicji składania $p_A\circ H=R$, $p_B\circ H=S$. Warunki te wyznaczają każdą parę należącą do $H$, ponieważ każdy element $P$ należy do dokładnie jednego składnika. To dowodzi własności iloczynu. Odwrotnie, dla $R':A\to Z$, $S':B\to Z$ określamy $H':P\to Z$ warunkami $i_A(a)\,H'\,z$ wtedy i tylko wtedy, gdy $a\,R'\,z$, oraz $i_B(b)\,H'\,z$ wtedy i tylko wtedy, gdy $b\,S'\,z$. Otrzymujemy jedyną relację spełniającą $H'i_A=R'$, $H'i_B=S'$, czyli własność koproduktu. Argument obejmuje również puste składniki. W szczególności iloczyn dwóch singletonów w $\Rel$ ma dwa elementy, a nie jeden jak w $\Set$.
\end{przyklad}

\begin{przyklad}[Klasyfikacja naturalnych operacji liniowych]\label{cat:ex:linear-natural-operations}
Ustalmy ciało $k$. Każda transformacja naturalna $\tau:V\oplus V\Rightarrow V$ (od funktora $V\mapsto V\oplus V$ do identyczności na $\Vect_k$) ma postać
\[
\tau_V(v,w)=av+bw
\]
dla jednoznacznie wyznaczonych skalarów $a,b\in k$. Istotnie, niech $a=\tau_k(1,0)$ i $b=\tau_k(0,1)$. Dla dowolnego $v\in V$ liniowy morfizm $f_v:k\to V$, $t\mapsto tv$, daje z naturalności
$\tau_V(v,0)=f_v(\tau_k(1,0))=av$. Analogicznie $\tau_V(0,w)=bw$, a liniowość $\tau_V$ daje żądany wzór. Odwrotnie, każda para skalarów definiuje liniowe składowe, a dla każdego liniowego $f:V\to W$ zachodzi $f(av+bw)=af(v)+bf(w)$. Jednoznaczność współczynników wynika z wartości na $(1,0)$ i $(0,1)$ w $k\oplus k$. Zatem w $\Vect_k$ naturalne operacje dwuargumentowe tworzą rodzinę parametryzowaną przez $k^2$, w przeciwieństwie do dwóch projekcji w $\Set$.
\end{przyklad}

\section{Słownik oznaczeń i krótkie podsumowanie}\label{sec:slownik}

Słownik zbiera najważniejsze oznaczenia stosowane w skrypcie. Ich znaczenie zależy od kategorii i ustalonych morfizmów.

\subsection{Kategorie i morfizmy}
\begin{center}
\renewcommand{\arraystretch}{1.15}
\begin{tabular}{@{}p{.34\linewidth}p{.60\linewidth}@{}}
\textbf{Oznaczenie} & \textbf{Znaczenie} \\
\hline
$\mathcal C,\mathcal D$ & kategorie \\
$\Ob(\mathcal C)$ & klasa obiektów kategorii \\
$\Hom_{\mathcal C}(A,B)$ & zbiór morfizmów z $A$ do $B$ (lokalna małość) \\
$g\circ f,\ \id_A$ & złożenie: najpierw $f$, następnie $g$; identyczność \\
$\mathcal C^{\op},\ f^{\op}$ & kategoria przeciwna i odwrócona strzałka \\
$A\cong B,\ f^{-1}$ & izomorfizm obiektów i odwrotność izomorfizmu \\
$\mathcal U$ & uniwersum Grothendiecka \\
$0_{A,B}$ & morfizm zerowy w kategorii z obiektem zerowym \\
\end{tabular}
\end{center}

\subsection{Przykładowe kategorie}
\begin{center}
\renewcommand{\arraystretch}{1.15}
\begin{tabular}{@{}p{.34\linewidth}p{.60\linewidth}@{}}
\textbf{Oznaczenie} & \textbf{Znaczenie} \\
\hline
$\Set,\ \Rel$ & zbiory i funkcje; zbiory i relacje \\
$\Grp,\ \Ab,\ \Mon$ & grupy, grupy abelowe i monoidy z homomorfizmami \\
$\Ring$ & pierścienie przemienne z jedynką i homomorfizmy zachowujące jedynkę \\
$\Vect_k,\ \Vect_k^{\mathrm{fd}}$ & przestrzenie wektorowe nad $k$ i pełna podkategoria przestrzeni skończenie wymiarowych \\
$\cat{Mod}_R^{\mathrm{l}}$ & lewe unitarne moduły nad pierścieniem $R$ z jedynką \\
$\Top,\ \Top_*$ & przestrzenie topologiczne i przestrzenie punktowane; odpowiednio odwzorowania ciągłe i ciągłe zachowujące punkt wyróżniony \\
$\Pos$ & zbiory częściowo uporządkowane z funkcjami monotonicznymi \\
$\mathbf 1$ & kategoria z jednym obiektem i tylko jego identycznością \\
\end{tabular}
\end{center}

\subsection{Funktory i transformacje}
\begin{center}
\renewcommand{\arraystretch}{1.15}
\begin{tabular}{@{}p{.34\linewidth}p{.60\linewidth}@{}}
\textbf{Oznaczenie} & \textbf{Znaczenie} \\
\hline
$F:\mathcal C\to\mathcal D$ & funktor kowariantny \\
$F:\mathcal C^{\op}\to\mathcal D$ & funktor kontrawariantny względem $\mathcal C$ \\
$\alpha:F\Rightarrow G$ & transformacja naturalna ze składowymi $\alpha_A:F(A)\to G(A)$ \\
$\Nat(F,G)$ & zbiór transformacji naturalnych (przy odpowiednich założeniach o wielkości) \\
$[\mathcal C,\mathcal D]$ & kategoria funktorów i transformacji naturalnych; lokalnie mała, gdy $\mathcal C$ jest mała i $\mathcal D$ lokalnie mała \\
$\beta\ast\alpha$ & składanie poziome; składanie pionowe oznaczamy $\circ$ \\
$h^A=\Hom(A,-)$ & kowariantny funktor reprezentowany przez $A$ \\
$h_A=\Hom(-,A)$ & kontrawariantny funktor reprezentowany przez $A$ \\
$\mathcal P,\ \mathcal P_+$ & funktor potęgowy kontrawariantny (przeciwobrazy) i kowariantny (obrazy); oba obejmują wszystkie podzbiory, także pusty \\
$\exists_f,\ \forall_f$ & obraz bezpośredni i zbiór punktów, których całe włókna są zawarte w zadanym podzbiorze (sekcja~\ref{sec:adjunctions}) \\
\end{tabular}
\end{center}

\subsection{Konstrukcje uniwersalne}
\begin{center}
\renewcommand{\arraystretch}{1.15}
\begin{tabular}{@{}p{.34\linewidth}p{.60\linewidth}@{}}
\textbf{Oznaczenie} & \textbf{Znaczenie} \\
\hline
$A\times B,\ A\sqcup B$ & iloczyn i koprodukt \\
$\prod_i A_i,\ \coprod_i A_i$ & iloczyn i koprodukt rodziny \\
$A\oplus B$ & suma prosta; w kategoriach addytywnych biprodukt \\
$A\times_C B,\ A\sqcup_C B$ & pullback (kwadrat kartezjański) i pushout, dla ustalonych strzałek do lub z $C$ \\
$\lim D,\ \colim D$ & granica i kogranica diagramu \\
$a\wedge b,\ a\vee b$ & kres dolny i górny w zbiorze częściowo uporządkowanym \\
$B^A,\ \ev$ & obiekt wykładniczy i ewaluacja $B^A\times A\to B$ \\
$\widetilde f$ & transpozycja $X\to B^A$ morfizmu $f:X\times A\to B$ w kontekście obiektów wykładniczych \\
$\int F$ & kategoria elementów funktora \\
$\mathcal C/B,\ B/\mathcal C$ & kategorie obiektów nad $B$ i pod $B$ \\
$(F\downarrow G)$ & kategoria przecinkowa \\
\end{tabular}
\end{center}

\subsection{Równoważności, sprzężenia i monady}
\begin{center}
\renewcommand{\arraystretch}{1.15}
\begin{tabular}{@{}p{.34\linewidth}p{.60\linewidth}@{}}
\textbf{Oznaczenie} & \textbf{Znaczenie} \\
\hline
$\mathcal C\simeq\mathcal D$ & równoważność kategorii \\
$F\cong G$ & izomorfizm naturalny funktorów \\
$F\dashv G$ & $F$ jest lewym sprzężeniem $G$ \\
$\eta,\ \varepsilon$ & jednostka $\id_{\mathcal C}\Rightarrow GF$ i kojednostka $FG\Rightarrow\id_{\mathcal D}$ \\
$(T,\eta,\mu)$ & monada: endofunktor, jednostka $\id\Rightarrow T$, mnożenie $T^2\Rightarrow T$ \\
$\mathcal C_T,\ g\star f$ & kategoria Kleisliego i jej składanie \\
$J,\ K$ & lewe i prawe sprzężenie Kleisliego, $KJ=T$ (stwierdzenie~\ref{cat:prop:kleisli-adj}) \\
$\mathcal C^T,\ (A,a)$ & kategoria Eilenberga--Moore'a i algebra monady $a:T(A)\to A$ \\
$F_T,\ U_T$ & funktory wolnej algebry i zapominania, $F_T\dashv U_T$ oraz $U_TF_T=T$ \\
$k^{(S)},\ G*H$ & wolna przestrzeń wektorowa na $S$ i iloczyn wolny grup \\
\end{tabular}
\end{center}

\subsection{Przykłady z innych działów}
\begin{center}
\renewcommand{\arraystretch}{1.15}
\begin{tabular}{@{}p{.34\linewidth}p{.60\linewidth}@{}}
\textbf{Oznaczenie} & \textbf{Znaczenie} \\
\hline
$\Ker f,\ \im f$ & jądro i obraz; w kategorii abelowej obraz to jądro kojądra \\
$\pi_1(X,x_0),\ H_n(X)$ & grupa podstawowa przestrzeni punktowanej i homologia singularna o współczynnikach całkowitych \\
$c(X),\ \pi_0(X)$ & zbiory składowych spójnych i składowych łukowej spójności \\
$\mathcal O(X),\ C(U,\mathbb R)$ & kategoria otwartych podzbiorów i zbiór funkcji ciągłych \\
$C_k(A,B)$ & funkcje ciągłe z topologią zwarto-otwartą \\
\end{tabular}
\end{center}

\subsection{Najważniejsze zależności}\label{cat:sec:summary-dependencies}

Kategoria opisuje obiekty i morfizmy (definicja~\ref{def:kategoria}), funktor zachowuje identyczności i składanie (definicja~\ref{def:funktor}), a transformacja naturalna porównuje funktory przez przemienne kwadraty (definicja~\ref{def:naturalna}). Własność uniwersalna daje istnienie i jednoznaczność odpowiedniego morfizmu.

\begin{itemize}
\item Granice i kogranice obejmują iloczyny, koprodukty, equalizatory, koequalizatory, pullbacki i pushouty (definicje~\ref{def:produkt}, \ref{def:koprodukt}, \ref{def:equalizer}, \ref{def:coeq}, \ref{def:pullback}, \ref{def:pushout}). Granica pustego diagramu jest obiektem końcowym, a kogranica początkowym (definicje~\ref{def:limit}, \ref{def:colimit}).
\item Reprezentowalność wiąże własność uniwersalną z funktorem Hom (definicja~\ref{def:representable}). Dla lokalnie małej $\mathcal C$ lemat Yonedy daje naturalne bijekcje $\Nat(h^A,F)\cong F(A)$ dla $F:\mathcal C\to\Set$ oraz $\Nat(h_A,G)\cong G(A)$ dla $G:\mathcal C^{\op}\to\Set$ (twierdzenia~\ref{tw:yoneda} i \ref{cat:thm:yoneda-contra}). Kontrawariantne funktory reprezentowalne tworzą włożenie Yonedy $A\mapsto h_A$, które jest w pełni wierne (twierdzenie~\ref{tw:yoneda-embedding}); dla funktorów kowariantnych dziedziną odpowiedniego włożenia jest $\mathcal C^{\op}$ (stwierdzenie~\ref{cat:prop:yoneda-cov-embedding}).
\item Kategoria obiektów nad bazą to szczególny przypadek kategorii przecinkowej (definicje~\ref{def:slice}, \ref{cat:def:comma}).
\item Funktor między małymi kategoriami jest równoważnością wtedy i tylko wtedy, gdy jest w pełni wierny i istotnie surjektywny na obiektach; sama pełna wierność nie wystarcza (twierdzenie~\ref{tw:equivalence}, z podanymi tam założeniami o wielkości i wyborze).
\item Sprzężenie $F\dashv G$ określa naturalne bijekcje $\Hom_{\mathcal D}(F(A),B)\cong\Hom_{\mathcal C}(A,G(B))$ (definicja~\ref{def:adjunction}). Lewe sprzężenie zachowuje istniejące małe kogranice, prawe istniejące małe granice (twierdzenie~\ref{tw:adjoints-preserve}). Konstrukcja swobodna i zapominanie to przykłady, nie definicja wszystkich sprzężeń.
\item Każde sprzężenie indukuje monadę $T=GF$ z mnożeniem $\mu_A=G(\varepsilon_{F(A)})$ (twierdzenie~\ref{cat:thm:adjunction-monad}). Każda monada na lokalnie małej kategorii pochodzi zarówno ze sprzężenia Kleisliego, jak i Eilenberga--Moore'a (stwierdzenia~\ref{cat:prop:kleisli-adj}, \ref{cat:prop:em-adj}).
\end{itemize}

\subsection{Rozróżnienia, o których warto pamiętać}\label{cat:sec:summary-pitfalls}
\begin{itemize}
\item Monomorfizm i epimorfizm definiujemy przez skracanie złożeń, nie przez injektywność i surjektywność (definicja~\ref{def:mono-epi}). Morfizm będący jednocześnie monomorfizmem i epimorfizmem nie musi być izomorfizmem (przykład~\ref{cat:ex:normed-nonabelian}).
\item Obiekt początkowy i końcowy nie muszą być tym samym obiektem. Obiekt zerowy spełnia obie własności (definicja~\ref{def:poczatkowy}).
\item Iloczyn i koprodukt mogą być różne, nawet gdy w innej kategorii realizują się przez tę samą konstrukcję. W $\Rel$ suma rozłączna jest obydwoma (przykład~\ref{cat:ex:rel-biproduct}).
\item Izomorfizm kategorii oznacza funktory wzajemnie odwrotne, a równoważność wymaga jedynie izomorfizmów naturalnych ich złożeń z identycznościami.
\item Funktor reprezentowany przez ustalony obiekt może być kowariantny lub kontrawariantny. Kierunek strzałek rozstrzyga, której postaci lematu Yonedy używamy.
\item Pełna wierność lewego sprzężenia jest równoważna odwracalności jednostki, a pełna wierność prawego sprzężenia odwracalności kojednostki (stwierdzenie~\ref{cat:prop:adj-full-faithful}). Same tożsamości trójkątne nie gwarantują odwracalności obu transformacji.
\item Istnienie iloczynów nie zapewnia istnienia obiektów wykładniczych. W $\Top$ topologia zwarto-otwarta daje obiekt wykładniczy przy wystarczającej hipotezie lokalnej zwartości i własności Hausdorffa przestrzeni $A$ (twierdzenie~\ref{cat:thm:exp-top}).
\item W kategorii abelowej jądra, kojądra i biprodukty współdziałają z addytywnością składania (definicja~\ref{cat:def:abelian-category}). Samo istnienie obiektu zerowego i skończonych iloczynów oraz koproduktów nie wystarcza do abelowości.
\end{itemize}

Najważniejsza praktyczna zasada brzmi: kiedy spotykamy nowy obiekt matematyczny, warto zapytać, \emph{jakie morfizmy do niego lub z niego wychodzą i czy można je opisać uniwersalnie}. W wielu sytuacjach taka charakterystyka jest trwalsza, prostsza i bardziej użyteczna niż konstrukcja na elementach.

\clearpage
\part*{Część II. Zaawansowana teoria kategorii}
\addcontentsline{toc}{part}{Część II. Zaawansowana teoria kategorii}
\section{Kategorie monoidalne, symetryczne i domknięte}\label{cat:sec:monoidal}

Dotychczas iloczyn obiektów był wyróżniany przez własność uniwersalną (definicja~\ref{def:produkt}). W wielu kategoriach potrzebujemy jednak innej operacji dwuargumentowej: na przykład iloczyn tensorowy modułów nie jest na ogół ich iloczynem kategorialnym. Przejdziemy od konkretnych izomorfizmów tensorowych do aksjomatów opisujących ich wzajemną zgodność.

\subsection{Struktura monoidalna i aksjomaty koherencji}
\begin{definicja}[Kategoria monoidalna]\label{cat:def:monoidal}
Kategorią monoidalną nazywamy kategorię $\mathcal C$ z bifunktorem $\otimes:\mathcal C\times\mathcal C\to\mathcal C$, obiektem $I$ oraz izomorfizmami naturalnymi
\[
\begin{aligned}
\alpha_{A,B,C}&:(A\otimes B)\otimes C\longrightarrow A\otimes(B\otimes C),\\
\lambda_A&:I\otimes A\longrightarrow A,\\
\rho_A&:A\otimes I\longrightarrow A.
\end{aligned}
\]
Wymagamy dwóch aksjomatów koherencji. Aksjomat pięciokąta ma postać
\begin{equation}\label{cat:eq:pentagon}
\alpha_{A,B,C\otimes D}\circ\alpha_{A\otimes B,C,D}
=(\id_A\otimes\alpha_{B,C,D})\circ\alpha_{A,B\otimes C,D}
\circ(\alpha_{A,B,C}\otimes\id_D).
\end{equation}
Obie strony są morfizmami $((A\otimes B)\otimes C)\otimes D\to A\otimes(B\otimes(C\otimes D))$.
Aksjomat trójkąta ma postać
\begin{equation}\label{cat:eq:triangle-monoidal}
(\id_A\otimes\lambda_B)\circ\alpha_{A,I,B}=\rho_A\otimes\id_B,
\end{equation}
gdzie obie strony prowadzą z $(A\otimes I)\otimes B$ do $A\otimes B$.
\end{definicja}

Naturalność $\alpha$ oznacza dla $f:A\to A'$, $g:B\to B'$, $h:C\to C'$ równość
\[
\alpha_{A',B',C'}\circ((f\otimes g)\otimes h)
=(f\otimes(g\otimes h))\circ\alpha_{A,B,C}.
\]
Analogicznie $\lambda_{A'}(\id_I\otimes f)=f\lambda_A$ i $\rho_{A'}(f\otimes\id_I)=f\rho_A$. Bifunktorialność (definicja~\ref{cat:def:bifunctor}) wymaga $\id_A\otimes\id_B=\id_{A\otimes B}$ oraz
\[
(f'\otimes g')\circ(f\otimes g)=(f'\circ f)\otimes(g'\circ g)
\]
dla składanych par morfizmów.

Pięciokąt można odczytać jako dwie drogi pomiędzy pięcioma nawiasowaniami czterech czynników. Trójkąt ma następujący diagram:
\[
\begin{CD}
(A\otimes I)\otimes B @>{\alpha_{A,I,B}}>> A\otimes(I\otimes B)\\
@V{\rho_A\otimes\id_B}VV @VV{\id_A\otimes\lambda_B}V\\
A\otimes B @>{\id_{A\otimes B}}>> A\otimes B .
\end{CD}
\]
Same izomorfizmy łączności i jednostki nie wystarczają: potrzebna jest zgodność ich złożeń.

\begin{definicja}[Ścisła kategoria monoidalna]\label{cat:def:strict-monoidal}
Kategoria monoidalna jest ścisła, gdy $\otimes$ jest łączny i ma jednostkę $I$ na poziomie obiektów i morfizmów, a $\alpha,\lambda,\rho$ są identycznościami.
\end{definicja}

\begin{przyklad}[Endofunktory i składanie]\label{cat:ex:endofunctor-monoidal}
Dla małej kategorii $\mathcal C$ kategoria $[\mathcal C,\mathcal C]$ endofunktorów jest ścisła monoidalna. Przyjmujemy $F\otimes G=F\circ G$, $I=\id_{\mathcal C}$. Dla transformacji $\eta:F\Rightarrow F'$ i $\theta:G\Rightarrow G'$ iloczynem jest złożenie poziome $\eta\ast\theta:FG\Rightarrow F'G'$ z definicji~\ref{cat:def:horizontal}. Prawo wymiany (twierdzenie~\ref{cat:thm:interchange}) daje bifunktorialność, a stwierdzenie~\ref{cat:prop:horizontal-assoc} daje ścisłą łączność i jednostki. Zatem pięciokąt i trójkąt składają się z identyczności.
\end{przyklad}

\begin{twierdzenie}[Struktura kartezjańska]\label{cat:thm:cartesian-monoidal}
Każda kategoria z wybranymi iloczynami binarnymi i obiektem końcowym $1$ ma strukturę monoidalną dla $\otimes=\times$ oraz $I=1$.
\end{twierdzenie}
\begin{proof}
Bifunktor iloczynu otrzymujemy ze stwierdzenia~\ref{cat:prop:product-functor}. Zarówno $(A\times B)\times C$, jak i $A\times(B\times C)$ reprezentują iloczyn trzech obiektów: morfizm do dowolnego z nich jest jednoznacznie wyznaczony przez trzy składowe. To daje kanoniczny izomorfizm $\alpha$ zachowujący wszystkie trzy rzuty (twierdzenie~\ref{tw:produkt-unique}). Rzuty $\lambda_A:1\times A\to A$ i $\rho_A:A\times1\to A$ są izomorfizmami: odwrotności to odpowiednio $\langle !_A,\id_A\rangle$ oraz $\langle\id_A,!_A\rangle$, gdzie $!_A:A\to1$ jest jedynym morfizmem. Złożenia z rzutami pokazują, że są to odwrotności z obu stron.

Obie strony równania naturalności $\alpha$ mają po złożeniu z odpowiednimi rzutami na $A',B',C'$ składowe $f p_A$, $g p_B$, $h p_C$, gdzie $p_A,p_B,p_C$ są rzutami z obiektu wyjściowego. Jednoznaczność morfizmu do iloczynu daje naturalność. Naturalność $\lambda,\rho$ jest warunkiem zachowania odpowiedniego rzutu przez $f\times\id$ lub $\id\times f$.

W pięciokącie każde zastosowanie $\alpha$ zachowuje rzuty na czynniki, więc obie drogi mają po złożeniu z czterema końcowymi rzutami odpowiednio $p_A,p_B,p_C,p_D$. Obiekt końcowy diagramu jest iterowanym iloczynem czterech obiektów; jednoznaczność morfizmu o tych składowych daje \eqref{cat:eq:pentagon}. W trójkącie obie drogi zachowują rzuty na $A,B$, a więc są równe przez własność uniwersalną $A\times B$. Otrzymujemy \eqref{cat:eq:triangle-monoidal}.
\end{proof}

\begin{twierdzenie}[Struktura tensorowa modułów]\label{cat:thm:module-monoidal}
Niech $R$ będzie przemiennym pierścieniem z jedynką. Kategoria unitarnych $R$-modułów z iloczynem $\otimes_R$ i jednostką $R$ jest monoidalna.
\end{twierdzenie}
\begin{proof}
Wykorzystujemy konstrukcję iloczynu tensorowego, w szczególności \atlasref{AA-18::aa:tw:tensor-universal}{AA-18: własność uniwersalna}, \atlasref{AA-18::aa:tw:tensor-sym-assoc}{AA-18: łączność i symetria} oraz \atlasref{AA-18::aa:tw:tensor-unit}{AA-18: jednostka}. Izomorfizmy $\alpha,\lambda,\rho$ działają na tensorach prostych przez
\[
(m\otimes n)\otimes p\mapsto m\otimes(n\otimes p),\quad
r\otimes m\mapsto rm,\quad m\otimes r\mapsto rm .
\]
Ich dobre określenie i odwracalność wynika z przytoczonych twierdzeń. Bifunktorialność wykazano w przykładzie~\ref{cat:ex:tensor-bifunctor}.

Sprawdzimy pozostałe aksjomaty. Dla homomorfizmów $f,g,h$ obie strony naturalności $\alpha$ wysyłają $(m\otimes n)\otimes p$ na $f(m)\otimes(g(n)\otimes h(p))$. Naturalność $\lambda$ wynika z $f(rm)=rf(m)$, a dla $\rho$ działa ten sam rachunek. Tensory proste generują odpowiednie moduły, więc te równości zachodzą dla wszystkich elementów.

W pięciokącie droga z dwoma zastosowaniami $\alpha$ prowadzi przez
\[
((m\otimes n)\otimes p)\otimes q
\longmapsto (m\otimes n)\otimes(p\otimes q)
\longmapsto m\otimes(n\otimes(p\otimes q)).
\]
Droga z trzema zastosowaniami ma kolejno wartości
\[
\begin{aligned}
((m\otimes n)\otimes p)\otimes q
&\longmapsto(m\otimes(n\otimes p))\otimes q\\
&\longmapsto m\otimes((n\otimes p)\otimes q)\\
&\longmapsto m\otimes(n\otimes(p\otimes q)).
\end{aligned}
\]
Iterowane tensory proste generują moduł wyjściowy: rozwijamy każdy czynnik kolejno jako skończoną sumę tensorów prostych i korzystamy z dwuliniowości. Obie drogi są liniowe i równe na tych generatorach, więc są równe. W trójkącie pierwsza droga daje $m\otimes(rn)$, a druga $(rm)\otimes n$. Równość jest relacją zbilansowania (por. \atlasref{AA-18::aa:def:tensor-balanced}{AA-18: odwzorowania zbilansowane}). Znów wystarczy równość na generatorach.
\end{proof}

\subsection{Plecenie i symetria}
\begin{definicja}[Kategoria monoidalna z pleceniem i symetryczna]\label{cat:def:braided-symmetric}
Plecenie jest izomorfizmem naturalnym $c_{A,B}:A\otimes B\to B\otimes A$ spełniającym dwa aksjomaty sześciokąta:
\begin{align*}
\alpha_{B,C,A}\,c_{A,B\otimes C}\,\alpha_{A,B,C}
&=(\id_B\otimes c_{A,C})\,\alpha_{B,A,C}\,(c_{A,B}\otimes\id_C),\\
\alpha^{-1}_{C,A,B}\,c_{A\otimes B,C}\,\alpha^{-1}_{A,B,C}
&=(c_{A,C}\otimes\id_B)\,\alpha^{-1}_{A,C,B}\,(\id_A\otimes c_{B,C}).
\end{align*}
W pierwszym równaniu obie strony prowadzą z $(A\otimes B)\otimes C$ do $B\otimes(C\otimes A)$, w drugim z $A\otimes(B\otimes C)$ do $(C\otimes A)\otimes B$. Składamy od prawej do lewej.
Kategoria monoidalna z pleceniem jest symetryczna, jeśli dodatkowo $c_{B,A}\circ c_{A,B}=\id_{A\otimes B}$.
\end{definicja}

\begin{stwierdzenie}[Symetria w przykładach kartezjańskich i tensorowych]\label{cat:prop:cartesian-tensor-symmetry}
Struktury z twierdzeń~\ref{cat:thm:cartesian-monoidal} i \ref{cat:thm:module-monoidal} są symetryczne.
\end{stwierdzenie}
\begin{proof}
W przypadku kartezjańskim bierzemy $c_{A,B}=\langle p_B,p_A\rangle$. Dwa złożenia zamian zachowują oba rzuty, więc są identycznościami. Naturalność wynika z tego, że obie drogi dla $f:A\to A'$ i $g:B\to B'$ mają składowe $g p_B$ oraz $f p_A$. Pierwszy sześciokąt daje po obu drogach uporządkowane składowe $B,C,A$, drugi $C,A,B$, z zachowaniem rzutów na poszczególne czynniki. Jednoznaczność morfizmów do iloczynu dowodzi obu równości.

W przypadku modułów stosujemy izomorfizm $c_{M,N}(m\otimes n)=n\otimes m$ z \atlasref{AA-18::aa:tw:tensor-sym-assoc}{AA-18}. Naturalność sprawdzamy na tensorach prostych: obie drogi dają $g(n)\otimes f(m)$. W pierwszym sześciokącie obie drogi wysyłają $(m\otimes n)\otimes p$ na $n\otimes(p\otimes m)$; w drugim obie wysyłają $m\otimes(n\otimes p)$ na $(p\otimes m)\otimes n$. Podwójna zamiana daje $m\otimes n$. Liniowość i generowanie przez tensory proste kończą dowód.
\end{proof}

\begin{przyklad}[Symetria przestrzeni z gradacją]\label{cat:ex:graded-symmetry}
Dla $\mathbb Z$-stopniowanych przestrzeni wektorowych nad ciałem $k$ i odwzorowań zachowujących stopień rozważamy standardowy iloczyn tensorowy z gradacją $(V\otimes W)_n=\bigoplus_{p+q=n}V_p\otimes W_q$. Jednostką jest $k$ w stopniu zero. Łączność i jednostki są odziedziczone z iloczynu tensorowego, zachowują stopień, a ich aksjomaty wynikają z twierdzenia~\ref{cat:thm:module-monoidal}.
Na jednorodnych wektorach $v\in V_p,w\in W_q$ definiujemy
\[
c(v\otimes w)=(-1)^{pq}w\otimes v.
\]
Wzór na każdej parze składowych stopniowanych jest dwuliniowy i daje liniowe odwzorowanie z sumy prostej; zachowuje stopień $p+q$. Naturalność wynika z zachowywania stopni przez morfizmy. Dwukrotna zamiana daje czynnik $(-1)^{2pq}=1$, więc $c$ jest izomorfizmem i spełnia symetrię. Dla trzeciego jednorodnego wektora stopnia $r$ pierwszy sześciokąt daje znaki $(-1)^{p(q+r)}$ i $(-1)^{pq+pr}$, drugi $(-1)^{(p+q)r}$ i $(-1)^{pr+qr}$. Znaki są równe, a kolejność tensorów po obu drogach taka sama. Jednorodne tensory proste generują rozpatrywane moduły, więc oba sześciokąty zachodzą. W charakterystyce $2$ wszystkie te znaki są równe $1$.
\end{przyklad}

\subsection{Domkniętość monoidalna}
\begin{definicja}[Prawa domkniętość monoidalna]\label{cat:def:monoidal-closed}
Kategoria monoidalna jest prawostronnie domknięta, gdy dla każdego $A$ funktor $-\otimes A$ ma prawe sprzężenie $[A,-]$. Oznacza to naturalną bijekcję
\[
\Hom_{\mathcal C}(X\otimes A,B)\cong\Hom_{\mathcal C}(X,[A,B]).
\]
Obiekt $[A,B]$ nazywamy wewnętrznym obiektem Hom. W kategorii symetrycznej domkniętość z jednej strony pociąga domkniętość z drugiej: symetria daje naturalny izomorfizm $A\otimes-\cong-\otimes A$, więc powyższe bijekcje można złożyć z bijekcjami indukowanymi przez ten izomorfizm.
\end{definicja}

\begin{przyklad}[Domkniętość kartezjańska i tensorowa]\label{cat:ex:closed-monoidal}
W kategorii kartezjańsko domkniętej $[A,B]=B^A$ (definicja~\ref{def:exponential}, twierdzenie~\ref{cat:thm:exp-naturality}). W kategorii modułów nad przemiennym pierścieniem $R$ przyjmujemy $[A,B]=\operatorname{Hom}_R(A,B)$ z punktowymi działaniami modułu. Naturalną bijekcję wyznaczają wzory
\[
\Phi(f)(x)(a)=f(x\otimes a),\qquad
\Psi(g)(x\otimes a)=g(x)(a).
\]
Ich dobre określenie i wzajemną odwrotność daje \atlasref{AA-18::aa:tw:tensor-hom-adjunction}{AA-18: sprzężenie tensor--Hom}. Dla $u:X'\to X$, $v:B\to B'$ naturalność sprawdzamy bezpośrednio:
\[
\Phi(vf(u\otimes\id_A))(x')(a)=v(f(u(x')\otimes a)).
\]
Jest to dokładnie wartość złożenia $\operatorname{Hom}_R(A,v)\circ\Phi(f)\circ u$ w $x'$ i $a$. W ten sposób uzyskujemy sprzężenie, nie tylko bijekcję poszczególnych zbiorów. Nie potrzeba skończonej wymiarowości ani skończonego generowania modułów.
\end{przyklad}

\begin{stwierdzenie}[Zachowywanie kogranic przez iloczyn monoidalny]\label{cat:prop:closed-colimits}
W prawostronnie domkniętej kategorii monoidalnej funktor $-\otimes A$ zachowuje istniejące małe kogranice. W przypadku symetrycznym dotyczy to także $A\otimes-$.
\end{stwierdzenie}
\begin{proof}
Z definicji~\ref{cat:def:monoidal-closed} jest to lewe sprzężenie. Stosujemy twierdzenie~\ref{tw:adjoints-preserve}. W przypadku symetrycznym drugi funktor jest naturalnie izomorficzny z pierwszym; izomorfizmy przenoszą własność uniwersalną kogranicy.
\end{proof}

\begin{uwaga}
Koherencja dla dowolnie długich wyrażeń monoidalnych, funktory monoidalne, obiekty monoidalne oraz obiekty posiadające dualność wymagają dalszych definicji i twierdzeń. Powyższe przykłady dowodzą konkretnych aksjomatów pięciokąta i trójkąta, a nie ogólnego twierdzenia Mac Lane'a o koherencji. W dalszej części omówimy je oddzielnie.
\end{uwaga}

\section{Koherencja Mac Lane'a i funktory monoidalne}\label{cat:sec:coherence}

Pracujemy z kategorią monoidalną w sensie definicji~\ref{cat:def:monoidal}. Najpierw wyprowadzimy tożsamości jednostki, następnie udowodnimy koherencję i skonstruujemy kategorię ściśle monoidalną równoważną wyjściowej.

\subsection{Dodatkowe tożsamości jednostki}
\begin{lemat}[Tożsamości jednostki]\label{cat:lem:unit-coherence}
W każdej kategorii monoidalnej zachodzą
\begin{align}
\lambda_{A\otimes B}\alpha_{I,A,B}&=\lambda_A\otimes\id_B,\label{cat:eq:left-unit-coherence}\\
\rho_{A\otimes B}&=(\id_A\otimes\rho_B)\alpha_{A,B,I},\label{cat:eq:right-unit-coherence}\\
\lambda_I&=\rho_I.\label{cat:eq:unitors-equal}
\end{align}
\end{lemat}
\begin{proof}
W dowodzie opuszczamy znak składania, zachowując kolejność od prawej do lewej. Funktory $I\otimes-$ oraz $-\otimes I$ są wierne, ponieważ są naturalnie izomorficzne z funktorem identycznościowym. Dokładniej, równość $I\otimes f=I\otimes g$ po złożeniu z $\lambda$ i wykorzystaniu naturalności daje $f\lambda=g\lambda$, skąd $f=g$.

Ustalmy $X=I$ (poniższe rachunki działają również dla dowolnego $X$). Pięciokąt dla $X,I,A,B$ daje
\[
\begin{aligned}
&\alpha_{X,I,A\otimes B}\alpha_{X\otimes I,A,B}\\
&=(\id_X\otimes\alpha_{I,A,B})
\alpha_{X,I\otimes A,B}(\alpha_{X,I,A}\otimes\id_B).
\end{aligned}
\]
Składamy obie strony z $\id_X\otimes\lambda_{A\otimes B}$. Lewa strona, dzięki trójkątowi i naturalności $\alpha$, staje się
\[
\begin{aligned}
&(\rho_X\otimes\id_{A\otimes B})\alpha_{X\otimes I,A,B}\\
&=\alpha_{X,A,B}((\rho_X\otimes\id_A)\otimes\id_B)\\
&=\alpha_{X,A,B}((\id_X\otimes\lambda_A)\otimes\id_B)
(\alpha_{X,I,A}\otimes\id_B)\\
&=(\id_X\otimes(\lambda_A\otimes\id_B))
\alpha_{X,I\otimes A,B}(\alpha_{X,I,A}\otimes\id_B).
\end{aligned}
\]
Druga równość to naturalność, trzecia wynika z trójkąta (po uwzględnieniu bifunktorialności), a ostatnia z naturalności $\alpha$. Prawa strona pięciokąta po złożeniu to
\[
(\id_X\otimes(\lambda_{A\otimes B}\alpha_{I,A,B}))
\alpha_{X,I\otimes A,B}(\alpha_{X,I,A}\otimes\id_B).
\]
Skracamy wspólne izomorfizmy po prawej, następnie korzystamy z wierności $I\otimes-$, otrzymując \eqref{cat:eq:left-unit-coherence}.

Aby otrzymać \eqref{cat:eq:right-unit-coherence}, stosujemy właśnie dowiedzioną tożsamość do kategorii z odwróconą kolejnością iloczynu: $A\odot B=B\otimes A$, $\alpha^\odot_{A,B,C}=\alpha^{-1}_{C,B,A}$, $\lambda^\odot_A=\rho_A$, $\rho^\odot_A=\lambda_A$. Pięciokąt dla $\odot$ jest odwróconym pięciokątem dla $\otimes$ z kolejnością czterech czynników odwróconą; trójkąt to równanie $(\rho_B\otimes\id_A)\alpha_{B,I,A}^{-1}=\id_B\otimes\lambda_A$, równoważne pierwotnemu. Pierwsza tożsamość w nowej strukturze brzmi $\rho_{B\otimes A}\alpha_{B,A,I}^{-1}=\id_B\otimes\rho_A$. Złożenie z $\alpha_{B,A,I}$ i zamiana nazw daje żądaną równość.

Naturalność $\lambda$ względem $\lambda_I:I\otimes I\to I$ daje $\lambda_I(\id_I\otimes\lambda_I)=\lambda_I\lambda_{I\otimes I}$. Skracamy izomorfizm $\lambda_I$, więc $\lambda_{I\otimes I}=\id_I\otimes\lambda_I$. Wstawiając do \eqref{cat:eq:left-unit-coherence} dla $A=B=I$ i stosując trójkąt, dostajemy $\lambda_I\otimes\id_I=\rho_I\otimes\id_I$. Wierność $-\otimes I$ daje \eqref{cat:eq:unitors-equal}.
\end{proof}

\subsection{Twierdzenie o koherencji}
Wyrażenia formalne budujemy z liter oznaczających obiekty, specjalnego symbolu jednostki $I$ oraz binarnego symbolu $\otimes$. Litera jest odróżniana od formalnego symbolu jednostki nawet wtedy, gdy przypisano jej obiekt jednostkowy. Słowo wyrażenia to uporządkowana lista liter po pominięciu formalnych jednostek.

\begin{twierdzenie}[Mac Lane'a o koherencji monoidalnej]\label{cat:thm:mac-lane-coherence}
Dla dwóch formalnych wyrażeń o tym samym słowie wszystkie morfizmy otrzymane z $\alpha,\lambda,\rho$, ich odwrotności, identyczności przez składanie i iloczyn monoidalny są równe. Istnieje co najmniej jeden taki izomorfizm. Porównujemy morfizmy o tych samych formalnych końcach, a nie dowolne morfizmy pomiędzy obiektami kategorii.
\end{twierdzenie}
\begin{proof}
Ustalamy orientację trzech operacji, które wolno wykonywać w dowolnym podwyrażeniu:
\[
((U\otimes V)\otimes W)\to U\otimes(V\otimes W),\qquad
I\otimes U\to U,\qquad U\otimes I\to U.
\]
Odpowiadają im odpowiednio $\alpha,\lambda,\rho$, z identycznościami w pozostałych czynnikach. Każdy krok jest izomorfizmem.

\textbf{Zakończenie redukcji.}
Niech $n(T)$ będzie liczbą wszystkich liści drzewa wyrażenia $T$ (łącznie z jednostkami), a $w(T)$ sumą, po wszystkich węzłach binarnych, liczby liści lewego poddrzewa. Porządkujemy pary $(n,w)$ leksykograficznie w $\mathbb N^2$. Usunięcie jednostki zmniejsza $n$. Zmiana nawiasowania nie zmienia $n$, ale zmniejsza $w$ o $n(U)$: wkład dwóch zmienionych węzłów wynosi przed krokiem $2n(U)+n(V)$, a po nim $n(U)+n(V)$. Wkłady węzłów poza podwyrażeniem pozostają takie same. Nie istnieje zatem nieskończony ciąg redukcji.

\textbf{Postać normalna.}
Każde wyrażenie nieredukowalne jest pojedynczą literą, samą jednostką albo wyrażeniem prawostronnie nawiasowanym złożonym z liter. Istotnie, gdy korzeń jest binarny, lewy czynnik nie może być binarny ani jednostką, prawy zaś nie może być jednostką. Stosujemy ten argument rekurencyjnie w prawym poddrzewie. Postać normalna $N(T)$ jest jednoznacznie wyznaczona przez słowo: dla pustego słowa jest to $I$. Każdy krok zachowuje słowo.

\textbf{Lokalne rozgałęzienia.}
Dwa różne kroki z tego samego wyrażenia można uzupełnić do diagramu przemiennego złożonego z redukcji. Jeżeli dotyczą rozłącznych poddrzew, wykonujemy oba w dowolnej kolejności; bifunktorialność iloczynu gwarantuje równość. Jeżeli jeden krok występuje wewnątrz zmiennej $U,V,W$ drugiej reguły, naturalność odpowiednio $\alpha,\lambda,\rho$ daje kwadrat przemienny. Pozostają nakładania w miejscach niebędących zmiennymi reguł. Są dokładnie następujące:
\[
\begin{array}{c|c}
\text{wyrażenie}&\text{równość uzupełniająca}\\ \hline
((U\otimes V)\otimes W)\otimes Z&\text{pięciokąt}\\
(I\otimes U)\otimes V&\eqref{cat:eq:left-unit-coherence}\\
(U\otimes I)\otimes V&\text{trójkąt}\\
(U\otimes V)\otimes I&\eqref{cat:eq:right-unit-coherence}\\
I\otimes I&\eqref{cat:eq:unitors-equal}
\end{array}
\]
Aby sprawdzić kompletność, zauważmy, że reguła $\alpha$ ma tylko jedno właściwe podwyrażenie binarne, które nie jest zmienną, mianowicie $U\otimes V$. Może ono rozpoczynać każdą z trzech reguł, co daje pierwsze trzy przypadki. Na korzeniu reguła $\alpha$ może zbiegać się z $\rho$, jeśli $W=I$, co daje czwarty przypadek. Reguły $\lambda,\rho$ mogą zbiegać się tylko w $I\otimes I$. Pozostałe właściwe podwyrażenia reguł jednostki są zmiennymi lub liśćmi. Jeżeli nakładanie występuje w większym kontekście, iloczyn z identycznościami zachowuje równość. Wszystkie wymienione diagramy są przemienne na mocy aksjomatów CAT-19 lub lematu~\ref{cat:lem:unit-coherence}.

\textbf{Globalna jednoznaczność.}
Indukujemy względem dobrze ufundowanej miary $(n,w)$. Dla wyrażenia normalnego jedyną drogą redukcji do postaci normalnej jest identyczność. Dla pozostałych wyrażeń rozważamy dwie drogi do postaci normalnej. Jeżeli zaczynają się tym samym krokiem, stosujemy założenie indukcyjne do jego celu. W przeciwnym razie oznaczamy pierwsze kroki przez $f:T\to U$, $g:T\to V$. Z lokalnej kontroli mamy drogi redukcji $p:U\to W$, $q:V\to W$ z $pf=qg$. Wybieramy drogę $r:W\to N(T)$, która istnieje dzięki zakończeniu redukcji i zachowaniu słowa. Miary $U,V$ są ściśle mniejsze od miary $T$ niezależnie od tego, czy $U=V$ jako wyrażenia formalne; także dwa różne kroki o wspólnym celu podlegają argumentowi lokalnemu (w szczególności $\lambda_I,\rho_I$). Wówczas $rp$ i $rq$ są drogami normalizującymi. Z założenia indukcyjnego są równe odpowiednim ogonom pierwotnych dróg. Równość $rpf=rqg$ dowodzi tezy.

Oznaczmy jedyny morfizm normalizujący przez $\nu_T:T\to N(T)$. Dla każdego elementarnego kroku $f:T\to U$ mamy $\nu_U f=\nu_T$, zatem $f=\nu_U^{-1}\nu_T$. Dla jego odwrotności zachodzi analogiczny wzór. Dowolny iloczyn dwóch morfizmów można rozpisać na kroki w podwyrażeniach, ponieważ $f\otimes g=(\id\otimes g)(f\otimes\id)$ z właściwymi identycznościami. W dowolnej drodze złożenia $\nu_U\nu_U^{-1}$ skracają się, pozostawiając $\nu_S^{-1}\nu_T$ dla końców $T,S$. Jest to również istniejący morfizm między tymi końcami. Dowód nie używa koherencji w przesłankach, tylko aksjomatów i lematu jednostki.
\end{proof}

\begin{uwaga}
Twierdzenie pozwala pomijać nawiasy w rachunkach strukturalnych, lecz nie utożsamia dowolnych morfizmów. Nie obejmuje samo z siebie permutacji czynników przez plecenie; dla kategorii symetrycznych udowodnimy odrębne twierdzenie o koherencji symetrii w ostatniej części rozdziału. Historyczny punkt odniesienia: S.~Mac Lane, \emph{Categories for the Working Mathematician}, rozdział VII; redukcja aksjomatów jednostki: G.~M.~Kelly, \emph{On MacLane's conditions for coherence of natural associativities, commutativities, etc.}, Journal of Algebra 1 (1964).
\end{uwaga}

\subsection{Funktory monoidalne}
\begin{definicja}[Funktor monoidalny typu lax]\label{cat:def:lax-monoidal}
Niech $\mathcal C,\mathcal D$ będą kategoriami monoidalnymi. Funktor monoidalny typu lax (o morfizmach strukturalnych skierowanych do obrazu iloczynu) składa się z funktora $F:\mathcal C\to\mathcal D$, naturalnej rodziny
\[
\phi_{A,B}:F(A)\otimes F(B)\to F(A\otimes B),\qquad
\phi_0:I_{\mathcal D}\to F(I_{\mathcal C})
\]
spełniających równości
\[
\begin{aligned}
&F(\alpha_{A,B,C})\phi_{A\otimes B,C}(\phi_{A,B}\otimes\id_{F(C)})\\
&=\phi_{A,B\otimes C}(\id_{F(A)}\otimes\phi_{B,C})\alpha_{F(A),F(B),F(C)},\\
&F(\lambda_A)\phi_{I_{\mathcal C},A}(\phi_0\otimes\id_{F(A)})=\lambda_{F(A)},\\
&F(\rho_A)\phi_{A,I_{\mathcal C}}(\id_{F(A)}\otimes\phi_0)=\rho_{F(A)}.
\end{aligned}
\]
W pierwszym równaniu dziedziną jest $(F(A)\otimes F(B))\otimes F(C)$, a przeciwdziedziną $F(A\otimes(B\otimes C))$; w dwóch ostatnich odpowiednio $I_{\mathcal D}\otimes F(A)$ oraz $F(A)\otimes I_{\mathcal D}$, z przeciwdziedziną $F(A)$. Funktor jest silny, jeśli wszystkie $\phi_{A,B}$ i $\phi_0$ są izomorfizmami; jest ścisły, jeśli są identycznościami (w szczególności wymagane są odpowiednie równości obiektów).
\end{definicja}

\begin{definicja}[Funktor komonoidalny]\label{cat:def:oplax-monoidal}
Funktor typu oplax (komonoidalny) ma naturalne morfizmy $\psi_{A,B}:F(A\otimes B)\to F(A)\otimes F(B)$ oraz $\psi_0:F(I_{\mathcal C})\to I_{\mathcal D}$, spełniające
\[
\begin{aligned}
&\alpha_{F(A),F(B),F(C)}(\psi_{A,B}\otimes\id_{F(C)})\psi_{A\otimes B,C}\\
&=(\id_{F(A)}\otimes\psi_{B,C})\psi_{A,B\otimes C}F(\alpha_{A,B,C}),\\
&\lambda_{F(A)}(\psi_0\otimes\id_{F(A)})\psi_{I_{\mathcal C},A}=F(\lambda_A),\\
&\rho_{F(A)}(\id_{F(A)}\otimes\psi_0)\psi_{A,I_{\mathcal C}}=F(\rho_A).
\end{aligned}
\]
Jeżeli struktura lax jest silna, jej odwrotne izomorfizmy definiują strukturę oplax: odwracamy obie drogi w każdym diagramie i w razie potrzeby składamy z izomorfizmem łączności lub jednostki, aby uzyskać powyższe równania.
\end{definicja}

\begin{definicja}[Monoidalna transformacja naturalna]\label{cat:def:monoidal-natural}
Dla funktorów lax $F,G$ o strukturach $\phi,\gamma$ transformacja naturalna $\eta:F\Rightarrow G$ jest monoidalna, jeśli
\[
\eta_{A\otimes B}\phi_{A,B}
=\gamma_{A,B}(\eta_A\otimes\eta_B),\qquad
\eta_{I_{\mathcal C}}\phi_0=\gamma_0.
\]
Dla funktorów silnych równoważne warunki otrzymujemy używając odwrotnych morfizmów strukturalnych.
\end{definicja}

\begin{stwierdzenie}[Składanie funktorów monoidalnych]\label{cat:prop:monoidal-composition}
Niech $F:\mathcal C\to\mathcal D$ oraz $G:\mathcal D\to\mathcal E$ będą funktorami typu lax o strukturach odpowiednio $\phi$ i $\gamma$. Wówczas $GF$ ma strukturę typu lax
\[
\theta_{A,B}=G(\phi_{A,B})\gamma_{F(A),F(B)},\qquad
\theta_0=G(\phi_0)\gamma_0.
\]
Jeśli oba funktory są silne (odpowiednio ścisłe), to ich złożenie również jest silne (odpowiednio ścisłe).
\end{stwierdzenie}
\begin{proof}
Typy morfizmów są następujące:
\[
GF(A)\otimes GF(B)\xrightarrow{\gamma_{F(A),F(B)}}
G(F(A)\otimes F(B))\xrightarrow{G(\phi_{A,B})}GF(A\otimes B).
\]
Morfizm jednostki prowadzi z $I_{\mathcal E}$ przez $G(I_{\mathcal D})$ do $GF(I_{\mathcal C})$. Naturalność $\theta$ wynika z naturalności $\phi,\gamma$: dla $a:A\to A'$, $b:B\to B'$ mamy
\[
\begin{aligned}
GF(a\otimes b)\theta_{A,B}
&=G(\phi_{A',B'})G(F(a)\otimes F(b))\gamma_{F(A),F(B)}\\
&=G(\phi_{A',B'})\gamma_{F(A'),F(B')}(GF(a)\otimes GF(b)).
\end{aligned}
\]
W rachunku łączności zapisujemy $X=F(A),Y=F(B),Z=F(C)$ oraz $P=F(A\otimes B)$, $Q=F(B\otimes C)$. Wstawiając definicję $\theta$ i stosując naturalność $\gamma$ względem $\phi_{A,B}:X\otimes Y\to P$, otrzymujemy
\[
\begin{aligned}
&GF(\alpha_{A,B,C})\theta_{A\otimes B,C}(\theta_{A,B}\otimes\id)\\
&=G(F(\alpha_{A,B,C})\phi_{A\otimes B,C}(\phi_{A,B}\otimes\id_Z))\\
&\qquad{}\cdot\gamma_{X\otimes Y,Z}(\gamma_{X,Y}\otimes\id_{G(Z)})\\
&=G(\phi_{A,B\otimes C}(\id_X\otimes\phi_{B,C}))\\
&\qquad{}\cdot G(\alpha_{X,Y,Z})\gamma_{X\otimes Y,Z}(\gamma_{X,Y}\otimes\id_{G(Z)})\\
&=G(\phi_{A,B\otimes C})G(\id_X\otimes\phi_{B,C})\\
&\qquad{}\cdot\gamma_{X,Y\otimes Z}(\id_{G(X)}\otimes\gamma_{Y,Z})
\alpha_{G(X),G(Y),G(Z)}\\
&=G(\phi_{A,B\otimes C})\gamma_{X,Q}
(\id_{G(X)}\otimes G(\phi_{B,C}))\\
&\qquad{}\cdot(\id_{G(X)}\otimes\gamma_{Y,Z})\alpha_{G(X),G(Y),G(Z)}\\
&=\theta_{A,B\otimes C}(\id_{GF(A)}\otimes\theta_{B,C})
\alpha_{GF(A),GF(B),GF(C)}.
\end{aligned}
\]
Kolejne równości wykorzystują naturalność $\gamma$, aksjomat łączności $F$, aksjomat łączności $G$, naturalność $\gamma$ i bifunktorialność. Wszystkie złożenia mają wspólne końce $(GF(A)\otimes GF(B))\otimes GF(C)$ i $GF(A\otimes(B\otimes C))$.

Dla lewej jednostki naturalność $\gamma$ względem $\phi_0$ daje
\[
\begin{aligned}
GF(\lambda_A)\theta_{I_{\mathcal C},A}(\theta_0\otimes\id)
&=G(F(\lambda_A)\phi_{I_{\mathcal C},A}(\phi_0\otimes\id_X))\\
&\quad{}\cdot\gamma_{I_{\mathcal D},X}(\gamma_0\otimes\id_{G(X)})\\
&=G(\lambda_X)\gamma_{I_{\mathcal D},X}(\gamma_0\otimes\id_{G(X)})\\
&=\lambda_{G(X)}.
\end{aligned}
\]
Analogiczny rachunek dla prawej jednostki, z podaniem odpowiednich czynników, to
\[
\begin{aligned}
GF(\rho_A)\theta_{A,I_{\mathcal C}}(\id\otimes\theta_0)
&=G(F(\rho_A)\phi_{A,I_{\mathcal C}}(\id_X\otimes\phi_0))\\
&\quad{}\cdot\gamma_{X,I_{\mathcal D}}(\id_{G(X)}\otimes\gamma_0)\\
&=G(\rho_X)\gamma_{X,I_{\mathcal D}}(\id_{G(X)}\otimes\gamma_0)\\
&=\rho_{G(X)}.
\end{aligned}
\]
Funktory zachowują izomorfizmy i identyczności, więc złożenia struktur silnych lub ścisłych mają tę samą własność.
\end{proof}

\begin{stwierdzenie}[Składanie monoidalnych transformacji]\label{cat:prop:monoidal-transform-composition}
Identyczności oraz złożenia pionowe i poziome monoidalnych transformacji naturalnych między funktorami typu lax są monoidalne, przy strukturach funktorów złożonych z poprzedniego stwierdzenia.
\end{stwierdzenie}
\begin{proof}
Dla identyczności warunki są bezpośrednimi równościami struktur. Niech $\eta:F\Rightarrow G$, $\zeta:G\Rightarrow H$ mają struktury odpowiednio $\phi,\gamma,\chi$. Wtedy
\[
(\zeta\eta)_{A\otimes B}\phi_{A,B}
=\zeta_{A\otimes B}\gamma_{A,B}(\eta_A\otimes\eta_B)
=\chi_{A,B}((\zeta_A\eta_A)\otimes(\zeta_B\eta_B)),
\]
a $(\zeta\eta)_I\phi_0=\zeta_I\gamma_0=\chi_0$.
Dla składania poziomego wystarczy sprawdzić oba rodzaje składania z identycznością funktora, ponieważ dowolne złożenie poziome jest ich złożeniem pionowym.
Jeżeli $K:\mathcal D\to\mathcal E$ ma strukturę $\kappa$, to
\[
\begin{aligned}
K(\eta_{A\otimes B})K(\phi_{A,B})\kappa_{F(A),F(B)}
&=K(\gamma_{A,B})K(\eta_A\otimes\eta_B)\kappa_{F(A),F(B)}\\
&=K(\gamma_{A,B})\kappa_{G(A),G(B)}(K(\eta_A)\otimes K(\eta_B)).
\end{aligned}
\]
Jednostka spełnia $K(\eta_I)K(\phi_0)\kappa_0=K(\gamma_0)\kappa_0$.
Jeżeli $L:\mathcal B\to\mathcal C$ ma strukturę $\ell$, to naturalność $\eta$ względem $\ell_{A,B}:L(A)\otimes L(B)\to L(A\otimes B)$ daje
\[
\begin{aligned}
\eta_{L(A\otimes B)}F(\ell_{A,B})\phi_{L(A),L(B)}
&=G(\ell_{A,B})\eta_{L(A)\otimes L(B)}\phi_{L(A),L(B)}\\
&=G(\ell_{A,B})\gamma_{L(A),L(B)}(\eta_{L(A)}\otimes\eta_{L(B)}).
\end{aligned}
\]
Podobnie $\eta_{L(I)}F(\ell_0)\phi_0=G(\ell_0)\eta_{I_{\mathcal C}}\phi_0=G(\ell_0)\gamma_0$.
Otrzymaliśmy wszystkie wymagane warunki.
\end{proof}

\begin{definicja}[Równoważność monoidalna]\label{cat:def:monoidal-equivalence}
Równoważnością monoidalną nazywamy parę silnych funktorów monoidalnych w przeciwnych kierunkach wraz z monoidalnymi izomorfizmami naturalnymi ich złożeń z odpowiednimi funktorami identycznościowymi. Funktory identycznościowe mają ścisłą strukturę monoidalną.
\end{definicja}

\subsection{Uściślanie kategorii monoidalnej}
\begin{twierdzenie}[Uściślanie Mac Lane'a]\label{cat:thm:strictification}
Każda lokalnie mała kategoria monoidalna jest silnie monoidalnie równoważna pewnej ścisłej kategorii monoidalnej, przy ustalonej konwencji wielkości.
\end{twierdzenie}
\begin{proof}
Niech obiektami $\mathcal S$ będą skończone listy obiektów $\mathcal C$, także lista pusta. Dla listy $u$ definiujemy $E(u)$ jako prawostronnie nawiasowany iloczyn jej elementów, przyjmując $E(())=I$ oraz $E((A))=A$. Definiujemy $\Hom_{\mathcal S}(u,v)=\Hom_{\mathcal C}(E(u),E(v))$ (morfizmy mają oznaczone końce $u,v$), składanie i identyczności odziedziczone z $\mathcal C$. Aksjomaty kategorii wynikają z aksjomatów $\mathcal C$.

Iloczyn obiektów to konkatenacja $uv$, jednostka to lista pusta. Z twierdzenia~\ref{cat:thm:mac-lane-coherence} otrzymujemy kanoniczne izomorfizmy $m_{u,v}:E(u)\otimes E(v)\to E(uv)$, stosując formalne litery dla pozycji listy. Dla morfizmów $f:u\to u'$, $g:v\to v'$ definiujemy
\[
f\otimes_{\mathcal S}g=m_{u',v'}(f\otimes g)m_{u,v}^{-1}.
\]
Jest to morfizm $uv\to u'v'$. Zachowanie identyczności wynika z $m_{u,v}m_{u,v}^{-1}=\id$. Dla dalszych morfizmów $f',g'$ środkowe izomorfizmy $m_{u',v'}^{-1}m_{u',v'}$ skracają się; bifunktorialność w $\mathcal C$ daje $(f'\otimes_{\mathcal S}g')(f\otimes_{\mathcal S}g)=(f'f)\otimes_{\mathcal S}(g'g)$.

Oznaczmy
\[
L_{u,v,w}=m_{uv,w}(m_{u,v}\otimes\id),\qquad
R_{u,v,w}=m_{u,vw}(\id\otimes m_{v,w}).
\]
Koherencja daje $L_{u,v,w}=R_{u,v,w}\alpha_{E(u),E(v),E(w)}$. Dla $h:w\to w'$ rozwinięcie definicji daje
\[
\begin{aligned}
(f\otimes_{\mathcal S}g)\otimes_{\mathcal S}h
&=L_{u',v',w'}((f\otimes g)\otimes h)L_{u,v,w}^{-1},\\
f\otimes_{\mathcal S}(g\otimes_{\mathcal S}h)
&=R_{u',v',w'}(f\otimes(g\otimes h))R_{u,v,w}^{-1}.
\end{aligned}
\]
Po podstawieniu $L=R\alpha$ pierwsze wyrażenie redukuje się do drugiego przez naturalność $\alpha$. Ponadto $m_{(),u}=\lambda_{E(u)}$, $m_{u,()}=\rho_{E(u)}$ z koherencji; naturalność jednostek daje $\id_{()}\otimes_{\mathcal S}f=f=f\otimes_{\mathcal S}\id_{()}$. Wobec tego $\mathcal S$ jest ściśle monoidalna.

Funktor $E:\mathcal S\to\mathcal C$ jest identycznością na odpowiednich zbiorach morfizmów, zatem jest w pełni wierny. Każdy $A$ jest obrazem listy $(A)$. Struktura silna to $m_{u,v}$ i $\id_I$. Naturalność $m$ wynika z definicji iloczynu morfizmów, aksjomat łączności to $L=R\alpha$, a aksjomaty jednostek wynikają z podanych wzorów dla pustej listy.

Jawny quasi-odwrotny funktor $J:\mathcal C\to\mathcal S$ wysyła $A$ na $(A)$ i $f$ na ten sam morfizm w odpowiednim zbiorze. Jego struktura $J(A)J(B)\to J(A\otimes B)$ oraz $()\to J(I)$ jest reprezentowana w $\mathcal C$ przez odpowiednie identyczności. Sprawdźmy jawnie aksjomaty $J$. Oznaczmy jego porównania przez $j_{A,B}$ i $j_0$. Z definicji $m_{(A),(B)}=\id_{A\otimes B}$. Naturalność $j$ po zastosowaniu wiernego $E$ sprowadza się więc do równości $a\otimes b=a\otimes b$ dla morfizmów $a,b$ w $\mathcal C$.

Dla list $u=(A),v=(B),w=(C)$ mamy $m_{uv,w}=\alpha_{A,B,C}$ i $m_{u,vw}=\id_{A\otimes(B\otimes C)}$; pozostałe porównania pojedynczych list są identycznościami. Dlatego
\[
\begin{aligned}
E(j_{A,B}\otimes_{\mathcal S}\id_{J(C)})&=\alpha_{A,B,C}^{-1},\\
E(\id_{J(A)}\otimes_{\mathcal S}j_{B,C})&=\id_{A\otimes(B\otimes C)}.
\end{aligned}
\]
Po zastosowaniu $E$ lewa droga aksjomatu łączności $J$ daje $\alpha_{A,B,C}\alpha_{A,B,C}^{-1}$, prawa identyczność (łączność $\mathcal S$ jest ścisła). Dla jednostki $E(j_0\otimes_{\mathcal S}\id)=\lambda_A^{-1}$, gdyż $m_{(),(A)}=\lambda_A$ i $m_{(I),(A)}=\id$. Podobnie $E(\id\otimes_{\mathcal S}j_0)=\rho_A^{-1}$. Złożenia z $E(J(\lambda_A))$ i $E(J(\rho_A))$ dają identyczności. Wierność $E$ dowodzi wszystkich aksjomatów; struktura $J$ jest silna, ponieważ jej składowe są reprezentowane przez identyczności. Mamy $EJ=\id_{\mathcal C}$ jako funktory silne: według stwierdzenia~\ref{cat:prop:monoidal-composition} struktura złożenia to $E(j_{A,B})m_{(A),(B)}=\id$ oraz $E(j_0)\id_I=\id_I$. Izomorfizm $JE\Rightarrow\id_{\mathcal S}$ w $u$ jest reprezentowany przez $\id_{E(u)}$, jest naturalny z definicji składania. Warunek monoidalności dla pary list po zastosowaniu $E$ ma po obu stronach morfizm $m_{u,v}:E(u)\otimes E(v)\to E(uv)$; warunek jednostki ma po obu stronach $\id_I$. Wierność $E$ daje wymagane równości. Otrzymaliśmy monoidalną równoważność.
\end{proof}

\begin{stwierdzenie}[Porównania kartezjańskie]\label{cat:prop:cartesian-oplax}
Każdy funktor między kategoriami z wybranymi skończonymi iloczynami, wyposażonymi w struktury kartezjańskie, ma kanoniczną strukturę oplax. Jest ona odwracalna dokładnie wtedy, gdy funktor zachowuje skończone iloczyny.
\end{stwierdzenie}
\begin{proof}
Definiujemy $\psi_{A,B}=\langle F(p_A),F(p_B)\rangle$ oraz $\psi_0$ jako jedyny morfizm do obiektu końcowego. Naturalność $\psi$ wynika z porównania obu rzutów i funktorialności $F$. Obie drogi w aksjomacie łączności mają składowe $F(p_A),F(p_B),F(p_C)$ dla złożonych rzutów z $(A\times B)\times C$, ponieważ kanoniczny izomorfizm łączności zachowuje te rzuty. Jednoznaczność morfizmu do iterowanego iloczynu daje równość. W aksjomatach jednostki po złożeniu z odpowiednim rzutem otrzymujemy $F(\lambda_A)$ lub $F(\rho_A)$, co daje wymaganą równość. Porównanie $\psi_{A,B}$ jest izomorfizmem wtedy i tylko wtedy, gdy $F(A\times B)$ z rzutami $F(p_A),F(p_B)$ jest iloczynem. Analogicznie $\psi_0$ jest izomorfizmem wtedy i tylko wtedy, gdy $F(1)$ jest końcowy. Skończone iloczyny otrzymujemy przez iterację binarnych oraz pusty iloczyn. To kończy dowód.
\end{proof}

\begin{definicja}[Zachowywanie plecenia]\label{cat:def:braided-functor}
Silny funktor monoidalny między kategoriami z pleceniem zachowuje plecenie, jeśli
\[
F(c^{\mathcal C}_{A,B})\phi_{A,B}
=\phi_{B,A}c^{\mathcal D}_{F(A),F(B)}.
\]
Obie strony są morfizmami $F(A)\otimes F(B)\to F(B\otimes A)$. Gdy obie kategorie są symetryczne, taki funktor nazywamy symetrycznym monoidalnym.
\end{definicja}

\begin{przyklad}
Funktor zachowujący skończone iloczyny z poprzedniego stwierdzenia, wyposażony w $\phi=\psi^{-1}$, jest symetryczny monoidalny: równoważna równość $\psi_{B,A}F(c_{A,B})=c_{F(A),F(B)}\psi_{A,B}$ zachodzi, ponieważ obie strony mają składowe $F(p_B),F(p_A)$.
\end{przyklad}

\subsection{Koherencja symetryczna}
\begin{lemat}[Przeniesienie symetrii do kategorii ścisłej]\label{cat:lem:strict-symmetry}
Jeżeli $\mathcal C$ jest symetryczna monoidalna, to jej uściślenie $\mathcal S$ z twierdzenia~\ref{cat:thm:strictification} jest ścisłe symetryczne monoidalne dla
\[
E(c^{\mathcal S}_{u,v})=m_{v,u}c_{E(u),E(v)}m_{u,v}^{-1}.
\]
Funktor $E$ zachowuje symetrię.
\end{lemat}
\begin{proof}
Prawa strona ma typ $E(uv)\to E(vu)$, więc określa dokładnie jeden morfizm w $\mathcal S$. Naturalność po podstawieniu wzoru na iloczyn morfizmów redukuje się do naturalności $c$ w $\mathcal C$: dla $f:u\to u'$, $g:v\to v'$ obie strony, po zastosowaniu $E$, są równe
\[
m_{v',u'}(g\otimes f)c_{E(u),E(v)}m_{u,v}^{-1}.
\]
Podwójna zamiana daje $m_{u,v}c_{E(v),E(u)}c_{E(u),E(v)}m_{u,v}^{-1}=\id$. Warunek zachowania symetrii przez $E$ jest dokładnie definicją $c^{\mathcal S}$.

Sprawdzimy sześciokąty. Zapiszmy $X=E(u)$, $Y=E(v)$, $Z=E(w)$ i użyjmy porównań $L_{u,v,w}$ z dowodu uściślania, identyfikujących lewostronnie nawiasowane trójki z wartościami konkatenacji. Po zastosowaniu $E$ i sprzężeniu morfizmów tymi porównaniami, tj. przez złożenie po prawej z $L$ dla źródła i po lewej z $L^{-1}$ dla celu, dostajemy dla pierwszego ścisłego sześciokąta:
\[
\begin{aligned}
c^{\mathcal S}_{u,vw}&\ \longmapsto\
 c_{X,Y\otimes Z}\alpha_{X,Y,Z},\\
(\id_v\otimes_{\mathcal S} c^{\mathcal S}_{u,w})
(c^{\mathcal S}_{u,v}\otimes_{\mathcal S}\id_w)
&\ \longmapsto\
 \alpha_{Y,Z,X}^{-1}(\id_Y\otimes c_{X,Z})\\
&\qquad{}\cdot\alpha_{Y,X,Z}(c_{X,Y}\otimes\id_Z).
\end{aligned}
\]
Istotnie, dla pierwszego morfizmu $L_{v,w,u}^{-1}m_{vw,u}=m_{v,w}^{-1}\otimes\id_X$ i $m_{u,vw}^{-1}L_{u,v,w}=(\id_X\otimes m_{v,w})\alpha_{X,Y,Z}$. Naturalność $c$ skraca $m_{v,w}$ z jego odwrotnością. W drugiej drodze pierwszy krok sprzęgamy przez $L_{v,u,w}^{-1}$ i $L_{u,v,w}$, co daje $c_{X,Y}\otimes\id_Z$. Dla drugiego kroku przechodzimy do prawostronnego nawiasowania przez $\alpha_{Y,X,Z}$, stosujemy $\id_Y\otimes c_{X,Z}$ i wracamy przez $\alpha_{Y,Z,X}^{-1}$; wynika to ze wzoru $L=R\alpha$. Równość otrzymanych wyrażeń jest pierwszym sześciokątem z CAT-19 po złożeniu z $\alpha_{Y,Z,X}^{-1}$.

Dla drugiego ścisłego sześciokąta ten sam wzór $L=R\alpha$ daje
\[
\begin{aligned}
c^{\mathcal S}_{uv,w}&\ \longmapsto\
 \alpha_{Z,X,Y}^{-1}c_{X\otimes Y,Z},\\
(c^{\mathcal S}_{u,w}\otimes_{\mathcal S}\id_v)
(\id_u\otimes_{\mathcal S} c^{\mathcal S}_{v,w})
&\ \longmapsto\
(c_{X,Z}\otimes\id_Y)\alpha_{X,Z,Y}^{-1}\\
&\qquad{}\cdot(\id_X\otimes c_{Y,Z})\alpha_{X,Y,Z}.
\end{aligned}
\]
W pierwszej drodze $L_{w,u,v}^{-1}m_{w,uv}=\alpha_{Z,X,Y}^{-1}(\id_Z\otimes m_{u,v}^{-1})$ i $m_{uv,w}^{-1}L_{u,v,w}=m_{u,v}\otimes\id_Z$, więc naturalność skraca porównania. W drugiej drodze pierwszy krok wymaga zmiany nawiasowania, a drugi działa na dwóch pierwszych czynnikach. Drugi sześciokąt z CAT-19, złożony po prawej z $\alpha_{X,Y,Z}$, daje równość. W obu przypadkach $E$ jest wierny, a porównania odwracalne, zatem ścisłe sześciokąty zachodzą również w $\mathcal S$. Rachunki obejmują puste listy, ponieważ wzór $L=R\alpha$ został dowiedziony także dla nich.
\end{proof}

\begin{lemat}[Relacje sąsiednich zamian]\label{cat:lem:coxeter-swaps}
W ścisłej symetrycznej kategorii monoidalnej zamiana czynników na pozycjach $i,i+1$, dana przez $c$ i identyczności na pozostałych pozycjach, spełnia relacje
\[
s_i^2=1,\qquad s_i s_j=s_j s_i\ (|i-j|>1),\qquad
s_i s_{i+1}s_i=s_{i+1}s_i s_{i+1}.
\]
W każdym kroku $s_i$ oznacza zamianę aktualnych czynników, a nie ustalony endomorfizm niezależny od słowa.
\end{lemat}
\begin{proof}
Pierwsza relacja jest symetrią, druga wynika z bifunktorialności dla rozłącznych par pozycji. Trzecia sprowadza się do trzech kolejnych czynników $A,B,C$. Naturalność $c$ względem morfizmu $c_{A,B}:A\otimes B\to B\otimes A$ i $\id_C$ daje
\[
c_{B\otimes A,C}(c_{A,B}\otimes\id_C)
=(\id_C\otimes c_{A,B})c_{A\otimes B,C}.
\]
Rozwinięcie obu zamian blokowych za pomocą drugiego ścisłego sześciokąta daje
\[
\begin{aligned}
&(c_{B,C}\otimes\id_A)(\id_B\otimes c_{A,C})(c_{A,B}\otimes\id_C)\\
&=(\id_C\otimes c_{A,B})(c_{A,C}\otimes\id_B)(\id_A\otimes c_{B,C}).
\end{aligned}
\]
To dokładnie trzecia relacja; obie drogi prowadzą z $A\otimes B\otimes C$ do $C\otimes B\otimes A$. Iloczyn z identycznościami przenosi tę równość na dowolne trzy kolejne pozycje.
\end{proof}

\begin{lemat}[Prezentacja grupy permutacji]\label{cat:lem:symmetric-presentation}
Dwa słowa w sąsiednich transpozycjach $s_1,\ldots,s_{n-1}$ wyznaczają tę samą permutację $n$ pozycji wtedy i tylko wtedy, gdy można przejść od jednego do drugiego przez relacje z lematu~\ref{cat:lem:coxeter-swaps}.
\end{lemat}
\begin{proof}
Dla $n=0,1$ twierdzenie jest bezpośrednie. Rozważmy grupę $G_n$ określoną przez podane generatory i relacje. Ponieważ $s_i^2=1$, odwrotności generatorów można zastąpić generatorami. Istnieje surjektywny homomorfizm $G_n\to S_n$, gdyż sąsiednie transpozycje spełniają relacje i generują $S_n$ (dowolną literę można przesunąć na żądane miejsce przez kolejne zamiany sąsiadów).

Niech $H$ będzie podgrupą generowaną przez $s_1,\ldots,s_{n-2}$. Jest obrazem $G_{n-1}$, zatem z indukcji ma co najwyżej $(n-1)!$ elementów. Definiujemy $r_j=s_{n-1}s_{n-2}\cdots s_j$ dla $1\le j<n$ oraz $r_n=1$. Pokażemy, że każdy element ma postać $hr_j$, $h\in H$. Zbiór tych postaci zawiera identyczność i jest zamknięty na mnożenie z prawej przez dowolny $s_i$. Dla $j<n$ mamy
\[
r_js_i=
\begin{cases}
s_i r_j,&i<j-1,\\
r_{j-1},&i=j-1,\\
r_{j+1},&i=j,\\
s_{i-1}r_j,&i>j.
\end{cases}
\]
Pierwszy przypadek wynika z przemienności rozłącznych zamian, drugi z definicji, trzeci z $s_j^2=1$. W ostatnim przypadku przenosimy prawy $s_i$ przez czynniki $s_j,\ldots,s_{i-2}$, otrzymując fragment $s_i s_{i-1}s_i$, zamieniamy go na $s_{i-1}s_i s_{i-1}$ i przenosimy pierwszy $s_{i-1}$ w lewo przez $s_{i+1},\ldots,s_{n-1}$. Dla $j=n$ mamy $r_ns_i=s_i r_n$ przy $i<n-1$ oraz $r_ns_{n-1}=r_{n-1}$. Wszystkie dodatkowe generatory po lewej należą do $H$. Indukcja po długości słowa dowodzi żądanej postaci.
Stąd $|G_n|\le n|H|\le n!$. Surjekcja na $S_n$ wymusza równość i jest izomorfizmem. Równość w grupie z prezentacji oznacza dokładnie możliwość stosowania relacji (w tym w dowolnym kontekście słowa), co daje tezę.
\end{proof}

\begin{twierdzenie}[Koherencja symetryczna]\label{cat:thm:symmetric-coherence}
W lokalnie małej symetrycznej kategorii monoidalnej dwa morfizmy strukturalne o tych samych formalnych końcach, zbudowane z $\alpha,\lambda,\rho,c$, ich odwrotności, identyczności, składania i iloczynu monoidalnego, są równe, jeśli wyznaczają tę samą permutację pozycji liter po pominięciu formalnych jednostek. Każda permutacja zgodna z literami źródła i celu jest realizowalna przez taki morfizm.
\end{twierdzenie}
\begin{proof}
Najpierw pracujemy w kategorii ścisłej. Ścisły pierwszy sześciokąt z dwoma czynnikami jednostkowymi daje $c_{A,I}=c_{A,I}c_{A,I}$, więc z odwracalności $c_{A,I}=\id_A$. Drugi sześciokąt daje analogicznie $c_{I,A}=\id_A$. Zamiana pustego bloku jest zatem identycznością.
Każdą zamianę dwóch niepustych bloków można rozłożyć na zamiany sąsiednich pojedynczych czynników: pierwszy sześciokąt rozkłada zamianę z iloczynem w drugim argumencie, a drugi rozkłada zamianę z iloczynem w pierwszym. Indukcja względem łącznej liczby czynników obu bloków (z przypadkiem pustym i parą jednoelementową jako bazami) daje rozkład. Odwrotność zamiany jest zamianą w przeciwnym kierunku przez symetrię. Iloczyn morfizmów rozkładamy na działanie kolejno w dwóch blokach przez bifunktorialność. W ścisłej kategorii pozostałe morfizmy strukturalne są identycznościami. Otrzymujemy słowo zamian sąsiadów, realizujące dokładnie permutację pierwotnego morfizmu. Dwa słowa o tej samej permutacji są połączone relacjami z lematu~\ref{cat:lem:symmetric-presentation}. Każda relacja daje równość morfizmów dla aktualnej kolejności czynników na mocy lematu~\ref{cat:lem:coxeter-swaps}. Złożenie z wcześniejszymi i późniejszymi zamianami zachowuje równość. Każdą permutację można uzyskać przez zamiany sąsiadów.

Dla ogólnej $\mathcal C$ stosujemy uściślenie $\mathcal S$ i lemat~\ref{cat:lem:strict-symmetry}. Dla formalnego wyrażenia $T$ oznaczmy przez $u_T$ listę obiektów odpowiadających literom i przez $\nu_T:T\to E(u_T)$ kanoniczny izomorfizm z twierdzenia~\ref{cat:thm:mac-lane-coherence}. Dla morfizmu strukturalnego $f:T\to T'$ przenosimy go do $\mathcal S$ jako morfizm reprezentowany przez $\nu_{T'}f\nu_T^{-1}$. Składanie przenosi się bezpośrednio przez skrócenie $\nu$ z odwrotnością. Z koherencji monoidalnej
\[
\nu_{T\otimes U}=m_{u_T,u_U}(\nu_T\otimes\nu_U).
\]
Dlatego przeniesienie iloczynu dwóch morfizmów jest iloczynem ich przeniesień w $\mathcal S$. Przeniesienia $\alpha,\lambda,\rho$ oraz ich odwrotności są identycznościami, gdyż normalizacja nie zależy od drogi. Dla zamiany $c_{T,U}$ naturalność $c$ i powyższa równość dają
\[
\nu_{U\otimes T}c_{T,U}\nu_{T\otimes U}^{-1}
=m_{u_U,u_T}c_{E(u_T),E(u_U)}m_{u_T,u_U}^{-1},
\]
czyli dokładnie reprezentanta $c^{\mathcal S}_{u_T,u_U}$. Indukcja po budowie morfizmu pokazuje, że jego przeniesienie jest morfizmem strukturalnym w ścisłej kategorii z tą samą permutacją. Tam wykazaliśmy jednoznaczność. Ponieważ $\nu_T,\nu_{T'}$ są izomorfizmami, równość przeniesień daje równość pierwotnych morfizmów. Istnienie otrzymujemy przez normalizację, realizację sąsiednich zamian z kanonicznymi zmianami nawiasowania i odwrotną normalizację. Gdy litery oznaczają ten sam obiekt, nadal rozróżniamy ich pozycje; sama równość słów obiektów nie wystarcza do stwierdzenia równości morfizmów.
\end{proof}

\section{Monoidy, komonoidy i dualność}\label{cat:sec:monoids-duality}
Korzystamy z twierdzeń~\ref{cat:thm:mac-lane-coherence} i \ref{cat:thm:strictification}. Rachunki prowadzimy w zapisie ścisłym: w ogólnej kategorii należy wstawić kanoniczne izomorfizmy łączności i jednostek. Kolejności czynników nie wolno zmieniać bez odpowiedniego morfizmu.

\subsection{Monoidy i komonoidy}
\begin{definicja}[Monoid]\label{cat:def:monoid-object}
Monoid w kategorii monoidalnej to obiekt $M$ z mnożeniem $m:M\otimes M\to M$ i jednostką $e:I\to M$, spełniającymi
\[
m(m\otimes\id_M)=m(\id_M\otimes m),\qquad
m(e\otimes\id_M)=\id_M=m(\id_M\otimes e).
\]
W zapisie nieścisłym pierwsza równość ma po prawej dodatkowy czynnik $\alpha_{M,M,M}$, a dwie równości jednostkowe mają po prawej odpowiednio $\lambda_M,\rho_M$. Morfizm monoidów $f:M\to N$ spełnia $fm_M=m_N(f\otimes f)$, $fe_M=e_N$. Identyczności spełniają te warunki, a dla kolejnego morfizmu $g$ mamy $gfm_M=gm_N(f\otimes f)=m_P(gf\otimes gf)$ i $gfe_M=e_P$. W kategorii z pleceniem monoid jest przemienny, gdy $mc_{M,M}=m$.
\end{definicja}
\begin{przyklad}[Monoidy i algebry]\label{cat:ex:monoids-algebras}
W $(\Set,\times,1)$ są to zwykłe monoidy. W $(\Vect_k,\otimes,k)$ własność uniwersalna iloczynu tensorowego daje odpowiednio dwuliniowe mnożenie i jednostkę $e(1)$. Aksjomaty sprawdzone na tensorach prostych dają dokładnie łączność i obie jednostki algebry; odwrotnie, równość na tensorach prostych daje równość liniowych morfizmów. Morfizmy monoidów są homomorfizmami algebr zachowującymi jedynkę.
\end{przyklad}
\begin{definicja}[Komonoid]\label{cat:def:comonoid-object}
Komonoid to obiekt $C$ z morfizmami $\Delta:C\to C\otimes C$ i $\epsilon:C\to I$ takimi, że
\[
(\Delta\otimes\id_C)\Delta=(\id_C\otimes\Delta)\Delta,\qquad
(\epsilon\otimes\id_C)\Delta=\id_C=(\id_C\otimes\epsilon)\Delta.
\]
W zapisie nieścisłym aksjomaty przyjmują postać
\[
\alpha_{C,C,C}(\Delta\otimes\id_C)\Delta=(\id_C\otimes\Delta)\Delta,\qquad
\lambda_C(\epsilon\otimes\id_C)\Delta=\id_C=\rho_C(\id_C\otimes\epsilon)\Delta.
\]
Morfizm komonoidów $f:C\to D$ spełnia $\Delta_Df=(f\otimes f)\Delta_C$ i $\epsilon_Df=\epsilon_C$. Identyczności spełniają te równości. Jeśli $g:D\to E$ jest kolejnym morfizmem komonoidów, to
\[
\Delta_Egf=(g\otimes g)\Delta_Df=(g\otimes g)(f\otimes f)\Delta_C
=(gf\otimes gf)\Delta_C,\qquad \epsilon_Egf=\epsilon_Df=\epsilon_C.
\]
W kategorii z pleceniem koprzemienność oznacza $c_{C,C}\Delta=\Delta$. Komonoid jest monoidem w kategorii przeciwnej z odwróconymi izomorfizmami strukturalnymi.
\end{definicja}
\begin{twierdzenie}[Komonoidy kartezjańskie]\label{cat:thm:cartesian-comonoid}
W kategorii z wybranymi skończonymi iloczynami każdy obiekt $X$ ma dokładnie jedną strukturę komonoidu dla iloczynu kartezjańskiego. Jest ona koprzemienna i zachowywana przez wszystkie morfizmy.
\end{twierdzenie}
\begin{proof}
Kojednostka jest wymuszona przez końcowość $1$. Dwie równości kojednostki, po wstawieniu izomorfizmów jednostkowych, dają $p_1\Delta=p_2\Delta=\id_X$. Własność uniwersalna iloczynu wymusza $\Delta=\langle\id_X,\id_X\rangle$. Obie drogi kołączności mają trzy składowe $\id_X$, a zamiana zachowuje obie składowe przekątnej. Dla $f:X\to Y$ oba morfizmy $\Delta_Y f,(f\times f)\Delta_X$ mają składowe $f,f$, a $\epsilon_Y f=\epsilon_X$ przez końcowość. To dowodzi istnienia, jednoznaczności i zachowania struktury.
\end{proof}
\begin{przyklad}
Dla zbioru $S$ na przestrzeni $k^{(S)}$ o bazie $(e_s)$ określamy liniowo $\Delta(e_s)=e_s\otimes e_s$, $\epsilon(e_s)=1$. Obie drogi kołączności dają $e_s\otimes e_s\otimes e_s$, obie kojednostki dają $e_s$. Zamiana zachowuje $\Delta(e_s)$. Równość morfizmów liniowych na bazie daje wszystkie aksjomaty, także dla nieskończonego $S$.
\end{przyklad}

\subsection{Moduły i monady}
\begin{definicja}[Moduł nad monoidem]\label{cat:def:monoid-module}
Lewy moduł nad $(M,m,e)$ to obiekt $X$ z działaniem $a:M\otimes X\to X$ takim, że
\[
a(m\otimes\id_X)=a(\id_M\otimes a),\qquad a(e\otimes\id_X)=\id_X.
\]
Prawy moduł ma działanie $b:X\otimes M\to X$ z $b(b\otimes\id_M)=b(\id_X\otimes m)$ i $b(\id_X\otimes e)=\id_X$. Bez założenia ścisłości warunki te mają postać
\[
\begin{aligned}
a(m\otimes\id_X)&=a(\id_M\otimes a)\alpha_{M,M,X},&
a(e\otimes\id_X)&=\lambda_X,\\
b(b\otimes\id_M)&=b(\id_X\otimes m)\alpha_{X,M,M},&
b(\id_X\otimes e)&=\rho_X.
\end{aligned}
\]
Morfizm lewych modułów $f:X\to Y$ spełnia $fa_X=a_Y(\id_M\otimes f)$, a prawych $fb_X=b_Y(f\otimes\id_M)$. Identyczności są morfizmami modułów. Dla lewych modułów z bifunktorialności mamy
$gfa_X=ga_Y(\id_M\otimes f)=a_Z(\id_M\otimes g)(\id_M\otimes f)=a_Z(\id_M\otimes gf)$,
a dla prawych analogicznie
$gfb_X=gb_Y(f\otimes\id_M)=b_Z(g\otimes\id_M)(f\otimes\id_M)=b_Z(gf\otimes\id_M)$.
Otrzymujemy więc kategorie lewych i prawych modułów. W kategorii zbiorów są to działania monoidu, w kategorii przestrzeni wektorowych moduły unitarne nad algebrą: dwuliniowość i aksjomaty na tensorach prostych dają dokładnie zwykłe aksjomaty.
\end{definicja}
\begin{stwierdzenie}[Monada działania]\label{cat:prop:monoid-module-monad}
Monoid $M$ wyznacza monadę $T=M\otimes-$, której algebry Eilenberga--Moore'a są dokładnie lewymi modułami nad $M$.
\end{stwierdzenie}
\begin{proof}
Na morfizmach $T(f)=\id_M\otimes f$. W zapisie ścisłym $\eta_X=e\otimes\id_X$, $\mu_X=m\otimes\id_X$. Naturalność wynika z bifunktorialności: dla $f:X\to Y$ obie drogi dają odpowiednio $e\otimes f$ i $m\otimes f$. Łączność monady redukuje się do
\[
(m(m\otimes\id_M))\otimes\id_X=(m(\id_M\otimes m))\otimes\id_X.
\]
Obie tożsamości jednostkowe monady sprawdzamy oddzielnie:
\[
\begin{aligned}
\mu_X\eta_{TX}&=(m\otimes\id_X)(e\otimes\id_M\otimes\id_X)
=[m(e\otimes\id_M)]\otimes\id_X=\id_{TX},\\
\mu_XT(\eta_X)&=(m\otimes\id_X)(\id_M\otimes e\otimes\id_X)
=[m(\id_M\otimes e)]\otimes\id_X=\id_{TX}.
\end{aligned}
\] W zapisie nieścisłym $\eta_X=(e\otimes\id_X)\lambda_X^{-1}$ oraz $\mu_X=(m\otimes\id_X)\alpha_{M,M,X}^{-1}$; koherencja uzasadnia rachunki. Aksjomaty algebry z definicji~\ref{cat:def:em-algebra}, $a\eta_X=\id_X$, $a\mu_X=aT(a)$, stają się dokładnie aksjomatami działania. Morfizmy algebr spełniają $fa_X=a_YT(f)$, czyli warunek morfizmu modułów, z tymi samymi identycznościami i składaniem.
\end{proof}
\begin{stwierdzenie}[Monady jako monoidy endofunktorów]\label{cat:prop:monads-as-monoids}
Dla małej kategorii $\mathcal C$ monady są dokładnie monoidami w $[\mathcal C,\mathcal C]$ z iloczynem danym przez składanie.
\end{stwierdzenie}
\begin{proof}
W przykładzie~\ref{cat:ex:endofunctor-monoidal} wykazano ścisłą monoidalność tej kategorii. Mnożenie to $\mu:T^2\Rightarrow T$, jednostka to $\eta:\id_{\mathcal C}\Rightarrow T$. Składowe iloczynów poziomych $\mu\ast\id_T$, $\id_T\ast\mu$ to $\mu_{TX},T(\mu_X)$; dla $\eta$ analogicznie $\eta_{TX},T(\eta_X)$. Aksjomaty monoidu są więc równościami
\[
\mu_X\mu_{TX}=\mu_XT(\mu_X),\qquad
\mu_X\eta_{TX}=\id_{TX}=\mu_XT(\eta_X),
\]
czyli aksjomatami monady (definicja~\ref{def:monad}). Warunki morfizmu monoidów $\tau:T\Rightarrow S$ to
$\tau_X\eta_X^T=\eta_X^S$ i $\tau_X\mu_X^T=\mu_X^S S(\tau_X)\tau_{TX}$.
\end{proof}

\subsection{Obiekty dualne}
\begin{definicja}[Dualność]\label{cat:def:dual-object}
Przyjmujemy konwencję: obiekt $D$ jest prawostronnie dualny do $A$, gdy istnieją koewaluacja $i_A:I\to A\otimes D$ i ewaluacja $e_A:D\otimes A\to I$ takie, że
\begin{equation}\label{cat:eq:snake}
(\id_A\otimes e_A)(i_A\otimes\id_A)=\id_A,\qquad
(e_A\otimes\id_D)(\id_D\otimes i_A)=\id_D.
\end{equation}
Lewostronna dualność ma typy $I\to L\otimes A$, $A\otimes L\to I$ i odwrócone kolejności w tożsamościach. Nazwy stron bywają zamieniane w literaturze, więc rozstrzygające są typy.
W zapisie nieścisłym tożsamości \eqref{cat:eq:snake} brzmią odpowiednio
\[
\begin{aligned}
\rho_A(\id_A\otimes e_A)\alpha_{A,D,A}(i_A\otimes\id_A)\lambda_A^{-1}&=\id_A,\\
\lambda_D(e_A\otimes\id_D)\alpha_{D,A,D}^{-1}(\id_D\otimes i_A)\rho_D^{-1}&=\id_D.
\end{aligned}
\]
\end{definicja}
\begin{twierdzenie}[Sprzężenie wynikające z dualności]\label{cat:thm:dual-adjunction}
Jeśli $D$ jest prawostronnie dualny do $A$, to $D\otimes-\dashv A\otimes-$.
\end{twierdzenie}
\begin{proof}
Dla $f:D\otimes X\to Y$ i $g:X\to A\otimes Y$ definiujemy
\[
\Phi(f)=(\id_A\otimes f)(i_A\otimes\id_X),\qquad
\Psi(g)=(e_A\otimes\id_Y)(\id_D\otimes g).
\]
Bifunktorialność i druga tożsamość trójkątna dają
\[
\begin{aligned}
\Psi\Phi(f)
&=(e_A\otimes\id_Y)(\id_D\otimes\id_A\otimes f)(\id_D\otimes i_A\otimes\id_X)\\
&=f(e_A\otimes\id_D\otimes\id_X)(\id_D\otimes i_A\otimes\id_X)=f.
\end{aligned}
\]
Analogicznie, używając pierwszej tożsamości,
\[
\begin{aligned}
\Phi\Psi(g)
&=(\id_A\otimes e_A\otimes\id_Y)(\id_A\otimes\id_D\otimes g)(i_A\otimes\id_X)\\
&=(\id_A\otimes e_A\otimes\id_Y)(i_A\otimes\id_A\otimes\id_Y)g=g.
\end{aligned}
\]
Dla $u:X'\to X$, $v:Y\to Y'$ bifunktorialność daje
$\Phi(vf(\id_D\otimes u))=(\id_A\otimes v)\Phi(f)u$.
Jest to naturalność bijekcji w obu zmiennych.
\end{proof}
\begin{stwierdzenie}[Jednoznaczność dualności]\label{cat:prop:dual-unique}
Dwie prawe dualności $(D,i,e)$ i $(D',i',e')$ dla $A$ są związane jedynym izomorfizmem $h:D\to D'$ spełniającym
$e'(h\otimes\id_A)=e$ i $(\id_A\otimes h)i=i'$.
\end{stwierdzenie}
\begin{proof}
Bijekcja $\Hom(X,D')\cong\Hom(X\otimes A,I)$ jest dana przez $h\mapsto e'(h\otimes\id_A)$, z odwrotnością $b\mapsto(b\otimes\id_{D'})(\id_X\otimes i')$. Złożenie w pierwszym kierunku upraszcza się do $(e'\otimes\id_{D'})(\id_{D'}\otimes i')h=h$; w drugim do $b(\id_X\otimes[(\id_A\otimes e')(i'\otimes\id_A)])=b$. Zatem jedyny kandydat to $h=(e\otimes\id_{D'})(\id_D\otimes i')$. Mamy
\[
\begin{aligned}
(\id_A\otimes h)i
&=(\id_A\otimes e\otimes\id_{D'})(i\otimes i')\\
&=(\id_A\otimes e\otimes\id_{D'})(i\otimes\id_A\otimes\id_{D'})i'=i'.
\end{aligned}
\]
Zamieniając dualności, otrzymujemy $k:D'\to D$ zachowujące ewaluację i koewaluację. Złożenia $kh,hk$ zachowują ewaluacje, więc przez jednoznaczność są identycznościami.
\end{proof}
\begin{stwierdzenie}[Kontrawariantna funktorialność dualności]\label{cat:prop:dual-functoriality}
Wybierzmy prawe dualności $(D_A,i_A,e_A)$ dla rozpatrywanych obiektów. Każdemu morfizmowi $f:A\to B$ odpowiada jedyny morfizm $f^*:D_B\to D_A$ taki, że
\begin{equation}\label{cat:eq:dual-morphism-eval}
e_A(f^*\otimes\id_A)=e_B(\id_{D_B}\otimes f).
\end{equation}
Morfizm ten spełnia również
\begin{equation}\label{cat:eq:dual-morphism-coeval}
(f\otimes\id_{D_A})i_A=(\id_B\otimes f^*)i_B.
\end{equation}
Ponadto $(\id_A)^*=\id_{D_A}$ oraz $(gf)^*=f^*g^*$ dla $g:B\to C$. Zatem po wyborze dualności otrzymujemy funktor kontrawariantny z pełnej podkategorii obiektów mających prawe dualności do $\mathcal C$.
\end{stwierdzenie}
\begin{proof}
Bijekcja z dowodu stwierdzenia~\ref{cat:prop:dual-unique}, zastosowana do $X=D_B$ i dualności obiektu $A$, daje istnienie i jednoznaczność oraz wzór
\[
f^*=(e_B\otimes\id_{D_A})(\id_{D_B}\otimes f\otimes\id_{D_A})(\id_{D_B}\otimes i_A).
\]
Aby wykazać \eqref{cat:eq:dual-morphism-coeval}, stosujemy bijekcję
\[
\Hom(I,B\otimes D_A)\longrightarrow\Hom(A,B),\qquad
t\longmapsto(\id_B\otimes e_A)(t\otimes\id_A).
\]
Jej odwrotnością jest $u\mapsto(u\otimes\id_{D_A})i_A$: pierwsze złożenie upraszcza się do $u(\id_A\otimes e_A)(i_A\otimes\id_A)=u$, a drugie do
\[
(\id_B\otimes e_A\otimes\id_{D_A})(t\otimes i_A)
=(\id_B\otimes[(e_A\otimes\id_{D_A})(\id_{D_A}\otimes i_A)])t=t.
\]
Obraz lewej strony \eqref{cat:eq:dual-morphism-coeval} to $f$. Obraz prawej strony obliczamy, stosując kolejno bifunktorialność, \eqref{cat:eq:dual-morphism-eval}, ponownie bifunktorialność oraz pierwszą tożsamość trójkątną dla $B$:
\[
\begin{aligned}
&(\id_B\otimes e_A)(\id_B\otimes f^*\otimes\id_A)(i_B\otimes\id_A)\\
&\quad=(\id_B\otimes e_B)(\id_B\otimes\id_{D_B}\otimes f)(i_B\otimes\id_A)\\
&\quad=(\id_B\otimes e_B)(i_B\otimes\id_B)f=f.
\end{aligned}
\]
Iniektywność bijekcji daje żądaną równość. Identyczność $\id_{D_A}$ spełnia \eqref{cat:eq:dual-morphism-eval} dla $f=\id_A$. Dla złożenia, korzystając dwukrotnie z \eqref{cat:eq:dual-morphism-eval} i z bifunktorialności, mamy
\[
\begin{aligned}
e_A((f^*g^*)\otimes\id_A)
&=e_B(\id_{D_B}\otimes f)(g^*\otimes\id_A)\\
&=e_B(g^*\otimes\id_B)(\id_{D_C}\otimes f)\\
&=e_C(\id_{D_C}\otimes g)(\id_{D_C}\otimes f)\\
&=e_C(\id_{D_C}\otimes gf).
\end{aligned}
\]
Jednoznaczność daje obie równości funktorialności. Nie potrzeba tu plecenia ani symetrii.
\end{proof}

\begin{stwierdzenie}[Dualność iloczynu]\label{cat:prop:dual-tensor}
Jeśli $D_A,D_B$ są prawostronnie dualne do $A,B$, to $D_B\otimes D_A$ jest prawostronnie dualny do $A\otimes B$ dla
\[
i_{AB}=(\id_A\otimes i_B\otimes\id_{D_A})i_A,\quad
e_{AB}=e_B(\id_{D_B}\otimes e_A\otimes\id_B).
\]
\end{stwierdzenie}
\begin{proof}
Pierwsza tożsamość: od $AB$ przechodzimy przez $AD_AAB$, następnie $ABD_BD_AAB$, po kontrakcji $D_AA$ do $ABD_BB$, po kontrakcji $D_BB$ do $AB$. Dzięki bifunktorialności wstawienie $BD_B$ można wykonać po kontrakcji $D_AA$: operacje dotyczą rozłącznych czynników. Złożenie redukuje się zatem do
\[
(\id_A\otimes[(\id_B\otimes e_B)(i_B\otimes\id_B)])
([(\id_A\otimes e_A)(i_A\otimes\id_A)]\otimes\id_B)=\id_{AB}.
\]
Druga tożsamość: od $D_BD_A$ przechodzimy przez $D_BD_AAD_A$, następnie $D_BD_AABD_BD_A$, po kontrakcji $D_AA$ do $D_BBD_BD_A$, po kontrakcji $D_BB$ do $D_BD_A$. Ponownie wstawienie $BD_B$ i pierwsza kontrakcja działają na rozłącznych czynnikach. Otrzymujemy
\[
([(e_B\otimes\id_{D_B})(\id_{D_B}\otimes i_B)]\otimes\id_{D_A})
(\id_{D_B}\otimes[(e_A\otimes\id_{D_A})(\id_{D_A}\otimes i_A)])=\id_{D_BD_A}.
\]
W obu rachunkach zapis $AB$ oznacza $A\otimes B$.
\end{proof}
\begin{stwierdzenie}[Dualność w kategorii symetrycznej]\label{cat:prop:symmetric-left-dual}
Prawa dualność $(D,i,e)$ dla $A$ daje lewą dualność $(D,c_{A,D}i,e c_{A,D})$.
\end{stwierdzenie}
\begin{proof}
Niech $i'=c_{A,D}i$ i $e'=e c_{A,D}$. W zapisie ścisłym naturalność symetrii i $c_{A,I}=\id_A$ dają
\[
c_{A,A\otimes D}(\id_A\otimes i)=i\otimes\id_A,\qquad
(e\otimes\id_A)c_{A,D\otimes A}=\id_A\otimes e.
\]
Z koherencji symetrycznej (twierdzenie~\ref{cat:thm:symmetric-coherence}) mamy
\[
(c_{A,D}\otimes\id_A)(\id_A\otimes c_{A,D})
=c_{A,D\otimes A}c_{A,A\otimes D},
\]
gdyż obie drogi prowadzą z $A\otimes A\otimes D$ do $D\otimes A\otimes A$ z kolejnością pozycji $3,1,2$. Stąd
\[
(e'\otimes\id_A)(\id_A\otimes i')=(\id_A\otimes e)(i\otimes\id_A)=\id_A.
\]
Dla drugiej tożsamości stosujemy ten sam dowód w kategorii z odwróconą kolejnością iloczynu (jak w lemacie~\ref{cat:lem:unit-coherence}), z symetrią $c^\odot_{X,Y}=c_{Y,X}$. Tam $A$ jest prawostronnie dualny do $D$ dla $i,e$, bo ich typy są $I\to D\odot A$, $A\odot D\to I$, a tożsamości trójkątne są zamienione. Pierwsza dowiedziona tożsamość w tej kategorii jest dokładnie $(\id_D\otimes e')(i'\otimes\id_D)=\id_D$.
\end{proof}
\begin{definicja}[Sztywność i zwartość domknięta]\label{cat:def:rigid-compact}
Kategoria monoidalna jest sztywna, jeśli każdy obiekt posiada dualność po obu stronach. Symetryczna kategoria monoidalna, w której każdy obiekt posiada prawą dualność, jest zwarta domknięta (ang.\ compact closed). Stwierdzenie~\ref{cat:prop:symmetric-left-dual} daje wtedy także lewą dualność. Nie utożsamiamy tego warunku z domkniętością monoidalną.
\end{definicja}
\begin{twierdzenie}[Dualność przestrzeni wektorowych]\label{cat:thm:vect-dualizable}
W kategorii wszystkich przestrzeni wektorowych nad ciałem $k$ obiekt ma prawą dualność wtedy i tylko wtedy, gdy jest skończenie wymiarowy.
\end{twierdzenie}
\begin{proof}
Dla skończenie wymiarowego $V$ bierzemy przestrzeń dualną $V^*=\Hom_k(V,k)$. Ewaluacja jest indukowana przez dwuliniowe odwzorowanie $(\varphi,v)\mapsto\varphi(v)$. Dla bazy $(v_j)$ i bazy dualnej $(v^j)$ określamy $i(1)=\sum_j v_j\otimes v^j$. Pierwsza tożsamość trójkątna wysyła $v$ na $\sum_jv^j(v)v_j=v$, druga wysyła $\varphi$ na $\sum_j\varphi(v_j)v^j=\varphi$. W wymiarze zero suma jest pusta, a identyczność przestrzeni zerowej jest morfizmem zerowym. Niezależność od bazy wynika ze stwierdzenia~\ref{cat:prop:dual-unique}: jedyny izomorfizm $V^*\to V^*$ zachowujący ewaluację jest identycznością, bo funkcjonały rozróżniamy przez wartości.
Odwrotnie, dla dowolnej prawej dualności $(D,i,e)$ zapisujemy $i(1)=\sum_{j=1}^n v_j\otimes d_j$, co jest skończoną sumą z definicji algebraicznego iloczynu tensorowego. Pierwsza tożsamość trójkątna wymusza
$v=\sum_j e(d_j\otimes v)v_j$ dla każdego $v$. Zatem $V$ jest skończenie wymiarowa. Nie zakładaliśmy $D=V^*$; kanoniczny izomorfizm wynika z jednoznaczności dualności.
\end{proof}

\subsection{Ślad kategorialny}
\begin{definicja}[Ślad]\label{cat:def:categorical-trace}
W kategorii symetrycznej monoidalnej dla endomorfizmu $f:A\to A$ obiektu z prawą dualnością $(D,i,e)$ określamy
\[
\operatorname{tr}_{\mathcal C}(f)=e c_{A,D}(f\otimes\id_D)i:I\to I.
\]
\end{definicja}
\[
\begin{CD}
I @>{i}>> A\otimes D @>{f\otimes\id_D}>> A\otimes D\\
@. @. @VV{c_{A,D}}V\\
@. I @<<{e}< D\otimes A
\end{CD}
\]
\begin{stwierdzenie}[Niezależność śladu od dualności]\label{cat:prop:trace-independent}
Ślad nie zależy od wyboru prawej dualności.
\end{stwierdzenie}
\begin{proof}
Dla drugiej dualności $(D',i',e')$ niech $h:D\to D'$ będzie izomorfizmem ze stwierdzenia~\ref{cat:prop:dual-unique}. Wtedy
\[
\begin{aligned}
e'c_{A,D'}(f\otimes\id_{D'})i'
&=e'c_{A,D'}(f\otimes h)i\\
&=e'(h\otimes\id_A)c_{A,D}(f\otimes\id_D)i\\
&=ec_{A,D}(f\otimes\id_D)i.
\end{aligned}
\]
Druga równość wynika z naturalności symetrii i bifunktorialności.
\end{proof}
\begin{stwierdzenie}[Cykliczność śladu kategorialnego]\label{cat:prop:trace-cyclic}
W kategorii symetrycznej monoidalnej niech $A,B$ mają prawe dualności. Dla $f:A\to B$ i $g:B\to A$ zachodzi
\[
\operatorname{tr}_{\mathcal C}(gf)=\operatorname{tr}_{\mathcal C}(fg).
\]
Obie strony są morfizmami $I\to I$, także gdy $A$ i $B$ nie są izomorficzne.
\end{stwierdzenie}
\begin{proof}
Wybieramy dualności $(D_A,i_A,e_A)$ i $(D_B,i_B,e_B)$ oraz morfizm $f^*:D_B\to D_A$ ze stwierdzenia~\ref{cat:prop:dual-functoriality}. Kolejno stosujemy \eqref{cat:eq:dual-morphism-coeval}, bifunktorialność, naturalność symetrii dla pary $(g,f^*)$ oraz \eqref{cat:eq:dual-morphism-eval}:
\[
\begin{aligned}
\operatorname{tr}_{\mathcal C}(gf)
&=e_Ac_{A,D_A}(g\otimes\id_{D_A})(f\otimes\id_{D_A})i_A\\
&=e_Ac_{A,D_A}(g\otimes\id_{D_A})(\id_B\otimes f^*)i_B\\
&=e_Ac_{A,D_A}(g\otimes f^*)i_B\\
&=e_A(f^*\otimes g)c_{B,D_B}i_B\\
&=e_A(f^*\otimes\id_A)(\id_{D_B}\otimes g)c_{B,D_B}i_B\\
&=e_B(\id_{D_B}\otimes f)(\id_{D_B}\otimes g)c_{B,D_B}i_B\\
&=e_B(\id_{D_B}\otimes fg)c_{B,D_B}i_B\\
&=e_Bc_{B,D_B}(fg\otimes\id_{D_B})i_B\\
&=\operatorname{tr}_{\mathcal C}(fg).
\end{aligned}
\]
Przedostatnia równość jest naturalnością symetrii dla pary $(fg,\id_{D_B})$. Niezależność wyniku od wyborów dualności zapewnia stwierdzenie~\ref{cat:prop:trace-independent}.
\end{proof}
\begin{przyklad}[Ślad iloczynów macierzy prostokątnych]\label{cat:ex:trace-rectangular}
Niech $F=(F_{ij})$ będzie macierzą $m\times n$, a $G=(G_{ji})$ macierzą $n\times m$ nad ciałem $k$. Traktując je jako odwzorowania $k^n\to k^m$ oraz $k^m\to k^n$, otrzymujemy
\[
\operatorname{tr}(GF)=\sum_{j=1}^n\sum_{i=1}^m G_{ji}F_{ij}
=\sum_{i=1}^m\sum_{j=1}^n F_{ij}G_{ji}=\operatorname{tr}(FG).
\]
Jest to szczególny przypadek cykliczności, potwierdzony bezpośrednim rachunkiem.
\end{przyklad}

\begin{przyklad}[Ślad liniowy]\label{cat:ex:trace-linear}
Dla skończenie wymiarowego $V$ ślad kategorialny jest odwzorowaniem $k\to k$ mnożącym przez $\sum_jv^j(f(v_j))$. Jest to suma elementów przekątnej macierzy $f$, czyli zwykły ślad. Porównaj \atlasref{DG-03::prop:kontrakcja-slad}{DG-03: kontrakcja i ślad}. Dla identyczności otrzymujemy $\dim V$ jako element ciała $k$; w dodatniej charakterystyce może on być zerem mimo $V\neq0$.
\end{przyklad}

\section{Kategorie addytywne i abelowe}\label{cat:sec:abelian-categories}
Punktem wyjścia są grupy abelowe i moduły nad pierścieniem z jedynką (por. \atlasref{AA-16::aa:def:modules-sub-hom}{AA-16: podmoduły i homomorfizmy} oraz \atlasref{AA-16::aa:def:modules-exact}{AA-16: ciągi dokładne}). Chcemy przenieść ich własności na kategorie, w których obiekty nie muszą być zbiorami z wyróżnionymi elementami. Wszystkie kategorie w tym rozdziale są lokalnie małe.

\subsection{Kategorie preaddytywne i biprodukty}
\begin{definicja}[Kategoria preaddytywna]\label{cat:def:preadditive}
Kategoria $\mathcal A$ jest preaddytywna, jeżeli każdy zbiór $\Hom(X,Y)$ jest grupą abelową, a składanie jest addytywne w każdej zmiennej. W szczególności $g0=0$ i $0f=0$: na przykład $g0=g(0+0)=g0+g0$, skąd po odjęciu $g0$ otrzymujemy żądaną równość. Funktor pomiędzy kategoriami preaddytywnymi jest addytywny, jeżeli jego działanie na każdym zbiorze morfizmów jest homomorfizmem grup abelowych.
\end{definicja}
\begin{przyklad}
Kategorie $\Ab$, $\Vect_k$ oraz lewych modułów nad pierścieniem $R$ są preaddytywne: dodawanie morfizmów jest punktowe, a obustronna addytywność składania wynika z addytywności homomorfizmów. Dla modułów por. \atlasref{AA-16::aa:st:modules-hom-abelian}{AA-16: grupa homomorfizmów}. Kategoria grup nieabelowych nie otrzymuje takiej struktury przez punktowe dodawanie: punktowy iloczyn dwóch homomorfizmów nie musi być homomorfizmem.
\end{przyklad}
\begin{definicja}[Biprodukt]\label{cat:def:biproduct}
W kategorii preaddytywnej biproduktem obiektów $X_1,\ldots,X_n$ nazywamy obiekt $B$ z morfizmami $i_j:X_j\to B$, $p_j:B\to X_j$ takimi, że
\[
 p_j i_k=\begin{cases}\id_{X_j}&j=k,\\0&j\ne k,\end{cases}
 \qquad \sum_{j=1}^n i_jp_j=\id_B.
\]
Dla pustej rodziny warunek przyjmuje postać $\id_B=0_{B,B}$.
\end{definicja}
\begin{twierdzenie}[Biprodukt jako iloczyn i koprodukt]\label{cat:thm:biproduct-universal}
Biprodukt jest jednocześnie iloczynem dla projekcji $p_j$ i koproduktem dla włożeń $i_j$. Odwrotnie, każdy skończony iloczyn w kategorii preaddytywnej ma kanoniczną strukturę biproduktu, podobnie każdy skończony koprodukt.
\end{twierdzenie}
\begin{proof}
Dla rodziny $u_j:T\to X_j$ definiujemy $u=\sum_j i_ju_j$. Wówczas $p_ku=\sum_j p_ki_ju_j=u_k$. Jeżeli $v:T\to B$ spełnia te same równania, to $v=\sum_j i_jp_jv=\sum_j i_ju_j=u$. Dla rodziny $v_j:X_j\to T$ definiujemy $v=\sum_j v_jp_j$ i otrzymujemy $vi_k=v_k$. Każde inne $w$ o tych samych składowych spełnia $w=\sum_j wi_jp_j=v$. To dowodzi obu własności uniwersalnych.
Odwrotnie, dla iloczynu $(B,p_j)$ definiujemy $i_k$ jednoznacznie przez warunki $p_ji_k=\delta_{jk}\id$ (dla $j\ne k$ oznacza to morfizm zerowy odpowiedniego typu). Ponieważ $p_k\sum_j i_jp_j=p_k$, jednoznaczność własności uniwersalnej daje $\sum_j i_jp_j=\id_B$. Dla koproduktu dowód otrzymujemy w kategorii przeciwnej, która jest preaddytywna z tymi samymi działaniami na odwróconych zbiorach morfizmów.
Dla $n=0$ warunek $\id_B=0$ pociąga $f=f\id_B=0$ i $g=\id_Bg=0$ dla każdego morfizmu z $B$ i do $B$. Zatem $B$ jest obiektem zerowym. Odwrotnie, obiekt początkowy lub końcowy ma dokładnie jeden endomorfizm, więc jego identyczność jest zerowa.
\end{proof}
\begin{definicja}[Kategoria addytywna]\label{cat:def:additive}
Kategoria addytywna jest kategorią preaddytywną posiadającą wszystkie skończone biprodukty, w tym obiekt zerowy. Biprodukt dwóch obiektów oznaczamy przez $X\oplus Y$.
\end{definicja}
\begin{stwierdzenie}[Rachunek macierzowy]\label{cat:prop:additive-matrices}
Dla biproduktów $X=\bigoplus_{j=1}^nX_j$, $Y=\bigoplus_{i=1}^mY_i$ morfizm $f:X\to Y$ jest jednoznacznie określony przez składowe $f_{ij}=p_i^Yfi_j^X$. Każda rodzina takich składowych wyznacza morfizm $f=\sum_{i,j}i_i^Y f_{ij}p_j^X$. Składanie odpowiada mnożeniu macierzy, w którym mnożenie wpisów oznacza składanie morfizmów.
\end{stwierdzenie}
\begin{proof}
Dla podanego wzoru $p_k^Yfi_l^X=\sum_{i,j}p_k^Yi_i^Y f_{ij}p_j^Xi_l^X=f_{kl}$. Z drugiej strony, wstawiając rozkłady obu identyczności, otrzymujemy
$f=\id_Y f\id_X=\sum_{i,j}i_i^Y(p_i^Yfi_j^X)p_j^X$.
Jeżeli $g:Y\to Z$, to
\[
 (gf)_{kj}=p_k^Zg\left(\sum_i i_i^Yp_i^Y\right)fi_j^X
 =\sum_i g_{ki}f_{ij}.
\]
Dodawanie jest składowe dzięki addytywności składania. Puste sumy są zerami w odpowiednich grupach morfizmów.
\end{proof}

\subsection{Jądra, kojądra, obraz i koobraz}
\begin{definicja}[Jądro i kojądro]\label{cat:def:kernel-cokernel}
Dla $f:X\to Y$ jądrem $\ker f:K\to X$ jest equalizator $f,0$, a kojądrem $\operatorname{coker}f:Y\to Q$ jest koequalizator $f,0$ (por. definicje~\ref{def:equalizer} i \ref{def:coeq}). W szczególności $\ker f$ jest monomorfizmem: jeśli $ku=kv$, to $u,v$ są dwoma rozłożeniami tego samego morfizmu przez jądro, więc $u=v$. Dualnie kojądro jest epimorfizmem. Jądra i kojądra są wyznaczone z dokładnością do jednoznacznego izomorfizmu zgodnego z morfizmem do $X$ lub z $Y$.
\end{definicja}
\begin{stwierdzenie}\label{cat:prop:mono-epi-zero-kernel}
W kategorii addytywnej posiadającej jądra i kojądra morfizm $f:X\to Y$ jest monomorfizmem wtedy i tylko wtedy, gdy jego jądro ma obiekt zerowy jako dziedzinę. Dualnie $f$ jest epimorfizmem wtedy i tylko wtedy, gdy jego kojądro ma obiekt zerowy jako przeciwdziedzinę.
\end{stwierdzenie}
\begin{proof}
Jeżeli $f$ jest monomorfizmem, to $fu=0=f0$ wymusza $u=0$. Morfizm $0\to X$ spełnia więc własność uniwersalną jądra. Odwrotnie, gdy jądro jest zerowe i $fu=fv$, addytywność daje $f(u-v)=0$. Morfizm $u-v$ rozkłada się przez obiekt zerowy, zatem $u-v=0$ i $u=v$. Drugi wniosek wynika z zastosowania pierwszego do kategorii przeciwnej.
\end{proof}
\begin{definicja}[Obraz i koobraz]\label{cat:def:image-coimage}
Zakładamy istnienie wszystkich jąder i kojąder. Dla $f:X\to Y$ przyjmujemy
\[
 k:K\to X=\ker f,\qquad q:Y\to Q=\operatorname{coker}f,
 \qquad p:X\to C=\operatorname{coker}k,\qquad j:J\to Y=\ker q,
\]
gdzie zapis oznacza odpowiednie morfizmy uniwersalne. Obiekty $C,J$ nazywamy odpowiednio koobrazem i obrazem $f$, a morfizmy $p,j$ ich strukturalnymi strzałkami.
\end{definicja}
\begin{stwierdzenie}[Kanoniczne porównanie obrazu i koobrazu]\label{cat:prop:coimage-image}
Istnieje jedyny morfizm $\bar f:C\to J$ taki, że $f=j\bar f p$.
\end{stwierdzenie}
\begin{proof}
Ponieważ $fk=0$, własność kojądra $p$ daje jedyny morfizm $a:C\to Y$ z $ap=f$. Mamy $qap=qf=0$; ponieważ $p$ jest epimorfizmem, $qa=0$. Własność jądra $j$ daje jedyny morfizm $\bar f:C\to J$ z $j\bar f=a$. Jeżeli $jbp=f$, to z epimorficzności $p$ wynika $jb=a$, a z monomorficzności $j$ wynika $b=\bar f$.
\end{proof}
\[
\begin{CD}
X @>{f}>> Y\\
@V{p}VV @AA{j}A\\
C @>{\bar f}>> J
\end{CD}
\]
\begin{definicja}[Kategoria abelowa]\label{cat:def:abelian}
Kategoria abelowa jest kategorią addytywną, w której istnieją wszystkie jądra i kojądra, a kanoniczny morfizm $\bar f$ ze stwierdzenia~\ref{cat:prop:coimage-image} jest izomorfizmem dla każdego $f$.
\end{definicja}
\begin{twierdzenie}[Rozkład morfizmu w kategorii abelowej]\label{cat:thm:abelian-factorization}
W kategorii abelowej każdy morfizm rozkłada się na epimorfizm, po którym następuje monomorfizm. Każdy monomorfizm jest jądrem swojego kojądra, a każdy epimorfizm jest kojądrem swojego jądra. Każdy morfizm będący jednocześnie monomorfizmem i epimorfizmem jest izomorfizmem.
\end{twierdzenie}
\begin{proof}
W rozkładzie $f=j\bar f p$ morfizm $\bar f p$ jest epimorfizmem jako złożenie epimorfizmu i izomorfizmu, a $j$ jest monomorfizmem.
Jeśli $f$ jest monomorfizmem, to $k:0\to X$ jest jego jądrem na mocy stwierdzenia~\ref{cat:prop:mono-epi-zero-kernel}. Identyczność $\id_X$ jest kojądrem $k$, ponieważ każdy morfizm z $X$ zeruje się po złożeniu z $k$ i jednoznacznie rozkłada się przez $\id_X$. Wybierając $p=\id_X$, mamy $f=j\bar f$, gdzie $\bar f$ jest izomorfizmem; stąd $f$ samo spełnia własność uniwersalną jądra $q$: rozłożenie przez $j$ składamy z $\bar f^{-1}$, a jednoznaczność wynika z monomorficzności $f$. Dualnie epimorfizm jest kojądrem swego jądra (kategoria przeciwna kategorii abelowej jest abelowa, gdyż porównanie obrazu i koobrazu staje się odwróconym izomorfizmem).
Jeśli $f$ jest także epimorfizmem, to jego kojądro jest zerowe. Jądrem morfizmu $Y\to0$ jest $\id_Y$, więc $f$ jest izomorfizmem przez jednoznaczność jąder.
\end{proof}
\begin{przyklad}[Moduły]\label{cat:ex:abelian-modules}
W kategorii lewych $R$-modułów biprodukty skończone są sumami prostymi. Jądro $f:M\to N$ jest włączeniem zwykłego jądra, a kojądro jest projekcją $N\to N/f(M)$. Własność uniwersalna jądra wynika z faktu, że $fu=0$ wtedy i tylko wtedy, gdy wartości $u$ należą do jądra; kojądra z własności uniwersalnej modułu ilorazowego. Stąd $C=M/\ker f$ i $J=f(M)$, zaś $\bar f([x])=f(x)$ jest izomorfizmem na mocy \atlasref{AA-16::aa:tw:modules-first-isomorphism}{AA-16: pierwsze twierdzenie o izomorfizmie}. To dowodzi abelowości, w szczególności dla $\Ab$ i $\Vect_k$.
\end{przyklad}
\begin{przyklad}[Kategoria addytywna, lecz nieabelowa]\label{cat:ex:additive-not-abelian}
W pełnej podkategorii skończenie generowanych wolnych grup abelowych istnieją skończone biprodukty. Morfizm $2:\mathbb Z\to\mathbb Z$ (mnożenie przez $2$) jest mono, bo jest iniektywny, i epi: jeśli $u,v:\mathbb Z\to F$ spełniają $u\circ2=v\circ2$, to $2(u(1)-v(1))=0$, skąd $u(1)=v(1)$, ponieważ $F$ jest bez torsji. Nie jest izomorfizmem, gdyż nie jest surjektywny. Stwierdzenie~\ref{cat:thm:abelian-factorization} wyklucza więc abelowość tej kategorii.
\end{przyklad}

\subsection{Ciągi dokładne i rozkładalność}
\begin{definicja}[Dokładność]\label{cat:def:abelian-exact}
W kategorii abelowej ciąg $X\xrightarrow{f}Y\xrightarrow{g}Z$ jest dokładny w $Y$, jeśli $gf=0$ i podobiekty $\operatorname{im}f$ oraz $\ker g$ są równe: istnieje izomorfizm między ich dziedzinami zgodny z monomorfizmami do $Y$. Krótki ciąg dokładny ma postać $0\to A\xrightarrow{i}B\xrightarrow{p}C\to0$, jest dokładny w każdym z trzech niezerowych miejsc (także gdy któryś z obiektów jest zerowy).
\end{definicja}
\begin{stwierdzenie}\label{cat:prop:ses-kernel-cokernel}
Krótki ciąg jest dokładny wtedy i tylko wtedy, gdy $i$ jest jądrem $p$, a $p$ kojądrem $i$.
\end{stwierdzenie}
\begin{proof}
Dokładność na końcach oznacza $\ker i=0$ i $\operatorname{im}p=C$ jako podobiekt $C$. Pierwszy warunek daje monomorficzność $i$. Drugi daje epimorficzność $p$, ponieważ $p$ rozkłada się na epimorfizm i izomorfizm obrazu z $C$. Dla monomorfizmu $i$ jego obraz jest reprezentowany przez samo $i$, więc dokładność w $B$ oznacza, że $i$ jest jądrem $p$. Ponieważ $p$ jest epimorfizmem, jest kojądrem swojego jądra, czyli $i$ (twierdzenie~\ref{cat:thm:abelian-factorization}). Odwrotnie, jądro jest mono, kojądro epi, ich złożenie jest zerowe, a obraz $i$ jest reprezentowany przez $i=\ker p$ jako podobiekt. Epimorfizm $p$ ma obraz całego $C$, gdyż jego kojądro jest zerowe. Wszystkie trzy warunki dokładności są spełnione.
\end{proof}
\begin{twierdzenie}[Rozkładalność krótkiego ciągu dokładnego]\label{cat:thm:abelian-split}
Dla krótkiego ciągu dokładnego $0\to A\xrightarrow{i}B\xrightarrow{p}C\to0$ następujące warunki są równoważne:
\begin{enumerate}
\item istnieje $s:C\to B$ z $ps=\id_C$;
\item istnieje $r:B\to A$ z $ri=\id_A$;
\item istnieje izomorfizm $u:A\oplus C\to B$ taki, że $ui_A=i$ i $pu=p_C$.
\end{enumerate}
\end{twierdzenie}
\begin{proof}
Z (1) wynika $p(\id_B-sp)=0$, więc własność jądra $i$ daje jedyny $r:B\to A$ z $ir=\id_B-sp$. Wówczas $iri=i-spi=i$, czyli $ri=\id_A$, oraz $irs=s-sps=0$, czyli $rs=0$. Obie równości wynikają z monomorficzności $i$. Morfizm $u=[i,s]:A\oplus C\to B$ ma odwrotność $v=\langle r,p\rangle:B\to A\oplus C$: iloczyn $vu$ ma macierz diagonalną z identycznościami, a $uv=ir+sp=\id_B$. Ponadto $ui_A=i$ oraz $pu=p_C$. To daje (3).
Z (2) wynika $(\id_B-ir)i=0$. Własność kojądra $p$ daje jedyny $s:C\to B$ z $sp=\id_B-ir$. Mamy $psp=p$, więc $ps=\id_C$ przez epimorficzność $p$. To daje (1).
Jeśli zachodzi (3), przyjmujemy $s=ui_C$ i $r=p_Au^{-1}$. Wtedy $ps=p_Ci_C=\id_C$ i $ri=p_Au^{-1}i=p_Ai_A=\id_A$.
\end{proof}
\begin{przyklad}
Ciąg $0\to\mathbb Z\xrightarrow{2}\mathbb Z\to\mathbb Z/2\mathbb Z\to0$ jest dokładny, ale nie jest rozkładalny: każdy homomorfizm $\mathbb Z/2\mathbb Z\to\mathbb Z$ jest zerowy, gdyż obraz klasy $1$ musi mieć rząd dzielący $2$. Nie może zatem być prawą odwrotnością projekcji. Por. \atlasref{AA-16::aa:tw:modules-splitting}{AA-16: rozkładalność ciągów modułów}.
\end{przyklad}
Dalsza teoria kompleksów i długich ciągów dokładnych korzysta z tych konstrukcji; w przypadku modułów ich szczegółowe dowody znajdują się w \atlasref{AT-03::def:kompleks-lancuchowy}{AT-03: kompleksy łańcuchowe} oraz \atlasref{AT-03::thm:les-homologia}{AT-03: długi ciąg dokładny homologii}.

\section{Obiekty rzutowe, injektywne i rezolucje}\label{cat:sec:projective-injective}
Pracujemy w lokalnie małej kategorii abelowej $\mathcal A$. Wykorzystujemy teorię funktorów dokładnych z rozdziału~\ref{cat:sec:abelian-categories}. W przypadku modułów odpowiednie pojęcie rzutowości oraz jego charakterystykę przez składniki sum prostych modułów wolnych znamy z \atlasref{AA-16::aa:tw:modules-projective-summand}{AA-16: charakterystyka modułów rzutowych}.

\subsection{Rzutowość i injektywność}
\begin{definicja}\label{cat:def:projective-injective}
Obiekt $P$ jest rzutowy, jeżeli dla każdego epimorfizmu $q:E\to N$ i każdego $f:P\to N$ istnieje $\widetilde f:P\to E$ z $q\widetilde f=f$.
Obiekt $I$ jest injektywny, jeżeli dla każdego monomorfizmu $j:M\to E$ i każdego $g:M\to I$ istnieje $\widetilde g:E\to I$ z $\widetilde g j=g$. Nie wymagamy jednoznaczności tych morfizmów.
\[
\begin{array}{ccc}
P&&\\
\downarrow\widetilde f&\searrow f&\\
E&\xrightarrow{q}&N
\end{array}
\qquad
\begin{CD}
M @>{j}>> E\\
@V{g}VV @VV{\widetilde g}V\\
I @>{\id_I}>> I
\end{CD}
\]
\end{definicja}
\begin{twierdzenie}[Charakteryzacje rzutowości i injektywności]\label{cat:thm:projective-characterization}
Dla obiektu $P$ równoważne są następujące warunki:
\begin{enumerate}
\item $P$ jest rzutowy;
\item funktor $\Hom(P,-):\mathcal A\to\Ab$ jest dokładny;
\item każdy epimorfizm $E\to P$ ma prawą odwrotność;
\item każdy krótki ciąg dokładny kończący się obiektem $P$ jest rozkładalny.
\end{enumerate}
Dualnie dla obiektu $I$ równoważne są injektywność, dokładność $\Hom(-,I):\mathcal A^{\op}\to\Ab$, istnienie lewej odwrotności każdego monomorfizmu $I\to E$ oraz rozkładalność każdego krótkiego ciągu dokładnego rozpoczynającego się obiektem $I$.
\end{twierdzenie}
\begin{proof}
Funktor $\Hom(P,-)$ jest lewostronnie dokładny na mocy przykładu~\ref{cat:ex:hom-left-exact}. Dla krótkiego ciągu dokładnego $0\to A\to E\xrightarrow{q}N\to0$ brakuje zatem jedynie surjektywności $\Hom(P,q)$. Jest ona równoważna istnieniu podniesienia każdego $f:P\to N$ wzdłuż $q$. Każdy epi $q$ występuje w takim ciągu, ponieważ jest kojądrem swego jądra. To daje równoważność (1) i (2).
Warunek (1) zastosowany do $f=\id_P$ daje (3). Jeśli zachodzi (3), to dla dowolnego epi $q:E\to N$ i $f:P\to N$ tworzymy iloczyn włóknisty $D=E\times_N P$ z projekcjami $a:D\to E$, $b:D\to P$. Z twierdzenia~\ref{cat:thm:abelian-base-change} morfizm $b$ jest epi. Istnieje więc $s:P\to D$ z $bs=\id_P$. Morfizm $\widetilde f=as$ spełnia $q\widetilde f=qas=fbs=f$. Otrzymujemy (1). Równoważność (3) i (4) wynika z twierdzenia~\ref{cat:thm:abelian-split} i z tego, że każdy epimorfizm $E\to P$ jest kojądrem swego jądra. Twierdzenie dla $I$ otrzymujemy w kategorii przeciwnej, która jest abelowa.
\end{proof}
\begin{stwierdzenie}[Retrakty i sumy proste]\label{cat:prop:projective-retracts}
Retrakt obiektu rzutowego jest rzutowy. Dowolny istniejący koprodukt rodziny obiektów rzutowych jest rzutowy. Dualnie retrakty obiektów injektywnych i dowolne istniejące iloczyny obiektów injektywnych są injektywne.
\end{stwierdzenie}
\begin{proof}
Niech $i:P\to Q$, $r:Q\to P$ spełniają $ri=\id_P$, gdzie $Q$ jest rzutowy. Dla epi $q:E\to N$ i $f:P\to N$ podnosimy $fr:Q\to N$ do $h:Q\to E$. Wtedy $qhi=fri=f$. Jeśli $P=\coprod_{s\in S}P_s$ z włożeniami $i_s$, wybieramy podniesienia $h_s:P_s\to E$ morfizmów $fi_s$. Dla nieskończonego zbioru $S$ korzystamy z aksjomatu wyboru. Własność koproduktu daje $h:P\to E$ z $hi_s=h_s$. Ponieważ $qhi_s=fi_s$ dla każdego $s$, mamy $qh=f$. Wnioski o injektywności są dualne.
\end{proof}
\begin{przyklad}
W $\Ab$ obiekt $\mathbb Z$ jest rzutowy, ponieważ homomorfizm z $\mathbb Z$ wyznacza obraz $1$ i można go podnieść wzdłuż każdej surjekcji. Obiekt $\mathbb Z/2\mathbb Z$ nie jest rzutowy, ponieważ ciąg $0\to\mathbb Z\xrightarrow{2}\mathbb Z\to\mathbb Z/2\mathbb Z\to0$ nie jest rozkładalny. Nie jest też injektywny: homomorfizm $2\mathbb Z\to\mathbb Z/2\mathbb Z$ wysyłający $2$ na klasę $1$ nie przedłuża się do $\mathbb Z$, gdyż każde przedłużenie $g$ musiałoby spełniać $g(2)=2g(1)=0$. Podobnie $\mathbb Z$ nie jest injektywna: homomorfizm $2\mathbb Z\to\mathbb Z$ wysyłający $2$ na $1$ nie może się przedłużyć, gdyż równanie $2g(1)=1$ nie ma rozwiązania w $\mathbb Z$.
\end{przyklad}

\subsection{Dostatecznie wiele obiektów i rezolucje}
\begin{definicja}\label{cat:def:enough-projectives-injectives}
Mówimy, że kategoria abelowa ma dostatecznie wiele obiektów rzutowych, gdy dla każdego $X$ istnieje epi $P\to X$ z $P$ rzutowym. Ma dostatecznie wiele obiektów injektywnych, gdy dla każdego $X$ istnieje mono $X\to I$ z $I$ injektywnym.
\end{definicja}
\begin{przyklad}[Moduły]\label{cat:ex:enough-projectives-mod}
Kategoria lewych modułów nad dowolnym pierścieniem $R$ z jedynką ma dostatecznie wiele obiektów rzutowych: moduł wolny $R^{(M)}$ na zbiorze elementów $M$ jest rzutowy, a homomorfizm wysyłający generator indeksowany przez $m$ na $m$ jest surjektywny. Rzutowość modułu wolnego oraz własność uniwersalną tej konstrukcji udowodniono w \atlasref{AA-16::aa:tw:modules-projective-summand}{AA-16: moduły rzutowe} oraz \atlasref{AA-16::aa:tw:modules-free-universal}{AA-16: moduły wolne}. W kategorii abelowej mającej dostatecznie wiele obiektów rzutowych każdy obiekt rzutowy jest retraktem każdego obiektu rzutowego, z którego istnieje na niego epimorfizm, na mocy twierdzenia~\ref{cat:thm:projective-characterization}.
\end{przyklad}
\begin{definicja}[Rezolucje]\label{cat:def:resolutions}
Rezolucją rzutową obiektu $X$ nazywamy dokładny ciąg
\[
 \cdots\xrightarrow{d_3}P_2\xrightarrow{d_2}P_1
 \xrightarrow{d_1}P_0\xrightarrow{\epsilon}X\to0,
\]
w którym wszystkie $P_n$ są rzutowe. Rezolucja injektywna ma postać dokładnego ciągu
\[
 0\to X\xrightarrow{\eta}I^0\xrightarrow{d^0}I^1
 \xrightarrow{d^1}I^2\to\cdots,
\]
w którym wszystkie $I^n$ są injektywne. Dokładność wymusza zerowanie złożeń dwóch kolejnych strzałek.
\end{definicja}
\begin{twierdzenie}[Istnienie rezolucji]\label{cat:thm:resolution-existence}
Jeżeli kategoria abelowa ma dostatecznie wiele obiektów rzutowych, to każdy obiekt ma rezolucję rzutową. Jeżeli ma dostatecznie wiele obiektów injektywnych, to każdy obiekt ma rezolucję injektywną. Przy nieskończonej konstrukcji przyjmujemy możliwość kolejnych wyborów obiektów i strzałek (na przykład pracujemy względem ustalonego uniwersum i stosujemy aksjomat wyboru zależnego).
\end{twierdzenie}
\begin{proof}
Przyjmijmy $K_0=X$. Wybieramy epi $e_0:P_0\to K_0$ z obiektem rzutowym $P_0$ i kładziemy $\epsilon=e_0$. Indukcyjnie, gdy mamy $e_n:P_n\to K_n$, wybieramy jądro $k_{n+1}:K_{n+1}\to P_n$ oraz epi $e_{n+1}:P_{n+1}\to K_{n+1}$ z obiektem rzutowym $P_{n+1}$. Definiujemy $d_{n+1}=k_{n+1}e_{n+1}$. Ponieważ $e_{n+1}$ jest epi, a $k_{n+1}$ mono, obraz $d_{n+1}$ jest reprezentowany przez $k_{n+1}$: ich kojądra mają tę samą własność uniwersalną, gdyż $hd_{n+1}=0$ wtedy i tylko wtedy, gdy $hk_{n+1}=0$.
Dla $n\ge1$ zachodzi $\ker d_n=\ker e_n$, bo $k_n$ jest mono; dla $n=0$ mamy $\ker\epsilon=\ker e_0$. W każdym przypadku $k_{n+1}$ jest więc jądrem strzałki wychodzącej z $P_n$. Tym samym otrzymujemy dokładność w każdym $P_n$, a dokładność w $X$ wynika z epimorficzności $\epsilon$. Rezolucję injektywną konstruujemy stosując ten dowód w $\mathcal A^{\op}$.
\end{proof}
\begin{przyklad}[Krótka rezolucja rzutowa]\label{cat:ex:cyclic-projective-resolution}
Dla $m\ge2$ ciąg $0\to\mathbb Z\xrightarrow{m}\mathbb Z\to\mathbb Z/m\mathbb Z\to0$ jest rezolucją rzutową grupy $\mathbb Z/m\mathbb Z$ (w stopniach $n\ge2$ przyjmujemy $P_n=0$). Nie jest to rozkładalny ciąg dokładny: gdyby projekcja miała prawą odwrotność, obraz klasy $1$ byłby elementem $\mathbb Z$ zabijanym przez $m$, czyli zerem, co jest niemożliwe.
\end{przyklad}

\subsection{Porównywanie rezolucji i homotopie}
Dla kompleksu łańcuchowego $P_\bullet$ zapisujemy różniczki jako $d_n^P:P_n\to P_{n-1}$, $n\ge1$, z $d_n^Pd_{n+1}^P=0$. Morfizm kompleksów $f:P_\bullet\to Q_\bullet$ jest rodziną $f_n:P_n\to Q_n$ spełniającą $d_n^Q f_n=f_{n-1}d_n^P$. W rezolucjach obiektów $X,Y$ morfizm ten podnosi $u:X\to Y$, gdy $\epsilon_Q f_0=u\epsilon_P$. W przypadku modułów definicje omówiono w \atlasref{TOP-03::def:kompleks-lancuchowy}{TOP-03: kompleksy łańcuchowe} i \atlasref{TOP-03::def:mapa-lancuchowa}{TOP-03: morfizmy łańcuchowe}.
\begin{definicja}[Homotopia łańcuchowa]\label{cat:def:resolution-homotopy}
Morfizmy kompleksów $f,g:P_\bullet\to Q_\bullet$ są homotopijne, jeśli istnieją $h_n:P_n\to Q_{n+1}$ takie, że
\begin{equation}\label{cat:eq:resolution-homotopy}
 f_0-g_0=d_1^Qh_0,\qquad
 f_n-g_n=d_{n+1}^Qh_n+h_{n-1}d_n^P\quad(n\ge1).
\end{equation}
Homotopia jest relacją równoważności: homotopia zerowa, zmiana znaku i dodawanie homotopii dają zwrotność, symetrię i przechodniość. Por. \atlasref{TOP-03::def:homotopia-lancuchowa}{TOP-03: homotopia łańcuchowa}.
\end{definicja}
\begin{twierdzenie}[Twierdzenie porównawcze]\label{cat:thm:resolution-comparison}
Niech $P_\bullet\to X$ będzie rezolucją rzutową, a $Q_\bullet\to Y$ dokładnym ciągiem augmentowanym postaci z definicji~\ref{cat:def:resolutions}, bez założenia rzutowości $Q_n$. Każdy $u:X\to Y$ podnosi się do morfizmu kompleksów $f:P_\bullet\to Q_\bullet$. Dowolne dwa podniesienia są homotopijne.
\end{twierdzenie}
\begin{proof}
W stopniu zerowym podnosimy $u\epsilon_P$ wzdłuż epimorfizmu $\epsilon_Q$ dzięki rzutowości $P_0$. Otrzymujemy $f_0$ z $\epsilon_Q f_0=u\epsilon_P$. Załóżmy, że skonstruowano $f_0,\ldots,f_{n-1}$ zgodne z różniczkami i augmentacją. Dla $n=1$ mamy $\epsilon_Q f_0d_1^P=u\epsilon_Pd_1^P=0$, a dla $n\ge2$
\[
 d_{n-1}^Q f_{n-1}d_n^P=f_{n-2}d_{n-1}^Pd_n^P=0.
\]
Niech $k:K\to Q_{n-1}$ będzie jądrem $\epsilon_Q$ dla $n=1$ lub $d_{n-1}^Q$ dla $n\ge2$. Istnieje jedyny $a:P_n\to K$ z $ka=f_{n-1}d_n^P$. Z dokładności $Q$ różniczka $d_n^Q$ rozkłada się jako $ke$ dla epi $e:Q_n\to K$. Rzutowość $P_n$ daje $f_n:P_n\to Q_n$ z $ef_n=a$, a stąd $d_n^Qf_n=ka=f_{n-1}d_n^P$. Indukcja daje istnienie.

Niech $f,g$ podnoszą ten sam $u$. Ponieważ $\epsilon_Q(f_0-g_0)=0$, morfizm $f_0-g_0$ rozkłada się przez jądro $\epsilon_Q$. Dokładność $Q$ i rzutowość $P_0$ pozwalają go podnieść do $h_0:P_0\to Q_1$ z $d_1^Qh_0=f_0-g_0$.
Jeśli skonstruowano $h_0,\ldots,h_{n-1}$ spełniające równania w niższych stopniach, kładziemy
\[
 a_n=f_n-g_n-h_{n-1}d_n^P:P_n\to Q_n.
\]
Dla $n=1$ otrzymujemy
$d_1^Qa_1=(f_0-g_0)d_1^P-d_1^Qh_0d_1^P=0$.
Dla $n\ge2$
\[
\begin{aligned}
 d_n^Qa_n
 &=(f_{n-1}-g_{n-1}-d_n^Qh_{n-1})d_n^P\\
 &=h_{n-2}d_{n-1}^Pd_n^P=0.
\end{aligned}
\]
Zatem $a_n$ rozkłada się przez jądro $d_n^Q$. Dokładność $Q$ dostarcza epimorfizmu z $Q_{n+1}$ na to jądro; rzutowość $P_n$ daje $h_n$ z $d_{n+1}^Qh_n=a_n$. Jest to równanie~\eqref{cat:eq:resolution-homotopy}. Indukcja kończy dowód (dla nieskończonych ciągów korzystamy z kolejnych wyborów jak przy istnieniu rezolucji).
\end{proof}
\begin{wniosek}[Niezależność rezolucji]\label{cat:cor:resolution-homotopy-equivalence}
Dwie rezolucje rzutowe tego samego obiektu są równoważne homotopijnie. Dla ustalonego wyboru rezolucji obiektów podniesienie morfizmu do klasy homotopii jest funktorialne.
\end{wniosek}
\begin{proof}
Stosujemy twierdzenie~\ref{cat:thm:resolution-comparison} do $\id_X$ w obu kierunkach, otrzymując $f:P_\bullet\to Q_\bullet$ i $g:Q_\bullet\to P_\bullet$. Złożenia $gf,fg$ podnoszą $\id_X$, podobnie jak identyczności kompleksów, więc $gf\simeq\id_P$ oraz $fg\simeq\id_Q$. Jeżeli $f,g$ podnoszą kolejno $u:X\to Y$, $v:Y\to Z$, to $\epsilon_Rg_0f_0=v\epsilon_Qf_0=vu\epsilon_P$. Stąd złożenie klas podniesień jest klasą podniesienia $vu$, a klasa podniesienia identyczności jest klasą identyczności.
Składanie jest dobrze określone na klasach homotopii: jeśli $f-f'=d^Qh+hd^P$ i $a:Q_\bullet\to R_\bullet$ jest morfizmem kompleksów, to
$af-af'=d^R(ah)+(ah)d^P$, gdzie $(ah)_n=a_{n+1}h_n$, dzięki przemienności $a$ z różniczkami.
Dla $b:S_\bullet\to P_\bullet$ mamy
$fb-f'b=d^Q(hb)+(hb)d^S$, gdzie $(hb)_n=h_nb_n$, dzięki przemienności $b$ z różniczkami.
Zatem zastąpienie obu czynników morfizmami homotopijnymi nie zmienia klasy złożenia.
\end{proof}
\begin{wniosek}[Porównywanie rezolucji injektywnych]\label{cat:cor:injective-comparison}
Niech $X\to I^\bullet$ będzie dokładnym ciągiem koaugmentowanym, a $Y\to J^\bullet$ rezolucją injektywną. Każdy $u:X\to Y$ przedłuża się do morfizmu kompleksów $f:I^\bullet\to J^\bullet$, jednoznacznego z dokładnością do homotopii kołańcuchowej. Dwie rezolucje injektywne tego samego obiektu są równoważne homotopijnie.
\end{wniosek}
\begin{proof}
Stosujemy twierdzenie~\ref{cat:thm:resolution-comparison} i wniosek~\ref{cat:cor:resolution-homotopy-equivalence} w kategorii przeciwnej. Tam $J^\bullet$ staje się rezolucją rzutową, a $u$ ma kierunek od $Y$ do $X$. Po powrocie do $\mathcal A$ otrzymujemy morfizm we wskazanym kierunku. Homotopia ma postać $s^n:I^n\to J^{n-1}$ dla $n\ge1$ i spełnia
\[
 f^n-g^n=d_J^{n-1}s^n+s^{n+1}d_I^n,
\]
gdzie pierwszy składnik dla $n=0$ jest zerem. Jest to odwrócenie złożeń z~\eqref{cat:eq:resolution-homotopy}.
\end{proof}

\subsection{Lemat Schanuela}
\begin{twierdzenie}[Lemat Schanuela]\label{cat:thm:schanuel}
Niech
\[
 0\to K\xrightarrow{i}P\xrightarrow{p}X\to0,\qquad
 0\to L\xrightarrow{j}Q\xrightarrow{q}X\to0
\]
będą krótkimi ciągami dokładnymi, w których $P,Q$ są rzutowe. Wówczas $K\oplus Q\cong L\oplus P$.
\end{twierdzenie}
\begin{proof}
Tworzymy iloczyn włóknisty $D=P\times_X Q$ z projekcjami $a:D\to P$, $b:D\to Q$. Obie projekcje są epi na mocy twierdzenia~\ref{cat:thm:abelian-base-change}. Para $(i,0)$ definiuje $k:K\to D$ z $ak=i$, $bk=0$. Jeśli $z:T\to D$ spełnia $bz=0$, to $paz=qbz=0$, więc istnieje jedyny $y:T\to K$ z $iy=az$. Morfizmy $ky$ i $z$ mają takie same złożenia z $a,b$, a zatem są równe przez własność iloczynu włóknistego. Jednoznaczność $y$ wynika z monomorficzności $i$. To dowodzi, że $k=\ker b$. Analogicznie para $(0,j)$ definiuje jądro $l:L\to D$ morfizmu $a$: dla $az=0$ mamy $qbz=0$, zatem $bz=jy$ dla jedynego $y:T\to L$ i $ly=z$.
Otrzymujemy krótkie ciągi dokładne
\[
 0\to K\xrightarrow{k}D\xrightarrow{b}Q\to0,\qquad
 0\to L\xrightarrow{l}D\xrightarrow{a}P\to0.
\]
Rzutowość $P,Q$ daje przekroje obu epimorfizmów, więc twierdzenie~\ref{cat:thm:abelian-split} daje $D\cong K\oplus Q$ oraz $D\cong L\oplus P$. Składając jeden izomorfizm z odwrotnością drugiego, otrzymujemy tezę. Wybór przekrojów nie musi być jednoznaczny, dlatego uzyskany izomorfizm nie jest na ogół kanoniczny.
\end{proof}
\begin{wniosek}\label{cat:cor:schanuel-projectivity}
Jeżeli $X$ ma krótką rezolucję rzutową $0\to K\to P\to X\to0$ z $K$ rzutowym, to jądro każdego epimorfizmu $Q\to X$ z $Q$ rzutowym jest rzutowe.
\end{wniosek}
\begin{proof}
Dla jądra $L$ tego epimorfizmu lemat Schanuela daje $L\oplus P\cong K\oplus Q$. Prawa strona jest rzutowa jako skończony koprodukt obiektów rzutowych. Obiekt $L$ jest retraktem lewej strony, zatem jest rzutowy na mocy stwierdzenia~\ref{cat:prop:projective-retracts}.
\end{proof}

\section{Funktory pochodne, Ext i Tor}\label{cat:sec:derived-functors}
Pracujemy w lokalnie małych kategoriach abelowych. Wykorzystujemy konstrukcje z rozdziałów~\ref{cat:sec:abelian-categories} i~\ref{cat:sec:projective-injective}. Wersję modułową teorii homologii i kohomologii omówiono w \atlasref{TOP-03::def:homologia-algebraiczna}{TOP-03: homologia} oraz \atlasref{TOP-03::def:kohomologia-algebraiczna}{TOP-03: kohomologia}.

\subsection{Homologia w kategorii abelowej}
\begin{definicja}[Cykle, brzegi i homologia]\label{cat:def:abelian-homology}
Niech $(C_\bullet,d)$ będzie kompleksem łańcuchowym, tj. rodziną obiektów $C_n$ i morfizmów $d_n:C_n\to C_{n-1}$ z $d_nd_{n+1}=0$. Dopuszczamy indeksy całkowite; dla kompleksu nieujemnie stopniowanego przyjmujemy $C_n=0$ dla $n<0$. Niech $z_n:Z_n\to C_n$ będzie jądrem $d_n$, a $b_n:B_n\to C_n$ obrazem $d_{n+1}$. Rozkład $d_{n+1}=b_ne_{n+1}$ z $e_{n+1}$ epi daje $d_nb_n=0$. Istnieje więc jedyny monomorfizm $j_n:B_n\to Z_n$ z $z_nj_n=b_n$; jest mono, ponieważ $b_n$ jest mono. Definiujemy
\[
 H_n(C_\bullet)=\operatorname{coker}(j_n)
\]
w sensie przeciwdziedziny kojądra $\pi_n:Z_n\to H_n(C_\bullet)$. Obiekty $Z_n,B_n$ nazywamy odpowiednio obiektami cykli i brzegów. Dla kompleksu kołańcuchowego stosujemy ten sam schemat z indeksami odwróconymi: $Z^n=\ker d^n$, $B^n=\operatorname{im}d^{n-1}$, $H^n=\operatorname{coker}(B^n\to Z^n)$, gdzie symbole oznaczają odpowiednie obiekty.
\end{definicja}
\begin{stwierdzenie}\label{cat:prop:homology-zero}
Kompleks jest dokładny w stopniu $n$ wtedy i tylko wtedy, gdy $H_n=0$ jako obiekt.
\end{stwierdzenie}
\begin{proof}
Kojądro $j_n$ jest zerowe dokładnie wtedy, gdy $j_n$ jest epi (stwierdzenie~\ref{cat:prop:mono-epi-zero-kernel}). Ponieważ $j_n$ jest mono, twierdzenie~\ref{cat:thm:abelian-factorization} daje równoważność z tym, że jest izomorfizmem, czyli że obraz i jądro określają ten sam podobiekt $C_n$. To jest dokładność.
\end{proof}
\begin{twierdzenie}[Funktorialność i niezmienniczość homotopijna]\label{cat:thm:abelian-homology-functor}
Dla ustalonego wyboru jąder, obrazów i kojąder $H_n$ jest funktorem addytywnym z kategorii kompleksów łańcuchowych do kategorii abelowej. Morfizmy homotopijne indukują ten sam morfizm homologii. Analogiczne twierdzenie obowiązuje dla kohomologii.
\end{twierdzenie}
\begin{proof}
Dla morfizmu kompleksów $f:C_\bullet\to D_\bullet$ mamy $d_n^D f_nz_n^C=f_{n-1}d_n^Cz_n^C=0$. Istnieje jedyny $f_Z:Z_n^C\to Z_n^D$ z $z_n^Df_Z=f_nz_n^C$. Ponadto
\[
 f_nb_n^Ce_{n+1}^C=f_nd_{n+1}^C=d_{n+1}^Df_{n+1}
 =b_n^De_{n+1}^Df_{n+1}.
\]
Jeśli $q$ jest kojądrem $b_n^D$, to $q f_nb_n^Ce_{n+1}^C=0$, a $e_{n+1}^C$ jest epi, więc $qf_nb_n^C=0$. Ponieważ $b_n^D=\ker q$ (twierdzenie~\ref{cat:thm:abelian-factorization}), istnieje jedyny $f_B:B_n^C\to B_n^D$ z $b_n^Df_B=f_nb_n^C$. Z monomorficzności $z_n^D$ wynika $f_Zj_n^C=j_n^Df_B$. Zatem $\pi_n^Df_Zj_n^C=0$, co daje jedyny morfizm $H_n(f)$ z $H_n(f)\pi_n^C=\pi_n^Df_Z$. Jednoznaczność konstrukcji na jądrach i kojądrach daje zachowywanie identyczności i złożeń. Addytywność wynika z równości $(f+g)_Z=f_Z+g_Z$ (składamy z mono $z_n^D$) i $(H_n(f)+H_n(g))\pi_n^C=\pi_n^D(f_Z+g_Z)$; skracamy epi $\pi_n^C$.

Jeżeli $f-g=d^Dh+hd^C$, to $(f_n-g_n)z_n^C=d_{n+1}^Dh_nz_n^C=b_n^De_{n+1}^Dh_nz_n^C$. Po skróceniu mono $z_n^D$ otrzymujemy $f_Z-g_Z=j_n^De_{n+1}^Dh_nz_n^C$. Składając z $\pi_n^D$, otrzymujemy zero, więc $(H_n(f)-H_n(g))\pi_n^C=0$. Epimorficzność $\pi_n^C$ kończy dowód. Dla kompleksów kołańcuchowych zamieniamy $C_n$ na $C^{-n}$ i stosujemy dowód bez zmian.
\end{proof}

\subsection{Lewe i prawe funktory pochodne}
\begin{definicja}[Funktory pochodne]\label{cat:def:derived-functors}
Niech $F:\mathcal A\to\mathcal B$ będzie funktorem addytywnym między kategoriami abelowymi. Jeżeli $\mathcal A$ ma dostatecznie wiele obiektów rzutowych, wybieramy rezolucję rzutową $P_\bullet(X)\to X$ i kładziemy
\[
 L_nF(X)=H_n(F(P_\bullet(X)))\quad(n\ge0).
\]
Jeżeli $\mathcal A$ ma dostatecznie wiele obiektów injektywnych, wybieramy rezolucję injektywną $X\to I^\bullet(X)$ i kładziemy
\[
 R^nF(X)=H^n(F(I^\bullet(X)))\quad(n\ge0).
\]
W obu wzorach pomijamy augmentację lub koaugmentację przed zastosowaniem homologii. Addytywność $F$ zapewnia zachowanie morfizmów zerowych, zatem obrazy kompleksów pozostają kompleksami. Przy wyborach dotyczących wszystkich obiektów pracujemy w ustalonym uniwersum, w którym można dokonać takich wyborów.
\end{definicja}
\begin{twierdzenie}[Dobrze określone funktory pochodne]\label{cat:thm:derived-well-defined}
Powyższe konstrukcje wyznaczają funktory addytywne, niezależne od wyboru rezolucji z dokładnością do kanonicznych izomorfizmów naturalnych zgodnych z kolejną zmianą wyborów.
\end{twierdzenie}
\begin{proof}
Dla $u:X\to Y$ wybieramy podniesienie $f:P_\bullet(X)\to P_\bullet(Y)$ z twierdzenia~\ref{cat:thm:resolution-comparison}. Definiujemy $L_nF(u)=H_n(Ff)$. Dwa podniesienia są homotopijne; zastosowanie addytywnego $F$ do równań homotopii zachowuje wszystkie sumy i złożenia, więc ich obrazy pozostają homotopijne. Twierdzenie~\ref{cat:thm:abelian-homology-functor} daje tę samą strzałkę homologii. Złożenie podniesień podnosi złożenie morfizmów, identyczność podnosi identyczność, a suma podniesień podnosi sumę; stąd funktorialność i addytywność.
Dla drugiego wyboru rezolucji $Q_\bullet(X)$ podniesienie $\id_X$ daje morfizm $P_\bullet(X)\to Q_\bullet(X)$. Podniesienie w drugą stronę jest jego odwrotnością z dokładnością do homotopii na mocy wniosku~\ref{cat:cor:resolution-homotopy-equivalence}. Po zastosowaniu $F$ i $H_n$ otrzymujemy izomorfizm $\alpha_X$. Dla $u:X\to Y$ oba złożenia w kwadracie naturalności przed zastosowaniem $F$ i $H_n$ podnoszą $u$ z $P_\bullet(X)$ do $Q_\bullet(Y)$, więc są homotopijne. Stąd kwadrat jest przemienny. Przy trzecim wyborze rezolucji złożenie dwóch izomorfizmów odpowiada kolejnemu podniesieniu $\id_X$, więc jest równe bezpośredniemu izomorfizmowi. Dla $R^nF$ stosujemy wniosek~\ref{cat:cor:injective-comparison} i kohomologiczną wersję twierdzenia~\ref{cat:thm:abelian-homology-functor}; wszystkie argumenty o podniesieniach zastępujemy przedłużeniami.
\end{proof}
\begin{stwierdzenie}[Stopień zerowy i zanikanie]\label{cat:prop:derived-degree-zero}
Jeżeli $F$ jest prawostronnie dokładny, to $L_0F\cong F$ naturalnie. Jeżeli jest lewostronnie dokładny, to $R^0F\cong F$ naturalnie. Dla obiektu rzutowego $P$ mamy $L_nF(P)=0$ dla $n>0$, a dla injektywnego $I$ mamy $R^nF(I)=0$ dla $n>0$. Jeśli $\operatorname{pd}X\le r$, to $L_nF(X)=0$ dla $n>r$; dualnie dla wymiaru injektywnego.
\end{stwierdzenie}
\begin{proof}
W rezolucji rzutowej $P_1\xrightarrow{d_1}P_0\xrightarrow{\epsilon}X\to0$ augmentacja jest kojądrem $d_1$: ich obraz i jądro są równe, a $\epsilon$ jest epi, więc własność uniwersalna kojądra $\ker\epsilon$ jest także własnością kojądra $d_1$. Prawostronna dokładność oznacza zachowywanie kojąder (twierdzenie~\ref{cat:thm:exact-functor-kernels}); zatem $F\epsilon$ jest kojądrem $Fd_1$. Ponieważ w kompleksie bez augmentacji $Z_0=FP_0$, mamy $L_0F(X)\cong FX$. Dla podniesienia $f$ morfizmu $u$ zachodzi $F\epsilon_Y Ff_0=Fu F\epsilon_X$; własność kojądra daje naturalność. Dualnie $F\eta$ jest jądrem $Fd^0$ dla lewostronnie dokładnego $F$, co daje naturalne $R^0F\cong F$.
Dla $P$ wybieramy rezolucję z identycznością w stopniu zerowym i zerami w dodatnich stopniach. Funktor addytywny zachowuje obiekty zerowe, więc wyższe homologie są zerowe. Dla skończonej rezolucji długości $r$ obiekty w stopniach $n>r$ są zerowe i otrzymujemy to samo zanikanie. Argument injektywny jest dualny.
\end{proof}
\begin{stwierdzenie}\label{cat:prop:exact-derived-zero}
Jeżeli $F$ jest dokładny, wszystkie jego istniejące funktory pochodne dodatniego stopnia znikają.
\end{stwierdzenie}
\begin{proof}
Funktor dokładny zachowuje dokładność ciągów, w szczególności rezolucji z augmentacją lub koaugmentacją. Po usunięciu tych strzałek dokładność w dodatnich stopniach pozostaje bez zmian. Stwierdzenie~\ref{cat:prop:homology-zero} daje tezę.
\end{proof}

\subsection{Długi ciąg dokładny homologii}
\begin{twierdzenie}[Długi ciąg dokładny]\label{cat:thm:abelian-homology-les}
Niech $0\to A_\bullet\xrightarrow{i}B_\bullet\xrightarrow{p}C_\bullet\to0$ będzie ciągiem kompleksów łańcuchowych dokładnym w każdym stopniu. Istnieją naturalne morfizmy łączące $\partial_n:H_n(C_\bullet)\to H_{n-1}(A_\bullet)$, dla których ciąg
\[
\cdots\to H_n(A_\bullet)\xrightarrow{H_n(i)}H_n(B_\bullet)
\xrightarrow{H_n(p)}H_n(C_\bullet)
\xrightarrow{\partial_n}H_{n-1}(A_\bullet)\to\cdots
\]
jest dokładny. Dla kompleksów kołańcuchowych morfizm łączący zwiększa stopień o jeden.
\end{twierdzenie}
\begin{proof}
Piszemy $z_n^X:Z_n^X\to X_n$ dla jądra różniczki, $\pi_n^X:Z_n^X\to H_n(X)$ dla projekcji na homologię, a $d_{n+1}^X=z_n^X j_n^X e_{n+1}^X$ dla rozkładu różniczki przez cykle i brzegi, gdzie $e_{n+1}^X$ jest epi.

\emph{Konstrukcja.}
Tworzymy iloczyn włóknisty $D=B_n\times_{C_n} Z_n^C$ z projekcjami $t:D\to B_n$, $r:D\to Z_n^C$. Morfizm $r$ jest epi. Ponieważ $p_{n-1}d_n^Bt=d_n^Cz_n^Cr=0$, istnieje jedyny $a:D\to A_{n-1}$ z $i_{n-1}a=d_n^Bt$. Mamy $i_{n-2}d_{n-1}^Aa=d_{n-1}^Bd_n^Bt=0$, więc $d_{n-1}^Aa=0$. Istnieje jedyny $v:D\to Z_{n-1}^A$ z $z_{n-1}^Av=a$. Jeśli $rh=0$, to $p_nth=0$, więc $th=i_nu$ dla pewnego $u$ i $z_{n-1}^Avh=d_n^Au$. Stąd $vh=j_{n-1}^Ae_n^Au$ i $\pi_{n-1}^Avh=0$. W szczególności $\pi_{n-1}^Av$ zeruje się na jądrze $r$ i rozkłada się jednoznacznie przez epi $r$ jako $\lambda r$, gdzie $\lambda:Z_n^C\to H_{n-1}(A)$.

Aby $\lambda$ rozkładało się przez $\pi_n^C$, wystarczy wykazać $\lambda j_n^C=0$. Tworzymy morfizm $w:B_{n+1}\to D$ o składowych $d_{n+1}^B$ i $j_n^C e_{n+1}^C p_{n+1}$; warunek iloczynu włóknistego wynika z przemienności $p$ z różniczkami. Wówczas $vw=0$, ponieważ $i_{n-1}z_{n-1}^Avw=d_n^Bd_{n+1}^B=0$. Zatem $\lambda j_n^C e_{n+1}^C p_{n+1}=0$. Ostatnie dwa morfizmy są epi, więc $\lambda j_n^C=0$. Otrzymujemy jedyny morfizm $\partial_n$ z
\begin{equation}\label{cat:eq:homology-connecting}
 \partial_n\pi_n^Cr=\pi_{n-1}^Av.
\end{equation}

\emph{Podnoszenie lokalne.}
W dowodzie dokładności stosujemy lemat~\ref{cat:lem:abelian-local-lift}: gdy chcemy podnieść morfizm przez epi, możemy zastąpić jego dziedzinę $T$ przez dziedzinę epi $T'\to T$. Wszystkie wcześniej uzyskane morfizmy składamy wówczas z tym epi. Piszemy równości bez oznaczania kolejnych zmian dziedziny; każda równość po zmianie ma wspólną nową dziedzinę. Kryterium~\ref{cat:lem:exactness-local-criterion} pozwala w ten sposób dowodzić dokładności. W szczególności każdy morfizm do $H_n(X)$ można lokalnie podnieść do $Z_n^X$, a cykl o zerowej klasie homologii można lokalnie zapisać jako różniczkę: jest to zastosowanie lematu do dokładnego ciągu $X_{n+1}\to Z_n^X\to H_n(X)\to0$, którego dokładność wynika z definicji obrazu i kojądra.

\emph{Zerowanie złożeń i dokładność.}
Przez cykl w poniższych rachunkach rozumiemy morfizm $x:T\to X_n$ z $d_n^Xx=0$, a przez $[x]$ złożenie jego jednoznacznego rozkładu przez $z_n^X$ z $\pi_n^X$. Każdy morfizm homologii jest lokalnie takiej postaci. Mamy $H_n(p)H_n(i)=0$ z $pi=0$. Dla cyklu $b$ w $B_n$ równanie~\eqref{cat:eq:homology-connecting} daje $\partial_n[p_nb]=0$, gdyż można użyć podniesienia $b$ o zerowej różniczce. Dla lokalnego podniesienia $b$ cyklu $c$ z $C_n$ mamy $p_nb=c$, $d_n^Bb=i_{n-1}a$ i $\partial_n[c]=[a]$. Zatem $H_{n-1}(i)\partial_n[c]=[d_n^Bb]=0$. Te złożenia zerują się na wszystkich morfizmach, gdyż epi użyte do zmian dziedziny można skrócić.

Jeśli $b$ jest cyklem w $B_n$ i $H_n(p)[b]=0$, to lokalnie $p_nb=d_{n+1}^Cc'$ dla $c':T'\to C_{n+1}$. Podnosimy $c'$ przez epi $p_{n+1}$ do $b'$. Wtedy $p_n(b-d_{n+1}^Bb')=0$, więc $b-d_{n+1}^Bb'=i_na$. Monomorficzność $i_{n-1}$ daje $d_n^Aa=0$. Stąd $[b]=H_n(i)[a]$.

Jeśli $c$ jest cyklem w $C_n$ i $\partial_n[c]=0$, podnosimy $c$ lokalnie do $b$ przez $p_n$. Mamy $d_n^Bb=i_{n-1}a$ z $a$ cyklem w $A_{n-1}$ i $[a]=0$. Lokalnie $a=d_n^Aa'$; wtedy $b-i_na'$ jest cyklem w $B_n$ i $p_n(b-i_na')=c$. Zatem $[c]=H_n(p)[b-i_na']$.

Jeśli $a$ jest cyklem w $A_{n-1}$ i $H_{n-1}(i)[a]=0$, to lokalnie $i_{n-1}a=d_n^Bb$ dla pewnego $b$. Wtedy $c=p_nb$ jest cyklem, gdyż $d_n^Cc=p_{n-1}i_{n-1}a=0$. Równanie~\eqref{cat:eq:homology-connecting} daje $\partial_n[c]=[a]$. Te trzy rachunki wraz z kryterium lokalnego podnoszenia dowodzą dokładności we wszystkich stopniach.

\emph{Naturalność.}
Dla morfizmu krótkich ciągów kompleksów $(f_A,f_B,f_C)$ lokalne podniesienie $b$ cyklu $c$ przechodzi na $f_{B,n}b$, podniesienie $f_{C,n}c$. Jeśli $d_n^Bb=i_{n-1}a$, to $d_n^{B'}f_{B,n}b=i'_{n-1}f_{A,n-1}a$. Z definicji morfizmu łączącego wynika zatem
$\partial'_nH_n(f_C)[c]=H_{n-1}(f_A)\partial_n[c]$.
Każda klasa ma lokalnego reprezentanta, więc po skróceniu epi otrzymujemy równość morfizmów. Pozostałe kwadraty są przemienne przez funktorialność homologii. Odwrócenie indeksów daje wersję kołańcuchową.
\end{proof}

\begin{twierdzenie}[Długie ciągi funktorów pochodnych]\label{cat:thm:derived-les}
Dla addytywnego funktora $F:\mathcal A\to\mathcal B$ oraz krótkiego ciągu dokładnego $0\to X\to Y\to Z\to0$ istnieją naturalne długie ciągi
\[
 \cdots\to L_nF(X)\to L_nF(Y)\to L_nF(Z)
 \to L_{n-1}F(X)\to\cdots\to L_0F(Z)\to0
\]
przy dostatecznej liczbie obiektów rzutowych, a przy dostatecznej liczbie injektywnych
\[
 0\to R^0F(X)\to R^0F(Y)\to R^0F(Z)
 \to R^1F(X)\to R^1F(Y)\to\cdots .
\]
\end{twierdzenie}
\begin{proof}
Lemat o podkowie (twierdzenie~\ref{cat:thm:horseshoe}) daje krótki ciąg rezolucji rozkładalny w każdym stopniu. Addytywny $F$ zachowuje biprodukty, więc po jego zastosowaniu ciąg pozostaje dokładny w każdym stopniu, niezależnie od dokładności $F$. Twierdzenie~\ref{cat:thm:abelian-homology-les} daje długi ciąg, a twierdzenie~\ref{cat:thm:derived-well-defined} identyfikuje jego obiekty i strzałki niełączące z ustalonymi funktorami pochodnymi.

Dla naturalności potrzebujemy zgodnych morfizmów między ciągami rezolucji z lematu o podkowie. Oznaczmy w stopniu $n$ składniki przez $P_n,T_n=P_n\oplus Q_n,Q_n$ i analogicznie z primem. W stopniu $0$ wybieramy podniesienia $f_P,f_Q$ morfizmów obiektów skrajnych przez augmentacje. Na środkowym obiekcie szukamy macierzy
\[
 f_T=\begin{pmatrix}f_P&h\\0&f_Q\end{pmatrix}.
\]
Jeśli na razie przyjmiemy $h=0$, różnica między wymaganym złożeniem z augmentacją a aktualnym złożeniem znika na $P_0$ i po projekcji do obiektu skrajnego prawego. Na $Q_0$ rozkłada się więc przez obiekt skrajny lewy docelowego ciągu. Rzutowość $Q_0$ pozwala podnieść tę różnicę przez augmentację $P'_0$; otrzymane $h:Q_0\to P'_0$ usuwa różnicę. Tak powstały przemienny diagram indukuje morfizm krótkich ciągów jąder. W następnym stopniu stosujemy ten sam argument do tych jąder i epimorfizmów użytych w konstrukcji podkowy. Indukcja daje morfizm krótkich ciągów rezolucji zgodny z augmentacjami. Twierdzenie o naturalności homologii daje naturalność strzałek łączących. Stosując konstrukcję do identyczności krótkiego ciągu i dwóch wyborów podków, otrzymujemy zgodność strzałek łączących z kanonicznymi izomorfizmami zmiany rezolucji. Wersję injektywną otrzymujemy odwracając strzałki w kategorii przeciwnej i stosując wersję kołańcuchową homologii.
\end{proof}

\subsection{Funktory Ext i Tor}
\begin{definicja}\label{cat:def:ext-tor}
Niech $R$ będzie pierścieniem z jedynką. Dla lewych $R$-modułów $M,N$ definiujemy
\[
 \operatorname{Ext}_R^n(M,N)=H^n(\Hom_R(P_\bullet(M),N)).
\]
Jest to prawy funktor pochodny kontrawariantnego funktora $\Hom_R(-,N)$ rozpatrywanego na kategorii przeciwnej; rezolucja rzutowa w kategorii modułów staje się tam rezolucją injektywną. Dla prawego $R$-modułu $U$ i lewego $R$-modułu $M$ definiujemy
\[
 \operatorname{Tor}_n^R(U,M)=H_n(U\otimes_R P_\bullet(M)).
\]
Jest to lewy funktor pochodny funktora $U\otimes_R-$ o wartościach w grupach abelowych. Nie zakładamy przemienności $R$. Obie konstrukcje są niezależne od wyboru rezolucji na mocy twierdzenia~\ref{cat:thm:derived-well-defined}.
\end{definicja}
\begin{stwierdzenie}\label{cat:prop:ext-tor-basic}
Mamy naturalne izomorfizmy $\operatorname{Ext}_R^0(M,N)\cong\Hom_R(M,N)$ oraz $\operatorname{Tor}_0^R(U,M)\cong U\otimes_R M$. Ext jest kontrawariantny w pierwszej i kowariantny w drugiej zmiennej, a Tor kowariantny w obu. Przy rzutowym $M$ oba funktory znikają w dodatnich stopniach. Dla injektywnego $N$ również $\operatorname{Ext}_R^n(M,N)=0$ dla $n>0$.
\end{stwierdzenie}
\begin{proof}
Funktor $\Hom_R(-,N)$ jest lewostronnie dokładny na kategorii przeciwnej, natomiast $U\otimes_R-$ jest prawostronnie dokładny, gdyż jest lewym sprzężonym funktora $\Hom_{\mathbb Z}(U,-)$ z lewym działaniem $(rf)(u)=f(ur)$. Własność sprzężenia wynika z bijekcji między homomorfizmami $U\otimes_R M\to D$ i $M\to\Hom_{\mathbb Z}(U,D)$, danej przez $g\mapsto(m\mapsto(u\mapsto g(u\otimes m)))$; relacja balansowania $ur\otimes m=u\otimes rm$ odpowiada $R$-liniowości. Twierdzenie~\ref{cat:thm:adjoint-exactness} i stwierdzenie~\ref{cat:prop:derived-degree-zero} dają izomorfizmy stopnia zerowego oraz zanikanie przy rzutowym $M$. Morfizmy w pozostałej zmiennej działają stopniowo przez postkompozycję w Hom lub przez iloczyn tensorowy z identycznością; są to morfizmy kompleksów. Komutują z morfizmami pochodzącymi z rezolucji, gdyż składanie homomorfizmów jest łączne, a iloczyn tensorowy jest bifunktorem. Stąd otrzymujemy bifunktory. Przy injektywnym $N$ funktor $\Hom_R(-,N)$ jest dokładny na mocy twierdzenia~\ref{cat:thm:projective-characterization}, więc wyższe funktory znikają.
\end{proof}
\begin{przyklad}[Grupy cykliczne]\label{cat:ex:ext-tor-cyclic}
Dla grupy abelowej $N$ i $m\ge2$ stosujemy rezolucję $0\to\mathbb Z\xrightarrow{m}\mathbb Z\to\mathbb Z/m\mathbb Z\to0$ z przykładu~\ref{cat:ex:cyclic-projective-resolution}. Utożsamienia $\Hom_{\mathbb Z}(\mathbb Z,N)\cong N$, $f\mapsto f(1)$, i $N\otimes_{\mathbb Z}\mathbb Z\cong N$, $x\otimes k\mapsto kx$, przekształcają odpowiednie różniczki na mnożenie przez $m$. Stąd
\[
\begin{aligned}
 \operatorname{Ext}_{\mathbb Z}^1(\mathbb Z/m\mathbb Z,N)&\cong N/mN,\\
 \operatorname{Tor}_1^{\mathbb Z}(N,\mathbb Z/m\mathbb Z)&\cong\{x\in N:mx=0\}.
\end{aligned}
\]
Wyższe stopnie $n\ge2$ znikają. Jeśli $N=\mathbb Z/k\mathbb Z$ i $g=\gcd(m,k)$, to jądro mnożenia przez $m$ jest generowane przez klasę $k/g$ o rzędzie $g$, a obraz przez klasę $g$: z tożsamości Bézouta wynika $g\in m\mathbb Z+k\mathbb Z$. Iloraz przez obraz ma więc rząd $g$ i jest cykliczny. Obie grupy są izomorficzne z $\mathbb Z/g\mathbb Z$, co nie oznacza, że Ext i Tor są tym samym funktorem.
\end{przyklad}

\subsection{Kompleksy podwójne i zrównoważoność}
\begin{lemat}[Porównanie z kompleksem totalnym]\label{cat:lem:total-acyclic}
Niech $D_{p,q}$, $p,q\ge0$, będzie podwójnym kompleksem grup abelowych z różniczkami poziomą $d_h$ i pionową $d_v$, które komutują i osobno mają kwadrat zero. Kompleks totalny ma $\operatorname{Tot}_n D=\bigoplus_{p+q=n}D_{p,q}$ i różniczkę $d=d_h+(-1)^p d_v$ na $D_{p,q}$.
Jeżeli wszystkie poziome wiersze mają dokładne augmentacje $D_{0,q}\to E_q\to0$, zgodne z różniczkami pionowymi, to naturalny morfizm $\operatorname{Tot}D\to E_\bullet$ indukuje izomorfizmy homologii. Analogicznie dla kompleksów podwójnych kołańcuchowych $D^{p,q}$ z różniczkami zwiększającymi indeksy, jeśli pionowe kolumny mają dokładne koaugmentacje $0\to E^p\to D^{p,0}$, zgodne z różniczkami poziomymi, to morfizm $E^\bullet\to\operatorname{Tot}D$ indukuje izomorfizmy kohomologii.
\end{lemat}
\begin{proof}
Z przemienności różniczek wynika $d^2=0$: dwa mieszane składniki na $D_{p,q}$ mają współczynniki $(-1)^p$ oraz $(-1)^{p-1}$; w wersji kołańcuchowej odpowiednio $(-1)^p$ i $(-1)^{p+1}$. W obu przypadkach się znoszą. Wszystkie sumy w ustalonym stopniu są skończone.

W wersji łańcuchowej projekcja na $p=0$ i augmentacja dają surjektywny w każdym stopniu morfizm kompleksów. Jego jądro jest totalizacją podwójnego kompleksu $K$ z $K_{p,q}=D_{p,q}$ dla $p>0$ oraz $K_{0,q}=\ker(D_{0,q}\to E_q)$. Wiersze $K$ są dokładne. Pokażemy, że każdy cykl $x$ stopnia $n$ w $\operatorname{Tot}K$ jest brzegiem. Niech $q$ będzie największym indeksem pionowym niezerowej składowej $x_{p,q}$, gdzie $p+q=n$. Jeśli $p>0$, równanie $dx=0$ w składowej $(p-1,q)$ daje $d_hx_{p,q}=0$, gdyż nie ma składowej o indeksie pionowym $q+1$. Dla $p=0$ ten sam warunek zachodzi przez przyjęcie $K_{-1,q}=0$. Dokładność wiersza pozwala wybrać $y_{p+1,q}$ z $d_hy_{p+1,q}=x_{p,q}$. Odejmując $dy_{p+1,q}$, usuwamy tę składową i zmieniamy co najwyżej składową o indeksie pionowym $q-1$. Powtarzamy aż do $q=0$, otrzymując $x=d(\sum y)$. Jądro jest zatem acykliczne. Długi ciąg dokładny z twierdzenia~\ref{cat:thm:abelian-homology-les} daje tezę.

W wersji kołańcuchowej koaugmentacja daje monomorfizm kompleksów. Jego iloraz jest totalizacją podwójnego kompleksu $K$ o dokładnych kolumnach pionowych: w stopniu $q=0$ zastępujemy $D^{p,0}$ ilorazem przez obraz $E^p$, a w wyższych stopniach nic nie zmieniamy. Dla cocyklu $x$ stopnia $n$ wybieramy najmniejszy indeks poziomy $p$ niezerowej składowej $x^{p,q}$, $p+q=n$. Równanie $dx=0$ w składowej $(p,q+1)$ daje $d_vx^{p,q}=0$, ponieważ nie ma składowej o indeksie $p-1$. Jeśli $q=0$, dokładność kolumny wymusza $x^{p,0}=0$. Jeśli $q>0$, istnieje $y^{p,q-1}$ z $d_vy^{p,q-1}=(-1)^p x^{p,q}$. Odejmując $dy^{p,q-1}$, usuwamy składową o indeksie $p$, zmieniając co najwyżej następną o indeksie $p+1$. Po skończenie wielu krokach otrzymujemy zero. Iloraz jest acykliczny, więc długi ciąg dokładny kohomologii daje izomorfizm. Obie konstrukcje są naturalne względem morfizmów podwójnych kompleksów zgodnych z augmentacjami.
\end{proof}
\begin{uwaga}
W lemacie można zamienić role wierszy i kolumn. Przy transpozycji $(p,q)\mapsto(q,p)$ mnożenie składowej przez $(-1)^{pq}$ daje izomorfizm kompleksów totalnych: dla różniczki poziomej zmiana znaku wynosi $(-1)^q$, a dla pionowej $(-1)^p$. Dotyczy to zarówno zwiększania, jak i zmniejszania indeksów.
\end{uwaga}

\begin{twierdzenie}[Zrównoważoność Ext]\label{cat:thm:ext-balanced}
Dla lewych $R$-modułów $M,N$ i dowolnej rezolucji injektywnej $N\to I^\bullet$ mamy naturalny izomorfizm
\[
 \operatorname{Ext}_R^n(M,N)\cong H^n(\Hom_R(M,I^\bullet)).
\]
\end{twierdzenie}
\begin{proof}
Wybieramy rezolucję rzutową $P_\bullet\to M$. Rozważamy $D^{p,q}=\Hom_R(P_p,I^q)$ z różniczkami $d_h(f)=f d_{p+1}^P$ i $d_v(f)=d_I^q f$. Mają one kwadrat zero i komutują przez łączność składania. Dla ustalonego $p$ funktor $\Hom_R(P_p,-)$ jest dokładny przez rzutowość $P_p$. Kolumna ma więc dokładną koaugmentację $\Hom_R(P_p,N)\to D^{p,0}$. Dla ustalonego $q$ dokładność $\Hom_R(-,I^q)$ przez injektywność $I^q$ daje dokładną koaugmentację wiersza $\Hom_R(M,I^q)\to D^{0,q}$. Obie koaugmentacje są zgodne z pozostałymi różniczkami przez przemienność złożeń. Lemat~\ref{cat:lem:total-acyclic} oraz transpozycja dają izomorfizmy
\[
 H^n(\Hom_R(P_\bullet,N))
 \longrightarrow H^n(\operatorname{Tot}D)
 \longleftarrow H^n(\Hom_R(M,I^\bullet)).
\]
Złożenie pierwszego z odwrotnością drugiego daje żądany izomorfizm. Dla morfizmu w pierwszej zmiennej wybieramy podniesienie między rezolucjami rzutowymi i stosujemy prekompozycję; dla morfizmu w drugiej zmiennej wybieramy przedłużenie między rezolucjami injektywnymi i stosujemy postkompozycję. Otrzymujemy morfizmy podwójnych kompleksów zgodne z obiema koaugmentacjami. Przemienność kwadratów indukowanych w kohomologii daje naturalność. Dla identyczności obiektów i zmiany rezolucji ta sama przemienność daje zgodność z kanonicznymi izomorfizmami z twierdzenia~\ref{cat:thm:derived-well-defined}. Zatem otrzymany izomorfizm nie zależy od wyboru rezolucji.
\end{proof}

\begin{definicja}[Moduł płaski]\label{cat:def:flat}
Lewy moduł $F$ jest płaski, jeśli funktor $-\otimes_R F$ na prawych modułach jest dokładny; prawy moduł $G$ jest płaski, jeśli $G\otimes_R-$ jest dokładny na lewych modułach. Por. \atlasref{AA-18::aa:def:flat-module}{AA-18: moduły płaskie}.
\end{definicja}
\begin{stwierdzenie}\label{cat:prop:projective-flat}
Każdy moduł rzutowy jest płaski.
\end{stwierdzenie}
\begin{proof}
Dla modułów wolnych $R^{(S)}$ funktor iloczynu tensorowego z nimi utożsamia się z sumą prostą $S$ kopii drugiego modułu, która zachowuje dokładność: cykle i obrazy wyznacza się współrzędnie, a element ma skończenie wiele niezerowych współrzędnych. Utożsamienie wynika z własności uniwersalnej iloczynu tensorowego; por. \atlasref{AA-18::aa:tw:tensor-direct-sum}{AA-18: iloczyn tensorowy i sumy proste} oraz \atlasref{AA-18::aa:tw:tensor-unit}{AA-18: iloczyn tensorowy z pierścieniem}. Moduł rzutowy jest składnikiem sumy prostej modułu wolnego (dla prawych modułów stosujemy twierdzenie o lewych modułach nad $R^{\op}$). Iloczyn tensorowy rozkłada się zgodnie z tą sumą prostą. Dokładność sumy dwóch ciągów implikuje dokładność każdego składnika, ponieważ jądra i obrazy rozkładają się na sumy prostych odpowiednich jąder i obrazów. Stąd płaskość; por. również \atlasref{AA-18::aa:tw:projective-flat}{AA-18: moduły rzutowe są płaskie}.
\end{proof}
\begin{twierdzenie}[Zrównoważoność Tor i rezolucje płaskie]\label{cat:thm:tor-balanced}
Niech $U$ będzie prawym, a $M$ lewym $R$-modułem. Jeżeli $G_\bullet\to U$ jest dokładną rezolucją przez płaskie prawe moduły, a $F_\bullet\to M$ dokładną rezolucją przez płaskie lewe moduły, to naturalnie
\[
 \operatorname{Tor}_n^R(U,M)\cong H_n(G_\bullet\otimes_R M)
 \cong H_n(U\otimes_R F_\bullet).
\]
W szczególności dowolna rezolucja rzutowa prawego modułu $U$ może służyć do obliczania Tor.
\end{twierdzenie}
\begin{proof}
Najpierw rozważamy dwie rezolucje rzutowe $Q_\bullet\to U$, $P_\bullet\to M$. Kompleks podwójny $D_{p,q}=Q_p\otimes_R P_q$ ma komutujące różniczki $d_h=d^Q\otimes1$, $d_v=1\otimes d^P$; komutują one dzięki bifunktorialności iloczynu tensorowego. Dokładność wierszy augmentowanych do $U\otimes_R P_q$ wynika z płaskości $P_q$, a kolumn augmentowanych do $Q_p\otimes_R M$ z płaskości $Q_p$. Lemat~\ref{cat:lem:total-acyclic} i transpozycja dają izomorfizmy
\[
 H_n(U\otimes_R P_\bullet)
 \longleftarrow H_n(\operatorname{Tot}D)
 \longrightarrow H_n(Q_\bullet\otimes_R M).
\]
Podniesienia morfizmów $U,M$ do rezolucji rzutowych indukują morfizmy kompleksów podwójnych zgodne z augmentacjami; naturalność lematu daje naturalność izomorfizmu. Niezależność od wyboru rezolucji wynika z zastosowania tej samej naturalności do podniesień identyczności.

Następnie stosujemy argument do $D_{p,q}=Q_p\otimes_R F_q$. Dokładność augmentowanych wierszy wynika z płaskości $F_q$, a kolumn z płaskości $Q_p$. Otrzymujemy izomorfizm $H_n(U\otimes_R F_\bullet)\cong H_n(Q_\bullet\otimes_R M)$, a ten drugi obiekt oblicza Tor na mocy pierwszej części dowodu. Zamiana $Q,P$ na $G,P$ daje analogiczny wynik dla $H_n(G_\bullet\otimes_R M)$. Nie zakładamy istnienia podniesień dowolnych morfizmów między rezolucjami płaskimi: do porównania służą podniesienia z rezolucji rzutowych (twierdzenie~\ref{cat:thm:resolution-comparison}). Dla dowolnego takiego podniesienia otrzymujemy przemienny diagram z augmentacjami w kompleksach podwójnych po użyciu rezolucji rzutowej w drugiej zmiennej. Stąd izomorfizm jest kanoniczny i zgodny z morfizmami rezolucji płaskich, gdy takie morfizmy są dane. Naturalne działanie na Tor jest zawsze określone przez rezolucje rzutowe.
\end{proof}
\begin{wniosek}\label{cat:cor:tor-flat}
Jeżeli $M$ lub $U$ jest płaski, to $\operatorname{Tor}_n^R(U,M)=0$ dla $n>0$.
\end{wniosek}
\begin{proof}
Dla płaskiego modułu można wybrać rezolucję płaską długości zero z identycznością jako augmentacją. Twierdzenie~\ref{cat:thm:tor-balanced} pozwala obliczyć Tor przez tę rezolucję. W dodatnich stopniach kompleks jest zerowy, więc jego homologia znika.
\end{proof}

\section{Kategorie triangulowane i pochodne}\label{cat:sec:triangulated-derived}
W tym rozdziale przechodzimy od pojedynczych grup homologii i funktorów pochodnych do kategorii, w których kompleksy są obiektami, a homotopie i quasi-izomorfizmy wpływają na pojęcie morfizmu. Korzystamy z rozdziałów~\ref{cat:sec:abelian-categories}, \ref{cat:sec:projective-injective} i~\ref{cat:sec:derived-functors}. Wersję łańcuchową homotopii omawia \atlasref{TOP-03::def:homotopia-lancuchowa}{TOP-03: homotopia łańcuchowa}. Przyjmujemy w całym rozdziale konwencję kołańcuchową.

\subsection{Kategoria homotopijna kompleksów}
\begin{definicja}\label{cat:def:complex-category}
Dla kategorii addytywnej $\mathcal A$ przez $\cat{C}(\mathcal A)$ oznaczamy kategorię kompleksów kołańcuchowych $X=(X^n,d_X^n)_{n\in\mathbb Z}$ z $d_X^{n+1}d_X^n=0$. Morfizm $f:X\to Y$ jest rodziną $f^n:X^n\to Y^n$ spełniającą $d_Y^nf^n=f^{n+1}d_X^n$. Homotopią między $f,g:X\to Y$ nazywamy rodzinę $h^n:X^n\to Y^{n-1}$ taką, że
\begin{equation}\label{cat:eq:cochain-homotopy}
 f^n-g^n=d_Y^{n-1}h^n+h^{n+1}d_X^n.
\end{equation}
Piszemy wówczas $f\simeq g$. Kategoria $\cat{K}(\mathcal A)$ ma te same obiekty co $\cat{C}(\mathcal A)$, a morfizmy są klasami homotopii. Symbole $\cat{C}^+,\cat{C}^-,\cat{C}^b$ i analogicznie $\cat{K}^+,\cat{K}^-,\cat{K}^b$ oznaczają pełne podkategorie kompleksów ograniczonych odpowiednio z dołu, z góry i z obu stron (w sensie zerowania się obiektów poza odpowiednim zakresem stopni).
\end{definicja}
\begin{stwierdzenie}\label{cat:prop:homotopy-category}
Kategoria $\cat{K}(\mathcal A)$ jest dobrze określoną kategorią addytywną, a funktor ilorazowy $\cat{C}(\mathcal A)\to\cat{K}(\mathcal A)$ jest addytywny.
\end{stwierdzenie}
\begin{proof}
Morfizmy kompleksów tworzą grupę abelową z dodawaniem stopniowym. Dla dowolnej rodziny $h^n:X^n\to Y^{n-1}$ morfizmy $u^n=d_Y^{n-1}h^n+h^{n+1}d_X^n$ komutują z różniczkami: obie strony równania $d_Y^nu^n=u^{n+1}d_X^n$ są równe $d_Y^nh^{n+1}d_X^n$. Takie morfizmy tworzą podgrupę, ponieważ suma i przeciwieństwo rodzin $h$ realizują odpowiednio sumę i przeciwieństwo $u$. Relacja homotopii jest więc relacją równoważności i iloraz zbioru morfizmów jest grupą abelową.
Jeśli $f-g=d_Yh+hd_X$, $a:W\to X$ i $b:Y\to Z$ są morfizmami kompleksów, to
\[
 b(f-g)a=d_Z(bha)+(bha)d_W.
\]
Wzór wynika z przemienności $a,b$ z różniczkami, z zachowaniem odpowiednich stopni. Homotopijność jest zatem zgodna ze składaniem i składanie klas jest dobrze określone, dwuliniowe i łączne. Identyczności reprezentują identyczności kompleksów. Obiekt zerowy i biprodukty w kategorii kompleksów otrzymujemy stopniowo; ich równania definiujące zachodzą także po przejściu do klas, więc są biproduktami w kategorii homotopijnej. Otrzymujemy kategorię addytywną oraz addytywny funktor ilorazowy.
\end{proof}
\begin{definicja}\label{cat:def:contractible-complex}
Kompleks $X$ nazywamy ściągalnym, jeśli $\id_X\simeq0$. Morfizm $f:X\to Y$ jest równoważnością homotopijną, jeśli istnieje $g:Y\to X$ z $gf\simeq\id_X$ i $fg\simeq\id_Y$.
\end{definicja}
\begin{stwierdzenie}\label{cat:prop:contractible-zero}
Kompleks jest ściągalny wtedy i tylko wtedy, gdy jest obiektem zerowym w $\cat{K}(\mathcal A)$. Równoważności homotopijne to dokładnie morfizmy, których klasy są izomorfizmami. Jeśli $\mathcal A$ jest abelowa, każdy ściągalny kompleks jest acykliczny, tj. $H^n(X)=0$ dla wszystkich $n$.
\end{stwierdzenie}
\begin{proof}
W kategorii addytywnej obiekt jest zerowy dokładnie wtedy, gdy jego identyczność jest zerowa: każdy morfizm do niego i z niego jest złożeniem z tą identycznością, a więc zerem. W $\cat{K}(\mathcal A)$ oznacza to $\id_X\simeq0$. Definicja izomorfizmu klas daje wprost oba równania homotopijne dla reprezentanta odwrotności. W kategorii abelowej kohomologia zachowuje identyczności i utożsamia morfizmy homotopijne na mocy twierdzenia~\ref{cat:thm:abelian-homology-functor}. Stąd $\id_{H^n(X)}=H^n(\id_X)=H^n(0)=0$, więc $H^n(X)$ jest zerowe.
\end{proof}
\begin{przyklad}\label{cat:ex:acyclic-not-contractible}
Kompleks grup abelowych skoncentrowany w stopniach $-1,0,1$,
\[
 \mathbb Z\xrightarrow{m}\mathbb Z\longrightarrow\mathbb Z/m\mathbb Z
 \qquad(m\ge2),
\]
gdzie druga strzałka jest projekcją, jest acykliczny, ponieważ wraz z zerami na końcach jest ciągiem dokładnym. Nie jest jednak ściągalny. Gdyby $\id_X=d_Xh+hd_X$, w stopniu $1$ mielibyśmy $\id_{\mathbb Z/m\mathbb Z}=d_X^0h^1$, ponieważ $d_X^1=0$. Każdy homomorfizm $h^1:\mathbb Z/m\mathbb Z\to\mathbb Z$ jest zerowy (obraz generatora jest elementem torsyjnym grupy bez torsji), co daje sprzeczność. Por. \atlasref{TOP-03::prz:acykliczny-niesciagalny}{TOP-03: kompleks acykliczny, który nie jest ściągalny}.
\end{przyklad}

\subsection{Przesunięcie kompleksu}
\begin{definicja}\label{cat:def:complex-shift}
Dla $r\in\mathbb Z$ definiujemy kompleks $X[r]$ i morfizm $f[r]$ przez
\[
 X[r]^n=X^{n+r},\qquad d_{X[r]}^n=(-1)^r d_X^{n+r},
 \qquad f[r]^n=f^{n+r}.
\]
W szczególności $X[1]$ ma różniczkę przeciwną do odpowiednio przesuniętej różniczki $X$.
\end{definicja}
\begin{stwierdzenie}\label{cat:prop:shift-functor}
Przesunięcie $[r]$ wyznacza addytywny automorfizm $\cat{C}(\mathcal A)$ oraz $\cat{K}(\mathcal A)$, z odwrotnością $[-r]$. Zachodzą równości $[r][s]=[r+s]$ i $[0]=\id$. Jeżeli $\mathcal A$ jest abelowa, to $H^n(X[r])\cong H^{n+r}(X)$ naturalnie w $X$.
\end{stwierdzenie}
\begin{proof}
Kwadrat różniczki przesuniętej jest zerem, ponieważ mnożymy złożenie dwóch kolejnych różniczek przez $(-1)^{2r}=1$. Równanie morfizmu kompleksów zachowuje się po pomnożeniu przez $(-1)^r$ i zmianie indeksu. Składanie, identyczność i dodawanie zachowują się stopniowo. Dla homotopii $h$ z~\eqref{cat:eq:cochain-homotopy} definiujemy $\widetilde h^n=(-1)^r h^{n+r}$. Wówczas
\[
 d_{Y[r]}^{n-1}\widetilde h^n+\widetilde h^{n+1}d_{X[r]}^n
 =d_Y^{n+r-1}h^{n+r}+h^{n+r+1}d_X^{n+r}
 =f[r]^n-g[r]^n.
\]
Przesunięcie przechodzi więc na kategorie homotopijne. Wzory na składanie przesunięć wynikają z dodawania indeksów oraz $(-1)^r(-1)^s=(-1)^{r+s}$. Jądro i obraz morfizmu w kategorii abelowej nie zmieniają się jako podobiekty po pomnożeniu go przez $-1$, ponieważ $-\id$ jest izomorfizmem. Stąd obiekty cykli, brzegów i ich ilorazy w $X[r]$ odpowiadają tym z $X$ o indeksie przesuniętym o $r$. Identyfikacje komutują z indukowanymi morfizmami, ponieważ $f[r]^n=f^{n+r}$.
\end{proof}

\subsection{Stożek morfizmu}
\begin{definicja}\label{cat:def:mapping-cone}
Dla morfizmu kompleksów $f:X\to Y$ definiujemy stożek $\operatorname{Cone}(f)$ przez
\[
 \operatorname{Cone}(f)^n=Y^n\oplus X^{n+1},\qquad
 d_{\operatorname{Cone}(f)}^n=
 \begin{pmatrix}d_Y^n&f^{n+1}\\0&-d_X^{n+1}\end{pmatrix}.
\]
Definiujemy $i:Y\to\operatorname{Cone}(f)$ jako włożenie pierwszego składnika i $p:\operatorname{Cone}(f)\to X[1]$ jako projekcję na drugi składnik.
\end{definicja}
\begin{stwierdzenie}\label{cat:prop:cone-complex}
Stożek jest kompleksem, a $i,p$ są morfizmami kompleksów. W kategorii abelowej ciąg
\begin{equation}\label{cat:eq:cone-ses}
 0\to Y\xrightarrow{i}\operatorname{Cone}(f)\xrightarrow{p}X[1]\to0
\end{equation}
jest dokładny i rozkładalny w każdym stopniu. Złożenia kolejnych strzałek $X\xrightarrow f Y\xrightarrow i\operatorname{Cone}(f)\xrightarrow p X[1]\xrightarrow{f[1]}Y[1]$ są zerowe w kategorii homotopijnej.
\end{stwierdzenie}
\begin{proof}
Iloczyn dwóch macierzy różniczek ma zerowe składniki diagonalne przez $d_X^2=d_Y^2=0$. Składnik górny prawy to $d_Y^{n+1}f^{n+1}-f^{n+2}d_X^{n+1}=0$, ponieważ $f$ jest morfizmem kompleksów. Włożenie i projekcja komutują z różniczkami przez postać macierzy, a dokładność stopniowa wynika z biproduktu. Mamy $pi=0$ już w kategorii kompleksów. Dla $if$ wybieramy homotopię $h^n:X^n\to\operatorname{Cone}(f)^{n-1}$, $h^n(x)=(0,x)$, gdzie zapis współrzędnych oznacza włożenie do biproduktu (nie zakładamy istnienia elementów obiektów). Rachunek macierzowy daje $d_{\operatorname{Cone}(f)}h+hd_X=(f,0)^{\mathsf T}=if$. Dla $f[1]p$ wybieramy homotopię $k^n:\operatorname{Cone}(f)^n\to Y[1]^{n-1}=Y^n$ równą projekcji na pierwszy składnik. Wówczas
\[
 d_{Y[1]}^{n-1}k^n+k^{n+1}d_{\operatorname{Cone}(f)}^n
 =[-d_Y^n,0]+[d_Y^n,f^{n+1}]=[0,f^{n+1}].
\]
To jest $f[1]^np^n$, więc wszystkie wskazane złożenia są zerowe w $\cat K(\mathcal A)$.
\end{proof}
\begin{stwierdzenie}\label{cat:prop:cone-homotopic}
Jeżeli $f,g:X\to Y$ są homotopijne, to istnieje izomorfizm kompleksów $\operatorname{Cone}(f)\to\operatorname{Cone}(g)$ komutujący z włożeniami $Y$ i projekcjami na $X[1]$.
\end{stwierdzenie}
\begin{proof}
Dla $f-g=d_Yh+hd_X$ wybieramy w stopniu $n$ macierz
\[
 T^n=\begin{pmatrix}1&h^{n+1}\\0&1\end{pmatrix}.
\]
Górne prawe składniki $d_{\operatorname{Cone}(g)}^nT^n$ i $T^{n+1}d_{\operatorname{Cone}(f)}^n$ to odpowiednio $d_Y^nh^{n+1}+g^{n+1}$ i $f^{n+1}-h^{n+2}d_X^{n+1}$. Są równe przez założenie. Pozostałe składniki są identyczne. Macierz odwrotna powstaje przez zastąpienie $h^{n+1}$ przez $-h^{n+1}$, a ponieważ $T$ komutuje z różniczkami, jego odwrotność również z nimi komutuje. Równania dotyczące włożenia i projekcji wynikają z postaci trójkątnej macierzy.
\end{proof}
\begin{twierdzenie}[Kohomologia stożka]\label{cat:thm:cone-les}
Jeżeli $\mathcal A$ jest abelowa, to istnieje naturalny długi ciąg dokładny
\[
 \cdots\to H^n(X)\xrightarrow{H^n(f)}H^n(Y)
 \xrightarrow{H^n(i)}H^n(\operatorname{Cone}(f))
 \xrightarrow{H^n(p)}H^{n+1}(X)\to\cdots.
\]
W konsekwencji $f$ indukuje izomorfizmy we wszystkich stopniach kohomologii wtedy i tylko wtedy, gdy $\operatorname{Cone}(f)$ jest acykliczny.
\end{twierdzenie}
\begin{proof}
Stosujemy długi ciąg z twierdzenia~\ref{cat:thm:abelian-homology-les} do~\eqref{cat:eq:cone-ses}. Identyfikujemy $H^n(X[1])$ z $H^{n+1}(X)$ na mocy stwierdzenia~\ref{cat:prop:shift-functor}. Morfizm łączący jest $H^{n+1}(f)$: dla cyklu $x:T\to X^{n+1}$ jego podniesienie do stożka ma składowe $(0,x)$, a różniczka podniesienia ma składowe $(f^{n+1}x,0)$. Konstrukcja morfizmu łączącego w~\eqref{cat:eq:homology-connecting}, po odwróceniu indeksów, daje klasę $f^{n+1}x$. Każdy morfizm do kohomologii można lokalnie podnieść do cyklu przez epi; skracając ten epi, uzyskujemy równość morfizmów w kategorii abelowej. Naturalność względem przemiennych kwadratów morfizmów kompleksów wynika z morfizmu ciągów~\eqref{cat:eq:cone-ses} o środkowej macierzy diagonalnej z odpowiednich morfizmów oraz naturalności długiego ciągu.

Jeżeli każde $H^n(f)$ jest izomorfizmem, dokładność daje $H^n(i)=0$ oraz $\ker H^{n+1}(f)=0$, więc $H^n(p)$ jest jednocześnie mono i zerowy. To wymusza $H^n(\operatorname{Cone}(f))=0$. Odwrotnie, przy zerowej kohomologii stożka dokładność po obu stronach $H^n(f)$ daje jego monomorficzność i epimorficzność. W kategorii abelowej oznacza to izomorfizm.
\end{proof}
\begin{definicja}\label{cat:def:quasi-isomorphism}
Morfizm kompleksów w kategorii abelowej nazywamy quasi-izomorfizmem, jeśli indukuje izomorfizmy $H^n(f)$ dla wszystkich $n\in\mathbb Z$. Własność ta zależy tylko od klasy homotopii na mocy twierdzenia~\ref{cat:thm:abelian-homology-functor}.
\end{definicja}
\begin{uwaga}
Równoważność homotopijna jest quasi-izomorfizmem, ponieważ kohomologia jest funktorem na kategorii homotopijnej. Odwrotność nie zachodzi: morfizm $0\to X$ dla kompleksu z przykładu~\ref{cat:ex:acyclic-not-contractible} jest quasi-izomorfizmem, ale nie jest równoważnością homotopijną przez stwierdzenie~\ref{cat:prop:contractible-zero}.
\end{uwaga}

\subsection{Trójkąty wyróżnione i pierwsze aksjomaty}
\begin{definicja}\label{cat:def:triangle}
Niech $\mathcal T$ będzie kategorią addytywną z ustalonym addytywnym autoizomorfizmem $[1]$ (ogólniej autoequiwalencją z ustaloną odwrotnością z dokładnością do izomorfizmu naturalnego). Trójkątem nazywamy diagram
\[
 X\xrightarrow f Y\xrightarrow u U\xrightarrow v X[1].
\]
Morfizmem trójkątów jest trójka $(a,b,c)$ morfizmów na obiektach $X,Y,U$, dla której $bf=f'a$, $cu=u'b$ oraz $a[1]v=v'c$. Jest to izomorfizm trójkątów wtedy i tylko wtedy, gdy $a,b,c$ są izomorfizmami: odwrotne morfizmy spełniają odpowiednie równania po złożeniu równań wyjściowych z odwrotnościami. Rotacją w przód nazywamy trójkąt
\[
 Y\xrightarrow u U\xrightarrow v X[1]\xrightarrow{-f[1]}Y[1],
\]
a rotacją wstecz trójkąt $U[-1]\xrightarrow{-v[-1]}X\xrightarrow f Y\xrightarrow u U$, gdzie stosujemy ustalone utożsamienia odwrotnych przesunięć.
\end{definicja}
\begin{definicja}\label{cat:def:cone-triangle}
W $\cat{K}(\mathcal A)$ trójkątem wyróżnionym nazywamy każdy trójkąt izomorficzny w kategorii trójkątów do trójkąta stożka
\begin{equation}\label{cat:eq:cone-triangle}
 X\xrightarrow{[f]}Y\xrightarrow{[i]}\operatorname{Cone}(f)
 \xrightarrow{[p]}X[1],
\end{equation}
dla pewnego morfizmu kompleksów $f$. Nawiasy kwadratowe oznaczają tu klasy homotopii morfizmów, w odróżnieniu od przesunięcia obiektu. Stwierdzenie~\ref{cat:prop:cone-homotopic} gwarantuje, że zmiana reprezentanta $[f]$ daje izomorficzny trójkąt stożka z identycznościami na $X,Y$.
\end{definicja}
\begin{stwierdzenie}[Warunki TR1--TR3]\label{cat:prop:tr-123}
Trójkąty wyróżnione w $\cat{K}(\mathcal A)$ spełniają następujące warunki.
\begin{enumerate}[label=\textup{(TR\arabic*)}]
 \item Klasa trójkątów wyróżnionych jest zamknięta na izomorfizmy. Trójkąt $X\xrightarrow{\id}X\to0\to X[1]$ jest wyróżniony. Każdy morfizm $X\to Y$ jest pierwszą strzałką pewnego trójkąta wyróżnionego.
 \item Trójkąt jest wyróżniony wtedy i tylko wtedy, gdy jego rotacja w przód jest wyróżniona.
 \item Dla dwóch trójkątów wyróżnionych $X\xrightarrow f Y\xrightarrow u U\xrightarrow v X[1]$ oraz $X'\xrightarrow{f'}Y'\xrightarrow{u'}U'\xrightarrow{v'}X'[1]$ każdy przemienny kwadrat $bf=f'a$ z morfizmami $a:X\to X'$ i $b:Y\to Y'$ można uzupełnić morfizmem $c:U\to U'$ do morfizmu trójkątów.
\end{enumerate}
Te same warunki zachodzą w $\cat{K}^+(\mathcal A)$, $\cat{K}^-(\mathcal A)$ i $\cat{K}^b(\mathcal A)$.
\end{stwierdzenie}
\begin{proof}
\emph{TR1.} Zamkniętość na izomorfizmy wynika z definicji. Dla stożka identyczności kontrakcję daje rodzina $h^n(x,y)=(0,x)$ z $\operatorname{Cone}(\id_X)^n=X^n\oplus X^{n+1}$ do stopnia $n-1$. Istotnie $dh(x,y)=(x,-d_X^nx)$ oraz $hd(x,y)=(0,d_X^nx+y)$, więc $dh+hd=\id$. Na mocy stwierdzenia~\ref{cat:prop:contractible-zero} stożek jest obiektem zerowym w kategorii homotopijnej; zastąpienie go przez $0$ daje żądany trójkąt. Dla dowolnej klasy morfizmu wybieramy reprezentanta i trójkąt~\eqref{cat:eq:cone-triangle}.

\emph{TR2, rotacja w przód.} Wystarczy rozważyć trójkąt stożka, gdyż rotacja przenosi izomorfizmy trójkątów na izomorfizmy. W rachunkach współrzędne oznaczają składowe morfizmu do biproduktu, a nie elementy obiektu kategorii. Oznaczmy $C=\operatorname{Cone}(f)$ i $B=\operatorname{Cone}(i)$ dla $i:Y\to C$. W stopniu $n$ piszemy
\[
 B^n=Y^n\oplus X^{n+1}\oplus Y^{n+1},\qquad
 d_B(y,x,z)=(d_Yy+fx+z,-d_Xx,-d_Yz).
\]
Pomijamy indeksy różniczek i morfizmów, gdy wynikają jednoznacznie z typów składowych. Morfizmy
\[
 r:B\to X[1],\quad r(y,x,z)=x,\qquad
 s:X[1]\to B,\quad s(x)=(0,x,-fx)
\]
komutują z różniczkami: dla $r$ jest to projekcja na $-d_Xx$, a dla $s$ mamy $d_Bs(x)=(0,-d_Xx,d_Yfx)=s(-d_Xx)$, ponieważ $f$ komutuje z różniczkami. Ponadto $rs=\id$. Rodzina $h^n:B^n\to B^{n-1}$ dana przez $h(y,x,z)=(0,0,y)$ spełnia
\[
 d_Bh+hd_B=(y,0,fx+z)=\id_B-sr.
\]
Stąd $r$ jest izomorfizmem w kategorii homotopijnej. Dla włożenia $j:C\to B$ zachodzi $rj=p$. Dla projekcji $q:B\to Y[1]$, $q(y,x,z)=z$, rodzina $k^n:B^n\to Y[1]^{n-1}$ z $k(y,x,z)=y$ daje
\[
 d_{Y[1]}k+kd_B=-d_Yy+(d_Yy+fx+z)=fx+z=q+f[1]r.
\]
A zatem $[q]=-[f[1]][r]$. Trójka $(\id_Y,\id_C,[r])$ jest izomorfizmem trójkąta stożka $i$ na rotację trójkąta stożka $f$.

\emph{TR2, rotacja wstecz.} Żeby uzyskać implikację odwrotną, pokażemy bezpośrednio zamkniętość na rotację wstecz. Dla $C=\operatorname{Cone}(f)$ bierzemy $B'=\operatorname{Cone}(-p[-1])$; jego składowe i różniczka są
\[
 B'^n=X^n\oplus Y^n\oplus X^{n+1},\qquad
 d_{B'}(x,y,z)=(d_Xx-z,d_Yy+fz,-d_Xz).
\]
Morfizmy $r':B'\to Y$, $r'(x,y,z)=fx+y$, oraz $s':Y\to B'$, $s'(y)=(0,y,0)$ są morfizmami kompleksów: $r'd_{B'}=fd_Xx-fz+d_Yy+fz=d_Yr'$ oraz $d_{B'}s'=s'd_Y$. Mamy $r's'=\id_Y$. Dla $h'^n:B'^n\to B'^{n-1}$, $h'(x,y,z)=(0,0,-x)$, rachunek daje
\[
 d_{B'}h'+h'd_{B'}=(x,-fx,z)=\id_{B'}-s'r'.
\]
Zatem $[r']$ jest izomorfizmem. Włożenie $j':X\to B'$ spełnia $r'j'=f$. Projekcja $q':B'\to C$ spełnia $q'(x,y,z)=(y,z)$. Różnica $q'-ir'$ ma składowe $(-fx,z)$ i jest homotopijna z zerem przez $k'^n:B'^n\to C^{n-1}$ dane wzorem $k'(x,y,z)=(0,-x)$:
\[
 d_Ck'+k'd_{B'}=(-fx,d_Xx)+(0,-d_Xx+z)=(-fx,z).
\]
Trójka $(\id_{C[-1]},\id_X,[r'])$ daje więc izomorfizm trójkąta stożka $-p[-1]$ na rotację wstecz trójkąta stożka $f$. Rotacje w przód i wstecz są wzajemnie odwrotne po identyfikacji przesunięć; zastosowanie rotacji wstecz do wyróżnionej rotacji w przód dowodzi drugiej implikacji.

\emph{TR3.} Przenosząc kwadrat przez izomorfizmy trójkątów, sprowadzamy się do trójkątów stożków $f,f'$. Wybieramy reprezentanty $a,b$ w kategorii kompleksów. Przemienność kwadratu w kategorii homotopijnej oznacza istnienie rodziny $h$ z $bf-f'a=d_{Y'}h+hd_X$. Definiujemy morfizm
\[
 c^n:\operatorname{Cone}(f)^n\to\operatorname{Cone}(f')^n,\qquad
 c^n=\begin{pmatrix}b^n&h^{n+1}\\0&a^{n+1}\end{pmatrix}.
\]
Składniki diagonalne równości $d_{\operatorname{Cone}(f')}^nc^n=c^{n+1}d_{\operatorname{Cone}(f)}^n$ są równe, gdyż $a,b$ komutują z różniczkami; dolne lewe są zerowe. Równość górnych prawych składników ma postać
\[
 d_{Y'}^nh^{n+1}+f'{}^{n+1}a^{n+1}
 =b^{n+1}f^{n+1}-h^{n+2}d_X^{n+1},
\]
co jest równaniem homotopii. Mamy $ci=i'b$ i $p'c=a[1]p$ już w kategorii kompleksów, więc $([a],[b],[c])$ jest morfizmem trójkątów. Po przeniesieniu z powrotem przez izomorfizmy dostajemy żądane uzupełnienie.

Wszystkie użyte konstrukcje (przesunięcia, skończone biprodukty i stożki) zachowują ograniczenie z dołu, z góry lub z obu stron; podkategorie są pełne, więc dowód działa bez zmian.
\end{proof}
\begin{uwaga}
Warunki TR1--TR3 nie są jeszcze pełną definicją kategorii triangulowanej. Potrzebny jest dodatkowo aksjomat ośmiościanu, który wiąże trójkąty odpowiadające dwóm składanym morfizmom i ich złożeniu. Uzupełnienie morfizmu trójkątów w TR3 nie jest na ogół jednoznaczne; dowód wymaga wyboru homotopii przemienności kwadratu, a nie tylko klas jego strzałek.
\end{uwaga}

\section{Rozszerzenia Kana}\label{cat:sec:kan-extensions}
Rozszerzenia Kana uogólniają konstrukcję granic, kogranic i sprzężeń: zamiast rozszerzać funkcję z podzbioru, poszukujemy uniwersalnego przedłużenia funktora wzdłuż innego funktora, który nie musi być włożeniem. Wykorzystamy kategorie przecinkowe z definicji~\ref{cat:def:comma}, granice i kogranice z rozdziału~\ref{sec:granice}, transformacje naturalne z rozdziału~\ref{sec:naturalne} oraz sprzężenia z rozdziału~\ref{sec:adjunctions}. W całym rozdziale $\mathcal A,\mathcal B$ są małymi kategoriami, $\mathcal C$ jest lokalnie mała. Kategorie funktorów oznaczamy przez $[\mathcal A,\mathcal C]$ i $[\mathcal B,\mathcal C]$; ich zbiory transformacji naturalnych są podzbiorami małych iloczynów zbiorów morfizmów.

\subsection{Własności uniwersalne}
\begin{definicja}[Lewe i prawe rozszerzenie Kana]\label{cat:def:kan-extension}
Niech $K:\mathcal A\to\mathcal B$ i $F:\mathcal A\to\mathcal C$. Lewym rozszerzeniem Kana funktora $F$ wzdłuż $K$ nazywamy parę $(L,\eta)$ z $L:\mathcal B\to\mathcal C$ i transformacją $\eta:F\Rightarrow LK$, dla której dla każdego $H:\mathcal B\to\mathcal C$ odwzorowanie
\[
 \Nat(L,H)\longrightarrow\Nat(F,HK),\qquad
 \alpha\longmapsto(\alpha K)\eta
\]
jest bijekcją. Piszemy $L=\operatorname{Lan}_K F$.
Prawym rozszerzeniem Kana nazywamy parę $(R,\epsilon)$ z $R:\mathcal B\to\mathcal C$ i $\epsilon:RK\Rightarrow F$, dla której dla każdego $H$ odwzorowanie
\[
 \Nat(H,R)\longrightarrow\Nat(HK,F),\qquad
 \beta\longmapsto\epsilon(\beta K)
\]
jest bijekcją. Piszemy $R=\operatorname{Ran}_K F$. Iloczyn transformacji oznacza ich składanie pionowe; $(\alpha K)_a=\alpha_{K(a)}$.
\end{definicja}
\begin{center}
\begin{tabular}{cc}
\(\begin{array}{ccc}
\mathcal A&\xrightarrow{K}&\mathcal B\\
&\searrow F&\downarrow L\\
&&\mathcal C
\end{array}\)
&
\(\begin{array}{ccc}
\mathcal A&\xrightarrow{K}&\mathcal B\\
&\searrow F&\downarrow R\\
&&\mathcal C
\end{array}\)\\
\(\eta:F\Rightarrow LK\)&\(\epsilon:RK\Rightarrow F\)
\end{tabular}
\end{center}
Nie wymagamy $LK=F$ ani nawet izomorficzności tych funktorów. Warunek uniwersalny dotyczy transformacji, nie równości funktorów.

\begin{stwierdzenie}[Jednoznaczność]\label{cat:prop:kan-unique}
Dwa lewe rozszerzenia Kana $(L,\eta),(L',\eta')$ są połączone jedynym izomorfizmem naturalnym $\theta:L\Rightarrow L'$ spełniającym $(\theta K)\eta=\eta'$. Dla prawych rozszerzeń analogicznie istnieje jedyny izomorfizm $\theta:R\Rightarrow R'$ z $\epsilon'(\theta K)=\epsilon$.
\end{stwierdzenie}
\begin{proof}
Z własności uniwersalnej $L$ otrzymujemy jedyną $\theta$ z podanym równaniem, a z własności $L'$ jedyną $\gamma:L'\Rightarrow L$ z $(\gamma K)\eta'=\eta$. Wtedy $((\gamma\theta)K)\eta=(\gamma K)\eta'=\eta$, więc jednoznaczność daje $\gamma\theta=\id_L$. Analogicznie $((\theta\gamma)K)\eta'=\eta'$ daje $\theta\gamma=\id_{L'}$. W przypadku prawym własność $R'$ daje $\theta$, a własność $R$ daje $\gamma:R'\Rightarrow R$ z $\epsilon(\gamma K)=\epsilon'$. Mamy $\epsilon((\gamma\theta)K)=\epsilon'(\theta K)=\epsilon$ i $\epsilon'((\theta\gamma)K)=\epsilon(\gamma K)=\epsilon'$, więc oba złożenia są identycznościami.
\end{proof}
\begin{twierdzenie}[Rozszerzenia jako sprzężenia]\label{cat:thm:kan-adjunction}
Niech $K^*:[\mathcal B,\mathcal C]\to[\mathcal A,\mathcal C]$ oznacza prekompozycję z $K$. Jeżeli lewe rozszerzenia istnieją dla wszystkich $F$, wyznaczają funktor $\operatorname{Lan}_K$ lewy sprzężony do $K^*$. Jeżeli prawe rozszerzenia istnieją dla wszystkich $F$, wyznaczają funktor $\operatorname{Ran}_K$ prawy sprzężony do $K^*$. Odwrotnie, istnienie odpowiedniego sprzężenia daje rozszerzenia Kana.
\end{twierdzenie}
\begin{proof}
Wybieramy rozszerzenia $(L_F,\eta_F)$. Dla $\mu:F\Rightarrow F'$ definiujemy $L_\mu:L_F\Rightarrow L_{F'}$ jako jedyną transformację z
\[
 (L_\mu K)\eta_F=\eta_{F'}\mu.
\]
Dla identyczności równanie spełnia $\id_{L_F}$. Dla $\nu:F'\Rightarrow F''$ mamy $((L_\nu L_\mu)K)\eta_F=(L_\nu K)\eta_{F'}\mu=\eta_{F''}\nu\mu$, więc jednoznaczność daje zachowanie składania. Własność uniwersalna daje bijekcję $\Nat(L_F,H)\cong\Nat(F,K^*H)$. Naturalność względem $H$ wynika z łączności składania, a względem $F$ z równania definiującego $L_\mu$: obraz $\alpha L_\mu$ jest równy $(\alpha K)\eta_{F'}\mu$. Jest to sprzężenie w sensie definicji~\ref{def:adjunction}.

Dla prawych rozszerzeń $(R_F,\epsilon_F)$ definiujemy $R_\mu:R_F\Rightarrow R_{F'}$ przez $\epsilon_{F'}(R_\mu K)=\mu\epsilon_F$. Identyczności spełniają to równanie, a dla $\nu$ mamy $\epsilon_{F''}((R_\nu R_\mu)K)=\nu\mu\epsilon_F$, co daje zachowanie składania. Bijekcja $\Nat(H,R_F)\cong\Nat(K^*H,F)$ jest naturalna względem $H$ przez prekompozycję, a względem $F$ przez równanie $\epsilon_{F'}(R_\mu K)=\mu\epsilon_F$. Odwrotnie jednostka sprzężenia $\operatorname{Lan}_K\dashv K^*$ daje $\eta_F$, a kojednostka $K^*\dashv\operatorname{Ran}_K$ daje $\epsilon_F$; bijekcje sprzężenia są dokładnie bijekcjami z definicji~\ref{cat:def:kan-extension}.
\end{proof}

\subsection{Wzory punktowe}
\begin{definicja}\label{cat:def:kan-comma}
Dla $b\in\Ob(\mathcal B)$ używamy dwóch kategorii przecinkowych. W $(K\downarrow b)$ obiektem jest $(a,u:K(a)\to b)$, a morfizmem do $(a',u')$ jest $f:a\to a'$ z $u'K(f)=u$. W $(b\downarrow K)$ obiektem jest $(a,v:b\to K(a))$, a morfizmem do $(a',v')$ jest $f:a\to a'$ z $K(f)v=v'$. Projekcje $P_b^-,P_b^+$ na $\mathcal A$ zapominają strzałkę $u$ lub $v$. Obie kategorie są małe: ich obiekty tworzą sumę rozłączną zbiorów morfizmów indeksowaną obiektami $\mathcal A$, a morfizmy są ograniczone warunkiem przemienności.
\end{definicja}
\begin{twierdzenie}[Punktowe lewe rozszerzenie Kana]\label{cat:thm:pointwise-lan}
Jeżeli dla każdego $b$ istnieje kogranica diagramu $FP_b^-:(K\downarrow b)\to\mathcal C$, to istnieje lewe rozszerzenie Kana i można je obliczyć przez
\[
 (\operatorname{Lan}_K F)(b)=\colim_{(a,u:K(a)\to b)} F(a).
\]
W szczególności warunek zachodzi, gdy $\mathcal C$ jest kozupełna.
\end{twierdzenie}
\begin{proof}
Wybieramy kogranicę $L(b)$ i morfizmy jej kostożka $\lambda^b_{a,u}:F(a)\to L(b)$, zatem $\lambda^b_{a',u'}F(f)=\lambda^b_{a,u}$ dla morfizmów kategorii przecinkowej. Dla $g:b\to c$ rodzina $\lambda^c_{a,gu}$ spełnia to samo równanie, gdyż z $u'K(f)=u$ wynika $(gu')K(f)=gu$. Własność uniwersalna kogranicy wyznacza jedyny $L(g)$ z
\begin{equation}\label{cat:eq:lan-action}
 L(g)\lambda^b_{a,u}=\lambda^c_{a,gu}.
\end{equation}
Dla $g=\id_b$ równanie spełnia $\id_{L(b)}$; dla $h:c\to d$ obie strony $L(hg)$ i $L(h)L(g)$ dają po prekompozycji z każdą $\lambda^b_{a,u}$ morfizm $\lambda^d_{a,hgu}$. Jednoznaczność w kogranicy daje funktorialność.
Kładziemy $\eta_a=\lambda^{K(a)}_{a,\id_{K(a)}}$. Dla $f:a\to a'$ równanie~\eqref{cat:eq:lan-action} daje
\[
 L(K(f))\eta_a=\lambda^{K(a')}_{a,K(f)}
 =\lambda^{K(a')}_{a',\id}F(f)=\eta_{a'}F(f),
\]
gdzie środkowa równość pochodzi z morfizmu $f$ w kategorii $(K\downarrow K(a'))$. Zatem $\eta$ jest naturalna.

Dla $H:\mathcal B\to\mathcal C$ i $\sigma:F\Rightarrow HK$ rodzina $H(u)\sigma_a:F(a)\to H(b)$ jest kostożkiem: jeśli $u'K(f)=u$, to
\[
 H(u')\sigma_{a'}F(f)=H(u')H(K(f))\sigma_a=H(u)\sigma_a.
\]
Istnieje jedyny $\alpha_b:L(b)\to H(b)$ z $\alpha_b\lambda^b_{a,u}=H(u)\sigma_a$. Dla $g:b\to c$ obie strony $H(g)\alpha_b$ i $\alpha_c L(g)$ po prekompozycji z $\lambda^b_{a,u}$ dają $H(gu)\sigma_a$, więc są równe. Mamy $\alpha_{K(a)}\eta_a=\sigma_a$ przez wybór $u=\id$. Jeśli $\alpha'$ również spełnia $(\alpha' K)\eta=\sigma$, to z~\eqref{cat:eq:lan-action} i naturalności $\alpha'$ wynika
\[
 \alpha'_b\lambda^b_{a,u}=\alpha'_b L(u)\eta_a
 =H(u)\alpha'_{K(a)}\eta_a=H(u)\sigma_a.
\]
Jednoznaczność w kogranicy daje $\alpha'=\alpha$. Dowodzi to własności uniwersalnej.
\end{proof}
\begin{twierdzenie}[Punktowe prawe rozszerzenie Kana]\label{cat:thm:pointwise-ran}
Jeżeli dla każdego $b$ istnieje granica diagramu $FP_b^+:(b\downarrow K)\to\mathcal C$, to istnieje prawe rozszerzenie Kana i
\[
 (\operatorname{Ran}_K F)(b)=\lim_{(a,v:b\to K(a))}F(a).
\]
W szczególności warunek zachodzi, gdy $\mathcal C$ jest zupełna.
\end{twierdzenie}
\begin{proof}
Wybieramy granicę $R(b)$ z projekcjami $\rho^b_{a,v}:R(b)\to F(a)$ spełniającymi $F(f)\rho^b_{a,v}=\rho^b_{a',v'}$ dla $K(f)v=v'$. Dla $g:b\to c$ rodzina $\rho^b_{a,wg}$ indeksowana obiektami $(a,w)$ w $(c\downarrow K)$ jest stożkiem, ponieważ $K(f)w=w'$ implikuje $K(f)wg=w'g$. Definiujemy jedyny $R(g):R(b)\to R(c)$ z
\begin{equation}\label{cat:eq:ran-action}
 \rho^c_{a,w}R(g)=\rho^b_{a,wg}.
\end{equation}
Identyczność spełnia to równanie dla $g=\id_b$. Dla $h:c\to d$ postkompozycja obu $R(hg),R(h)R(g)$ z każdą projekcją $\rho^d_{a,z}$ daje $\rho^b_{a,zhg}$, więc z jednoznaczności w granicy otrzymujemy zachowanie składania.
Definiujemy $\epsilon_a=\rho^{K(a)}_{a,\id}$. Dla $f:a\to a'$ mamy
\[
 \epsilon_{a'}R(K(f))=\rho^{K(a)}_{a',K(f)}
 =F(f)\rho^{K(a)}_{a,\id}=F(f)\epsilon_a.
\]
Środkowa równość to warunek stożka w $(K(a)\downarrow K)$.

Dla $\sigma:HK\Rightarrow F$ rodzina $\sigma_a H(v):H(b)\to F(a)$ jest stożkiem, gdyż $F(f)\sigma_a H(v)=\sigma_{a'}H(K(f)v)=\sigma_{a'}H(v')$. Istnieje jedyny $\beta_b:H(b)\to R(b)$ z $\rho^b_{a,v}\beta_b=\sigma_a H(v)$. Dla $g:b\to c$ postkompozycja $R(g)\beta_b$ oraz $\beta_c H(g)$ z każdą projekcją $\rho^c_{a,w}$ daje $\sigma_a H(wg)$, zatem $\beta$ jest naturalna. Mamy $\epsilon_a\beta_{K(a)}=\sigma_a$. Jeżeli $\beta'$ spełnia to samo równanie, to dla każdego $(a,v)$
\[
 \rho^b_{a,v}\beta'_b=\epsilon_a R(v)\beta'_b
 =\epsilon_a\beta'_{K(a)} H(v)=\sigma_a H(v).
\]
Jednoznaczność w granicy daje $\beta'=\beta$.
\end{proof}
\begin{uwaga}
Powyższe dowody podają warunki wystarczające istnienia rozszerzeń. Samo istnienie rozszerzenia w sensie własności uniwersalnej nie jest definicją rozszerzenia punktowego. Rozszerzeniami punktowymi nazywamy tu rozszerzenia obliczane za pomocą wskazanych granic lub kogranic i ich kanonicznych transformacji.
\end{uwaga}
\begin{przyklad}[Pusta kategoria przecinkowa]\label{cat:ex:kan-empty}
Jeżeli $(K\downarrow b)$ jest pusta, wartość punktowego lewego rozszerzenia w $b$ jest obiektem początkowym $\mathcal C$, o ile taki obiekt istnieje. Jeżeli $(b\downarrow K)$ jest pusta, wartość prawego rozszerzenia jest obiektem końcowym. Wynika to z definicji kogranicy i granicy pustego diagramu: rodzina składowych kostożka lub stożka jest pusta, a własność uniwersalna wymaga dokładnie jednego morfizmu w odpowiednim kierunku.
\end{przyklad}

\subsection{Rozszerzenia wzdłuż funktorów w pełni wiernych}
\begin{twierdzenie}\label{cat:thm:kan-fully-faithful}
Jeżeli $K:\mathcal A\to\mathcal B$ jest w pełni wierny, to jednostka $\eta:F\Rightarrow LK$ punktowego lewego rozszerzenia oraz kojednostka $\epsilon:RK\Rightarrow F$ punktowego prawego rozszerzenia (o ile odpowiednie rozszerzenie istnieje) są izomorfizmami naturalnymi.
\end{twierdzenie}
\begin{proof}
Obiekt $(a,\id_{K(a)})$ jest końcowy w $(K\downarrow K(a))$: morfizm z $(x,u)$ do niego to strzałka $f:x\to a$ z $K(f)=u$, istniejąca przez pełność i jedyna przez wierność. Przypomnijmy, że dla diagramu $D:\mathcal J\to\mathcal C$ z obiektem końcowym $t$ w $\mathcal J$ kogranicą jest $D(t)$. Jedynym strzałkom $j_x:x\to t$ odpowiada kostożek $D(j_x)$, ponieważ $j_yh=j_x$ dla $h:x\to y$. Każdy kostożek $(k_x:D(x)\to W)_x$ spełnia $k_x=k_tD(j_x)$, a $D(j_t)=\id$ wymusza jednoznaczność $k_t$. Stosując to do $FP^-_{K(a)}$, otrzymujemy izomorfizm $\eta_a=\lambda^{K(a)}_{a,\id}:F(a)\to L(K(a))$ z jednoznaczności kogranicy.

Dla prawego rozszerzenia $(a,\id_{K(a)})$ jest początkowy w $(K(a)\downarrow K)$: strzałka do $(x,v)$ odpowiada jedynemu $f:a\to x$ z $K(f)=v$. Gdy $s$ jest początkowy w $\mathcal J$, jedyne strzałki $j_x:s\to x$ dają stożek $D(j_x)$, a dowolny stożek $(k_x:W\to D(x))_x$ spełnia $k_x=D(j_x)k_s$. Zatem $D(s)$ jest granicą. Dostajemy izomorfizm $\epsilon_a=\rho^{K(a)}_{a,\id}$. Naturalność obu transformacji wynika z twierdzeń~\ref{cat:thm:pointwise-lan} i~\ref{cat:thm:pointwise-ran}; odwrotności składowych są naturalne przez twierdzenie~\ref{tw:nat-iso-components}.
\end{proof}
\begin{wniosek}\label{cat:cor:kan-fully-faithful}
Jeżeli $K$ jest w pełni wierny, to istniejące dla wszystkich $F$ punktowe funktory $\operatorname{Lan}_K$ i $\operatorname{Ran}_K$ są w pełni wierne (każdy przy odpowiednim założeniu istnienia).
\end{wniosek}
\begin{proof}
Jednostka sprzężenia $\operatorname{Lan}_K\dashv K^*$ i kojednostka $K^*\dashv\operatorname{Ran}_K$ są odpowiednio rodzinami $\eta_F$ i $\epsilon_F$, gdyż odpowiadają identycznościom w bijekcjach z definicji~\ref{cat:def:kan-extension}. Są izomorfizmami przez twierdzenie~\ref{cat:thm:kan-fully-faithful}. Stosujemy kryterium pełnej wierności funktora sprzężonego ze stwierdzenia~\ref{cat:prop:adj-full-faithful}.
\end{proof}
\begin{przyklad}[Porządki]\label{cat:ex:kan-poset}
Dla podzbioru $S$ zbioru uporządkowanego $P$, włożenia $K:S\hookrightarrow P$ z porządkiem indukowanym i monotonicznego $F:S\to Q$ do kraty zupełnej mamy
\[
 L(p)=\bigvee_{s\in S,\ s\le p}F(s),\qquad
 R(p)=\bigwedge_{s\in S,\ p\le s}F(s).
\]
Kategorie przecinkowe odpowiadają wskazanym podzbiorom; kostożki w $Q$ to ograniczenia górne, a stożki to ograniczenia dolne. Stąd wzory punktowe. Dla $s\in S$ obie rodziny zawierają $F(s)$, odpowiednio jako największy i najmniejszy element, więc $L(s)=F(s)=R(s)$. Jeśli monotoniczny $H:P\to Q$ przedłuża $F$, to $F(s)\le H(p)$ dla $s\le p$ oraz $H(p)\le F(s)$ dla $p\le s$. Z własności kresów $L(p)\le H(p)\le R(p)$. Dla $P=\{0<1<2\}$, $S=\{0,2\}$, $F(0)=0,F(2)=2$ i $Q=P$ otrzymujemy $L(1)=0$, $R(1)=2$. Oba rozszerzenia mogą więc być różne.
\end{przyklad}
\begin{przyklad}[Brak pełnej wierności]\label{cat:ex:kan-not-fully-faithful}
Niech $\mathcal A$ będzie dyskretną kategorią dwóch obiektów, $K:\mathcal A\to\mathbf1$ jedynym funktorem, a $F$ przypisuje obu obiektom dwuelementowe zbiory. Lewe rozszerzenie jest ich sumą rozłączną, prawe iloczynem. Jednostka składa się z włożeń, kojednostka z projekcji; żadne nie są bijekcjami. Funktor $K$ jest wierny, ale nie pełny, bo nie podnosi identyczności w $\mathbf1$ do strzałki między różnymi obiektami.
\end{przyklad}

\subsection{Rozszerzenia wzdłuż funktorów sprzężonych}
\begin{twierdzenie}\label{cat:thm:kan-along-adjoint}
Jeżeli $K\dashv T$ z jednostką $u:\id_{\mathcal A}\Rightarrow TK$ i kojednostką $v:KT\Rightarrow\id_{\mathcal B}$, to $(FT,Fu)$ jest punktowym lewym rozszerzeniem Kana $F$ wzdłuż $K$. Jeżeli $S\dashv K$ z jednostką $u':\id_{\mathcal B}\Rightarrow KS$ i kojednostką $v':SK\Rightarrow\id_{\mathcal A}$, to $(FS,Fv')$ jest punktowym prawym rozszerzeniem. Nie potrzeba zupełności ani kozupełności kategorii wartości.
\end{twierdzenie}
\begin{proof}
Dla $K\dashv T$ obiekt $(T(b),v_b)$ jest końcowy w $(K\downarrow b)$, gdyż dla $f:K(a)\to b$ sprzężenie daje dokładnie jedną $h:a\to T(b)$ z $v_bK(h)=f$. Kogranica istnieje i jest $F(T(b))$, z kostożkiem $F(h)$, przez argument z dowodu twierdzenia~\ref{cat:thm:kan-fully-faithful}. Dla $g:b\to c$ transpozycją $gf$ jest $T(g)h$, gdyż $v_cK(T(g))K(h)=gv_bK(h)=gf$. Wzór~\eqref{cat:eq:lan-action} daje zatem $L(g)=F(T(g))$. Transpozycją $\id_{K(a)}$ jest $u_a$ przez równanie trójkątne $v_{K(a)}K(u_a)=\id_{K(a)}$, więc jednostką rozszerzenia jest $Fu$.

Dla $S\dashv K$ obiekt $(S(b),u'_b)$ jest początkowy w $(b\downarrow K)$: dla $f:b\to K(a)$ istnieje jedyna $h:S(b)\to a$ z $K(h)u'_b=f$. Granica jest $F(S(b))$, z projekcjami $F(h)$. Dla $g:b\to c$ i $f:c\to K(a)$ transpozycją $fg$ jest $hS(g)$, bo $K(h)K(S(g))u'_b=K(h)u'_cg=fg$. Wzór~\eqref{cat:eq:ran-action} daje $R(g)=F(S(g))$. Transpozycją $\id_{K(a)}$ jest $v'_a$ przez równanie $K(v'_a)u'_{K(a)}=\id_{K(a)}$. Stąd kojednostka $Fv'$.
\end{proof}
\begin{przyklad}
Dla równoważności kategorii $K$ z funktorem odwrotnym $T$ z dokładnością do izomorfizmu naturalnego oba rozszerzenia są kanonicznie izomorficzne z $FT$. Istotnie równoważność można wyposażyć w oba sprzężenia (rozdział~\ref{sec:adjunctions}), a twierdzenie~\ref{cat:thm:kan-along-adjoint} stosuje się do każdego z nich. Dla włożenia $K:\mathbf1\to\mathcal B$ wskazującego obiekt początkowy $\mathcal B$ mamy $K\dashv T$ z jedynym $T:\mathcal B\to\mathbf1$, więc lewe rozszerzenie dowolnego obiektu jest funktorem stałym; dla obiektu końcowego analogicznie prawe rozszerzenie jest stałe.
\end{przyklad}

\subsection{Zachowywanie, absolutność i składanie}
\begin{definicja}\label{cat:def:kan-preserved}
Funktor $J:\mathcal C\to\mathcal D$ zachowuje lewe rozszerzenie $(L,\eta)$, jeśli $(JL,J\eta)$ jest lewym rozszerzeniem $JF$ wzdłuż $K$. Dla prawego rozszerzenia analogicznie wymagamy, aby $(JR,J\epsilon)$ było prawym rozszerzeniem. Rozszerzenie nazywamy absolutnym, gdy zachowuje je każdy funktor wychodzący z $\mathcal C$ (w rozpatrywanym uniwersum).
\end{definicja}
\begin{stwierdzenie}\label{cat:prop:kan-preservation}
Jeśli $J$ zachowuje kogranice wszystkich diagramów $FP_b^-$, to zachowuje punktowe lewe rozszerzenie Kana. Jeśli zachowuje granice wszystkich $FP_b^+$, to zachowuje punktowe prawe rozszerzenie. W szczególności lewe funktory sprzężone zachowują punktowe lewe rozszerzenia, a prawe funktory sprzężone punktowe prawe rozszerzenia.
\end{stwierdzenie}
\begin{proof}
W przypadku lewym $JL(b)$ z kostożkiem $J\lambda^b_{a,u}$ jest kogranicą $JFP_b^-$. Równanie~\eqref{cat:eq:lan-action} po zastosowaniu $J$ pokazuje, że $JL(g)$ jest dokładnie morfizmem wyznaczonym przez te kostożki; jednostką konstrukcji jest $J\eta$. Twierdzenie~\ref{cat:thm:pointwise-lan} daje tezę. W przypadku prawym $JR(b)$ z projekcjami $J\rho^b_{a,v}$ jest granicą $JFP_b^+$, a równanie~\eqref{cat:eq:ran-action} wyznacza $JR(g)$ i kojednostkę $J\epsilon$. Stosujemy twierdzenie~\ref{cat:thm:pointwise-ran}. Ostatnie zdanie wynika z twierdzenia~\ref{tw:adjoints-preserve}.
\end{proof}
\begin{wniosek}\label{cat:cor:kan-absolute}
Rozszerzenia skonstruowane w twierdzeniu~\ref{cat:thm:kan-along-adjoint} są absolutne.
\end{wniosek}
\begin{proof}
Dla dowolnego $J:\mathcal C\to\mathcal D$ to samo sprzężenie $K\dashv T$ daje lewe rozszerzenie $(JF)T$ z jednostką $(JF)u=J(Fu)$. Jest to dokładnie obraz $(FT,Fu)$. Dla $S\dashv K$ analogicznie $(JF)S$ z kojednostką $(JF)v'=J(Fv')$ jest prawym rozszerzeniem.
\end{proof}
\begin{przyklad}[Rozszerzenie nieabsolutne]\label{cat:ex:kan-not-absolute}
Dla $K$ z dyskretnej kategorii dwóch obiektów do $\mathbf1$ i $F$ o wartościach singletonowych lewe rozszerzenie ma dwa elementy. Funktor $J:\Set\to\Set$, $J(X)=X\times X$, daje zbiór czteroelementowy, natomiast lewe rozszerzenie $JF$ ma dwa elementy. Nie mogą być izomorficzne. Dla prawego rozszerzenia tego samego $F$ funktor $J'(X)=X\sqcup\{*\}$ daje dwa elementy na obrazie iloczynu, a cztery elementy w iloczynie obrazów. Zatem również prawe rozszerzenia nie zawsze są absolutne.
\end{przyklad}
\begin{twierdzenie}[Składanie rozszerzeń]\label{cat:thm:kan-composition}
Niech $K:\mathcal A\to\mathcal B$ i $M:\mathcal B\to\mathcal E$, gdzie również $\mathcal E$ jest mała. Jeśli istnieją kolejne lewe rozszerzenia $(L,\eta)$ wzdłuż $K$ i $(P,\theta)$ funktora $L$ wzdłuż $M$, to $(P,(\theta K)\eta)$ jest lewym rozszerzeniem $F$ wzdłuż $MK$. Dla kolejnych prawych rozszerzeń $(R,\epsilon)$ i $(Q,\zeta)$ otrzymujemy $(Q,\epsilon(\zeta K))$ jako prawe rozszerzenie. W szczególności, gdy istnieją, mamy kanoniczne izomorfizmy
\[
 \operatorname{Lan}_{MK}F\cong\operatorname{Lan}_M(\operatorname{Lan}_K F),
 \qquad
 \operatorname{Ran}_{MK}F\cong\operatorname{Ran}_M(\operatorname{Ran}_K F).
\]
\end{twierdzenie}
\begin{proof}
Dla $H:\mathcal E\to\mathcal C$ mamy złożenie bijekcji
\[
 \Nat(P,H)\longrightarrow\Nat(L,HM)
 \longrightarrow\Nat(F,HMK).
\]
Transformacja $\alpha$ przechodzi kolejno na $(\alpha M)\theta$ i $((\alpha M)\theta K)\eta=(\alpha MK)(\theta K)\eta$. Jest to wymagana bijekcja dla rozszerzenia wzdłuż $MK$. Dla prawego rozszerzenia składamy bijekcje $\Nat(H,Q)\to\Nat(HM,R)\to\Nat(HMK,F)$, a obraz $\beta$ to $\epsilon(\zeta K)(\beta MK)$. Jednoznaczność wynika z~\ref{cat:prop:kan-unique}. Dla trzech kolejnych funktorów oba sposoby grupowania dają tę samą złożoną transformację dzięki łączności składania; jednoznaczność porównań daje zgodność kanonicznych izomorfizmów. W przypadku funktorów rozszerzeń naturalność porównań względem $F$ wynika z równania definiującego działanie na transformacjach w dowodzie twierdzenia~\ref{cat:thm:kan-adjunction}: oba złożenia w kwadracie porównawczym mają tę samą transformację po prekompozycji z $MK$ i jednostką (odpowiednio po złożeniu z kojednostką).
\end{proof}
\begin{wniosek}\label{cat:cor:kan-limit}
Dla jedynego funktora $!:\mathcal A\to\mathbf1$ lewe rozszerzenie jest kogranicą $F$, a prawe granicą $F$; równoważność dotyczy także istnienia tych obiektów.
\end{wniosek}
\begin{proof}
Funktor $\mathbf1\to\mathcal C$ wybiera obiekt $X$, a transformacja $F\Rightarrow X!$ jest kostożkiem. Własność uniwersalna lewego rozszerzenia mówi dokładnie, że każdy kostożek rozkłada się jednoznacznie przez morfizm z jego wierzchołka. Dla transformacji $X!\Rightarrow F$ i prawego rozszerzenia otrzymujemy własność uniwersalną granicy.
\end{proof}

\subsection{Jawne modele w kategorii zbiorów}
\begin{twierdzenie}\label{cat:thm:kan-set}
Dla $F:\mathcal A\to\Set$ lewy wzór punktowy ma postać
\[
 L(b)=\left(\coprod_{a\in\Ob(\mathcal A)}
       \Hom_{\mathcal B}(K(a),b)\times F(a)\right)/{\sim},
\]
gdzie $\sim$ jest najmniejszą relacją równoważności zawierającą
\[
 (a,u'K(f),x)\sim(a',u',F(f)(x))
 \quad(f:a\to a',\ u':K(a')\to b,\ x\in F(a)).
\]
Prawy wzór punktowy ma postać
\[
 R(b)=\left\{(x_{a,v})\in\prod_{(a,v)\in(b\downarrow K)}F(a):
 F(f)(x_{a,v})=x_{a',K(f)v}\ \text{dla wszystkich }f:a\to a'\right\}.
\]
Dla $g:b\to c$ mamy $L(g)[a,u,x]=[a,gu,x]$ oraz
$(R(g)x)_{a,w}=x_{a,wg}$.
\end{twierdzenie}
\begin{proof}
Iloraz jest konstrukcją kogranicy zbiorów z twierdzenia~\ref{tw:colimits-set}, zastosowaną do $(K\downarrow b)$: obiekty indeksujące to pary $(a,u)$, a każda strzałka $f$ narzuca dokładnie wskazaną relację. Włożenia $x\mapsto[a,u,x]$ tworzą kostożek. Dla dowolnego kostożka $k_{a,u}$ funkcja $[a,u,x]\mapsto k_{a,u}(x)$ jest dobrze określona, gdyż zachowuje relacje generujące, a jednoznaczność wynika z tego, że każda klasa ma reprezentanta. Działanie $L(g)$ zachowuje relacje, bo $g(u'K(f))=(gu')K(f)$, i spełnia~\eqref{cat:eq:lan-action}. Analogicznie zbiór zgodnych rodzin jest konstrukcją granicy z twierdzenia~\ref{tw:limits-set}. Dla dowolnego stożka $k_{a,v}:W\to F(a)$ funkcja $w\mapsto(k_{a,v}(w))$ ma wartości w zbiorze zgodnych rodzin, a projekcje wyznaczają ją jednoznacznie. Działanie $R(g)$ zachowuje zgodność, bo $K(f)(wg)=(K(f)w)g$, i spełnia~\eqref{cat:eq:ran-action}. Dla pustego indeksu iloraz jest pusty, a iloczyn singletonem, zgodnie z przykładem~\ref{cat:ex:kan-empty}.
\end{proof}
\begin{stwierdzenie}\label{cat:prop:kan-representables}
Dla $a_0\in\Ob(\mathcal A)$ istnieje naturalny izomorfizm
\[
 \operatorname{Lan}_K\Hom_{\mathcal A}(a_0,-)
 \cong\Hom_{\mathcal B}(K(a_0),-).
\]
Dla presnopów (funktorów kontrawariantnych do $\Set$) analogicznie
\[
 \operatorname{Lan}_{K^{\op}}\Hom_{\mathcal A}(-,a_0)
 \cong\Hom_{\mathcal B}(-,K(a_0)).
\]
Nie trzeba zakładać pełnej wierności $K$.
\end{stwierdzenie}
\begin{proof}
W lewym wzorze z twierdzenia~\ref{cat:thm:kan-set} przypisujemy klasie $[a,u,h:a_0\to a]$ strzałkę $uK(h)$. Funkcja jest dobrze określona, ponieważ $u'K(f)K(h)=u'K(fh)$. Odwrotność wysyła $w:K(a_0)\to b$ na $[a_0,w,\id_{a_0}]$. Jedno złożenie jest identycznością, a drugie również, gdyż relacja generowana przez $h$ daje $[a_0,uK(h),\id]=[a,u,h]$. Naturalność względem $g:b\to c$ wynika z $g(uK(h))=(gu)K(h)$. Stosując wynik do $K^{\op}$, otrzymujemy wersję kontrawariantną. Jest to także konsekwencja lematu Yonedy: dla $H:\mathcal B\to\Set$ mamy $\Nat(\Hom_{\mathcal A}(a_0,-),HK)\cong H(K(a_0))\cong\Nat(\Hom_{\mathcal B}(K(a_0),-),H)$ przez twierdzenie~\ref{tw:yoneda}. Obie bijekcje są naturalne w $H$.
\end{proof}

\subsection{Przykłady algebraiczne: indukcja i koindukcja}
\begin{twierdzenie}\label{cat:thm:kan-group-modules}
Niech $\varphi:G\to H$ będzie homomorfizmem grup, $k$ ciałem, $BG,BH$ kategoriami jednoobiektowymi o morfizmach odpowiednio $G,H$ i składaniu danym przez mnożenie. Funktory $BG\to\Vect_k$ odpowiadają lewym $kG$-modułom; prekompozycja z $B\varphi$ odpowiada ograniczeniu skalarów wzdłuż $kG\to kH$. Lewe i prawe rozszerzenie Kana mają postać
\[
 \operatorname{Lan}_{B\varphi}V\cong kH\otimes_{kG}V,\qquad
 \operatorname{Ran}_{B\varphi}V\cong\Hom_{kG}(kH,V).
\]
W pierwszym wzorze $kH$ jest prawym $kG$-modułem przez $\varphi$ i lewym $kH$-modułem przez mnożenie. W drugim wzorze $kH$ jest lewym $kG$-modułem przez $\varphi$, a lewe działanie $H$ na Hom określamy przez $(h\cdot f)(x)=f(xh)$.
\end{twierdzenie}
\begin{proof}
Funktor z $BG$ wybiera przestrzeń wektorową $V$ i homomorfizm $G\to\Aut(V)$; przedłużenie liniowe daje działanie $kG$. Transformacje naturalne to dokładnie liniowe morfizmy komutujące z działaniami. To identyfikuje kategorię funktorów z kategorią modułów.

Dla $kH$-modułu $W$ i $kG$-liniowego $a:V\to W$ (z ograniczonym działaniem) definiujemy $\bar a(x\otimes v)=xa(v)$. Jest to dobrze określone, gdyż $x\varphi(g)a(v)=xa(gv)$ dla $g\in G$; przez liniowość to samo zachodzi dla $kG$. Morfizm jest $kH$-liniowy. Odwrotnie, $b:kH\otimes_{kG}V\to W$ wyznacza $a(v)=b(1\otimes v)$, które jest $kG$-liniowe, bo $1\otimes gv=\varphi(g)\otimes v$. Obie konstrukcje są odwrotne, a naturalność wynika z ich wzorów. Jednostką jest $v\mapsto1\otimes v$. Zatem definicja~\ref{cat:def:kan-extension} daje lewe rozszerzenie.

W przypadku prawym $f:kH\to V$ jest $kG$-liniowe; $h\cdot f$ również, bo $(\varphi(g)x)h=\varphi(g)(xh)$. Ponadto $(h_1\cdot(h_2\cdot f))(x)=f(xh_1h_2)$, więc wzór definiuje lewe działanie. Dla $kG$-liniowego $a:W\to V$ definiujemy $\bar a(w)(x)=a(xw)$. Jest to $kG$-liniowe w $x$, a $\bar a(hw)(x)=a(xhw)=(h\cdot\bar a(w))(x)$, więc $\bar a$ jest $kH$-liniowe. Odwrotnie, $b:W\to\Hom_{kG}(kH,V)$ daje $a(w)=b(w)(1)$; mamy $a(\varphi(g)w)=b(w)(\varphi(g))=ga(w)$. Wzory są odwrotne, ponieważ $b(xw)(1)=b(w)(x)$, i naturalne względem morfizmów $V,W$. Kojednostką jest ewaluacja w $1$. Stąd prawe rozszerzenie. Ponieważ $\Vect_k$ jest zupełna i kozupełna (iloczyny i sumy przestrzeni, jądra i kojądra różnic morfizmów, twierdzenie~\ref{cat:thm:limits-products}), oba rozszerzenia są punktowe przez jednoznaczność.
\end{proof}
\begin{przyklad}
Dla homomorfizmu $G\to\{1\}$ powyższe wzory dają odpowiednio koinwarianty $V_G=V/\operatorname{span}_k\{gv-v:g\in G,v\in V\}$ oraz inwarianty $V^G=\{v:gv=v\ \forall g\in G\}$. Istotnie w iloczynie tensorowym $1\otimes gv=1\otimes v$, a funkcja $1\otimes v\mapsto[v]$ i $[v]\mapsto1\otimes v$ są odwrotne. Morfizm $kG$-liniowy $k\to V$ jest wyznaczony przez wartość w $1$, która musi być niezmiennicza; każdy niezmienniczy wektor definiuje taki morfizm. Są to odpowiednio kogranica i granica reprezentacji jako diagramu.
\end{przyklad}
\begin{uwaga}
Wzory z twierdzenia~\ref{cat:thm:kan-group-modules} działają również dla dowolnego homomorfizmu pierścieni z jedynką $R\to S$, jako sprzężenia z funktorem ograniczenia skalarów między kategoriami modułów. Wtedy $S\otimes_R-$ i $\Hom_R(S,-)$ mają te same dowody, z $s\in S$ zamiast $h\in H$ i $r\in R$ zamiast $g\in G$. Aby traktować je jako rozszerzenia Kana w języku funktorów z jednoobiektowych kategorii pierścieni, należy wymagać addytywności: sam zwykły funktor z multiplikatywnego monoidu nie musi respektować dodawania pierścienia. Wersję rozszerzeń w kategoriach wzbogaconych omówimy osobno.
\end{uwaga}

\end{document}
